%% SPDX-License-Identifier: GPL-2.0-or-later
%% Copyright (C) 2026 Nicolas Rooney
%%%%%%%%%%%%%%%%%%%%%%%%%%%%%%%%%%%%%%%%%%%%%%%%%%%%%%%%%%%%%%%%%%%%%%%%%%%%%%
%%% doom.tex -- Doom, implemented entirely in TeX macros.
%%%
%%% Pure pdfTeX.  No shell escape, no LuaTeX, no external programs, no
%%% packages.  Arithmetic is 16.16 fixed point on \numexpr integers -- the
%%% same fixed_t Doom uses -- and every pixel is written by macro expansion
%%% into the PDF content stream as an indexed inline image.
%%%
%%% The renderer follows Doom's own: the BSP is walked front to back,
%%% subsector by subsector, and each seg narrows a pair of per-column clip
%%% arrays.  Two-sided linedefs act as portals, contributing upper and lower
%%% textures and leaving the middle open for what lies beyond.  The game on
%%% top of it is Doom's too -- collision, doors and lifts, pickups, monsters,
%%% missiles, weapons -- each in a section of its own further down.
%%%
%%%   pdftex doom.tex                 sixteen frames of the built-in map
%%%
%%% This file reads top to bottom: configuration, fixed point, trigonometry,
%%% palette, map arrays, renderer, movement, things, the status bar, the
%%% sector machinery, enemies, missiles, sound, firing, and the main loop.
%%% wad.tex loads a real WAD into the same arrays and play.tex is the
%%% interactive loop; README.md says how to run all three and INTERNALS.md
%%% why each part is shaped the way it is.
%%%
%%% Verified with pdfTeX 3.141592653-2.6-1.40.25 (TeX Live) and MiKTeX 25.12.
%%% It needs an eTeX-mode pdftex format -- see the README if \newcount runs
%%% out at 255.
%%%%%%%%%%%%%%%%%%%%%%%%%%%%%%%%%%%%%%%%%%%%%%%%%%%%%%%%%%%%%%%%%%%%%%%%%%%%%%
\catcode`\@=11
\newlinechar=`\^^J

%%% ===========================================================  CONFIGURATION
%%% Override by defining \DTW \DTH \DTPIX \DTFRAMES \DTWAD \DTMAP \DTPLAY
%%% *before* inputting this file.  The README lists every such macro; the
%%% \DTNO... family each switch one feature off and are read where that
%%% feature is declared, or just before the main loop at the foot of the file.
%%%
%%% \DTW x \DTH is the RENDER resolution and therefore the PDF's.  In
%%% interactive mode \DTASC{n} shrinks the terminal viewport to n x 2n pixels
%%% a cell without shrinking the render, which is how the recording gets a
%%% resolution the terminal could never show; \DTPIPE hands the frame to a
%%% front-end instead of drawing it.  Both live in play.tex.
\newcount\SCRW    \SCRW=200
\newcount\SCRH    \SCRH=120
\newcount\PIXSZ   \PIXSZ=3
\newcount\NFRAMES \NFRAMES=16
\expandafter\ifx\csname DTW\endcsname\relax\else\SCRW=\DTW\relax\fi
\expandafter\ifx\csname DTH\endcsname\relax\else\SCRH=\DTH\relax\fi
\expandafter\ifx\csname DTPIX\endcsname\relax\else\PIXSZ=\DTPIX\relax\fi
\expandafter\ifx\csname DTFRAMES\endcsname\relax\else\NFRAMES=\DTFRAMES\relax\fi

%%% Plain TeX's \loop cannot nest: every \loop redefines the same \body.  This
%%% tail-recursive version consumes its \fi before recursing, so it nests
%%% freely and does not grow the conditional stack.
\long\def\dtWhile#1#2{\ifnum#1 \expandafter\dt@wY\else\expandafter\dt@wN\fi
  {#1}{#2}}
\long\def\dt@wY#1#2{#2\dtWhile{#1}{#2}}
\long\def\dt@wN#1#2{}

%%% =============================================================  FIXED POINT
%%% eTeX evaluates the scaling operation a*b/c with a 64-bit intermediate,
%%% which is exactly what FixedMul and FixedDiv need.  (\dimexpr would be the
%%% obvious home for 16.16, but dimensions cap at 16383.99998pt and cannot
%%% hold Doom's map coordinates.)
\def\fmul#1#2{\number\numexpr(#1)*(#2)/65536\relax}
\def\fdiv#1#2{\number\numexpr(#1)*65536/(#2)\relax}

\newcount\tmpa \newcount\tmpb \newcount\tmpc \newcount\tmpd \newcount\tmpi
\newcount\tmpj \newcount\tmpk

\newcount\isq@r \newcount\isq@n
\def\dtIsqrt#1#2{%
  \isq@n=\numexpr#2\relax
  \ifnum\isq@n<2 #1=\isq@n
  \else
    \isq@r=\isq@n \divide\isq@r by2 \tmpk=0
    \dtWhile{\tmpk=0}{%
      \tmpa=\isq@n \divide\tmpa by\isq@r \advance\tmpa by\isq@r
      \divide\tmpa by2
      \ifnum\tmpa<\isq@r \isq@r=\tmpa \else\tmpk=1 \fi}%
    #1=\isq@r
  \fi}

%%% ====================================================================  TRIG
%%% 8192 units to a full turn, as in Doom's FINEANGLES.
\def\dt@taylorsin#1{%
  \tmpa=#1\relax
  \tmpb=\fmul{\the\tmpa}{\the\tmpa}\relax
  \tmpc=\fmul{\the\tmpa}{\the\tmpb}\relax
  \tmpd=\the\tmpa
  \advance\tmpd by-\numexpr\the\tmpc/6\relax
  \tmpc=\fmul{\the\tmpc}{\the\tmpb}\relax
  \advance\tmpd by\numexpr\the\tmpc/120\relax
  \tmpc=\fmul{\the\tmpc}{\the\tmpb}\relax
  \advance\tmpd by-\numexpr\the\tmpc/5040\relax
  \tmpc=\fmul{\the\tmpc}{\the\tmpb}\relax
  \advance\tmpd by\numexpr\the\tmpc/362880\relax
  \tmpc=\fmul{\the\tmpc}{\the\tmpb}\relax
  \advance\tmpd by-\numexpr\the\tmpc/39916800\relax}

\message{[doom.tex] sine table...}
\tmpi=0
\dtWhile{\tmpi<2049}{%
  \dt@taylorsin{\numexpr\the\tmpi*102944/2048\relax}%
  \expandafter\xdef\csname dt@s\the\tmpi\endcsname{\the\tmpd}%
  \advance\tmpi by1\relax }
\expandafter\xdef\csname dt@s2048\endcsname{65536}

%%% \numexpr DIVISION ROUNDS, so modular arithmetic must use \divide.
\def\dtnorm#1#2{#1=\numexpr#2\relax
  \tmpa=#1\divide\tmpa by8192 \multiply\tmpa by8192 \advance#1 by-\tmpa
  \ifnum#1<0 \advance#1 by8192 \fi}
\def\dtsin#1{\dt@sin{\number\numexpr#1\relax}}
\def\dt@sin#1{%
  \ifnum#1>4096
    -\csname dt@s\ifnum#1>6144 \number\numexpr8192-(#1)\relax
       \else\number\numexpr(#1)-4096\relax\fi\endcsname
  \else
    \csname dt@s\ifnum#1>2048 \number\numexpr4096-(#1)\relax
       \else\number\numexpr#1\relax\fi\endcsname
  \fi}
\def\dtcos#1{\dt@sin{\number\numexpr(#1)+2048
   -\ifnum\numexpr(#1)+2048\relax>8191 8192\else0\fi\relax}}

%%% ---- arctangent ----------------------------------------------------------
%%% Only needed to pick a sprite rotation, so it is built from the sine table
%%% rather than by a second series: walk the eighth-turn 0..1024, take each
%%% angle's tangent scaled to 1024, and fill every slot up to it.  The
%%% tangent is monotonic there, so one pass fills the table exactly and no
%%% division by a vanishing cosine ever happens at render time.
\newcount\dtATA \newcount\dtATT \newcount\dtATP
\message{[doom.tex] arctangent table...}
\dtATA=0 \dtATP=0
\dtWhile{\dtATA<1025}{%
  \dtATT=\numexpr\dtsin{\the\dtATA}*1024/\dtcos{\the\dtATA}\relax
  \ifnum\dtATT>1024 \dtATT=1024 \fi
  \dtWhile{\dtATP<\numexpr\dtATT+1\relax}{%
    \expandafter\xdef\csname dtAT\the\dtATP\endcsname{\the\dtATA}%
    \advance\dtATP by1\relax }%
  \advance\dtATA by1\relax }

%%% Angle from the origin to (#1,#2) in 0..8191, by folding onto that eighth
%%% turn: the table covers 0..45 degrees and everything else is a reflection.
%%
%% \dtATX and \dtATY are this macro's own.  They were \dtAX and \dtAY, which
%% the collision scan also declares and uses for a seg's first vertex -- two
%% \newcounts of one name, so the later one quietly won and every caller of
%% \dtAtan had to take its bearing before going anywhere near a move test.
%% Four comments in this file said so.  A register each costs nothing.
\newcount\dtATX \newcount\dtATY \newcount\dtANG \newcount\dtAQ
\def\dtAtan#1#2{%
  \dtATX=\numexpr#1\relax \dtATY=\numexpr#2\relax
  \dtAQ=0
  \ifnum\dtATX<0 \dtATX=-\dtATX \advance\dtAQ by1 \fi
  \ifnum\dtATY<0 \dtATY=-\dtATY \advance\dtAQ by2 \fi
  \ifnum\dtATX=0 \ifnum\dtATY=0 \dtATX=1 \fi\fi
  \ifnum\dtATY>\dtATX
    \dtANG=\numexpr2048-\csname dtAT\number\numexpr\dtATX*1024/\dtATY\relax
      \endcsname\relax
  \else
    \dtANG=\csname dtAT\number\numexpr\dtATY*1024/\dtATX\relax\endcsname\relax
  \fi
  \ifcase\dtAQ
  \or \dtANG=\numexpr4096-\dtANG\relax
  \or \dtANG=\numexpr8192-\dtANG\relax
  \or \dtANG=\numexpr4096+\dtANG\relax
  \fi
  \ifnum\dtANG>8191 \advance\dtANG by-8192 \fi}

%%% ==============================================================  HEX OUTPUT
%%% Built by construction (q*16+r) so no truncating division is needed.
\def\dt@hd#1{\ifcase\numexpr#1\relax 0\or1\or2\or3\or4\or5\or6\or7\or8\or9\or
  A\or B\or C\or D\or E\or F\fi}
\newcount\dt@q \newcount\dt@r
\dt@q=0 \dt@r=0 \tmpi=0
\dtWhile{\tmpi<256}{%
  \expandafter\xdef\csname dt@h\the\tmpi\endcsname
    {\dt@hd{\the\dt@q}\dt@hd{\the\dt@r}}%
  \advance\dt@r by1
  \ifnum\dt@r>15 \dt@r=0 \advance\dt@q by1 \fi
  \advance\tmpi by1\relax }
\def\hx#1{\csname dt@h\number\numexpr#1\relax\endcsname}

%%% Only the first argument needs pre-expansion here, so the single
%%% \expandafter idiom is safe.
\def\Rep#1#2{\ifnum\numexpr#1\relax>0 #2\expandafter\dt@repN\else
  \expandafter\dt@repS\fi{#1}{#2}}
\def\dt@repN#1#2{\expandafter\Rep\expandafter{\number\numexpr(#1)-1\relax}{#2}}
\def\dt@repS#1#2{}

%%% ==================================================================  PALETTE
\newcount\pj@
\gdef\DTPAL{}
\def\dt@pe#1#2#3{\xdef\DTPAL{\DTPAL\hx{#1}\hx{#2}\hx{#3}}}
\def\dt@ramp#1#2#3#4#5#6{\pj@=0
  \dtWhile{\pj@<16}{%
    \dt@pe{\numexpr#1+((#4)-(#1))*\the\pj@/15\relax}%
          {\numexpr#2+((#5)-(#2))*\the\pj@/15\relax}%
          {\numexpr#3+((#6)-(#3))*\the\pj@/15\relax}%
    \advance\pj@ by1\relax }}
\message{[doom.tex] palette...}
\dt@ramp{255}{255}{255}{15}{15}{15}     %   0 grey
\dt@ramp{239}{223}{207}{31}{23}{11}     %  16 bone
\dt@ramp{255}{183}{119}{35}{19}{7}      %  32 tan
\dt@ramp{255}{31}{31}{27}{0}{0}         %  48 red
\dt@ramp{175}{127}{87}{23}{15}{9}       %  64 brown
\dt@ramp{119}{255}{111}{7}{31}{7}       %  80 green
\dt@ramp{143}{175}{255}{7}{11}{35}      %  96 blue
\dt@ramp{255}{235}{155}{43}{31}{7}      % 112 gold
\dt@ramp{207}{155}{99}{27}{19}{11}      % 128 wood
\dt@ramp{163}{163}{171}{17}{17}{21}     % 144 blue-grey
\dt@ramp{255}{119}{0}{31}{13}{0}        % 160 orange
\dt@ramp{171}{139}{163}{21}{17}{21}     % 176 mauve
\dt@ramp{111}{131}{147}{13}{17}{21}     % 192 steel
\dt@ramp{255}{255}{115}{35}{35}{11}     % 208 yellow
\dt@ramp{95}{95}{95}{9}{9}{9}           % 224 dark grey
\dt@ramp{63}{47}{35}{5}{4}{3}           % 240 near black

%%% ===========================================================  WALL TEXTURE
%%% 64 x 65 shade offsets 0..7; row 64 duplicates row 0 to absorb rounding.
\newcount\tu \newcount\tv \newcount\tsh \newcount\tq \newcount\tr
\message{[doom.tex] wall texture...}
\tu=0 \tv=0 \tmpi=0
\dtWhile{\tmpi<4160}{%
  \tq=\tv \divide\tq by16
  \tr=\numexpr\tv-\tq*16\relax
  \tmpj=\numexpr\tu+\ifodd\tq 16\else0\fi\relax
  \tmpk=\tmpj \divide\tmpk by32
  \tmpk=\numexpr\tmpj-\tmpk*32\relax
  \ifnum\tr<2 \tsh=6 \else
    \ifnum\tmpk<2 \tsh=6 \else
      \tsh=\numexpr1+\tr/6\relax
      \ifodd\tmpj \advance\tsh by1 \fi
    \fi
  \fi
  \expandafter\xdef\csname dtT\the\tu @\the\tv\endcsname{\the\tsh}%
  \advance\tu by1
  \ifnum\tu>63 \tu=0 \advance\tv by1 \fi
  \advance\tmpi by1\relax }

%%% ------------------------------------------------------- DOOR TEXTURE
%%% A door has to read as a door at a glance.  Without TEXTURE1 composition
%%% there is no BIGDOOR2 to sample, so the cue has to be structural: the
%%% brick table above bands horizontally, and this one bands vertically --
%%% ribbed panels, a centre seam, a frame edge, and a single rail where the
%%% 64-unit tile repeats so the door has a top and a bottom rather than
%%% being infinitely corrugated.  Same 64 x 65 shape and the same 0..7 range
%%% of shade offsets, so selecting between the two is a change of name and
%%% nothing else.
\message{[doom.tex] door texture...}
\tu=0 \tv=0 \tmpi=0
\dtWhile{\tmpi<4160}{%
  \tq=\tu \divide\tq by8 \multiply\tq by8
  \tr=\numexpr\tu-\tq\relax
  \ifcase\tr\relax
    \tsh=6 \or \tsh=4 \or \tsh=2 \or \tsh=1 \or
    \tsh=1 \or \tsh=2 \or \tsh=3 \else \tsh=5
  \fi
  \ifnum\tu<2 \tsh=6 \fi
  \ifnum\tu>61 \tsh=6 \fi
  \ifnum\tu>30 \ifnum\tu<33 \tsh=6 \fi\fi
  \ifnum\tv<2 \tsh=5 \fi
  \ifnum\tv>62 \tsh=5 \fi
  \expandafter\xdef\csname dtD\the\tu @\the\tv\endcsname{\the\tsh}%
  \advance\tu by1
  \ifnum\tu>63 \tu=0 \advance\tv by1 \fi
  \advance\tmpi by1\relax }

%%% =====================================================================  MAP
%%% Arrays in csnames, so map size is bounded by hash and main memory, not by
%%% register count.  wad.tex fills exactly these from a real WAD.
\newcount\NVERT \newcount\NLINE \newcount\NSEC
\newcount\NSEG  \newcount\NSSEC \newcount\NNODE \newcount\DTROOT
\NVERT=0 \NLINE=0 \NSEC=0 \NSEG=0 \NSSEC=0 \NNODE=0

\def\dtAddVert#1#2{%
  \expandafter\xdef\csname dtVx\the\NVERT\endcsname{#1}%
  \expandafter\xdef\csname dtVy\the\NVERT\endcsname{#2}%
  \global\advance\NVERT by1\relax}
%% front sector, back sector (-1 for a one-sided wall)
\def\dtAddLine#1#2{%
  \expandafter\xdef\csname dtLF\the\NLINE\endcsname{#1}%
  \expandafter\xdef\csname dtLB\the\NLINE\endcsname{#2}%
  \global\advance\NLINE by1\relax}
%%% ---- sidedefs and the texture references that hang off them --------------
%%% A seg reaches its textures the way Doom does: seg -> linedef -> side ->
%%% sidedef.  The built-in map has no sidedefs at all and never calls these,
%%% which is what keeps it on the procedural shade tables.
%%
%% sector, upper / lower / middle texture index (-1 = none), x and y offset
\newcount\NSIDE \NSIDE=0
\def\dtAddSide#1#2#3#4#5#6{%
  \expandafter\xdef\csname dtYs\the\NSIDE\endcsname{#1}%
  \expandafter\xdef\csname dtYu\the\NSIDE\endcsname{#2}%
  \expandafter\xdef\csname dtYl\the\NSIDE\endcsname{#3}%
  \expandafter\xdef\csname dtYm\the\NSIDE\endcsname{#4}%
  \expandafter\xdef\csname dtYx\the\NSIDE\endcsname{#5}%
  \expandafter\xdef\csname dtYy\the\NSIDE\endcsname{#6}%
  \global\advance\NSIDE by1\relax}
%% linedef -> its two sidedefs, and the two pegging bits that decide where a
%% texture's v origin sits: 1 = DONTPEGTOP, 2 = DONTPEGBOTTOM.  Kept apart
%% from \dtAddLine so the built-in map's two-argument calls still work.
\def\dtAddLineSide#1#2#3#4{%
  \expandafter\xdef\csname dtLR#1\endcsname{#2}%
  \expandafter\xdef\csname dtLL#1\endcsname{#3}%
  \expandafter\xdef\csname dtLG#1\endcsname{#4}}
%% Distance from the linedef's first vertex to this seg's first vertex.  Two
%% segs of one linedef must continue the same texture, not restart it.
\def\dtAddSegOff#1#2{\expandafter\xdef\csname dtGo#1\endcsname{#2}}
%% Byte 8 of a linedef: the tag naming which sectors it acts on.  A manual
%% door acts on the line's own back sector and carries tag 0; everything else
%% acts on *every* sector whose tag matches, which is why this cannot be
%% inferred from the geometry.  Only nonzero tags are stored, so \dtLineTag
%% answers 0 for the built-in map without it needing any changes.
\def\dtAddLineTag#1#2{\expandafter\xdef\csname dtLT#1\endcsname{#2}}
\def\dtLineTag#1{\ifcsname dtLT#1\endcsname\csname dtLT#1\endcsname\else0\fi}

%% floor, ceiling, light 0..255, wall ramp, floor ramp, ceiling ramp
\def\dtAddSector#1#2#3#4#5#6{%
  \expandafter\xdef\csname dtSf\the\NSEC\endcsname{#1}%
  \expandafter\xdef\csname dtSc\the\NSEC\endcsname{#2}%
  \expandafter\xdef\csname dtSl\the\NSEC\endcsname{#3}%
  \expandafter\xdef\csname dtSw\the\NSEC\endcsname{#4}%
  \expandafter\xdef\csname dtSp\the\NSEC\endcsname{#5}%
  \expandafter\xdef\csname dtSq\the\NSEC\endcsname{#6}%
  \global\advance\NSEC by1\relax}
%% Byte 24 of a sector: the tag a tagged linedef names it by.  Stored only
%% when nonzero, and set after \dtAddSector because the sector index it
%% belongs to is the one that call just consumed.
\def\dtAddSectorTag#1#2{\expandafter\xdef\csname dtSt#1\endcsname{#2}}
\def\dtSecTag#1{\ifcsname dtSt#1\endcsname\csname dtSt#1\endcsname\else0\fi}
%% Byte 22 of a sector: its special -- a light that strobes, a floor that
%% burns, a secret.  Stored only when nonzero, on the tag's own argument.
\def\dtAddSectorSpec#1#2{\expandafter\xdef\csname dtSs#1\endcsname{#2}}
\def\dtSecSpec#1{\ifcsname dtSs#1\endcsname\csname dtSs#1\endcsname\else0\fi}
%% v1, v2, linedef, side
\def\dtAddSeg#1#2#3#4{%
  \expandafter\xdef\csname dtGa\the\NSEG\endcsname{#1}%
  \expandafter\xdef\csname dtGb\the\NSEG\endcsname{#2}%
  \expandafter\xdef\csname dtGl\the\NSEG\endcsname{#3}%
  \expandafter\xdef\csname dtGs\the\NSEG\endcsname{#4}%
  \dtIsqrt\tmpa{(\csname dtVx#2\endcsname-(\csname dtVx#1\endcsname))*%
                (\csname dtVx#2\endcsname-(\csname dtVx#1\endcsname))+%
                (\csname dtVy#2\endcsname-(\csname dtVy#1\endcsname))*%
                (\csname dtVy#2\endcsname-(\csname dtVy#1\endcsname))}%
  \expandafter\xdef\csname dtGn\the\NSEG\endcsname{\the\tmpa}%
  \global\advance\NSEG by1\relax}
%% seg count, first seg
\def\dtAddSSec#1#2{%
  \expandafter\xdef\csname dtCn\the\NSSEC\endcsname{#1}%
  \expandafter\xdef\csname dtCf\the\NSSEC\endcsname{#2}%
  \global\advance\NSSEC by1\relax}
%% x, y, dx, dy, right child, left child (child >= 32768 means subsector)
\def\dtAddNode#1#2#3#4#5#6{%
  \expandafter\xdef\csname dtNx\the\NNODE\endcsname{#1}%
  \expandafter\xdef\csname dtNy\the\NNODE\endcsname{#2}%
  \expandafter\xdef\csname dtNu\the\NNODE\endcsname{#3}%
  \expandafter\xdef\csname dtNv\the\NNODE\endcsname{#4}%
  \expandafter\xdef\csname dtNr\the\NNODE\endcsname{#5}%
  \expandafter\xdef\csname dtNl\the\NNODE\endcsname{#6}%
  \global\advance\NNODE by1\relax}

%% x, y, type, floor height of the sector it stands in
\newcount\NTHING \NTHING=0
\def\dtAddThing#1#2#3#4{%
  \expandafter\xdef\csname dtHx\the\NTHING\endcsname{#1}%
  \expandafter\xdef\csname dtHy\the\NTHING\endcsname{#2}%
  \expandafter\xdef\csname dtHt\the\NTHING\endcsname{#3}%
  \expandafter\xdef\csname dtHz\the\NTHING\endcsname{#4}%
  \global\advance\NTHING by1\relax}

%%% Built-in level: three rooms in a row joined by portals that step both up
%%% and down, so upper and lower textures both appear.  Each room is convex,
%%% so the BSP is two nodes.  A WAD replaces all of this.
\def\dtBuiltinMap{%
  \dtAddVert{0}{0}\dtAddVert{0}{512}\dtAddVert{384}{512}\dtAddVert{384}{352}%
  \dtAddVert{384}{160}\dtAddVert{384}{0}\dtAddVert{640}{512}%
  \dtAddVert{640}{352}\dtAddVert{640}{160}\dtAddVert{640}{0}%
  \dtAddVert{1024}{512}\dtAddVert{1024}{0}%
  \dtAddSector{0}{128}{224}{64}{32}{176}%
  \dtAddSector{24}{96}{176}{128}{112}{192}%
  \dtAddSector{-16}{144}{128}{192}{80}{96}%
  \dtAddLine{0}{-1}\dtAddLine{0}{-1}\dtAddLine{0}{-1}\dtAddLine{0}{1}%
  \dtAddLine{0}{-1}\dtAddLine{0}{-1}\dtAddLine{1}{-1}\dtAddLine{1}{-1}%
  \dtAddLine{1}{-1}\dtAddLine{1}{-1}\dtAddLine{1}{2}\dtAddLine{1}{-1}%
  \dtAddLine{1}{-1}\dtAddLine{2}{-1}\dtAddLine{2}{-1}\dtAddLine{2}{-1}%
  \dtAddLine{2}{-1}\dtAddLine{2}{-1}%
  %% subsector 0: west room
  \dtAddSeg{0}{1}{0}{0}\dtAddSeg{1}{2}{1}{0}\dtAddSeg{2}{3}{2}{0}%
  \dtAddSeg{3}{4}{3}{0}\dtAddSeg{4}{5}{4}{0}\dtAddSeg{5}{0}{5}{0}%
  \dtAddSSec{6}{0}%
  %% subsector 1: middle room; its second seg is the west portal, back side
  \dtAddSeg{5}{4}{6}{0}\dtAddSeg{4}{3}{3}{1}\dtAddSeg{3}{2}{7}{0}%
  \dtAddSeg{2}{6}{8}{0}\dtAddSeg{6}{7}{9}{0}\dtAddSeg{7}{8}{10}{0}%
  \dtAddSeg{8}{9}{11}{0}\dtAddSeg{9}{5}{12}{0}%
  \dtAddSSec{8}{6}%
  %% subsector 2: east room
  \dtAddSeg{9}{8}{13}{0}\dtAddSeg{8}{7}{10}{1}\dtAddSeg{7}{6}{14}{0}%
  \dtAddSeg{6}{10}{15}{0}\dtAddSeg{10}{11}{16}{0}\dtAddSeg{11}{9}{17}{0}%
  \dtAddSSec{6}{14}%
  %% node 0 splits west from middle; node 1, the root, splits off the east
  \dtAddNode{384}{0}{0}{512}{32769}{32768}%
  \dtAddNode{640}{0}{0}{512}{32770}{0}%
  \DTROOT=1 }

%%% ================================================================  RENDERER
\newcount\PPLANE \newcount\HALFH
\newcount\NEARZ  \NEARZ=327680
\newcount\PX \newcount\PY \newcount\PZ \newcount\PA
\newcount\vsin \newcount\vcos
\newcount\ra \newcount\rb \newcount\rc \newcount\rd \newcount\re \newcount\rf
\newcount\d@one \newcount\d@two \newcount\l@one \newcount\l@two
\newcount\x@one \newcount\x@two \newcount\s@one \newcount\s@two
\newcount\u@one \newcount\u@two \newcount\cx \newcount\wid
\newcount\dtOPEN \newcount\dtSIDE \newcount\dtSA \newcount\dtSB
\newcount\dtG \newcount\dtLD \newcount\dtSDE \newcount\dtFS \newcount\dtBS
\newcount\ct \newcount\cb \newcount\yt \newcount\yb \newcount\ut \newcount\lb
\newcount\ytr \newcount\ybr \newcount\utr \newcount\lbr
\newcount\vstep \newcount\vstepr \newcount\vstart \newcount\wbase
\newcount\wlight
%% Real wall textures.  \dtCMN is how many COLORMAP light levels are loaded;
%% 0 means no WAD at all, in which case nothing below is ever consulted.
\newcount\dtCMN \dtCMN=0
\newcount\dtSDI \newcount\dtWTU \newcount\dtWTL \newcount\dtWTM
\newcount\dtWT  \newcount\dtTU  \newcount\dtTV  \newcount\dtTW
\newcount\dtTG
\newcount\dtUOFF \newcount\dtVOFF \newcount\dtPEGF \newcount\dtFAKE
\newcount\dtWCM
\newif\ifdt@realtex
%% Dedicated to seg drawing.  These used to be \tmpi/\tmpj/\tmpk, which the
%% frame loop also used -- so drawing a seg reset the frame counter and the
%% camera walked off to infinity.
\newcount\dtVA \newcount\dtVB \newcount\dtSKIP \newcount\dtFRAME

%%% Per-column clip arrays: rows [CT,CB) are still open.  eTeX gives 32768
%%% registers, addressed directly (plain TeX's \newcount caps at 255).
\def\CT#1{\count\numexpr1000+(#1)\relax}
\def\CB#1{\count\numexpr2000+(#1)\relax}
%%% Bank 3000: the scale at which column x was closed by a solid wall, which
%%% is the depth a sprite must beat to be visible in that column.  Scale is
%%% inverse depth, so larger is nearer and 0 means nothing has closed it.
\def\CD#1{\count\numexpr3000+(#1)\relax}
%%% Banks 4000/5000: the world direction of column x's ray, scaled by 4096,
%%% which is what turns a plane pixel's depth into its world position.  Bank
%%% 6000: the same column's sky texture column.  All three are rebuilt once a
%%% frame by \dtPlaneFrame, because all three depend on the facing.
\def\PU#1{\count\numexpr4000+(#1)\relax}
\def\PV#1{\count\numexpr5000+(#1)\relax}
\def\PK#1{\count\numexpr6000+(#1)\relax}
%%% Bank 7000: which of the eight texels in that group, since the sky texture
%%% is packed along u exactly as every other texture is.
\def\PJ#1{\count\numexpr7000+(#1)\relax}

%%% ---- silhouettes ---------------------------------------------------------
%%% \CD answers one question -- "is this sprite behind the wall that closed
%%% this column" -- and a portal never closes a column, so until a door had
%%% shut completely or a step reached the ceiling neither hid anything at all.
%%% A door part way open drew its whole slab and every sprite behind it came
%%% through, and a sprite standing behind a raised platform showed through the
%%% platform.
%%%
%%% Doom keeps a drawseg per wall carrying its own clip arrays and a
%%% silhouette flag, and clips each sprite against every drawseg in front of
%%% it.  The equivalent here is to record, per column, every narrowing of CT
%%% and CB along with the scale at which it happened.  Traversal is front to
%%% back, so within a column those records already run near to far: the clip
%%% in force at a given depth is simply the last record nearer than it, and
%%% the lookup is one forward walk with no sorting and no search.
%%%
%%% The lists are pairs {scale}{row}, and only an actual narrowing is
%%% recorded -- a portal that changes nothing in a column adds nothing, so on
%%% ordinary geometry these stay one or two entries long.
\def\dt@silT#1#2#3{%
  \expandafter\xdef\csname dtZT#1\endcsname{\csname dtZT#1\endcsname{#2}{#3}}}
\def\dt@silB#1#2#3{%
  \expandafter\xdef\csname dtZB#1\endcsname{\csname dtZB#1\endcsname{#2}{#3}}}

%%% Per-column pixel spans.  Front-to-back rendering fills a column from the
%%% outside in, so the top half only ever grows downward and the bottom half
%%% only ever grows upward: append to one, prepend to the other, and nothing
%%% ever needs random access into the middle of a column.
\def\dtAppTop#1#2{%
  \expandafter\xdef\csname dtTS#1\endcsname{\csname dtTS#1\endcsname#2}}
\def\dtPreBot#1#2{%
  \expandafter\xdef\csname dtBS#1\endcsname{#2\csname dtBS#1\endcsname}}

\def\dtClearClip{%
  \cx=0
  \dtWhile{\cx<\SCRW}{%
    \CT\cx=0 \CB\cx=\SCRH \CD\cx=0
    \expandafter\gdef\csname dtTS\the\cx\endcsname{}%
    \expandafter\gdef\csname dtBS\the\cx\endcsname{}%
    \expandafter\gdef\csname dtZT\the\cx\endcsname{}%
    \expandafter\gdef\csname dtZB\the\cx\endcsname{}%
    \advance\cx by1\relax }%
  \dtOPEN=\SCRW}

%%% Row shading for floors and ceilings: distance from the horizon stands in
%%% for depth, which is what makes a flat plane fade into the distance.
%%
%% dtYR, and it was dtRS until the ceiling movers arrived and took the same
%% name for a sector's state: marking a crusher on sector 3 then repainted
%% screen row 3 of every procedural plane, one shade wrong, for the rest of
%% the run.  Register aliasing again, in a \csname this time, where no
%% \newcount is there to object.
\def\dtPlaneShade{%
  \tmpi=0
  \dtWhile{\tmpi<\SCRH}{%
    \tmpa=\numexpr\tmpi-\HALFH\relax
    \ifnum\tmpa<0 \tmpa=-\tmpa \fi
    \tmpb=\numexpr13-\tmpa*30/\HALFH\relax
    \ifnum\tmpb<0 \tmpb=0 \fi \ifnum\tmpb>13 \tmpb=13 \fi
    \expandafter\xdef\csname dtYR\the\tmpi\endcsname{\the\tmpb}%
    \advance\tmpi by1\relax }}

%%% Expandable plane span: rows #1..#2-1 at ramp base #3.
\def\PlaneRun#1#2#3{%
  \ifnum\numexpr#1\relax<\numexpr#2\relax
    \hx{\numexpr(#3)+\csname dtYR\number\numexpr#1\relax\endcsname\relax}%
    \expandafter\dt@prN\else\expandafter\dt@prS\fi{#1}{#2}{#3}}
\def\dt@prN#1#2#3{\expandafter\PlaneRun\expandafter
  {\number\numexpr(#1)+1\relax}{#2}{#3}}
\def\dt@prS#1#2#3{}

%%% ---- real flats ----------------------------------------------------------
%%% The same span over a 64x64 flat read from the WAD.  A plane pixel needs
%%% its world position, which a wall pixel does not, and that is the whole
%%% difference: Doom gets it by walking horizontal spans, where depth is
%%% constant and (x,y) step linearly, but this renderer is column-wise and
%%% cannot.  What makes the column-wise form affordable is that along ONE
%%% column the world point is a scalar multiple of the depth:
%%%
%%%   lateral/depth = (PPLANE-x)/PPLANE  =: k     -- fixed for the column
%%%   world = (PX,PY) + depth * (cos - k sin, sin + k cos)
%%%
%%% so the two direction terms are per-column constants (banks 4000/5000,
%%% built once a frame) and each pixel costs one multiply-add per axis.
%%%
%%% Depth itself is planeheight*PPLANE/(y-HALFH), and the second factor
%%% depends on nothing but the row -- so dtYS<y> holds it for the whole run
%%% and no pixel ever divides.  Doom's own half-pixel bias is in there too,
%%% which is also what removes the singularity at the horizon row.
%%%
%%% Units: everything is 12.4 (sixteenths of a map unit).  16.16 would
%%% overflow on the near-horizon rows, and whole units make the texture jump
%%% a texel as the camera moves.  \dt@uxq is the direction scaled by 4096, so
%%% D*UX/4096 is again 12.4 -- and that product needs \numexpr's 64-bit
%%% intermediate, which is exactly what the a*b/c form gets.
\def\dt@fkey{dtFL0@}
\def\dt@phq{0}
\def\dt@psm{0}
\def\dt@pxq{0}\def\dt@pyq{0}
\def\dt@uxq{0}\def\dt@uyq{0}
\def\PlaneTex#1#2{%
  \ifnum\numexpr#1\relax<\numexpr#2\relax
    \expandafter\dt@ptD\expandafter{\number\numexpr
      \dt@phq*\csname dtYS\number\numexpr#1\relax\endcsname/256\relax}%
    \expandafter\dt@ptN\else\expandafter\dt@ptS\fi{#1}{#2}}
\def\dt@ptN#1#2{\expandafter\PlaneTex\expandafter
  {\number\numexpr(#1)+1\relax}{#2}}
\def\dt@ptS#1#2{}
%% #1 is the depth, wanted three times over -- for u, for v and for the light
%% -- so each stage hands it on rather than recomputing it.  A macro's
%% arguments are the only expandable place to put a temporary.
\def\dt@ptD#1{\expandafter\dt@ptU\expandafter{\number\numexpr
  \dt@pxq+#1*\dt@uxq/4096\relax}{#1}}
\def\dt@ptU#1#2{\expandafter\dt@ptV\expandafter{\number\numexpr
  \dt@pyq+#2*\dt@uyq/4096\relax}{#1}{#2}}
%% Doom's zlight, near enough: level = startmap - 80/(j+1) with j the depth
%% in sixteenths of a unit, which in 12.4 is 20480/(D+256).  It runs from
%% startmap-80 up against startmap, so the result needs clamping into the
%% table -- and \ifnum inside a \csname is a trap, so dtCC does it as a
%% lookup instead, biased by 128 to keep the key non-negative.
%%
%% The 20480 lives in a macro for the same reason \DOORRAMP is a count: a
%% fixture can set it to zero and take the distance term out, and comparing
%% that render against the real one is the only thing that says the term is
%% doing any work.  Counting how many colours a plane pass reaches does not
%% -- 30 flats across 59 sector lights reach plenty with the depth term
%% deleted.
%%
%% The texel arrives packed eight to a control sequence (see wad.tex on why),
%% so u splits into a group and an offset within it.  THREE \expandafter, not
%% one: \csname yields the control sequence \dtFL7@3@2, not its body, and a
%% macro's arguments are not expanded while being grabbed -- so one gets
%% \dt@fpick the token itself as its first pixel and it eats \endcsname and
%% everything after it looking for the other seven.  The first expansion
%% builds the name, the second expands what it named.
%% Every line end in here is either after a control word or commented out.  A
%% space after a CHARACTER is a real token and would land inside the key --
%% "dtCM12@ 4F" is undefined, \csname quietly makes it \relax, and \relax
%% then travels into the column as if it were a pixel.
\def\dt@fpick#1#2#3#4#5#6#7#8#9{%
  \ifcase#9 #1\or#2\or#3\or#4\or#5\or#6\or#7\else#8\fi}
\def\dt@plit{20480}
\def\dt@ptV#1#2#3{%
  \csname dtCM\csname dtCC\number\numexpr\dt@psm-\dt@plit/(#3+256)+128\relax
    \endcsname
    @\expandafter\expandafter\expandafter\dt@fpick\csname\dt@fkey
    \number\numexpr(#1-8)/16-64*((#1-512)/1024)\relax
    @\number\numexpr(#2-64)/128-8*((#2-512)/1024)\relax\endcsname
    {\number\numexpr(#2-8)/16-8*((#2-64)/128)\relax}\endcsname}

%%% ---- sky -----------------------------------------------------------------
%%% A sky plane is not a plane at all: Doom draws it as a vertical strip of a
%%% wall texture picked by view angle, with no perspective and no shading, so
%%% that it never moves as the player walks and never gets darker.
%%%
%%% The sky lives in its own packed array (wad.tex says why) keyed by texture
%%% column, so \dt@skyk -- the column, and the offset within its group -- is a
%%% per-column constant and only the row moves.  That makes this the cheapest
%%% span in the renderer: two table lookups and a pick, no arithmetic at all
%%% beyond the row counter.
%% Its own pick, not \dt@pick: a portal's ceiling run and its upper wall run
%% land in the SAME \edef, and \dt@usetex has already pointed \dt@pick at the
%% wall's texture column by then.
\def\dt@skyk{dtX0@0@}
\let\dt@spick\dt@pick
\def\SkyRun#1#2{%
  \ifnum\numexpr#1\relax<\numexpr#2\relax
    \expandafter\expandafter\expandafter\dt@spick
      \csname\dt@skyk\csname dtSV\number\numexpr#1\relax\endcsname\endcsname
    \expandafter\dt@skN\else\expandafter\dt@skS\fi{#1}{#2}}
\def\dt@skN#1#2{\expandafter\SkyRun\expandafter
  {\number\numexpr(#1)+1\relax}{#2}}
\def\dt@skS#1#2{}

%%% ---- the tables the two runs above read -----------------------------------
%%% dtYS<y>  PPLANE*512/(2|y-HALFH|+1) -- the row half of a plane's depth.
%%%          Doom's own half-pixel bias is the +1, and it is also what stops
%%%          the horizon row dividing by zero.
%%% dtCC<n>  clamp n-128 into the loaded COLORMAP.  A lookup rather than two
%%%          \ifnum because this is read from inside a \csname, where a
%%%          conditional's terminating space would land in the key.
%%% dtKA<x>  the angle of column x's ray away from the view centre, for the
%%%          sky.  Depends on nothing but the projection, so it is built here
%%%          rather than per frame.
%%% dtSV<y>  the sky texture row for screen row y.  Doom pins texture row 100
%%%          to the horizon and steps one texture row per screen row of a
%%%          200-row view, which is why the sky neither scrolls vertically nor
%%%          scales with the window.
\newcount\dtSKYT \dtSKYT=-1
\newcount\dtSKYW \newcount\dtPLH \newcount\dtPSM \newcount\dtPFI
%%% Doom's fixedcolormap: -1 for none, otherwise the one COLORMAP table the
%%% whole view is drawn through.  It is set by \dt@powerview from the
%%% player's powers and read at the four places a light level is settled --
%%% a wall column, a plane seg, a sprite and a masked seg -- and at none of
%%% the places a pixel is drawn.  The sky is exempt, as it is in Doom.
\newcount\dtFIXCM \dtFIXCM=-1
\newif\ifdt@pcol \newif\ifdt@psky
\def\dtPlaneTables{%
  \tmpi=0
  \dtWhile{\tmpi<\SCRH}{%
    \tmpa=\numexpr\tmpi-\HALFH\relax
    \ifnum\tmpa<0 \tmpa=-\tmpa \fi
    \expandafter\xdef\csname dtYS\the\tmpi\endcsname
      {\number\numexpr\PPLANE*512/(2*\tmpa+1)\relax}%
    \advance\tmpi by1\relax }%
  \tmpi=0
  \dtWhile{\tmpi<256}{%
    \tmpb=\numexpr\tmpi-128\relax
    \ifnum\tmpb>\numexpr\dtCMN-1\relax \tmpb=\numexpr\dtCMN-1\relax \fi
    \ifnum\tmpb<0 \tmpb=0 \fi
    \expandafter\xdef\csname dtCC\the\tmpi\endcsname{\the\tmpb}%
    \advance\tmpi by1\relax }%
  %% ...and thirty-three more that do not clamp at all: 300+n is table n,
  %% which is how a fixed colormap -- 1 for the visor, 32 for invulnerability
  %% -- gets through a lookup built to keep everything under 32.
  \tmpi=0
  \dtWhile{\tmpi<33}{%
    \expandafter\xdef\csname dtCC\number\numexpr300+\tmpi\relax\endcsname
      {\the\tmpi}%
    \advance\tmpi by1\relax }%
  \tmpi=0
  \dtWhile{\tmpi<\SCRW}{%
    \dtAtan{\PPLANE}{\numexpr\PPLANE-\tmpi\relax}%
    \expandafter\xdef\csname dtKA\the\tmpi\endcsname{\the\dtANG}%
    \advance\tmpi by1\relax }%
  \ifnum\dtSKYT<0\else
    \dtSKYW=\csname dtXw\the\dtSKYT\endcsname\relax
    \ifnum\dtSKYW=0
      \dtSKYT=-1
      \message{[doom.tex] no sky texture in this WAD -- sky stays a ramp.}%
    \else
      \tmpc=\csname dtXh\the\dtSKYT\endcsname\relax
      \tmpi=0
      \dtWhile{\tmpi<\SCRH}{%
        \tmpb=\numexpr100+(\tmpi-\HALFH)*200/\SCRH\relax
        \ifnum\tmpb<0 \tmpb=0 \fi
        \ifnum\tmpb>\tmpc \tmpb=\tmpc \fi
        \expandafter\xdef\csname dtSV\the\tmpi\endcsname{\the\tmpb}%
        \advance\tmpi by1\relax }%
      \message{[doom.tex] sky: texture \the\dtSKYT, \the\dtSKYW x\the\tmpc}%
    \fi
  \fi}

%%% Once a frame: the camera in 12.4 with the bias that keeps every world
%%% coordinate positive (the modulus in \dt@ptV is a rounding \numexpr and
%%% only floors correctly above zero), and the per-column ray directions.
%%% The bias is 2^28, a whole number of 64-unit tiles, so it shifts nothing.
\def\dtPlaneFrame{%
  \edef\dt@pxq{\number\numexpr\PX/4096+268435456\relax}%
  \edef\dt@pyq{\number\numexpr\PY/4096+268435456\relax}%
  \cx=0
  \dtWhile{\cx<\SCRW}{%
    \PU\cx=\numexpr(\vcos-(\PPLANE-\cx)*\vsin/\PPLANE)*4096/65536\relax
    \PV\cx=\numexpr(\vsin+(\PPLANE-\cx)*\vcos/\PPLANE)*4096/65536\relax
    \ifnum\dtSKYT<0\else
      \tmpa=\numexpr\PA+\csname dtKA\the\cx\endcsname\relax
      \divide\tmpa by8
      \tmpb=\tmpa \divide\tmpb by\dtSKYW \multiply\tmpb by\dtSKYW
      \tmpa=\numexpr\tmpa-\tmpb\relax
      \tmpb=\tmpa \divide\tmpb by8
      \PK\cx=\tmpb
      \multiply\tmpb by8
      \PJ\cx=\numexpr\tmpa-\tmpb\relax
    \fi
    \advance\cx by1\relax }}

%%% ---- picking a plane's surface, once per seg -----------------------------
%%% Which flat a plane wears, how far its height is from the eye and how dark
%%% its sector is are all constant along a seg, so they are settled here and
%%% the column loop only switches between the two answers.  A flat that never
%%% resolved, or a sky with no sky texture, falls back to the sector's
%%% procedural ramp exactly as an unresolved wall texture does.
%%%
%%% #1 is C or F, #2 the plane height, #3 the flat array, #4 the ramp base.
\def\dt@planeone#1#2#3#4{%
  \dtPFI=-1
  \ifcsname #3\the\dtFS\endcsname \dtPFI=\csname #3\the\dtFS\endcsname\relax \fi
  \ifnum\dtPFI=-2
    \ifnum\dtSKYT<0 \dtPFI=-1 \fi
  \else
    \ifnum\dtPFI<0\else
      \ifnum\csname dtFw\the\dtPFI\endcsname=0 \dtPFI=-1 \fi
    \fi
  \fi
  \ifnum\dtPFI=-2
    \global\dt@pskytrue
    \expandafter\let\csname dt@#1run\endcsname\SkyRun
  \else\ifnum\dtPFI<0
    \expandafter\edef\csname dt@#1run\endcsname##1##2{%
      \noexpand\PlaneRun{##1}{##2}{\number\numexpr#4\relax}}%
  \else
    \global\dt@pcoltrue
    \expandafter\let\csname dt@#1run\endcsname\PlaneTex
    %% planeheight in 12.4, capped at 4096 units: past that the plane is a
    %% row or two of horizon and the cap keeps the world coordinate inside
    %% the bias \dtPlaneFrame allows for.
    \dtPLH=\numexpr(\PZ-(#2)*65536)/4096\relax
    \ifnum\dtPLH<0 \dtPLH=-\dtPLH \fi
    \ifnum\dtPLH>65536 \dtPLH=65536 \fi
    \expandafter\xdef\csname dt@#1key\endcsname{dtFL\the\dtPFI @}%
    \expandafter\xdef\csname dt@#1phq\endcsname{\the\dtPLH}%
    \expandafter\xdef\csname dt@#1psm\endcsname{\the\dtPSM}%
  \fi\fi}

\def\dtPlaneSeg{%
  \global\dt@pcolfalse \global\dt@pskyfalse
  \dtPSM=\csname dtSl\the\dtFS\endcsname\relax
  \divide\dtPSM by16
  \dtPSM=\numexpr(15-\dtPSM)*4\relax
  %% A fixed colormap reaches a plane through the clamp table rather than
  %% past it: dtCC<300+n> answers n, \dt@setfixcm has taken the distance
  %% term out, and 172+128 is 300.  No plane pixel pays for the feature.
  \ifnum\dtFIXCM<0 \else \dtPSM=\numexpr172+\dtFIXCM\relax \fi
  \dt@planeone{C}{\csname dtSc\the\dtFS\endcsname}{dtAc}%
    {\csname dtSq\the\dtFS\endcsname}%
  \dt@planeone{F}{\csname dtSf\the\dtFS\endcsname}{dtAf}%
    {\csname dtSp\the\dtFS\endcsname}}

%% Four \let per surface per column, against two \ifcsname and a division if
%% this were decided per column instead.
\def\dt@Ckey{dtFL0@}\def\dt@Cphq{0}\def\dt@Cpsm{0}
\def\dt@Fkey{dtFL0@}\def\dt@Fphq{0}\def\dt@Fpsm{0}
\def\dt@Crun#1#2{\PlaneRun{#1}{#2}{0}}
\def\dt@Frun#1#2{\PlaneRun{#1}{#2}{0}}
\def\dt@setceil{\let\dt@prun\dt@Crun \let\dt@fkey\dt@Ckey
  \let\dt@phq\dt@Cphq \let\dt@psm\dt@Cpsm}
\def\dt@setfloor{\let\dt@prun\dt@Frun \let\dt@fkey\dt@Fkey
  \let\dt@phq\dt@Fphq \let\dt@psm\dt@Fpsm}

%%% Expandable textured wall span: rows, v (16.16), vstep (16.16), ramp base.
%%% v wraps by one subtraction, which stays inside \numexpr and therefore
%%% inside \edef -- so a whole span is one expansion with no assignments.
\def\dt@uu{0}
%%% Which shade table this span samples: dtT (brick) or dtD (door).  Set
%%% before the \edef that expands \WallRun -- it is read during that
%%% expansion and nowhere else, so it costs nothing per pixel.
\def\dt@tx{dtT}
\def\WallRun#1#2#3#4{%
  \ifnum\numexpr#1\relax>0
    \hx{\numexpr(#4)+\csname \dt@tx\dt@uu @\number\numexpr(#2)/65536\relax
        \endcsname\relax}%
    \expandafter\dt@wrN\else\expandafter\dt@wrS\fi{#1}{#2}{#3}{#4}}
%% Two arguments need pre-expansion, and a macro grabs its arguments before a
%% later \expandafter can fire -- so expand them one at a time.
\def\dt@wrN#1#2#3#4{%
  \expandafter\dt@wrA\expandafter{\number\numexpr(#2)+(#3)%
    -\ifnum\numexpr(#2)+(#3)\relax>4194303 4194304\else0\fi\relax}{#1}{#3}{#4}}
\def\dt@wrA#1#2#3#4{%
  \expandafter\dt@wrB\expandafter{\number\numexpr(#2)-1\relax}{#1}{#3}{#4}}
\def\dt@wrB#1#2#3#4{\WallRun{#1}{#2}{#3}{#4}}
\def\dt@wrS#1#2#3#4{}

%%% The same span over a real composited texture.  Two differences: the texel
%%% IS a palette byte rather than a shade offset, so it goes through COLORMAP
%%% exactly as a sprite pixel does and needs no \hx on the way out; and v
%%% wraps at the texture's own height rather than at a fixed 64 rows, so the
%%% modulus arrives in \dt@vmod instead of being written into the macro.
%%%
%%% The nested \csname is the whole trick: scanning the outer name expands
%%% the inner one, which yields the two hex characters that finish the
%%% COLORMAP key.  It also means an undefined texel is a "Missing
%%% \endcsname", not a wrong pixel -- which is why the loader fills holes.
%%%
%%% #4 is the ramp base, unused here.  Both runs take it so that one \let
%%% chooses between them and the \edef below does not have to know which.
\def\dt@ux{}
\def\dt@cml{dtCM0@}
\def\dt@vmod{4194304}
%%% A composited texel arrives eight to a control sequence, packed along u --
%%% see wad.tex for why the string pool leaves no choice.  Because u is fixed
%%% for a whole column span, WHICH of the eight is fixed too, so \dt@usetex
%%% \let\dt@pick at one of these once per column and the span expansion below
%%% pays a single macro call and no arithmetic.
%%%
%%% Defined through \csname because a digit is not a letter: \dt@pk0 written
%%% out is \dt@pg followed by the character 0, which would silently define
%%% \dt@pg eight times over with conflicting parameter texts.  \csname names
%%% may use digits; macro names may not.
%%%
%%% Not \dt@pk either: play.tex holds the pressed key in \dt@pk, so the flat
%%% picker under that name became the letter `w' the moment anyone walked
%%% forward, and the braces of the group it should have taken apart travelled
%%% into the column for \dt@spDN to trip over.  Batch mode never sees it.
\expandafter\def\csname dt@pg0\endcsname#1#2#3#4#5#6#7#8{#1}
\expandafter\def\csname dt@pg1\endcsname#1#2#3#4#5#6#7#8{#2}
\expandafter\def\csname dt@pg2\endcsname#1#2#3#4#5#6#7#8{#3}
\expandafter\def\csname dt@pg3\endcsname#1#2#3#4#5#6#7#8{#4}
\expandafter\def\csname dt@pg4\endcsname#1#2#3#4#5#6#7#8{#5}
\expandafter\def\csname dt@pg5\endcsname#1#2#3#4#5#6#7#8{#6}
\expandafter\def\csname dt@pg6\endcsname#1#2#3#4#5#6#7#8{#7}
\expandafter\def\csname dt@pg7\endcsname#1#2#3#4#5#6#7#8{#8}
\expandafter\let\expandafter\dt@pick\csname dt@pg0\endcsname
%% THREE \expandafter: \csname yields the control sequence naming the group,
%% not the group, and a macro's arguments are not expanded while being
%% grabbed -- with one, \dt@pick takes that token as its first texel and eats
%% \endcsname looking for the other seven.
\def\WallTex#1#2#3#4{%
  \ifnum\numexpr#1\relax>0
    \csname\dt@cml\expandafter\expandafter\expandafter\dt@pick
      \csname\dt@ux\number\numexpr(#2)/65536\relax\endcsname\endcsname
    \expandafter\dt@wtN\else\expandafter\dt@wtS\fi{#1}{#2}{#3}{#4}}
%% \dt@vmod sits inside a \numexpr that \relax terminates, which is the only
%% safe way to read a macro as a number here: a bare \dt@vmod followed by a
%% space produces no space token and the scan runs straight into the digits
%% that follow it.
\def\dt@wtN#1#2#3#4{%
  \expandafter\dt@wtA\expandafter{\number\numexpr(#2)+(#3)%
    -\ifnum\numexpr(#2)+(#3)-\dt@vmod\relax<0 0\else\dt@vmod\fi\relax}%
    {#1}{#3}{#4}}
\def\dt@wtA#1#2#3#4{%
  \expandafter\dt@wtB\expandafter{\number\numexpr(#2)-1\relax}{#1}{#3}{#4}}
\def\dt@wtB#1#2#3#4{\WallTex{#1}{#2}{#3}{#4}}
\def\dt@wtS#1#2#3#4{}

%%% View-space transform.  Every \fmul here has one operand bounded by
%%% FRACUNIT, so the 64-bit intermediate cannot overflow on the way out.
\def\dtXform#1#2{%
  \rc=\numexpr(#1)*65536-\PX\relax
  \rd=\numexpr(#2)*65536-\PY\relax
  \ra=\numexpr\fmul{\the\rc}{\the\vcos}+\fmul{\the\rd}{\the\vsin}\relax
  \rb=\numexpr\fmul{\the\rd}{\the\vcos}-\fmul{\the\rc}{\the\vsin}\relax}

\def\dtScale#1{\number\numexpr\PPLANE*65536*65536/(#1)\relax}

%%% ---- BSP traversal, front to back ----------------------------------------
%%% R_PointOnSide: the sign of the cross product of the partition delta with
%%% the vector from the partition origin to the query point.
%%%
%%% The query point is (\dtPSX,\dtPSY) in *map units*, not the camera, so that
%%% point location (\dtPointInSSec) can reuse this untouched.  \dtViewSetup
%%% loads it from the camera once per frame, which also removes a division
%%% from the inner traversal.
\newcount\dtPSX \newcount\dtPSY
\def\dtPointSide#1{%
  \tmpa=\numexpr\dtPSX-(\csname dtNx#1\endcsname)\relax
  \tmpb=\numexpr\dtPSY-(\csname dtNy#1\endcsname)\relax
  \tmpc=\numexpr(\csname dtNu#1\endcsname)*\tmpb
        -(\csname dtNv#1\endcsname)*\tmpa\relax
  \ifnum\tmpc<0 \dtSIDE=0 \else\dtSIDE=1 \fi}

%%% Stops as soon as every column is closed, which is what makes walking a
%%% real level cheap: geometry behind a solid wall is never visited.
\def\dtRenderBSP#1{%
  \ifnum\dtOPEN>0
    \ifnum\numexpr#1\relax>32767
      \dtSubsector{\number\numexpr(#1)-32768\relax}%
    \else
      \dtPointSide{#1}%
      \ifnum\dtSIDE=0
        \dtRenderBSP{\csname dtNr#1\endcsname}%
        \dtRenderBSP{\csname dtNl#1\endcsname}%
      \else
        \dtRenderBSP{\csname dtNl#1\endcsname}%
        \dtRenderBSP{\csname dtNr#1\endcsname}%
      \fi
    \fi
  \fi}

%%% Subsectors are leaves, so this loop never nests inside itself.
\def\dtSubsector#1{%
  \dtSA=\csname dtCf#1\endcsname\relax
  \dtSB=\numexpr\dtSA+\csname dtCn#1\endcsname\relax
  \dtWhile{\dtSA<\dtSB}{\dtSegDraw{\the\dtSA}\advance\dtSA by1\relax }}

%%% ==================================  POINT LOCATION / FLOOR / COLLISION (M1-M3)
%%% All of this runs *between* frames, never inside the render path.  It still
%%% avoids \tmpa-\tmpk except where \dtPointSide already uses them, because
%%% \dtFollowFloor is one edit away from being called from inside a frame.

\newcount\dtPI \newcount\dtPJ \newcount\dtPK
\newcount\dtSSEC \newcount\dtSECN     % subsector / sector containing the point

%%% ---- M1: R_PointInSubsector ---------------------------------------------
%%% Descend from the root taking the near child until a leaf falls out.  A
%%% child index >= 32768 is a subsector, exactly as in the WAD.
\def\dtPointInSSec{%
  \dtPI=\DTROOT
  \dtWhile{\dtPI<32768}{%
    \dtPointSide{\the\dtPI}%
    \ifnum\dtSIDE=0
      \dtPI=\csname dtNr\the\dtPI\endcsname\relax
    \else
      \dtPI=\csname dtNl\the\dtPI\endcsname\relax
    \fi}%
  \dtSSEC=\numexpr\dtPI-32768\relax
  \dtSSecSector{\the\dtSSEC}}

%%% Subsectors do not record their sector.  Recover it the way Doom does, from
%%% the first seg: the seg's side selects the linedef's front or back sector.
\def\dtSSecSector#1{%
  \dtPJ=\csname dtCf#1\endcsname\relax
  \dtPK=\csname dtGl\the\dtPJ\endcsname\relax
  \ifnum\csname dtGs\the\dtPJ\endcsname=0
    \dtSECN=\csname dtLF\the\dtPK\endcsname\relax
  \else
    \dtSECN=\csname dtLB\the\dtPK\endcsname\relax
  \fi
  %% a malformed map can give a side-1 seg on a one-sided linedef
  \ifnum\dtSECN<0 \dtSECN=\csname dtLF\the\dtPK\endcsname\relax\fi}

\def\dtLocatePlayer{%
  \dtPSX=\numexpr\PX/65536\relax
  \dtPSY=\numexpr\PY/65536\relax
  \dtPointInSSec}

%%% ---- M2: floor following -------------------------------------------------
%%% Doom moves viewz *toward* the target rather than snapping, which is what
%%% makes a staircase read as a staircase instead of a lift.
\newcount\VIEWH \VIEWH=41             % Doom's VIEWHEIGHT
\newcount\dtTZ \newcount\dtDZ
\def\dtFloorZ{\csname dtSf\the\dtSECN\endcsname}
\def\dtSnapFloor{\PZ=\numexpr(\dtFloorZ+\VIEWH)*65536\relax}
%%% Halving the gap converges in two or three frames, which reads as a step.
%%% Dividing by three took seven and read as floating up a ramp; the snap
%%% threshold has to be on the *remaining gap*, or the last unit never closes
%%% because \numexpr rounds a small quotient to zero.
\def\dtFollowFloor{%
  \dtTZ=\numexpr(\dtFloorZ+\VIEWH)*65536\relax
  \dtDZ=\numexpr\dtTZ-\PZ\relax
  \ifnum\dtDZ<0 \dtPJ=-\dtDZ \else\dtPJ=\dtDZ \fi
  \ifnum\dtPJ<1048576
    \PZ=\dtTZ
  \else
    \advance\PZ by\numexpr\dtDZ/2\relax
  \fi}

%%% ---- M3: collision -------------------------------------------------------
%%% Doom's constants.  A move is blocked by a one-sided line, by an opening too
%%% short to fit through, or by a step too high to climb.
\newcount\PRAD  \PRAD=16
\newcount\PSTEP \PSTEP=24
\newcount\PHEAD \PHEAD=56

%%% The radius and the headroom THIS move is tested with.  The scan below
%%% used to read \PRAD and \PHEAD directly, which was fine while the player
%%% was the only thing that moved; a monster has its own radius -- a demon is
%%% thirty units wide against the player's sixteen -- and borrows the same
%%% scan.  \dtTryMove sets both back to the player's on every call, so the
%%% player's route through here is exactly the route it was.
\newcount\dtMRAD \dtMRAD=16
\newcount\dtMHGT \dtMHGT=56
%%% Which thing is doing the moving, so the scan below does not stop it dead
%%% against itself.  -1 is the player.
\newcount\dtMSELF \dtMSELF=-1
%%% \DTNOMON leaves the map exactly as it was before there were monsters in
%%% it: nothing ticks and nothing blocks.  Declared here rather than with the
%%% rest of M15 because the collision scan below is the first thing to ask.
\newif\ifdt@mon \dt@montrue
%%% ...and \DTNOMIS the same for M18, for the same reason: \dt@segopen below
%%% is the first thing to ask whether the mover it is testing is a missile.
%%% A missile does not climb steps and does not duck -- it is stopped by the
%%% geometry at its own height -- so the one branch here is the whole of the
%%% departure.  \dtMISZ is the z it is flying at while that test runs.
\newif\ifdt@mis \dt@misfalse
\newcount\dtMISZ
%%% Which thing stopped the move, or -1.  The player and a monster only ever
%%% ask WHETHER they were blocked; a missile has to damage what it hit, and
%%% the scan below is already standing on the answer.
\newcount\dtMHIT \dtMHIT=-1
\newcount\dtBKI \newcount\dtBKA \newcount\dtBKB \newcount\dtBKR
\newcount\dtBKS

\newcount\dtCX \newcount\dtCY         % candidate position, map units
\newcount\dtPX \newcount\dtPY         % move origin, map units
\newcount\dtOX \newcount\dtOY         % move origin RELATIVE to the candidate
\newcount\dtBXL \newcount\dtBXH \newcount\dtBYL \newcount\dtBYH
\newcount\dtDA \newcount\dtDB \newcount\dtDC \newcount\dtDD \newcount\dtHIT
\newcount\dtBLK \newcount\dtSGI \newcount\dtRJ
\newcount\dtCA \newcount\dtCB
\newcount\dtAX \newcount\dtAY \newcount\dtBX \newcount\dtBY
\newcount\dtABX \newcount\dtABY \newcount\dtDOT \newcount\dtLSQ \newcount\dtCR
\newcount\dtCL \newcount\dtCF \newcount\dtCK \newcount\dtOT \newcount\dtOB
\newcount\dtFLZ                       % floor the player is standing on now

%%% Collision runs against SEGS, not LINEDEFS.  Segs already carry vertex
%%% indices and a linedef index, so no new map arrays are needed, and segs
%%% partition the linedefs exactly -- the same geometry, already loaded.
%%%
%%% This tests a *move*, not a position.  Testing only where the move lands
%%% tunnels: a 40-unit step that starts 20 units in front of a wall ends 20
%%% units behind it, both outside the 16-unit radius, and the player passes
%%% straight through into the void.  A sealed map cannot be left without
%%% crossing some seg, so the crossing test below is what actually keeps the
%%% player inside the level; the radius test is what stops them resting
%%% inside a wall.
\def\dtBlocked{%
  \dtBLK=0 \dtSGI=0
  %% the swept box, in coordinates relative to the destination
  \dtBXL=\ifnum\dtOX<0 \dtOX\else0\fi\relax \advance\dtBXL by-\dtMRAD
  \dtBXH=\ifnum\dtOX>0 \dtOX\else0\fi\relax \advance\dtBXH by\dtMRAD
  \dtBYL=\ifnum\dtOY<0 \dtOY\else0\fi\relax \advance\dtBYL by-\dtMRAD
  \dtBYH=\ifnum\dtOY>0 \dtOY\else0\fi\relax \advance\dtBYH by\dtMRAD
  \dtWhile{\dtSGI<\NSEG}{\dt@segblock\advance\dtSGI by1\relax }%
  \ifnum\dtBLK=0 \dt@mobjblock \fi}

%%% P_CheckPosition's other half: a live monster is solid.  Doom compares the
%%% sum of the two radii on each axis, which is the box test \dt@touch
%%% already uses for items -- an item is reached with it and a monster is
%%% stopped by it.
%%%
%%% This runs on the destination only, not swept.  A monster is not a wall:
%%% tunnelling past one at speed is a missed shove rather than an escape from
%%% the level, and the swept form would cost three more scans a move for it.
%%%
%%% Dying and dead do not block, which is Doom clearing MF_SOLID as the
%%% death animation starts -- you can walk over a corpse.
\def\dt@mobjblock{%
  \ifdt@mon
    \ifnum\dtNMON>0
      \dtBKI=0
      \dtWhile{\dtBKI<\NTHING}{%
        \ifnum\dtBKI=\dtMSELF \else
          \ifcsname dtOs\the\dtBKI\endcsname \dt@mobjone \fi
        \fi
        \advance\dtBKI by1\relax }%
    \fi
  \fi}
\def\dt@mobjone{%
  \ifnum\csname dtOs\the\dtBKI\endcsname>3 \else
    %% PIT_CheckThing's species rule (M20): a missile passes straight through
    %% a thing of the same type as whatever fired it.  Without it the first
    %% fireball an imp throws across a room of imps starts a brawl that Doom
    %% would never have started, and infighting stops being a thing the
    %% player can provoke and becomes the default state of every map.
    \dtBKS=1
    \ifdt@mis \ifnum\dtMSELF>-1
      \ifnum\csname dtHt\the\dtBKI\endcsname=\csname dtHt\the\dtMSELF\endcsname
        \relax \dtBKS=0
      \fi
    \fi\fi
    \ifnum\dtBKS=1
    \dtBKR=\numexpr\dtMRAD
      +\csname dtOr\csname dtHt\the\dtBKI\endcsname\endcsname\relax
    \dtBKA=\numexpr\csname dtHx\the\dtBKI\endcsname-\dtCX\relax
    \ifnum\dtBKA<0 \dtBKA=-\dtBKA \fi
    \ifnum\dtBKA<\dtBKR
      \dtBKB=\numexpr\csname dtHy\the\dtBKI\endcsname-\dtCY\relax
      \ifnum\dtBKB<0 \dtBKB=-\dtBKB \fi
      \ifnum\dtBKB<\dtBKR \dtBLK=1 \dtMHIT=\dtBKI \dtBKI=\NTHING \fi
    \fi
    \fi
  \fi}

%%% Set up a move from (#1,#2) to (#3,#4), all in map units, and test it.
\def\dtMoveTest#1#2#3#4{%
  \dtCX=\numexpr#3\relax \dtCY=\numexpr#4\relax
  \dtOX=\numexpr(#1)-\dtCX\relax \dtOY=\numexpr(#2)-\dtCY\relax
  \dtBlocked}

%%% Vertices relative to the candidate point.  The bounding-box reject that
%%% follows is what keeps the products below 2^31: past it, the interesting
%%% coordinates are bounded by the seg's extent plus the swept box.
\def\dt@segblock{%
  \dtCA=\csname dtGa\the\dtSGI\endcsname\relax
  \dtCB=\csname dtGb\the\dtSGI\endcsname\relax
  \dtAX=\numexpr\csname dtVx\the\dtCA\endcsname-\dtCX\relax
  \dtAY=\numexpr\csname dtVy\the\dtCA\endcsname-\dtCY\relax
  \dtBX=\numexpr\csname dtVx\the\dtCB\endcsname-\dtCX\relax
  \dtBY=\numexpr\csname dtVy\the\dtCB\endcsname-\dtCY\relax
  \dtRJ=0
  \ifnum\dtAX>\dtBXH \ifnum\dtBX>\dtBXH \dtRJ=1 \fi\fi
  \ifnum\dtAX<\dtBXL \ifnum\dtBX<\dtBXL \dtRJ=1 \fi\fi
  \ifnum\dtAY>\dtBYH \ifnum\dtBY>\dtBYH \dtRJ=1 \fi\fi
  \ifnum\dtAY<\dtBYL \ifnum\dtBY<\dtBYL \dtRJ=1 \fi\fi
  \ifnum\dtRJ=0
    \dtHIT=0
    \dt@segcross
    \ifnum\dtHIT=0 \dt@segnear \fi
    \ifnum\dtHIT=1 \dt@segopen \fi
  \fi}

%%% Does the move cross this seg?  The move runs from (\dtOX,\dtOY) to the
%%% origin, so the two endpoint tests against AB are cheap.  Collinear and
%%% exactly-touching cases fall through to the radius test rather than being
%%% special-cased here.
\def\dt@segcross{%
  \dtABX=\numexpr\dtBX-\dtAX\relax
  \dtABY=\numexpr\dtBY-\dtAY\relax
  \dtDA=\numexpr\dtABX*(\dtOY-\dtAY)-\dtABY*(\dtOX-\dtAX)\relax
  \dtDB=\numexpr\dtABY*\dtAX-\dtABX*\dtAY\relax
  \ifnum\dtDA=0 \else\ifnum\dtDB=0 \else
    \ifnum\dtDA<0
      \ifnum\dtDB>0 \dt@segcrossB \fi
    \else
      \ifnum\dtDB<0 \dt@segcrossB \fi
    \fi
  \fi\fi}
\def\dt@segcrossB{%
  \dtDC=\numexpr-\dtOX*(\dtAY-\dtOY)+\dtOY*(\dtAX-\dtOX)\relax
  \dtDD=\numexpr-\dtOX*(\dtBY-\dtOY)+\dtOY*(\dtBX-\dtOX)\relax
  \ifnum\dtDC=0 \else\ifnum\dtDD=0 \else
    \ifnum\dtDC<0
      \ifnum\dtDD>0 \dtHIT=1 \fi
    \else
      \ifnum\dtDD<0 \dtHIT=1 \fi
    \fi
  \fi\fi}

%%% Distance from the destination to the segment, without a square root:
%%% compare against an endpoint when the projection falls outside the seg, and
%%% against |cross| < radius * length when it falls inside.  \dtGn is that
%%% length, already computed at load time for texture mapping.  \dtABX/\dtABY
%%% were set by \dt@segcross, which always runs first.
\def\dt@segnear{%
  \dtLSQ=\numexpr\dtABX*\dtABX+\dtABY*\dtABY\relax
  \ifnum\dtLSQ>0
    \dtDOT=\numexpr-\dtAX*\dtABX-\dtAY*\dtABY\relax
    \ifnum\dtDOT<1
      \ifnum\numexpr\dtAX*\dtAX+\dtAY*\dtAY\relax<\numexpr\dtMRAD*\dtMRAD\relax
        \dtHIT=1
      \fi
    \else
      \ifnum\dtDOT>\dtLSQ
        \ifnum\numexpr\dtBX*\dtBX+\dtBY*\dtBY\relax<\numexpr\dtMRAD*\dtMRAD\relax
          \dtHIT=1
        \fi
      \else
        \dtCR=\numexpr\dtAX*\dtBY-\dtAY*\dtBX\relax
        \ifnum\dtCR<0 \dtCR=-\dtCR \fi
        \ifnum\dtCR<\numexpr\dtMRAD*(\csname dtGn\the\dtSGI\endcsname)\relax
          \dtHIT=1
        \fi
      \fi
    \fi
  \fi}

%%% P_LineOpening: the gap a two-sided line leaves is between the lower of the
%%% two ceilings and the higher of the two floors.
\def\dt@segopen{%
  \dtCL=\csname dtGl\the\dtSGI\endcsname\relax
  \dtCF=\csname dtLF\the\dtCL\endcsname\relax
  \dtCK=\csname dtLB\the\dtCL\endcsname\relax
  \ifnum\dtCK<0
    \dtBLK=1
  \else
    \ifnum\dtCF<0
      \dtBLK=1
    \else
      \dtOT=\csname dtSc\the\dtCF\endcsname\relax
      \ifnum\csname dtSc\the\dtCK\endcsname<\dtOT
        \dtOT=\csname dtSc\the\dtCK\endcsname\relax
      \fi
      \dtOB=\csname dtSf\the\dtCF\endcsname\relax
      \ifnum\csname dtSf\the\dtCK\endcsname>\dtOB
        \dtOB=\csname dtSf\the\dtCK\endcsname\relax
      \fi
      \ifnum\numexpr\dtOT-\dtOB\relax<\dtMHGT \dtBLK=1 \fi
      %% A missile is stopped by the opening it actually arrives at rather
      %% than by a step it cannot climb -- P_TryMove skips the step test for
      %% MF_MISSILE and P_ZMovement explodes it against the floor instead.
      %% That is what lets a fireball fly over a parapet the imp behind it
      %% could not walk up, and stop dead against the wall above a doorway.
      \ifdt@mis
        \ifnum\dtMISZ<\dtOB \dtBLK=1 \fi
        \ifnum\numexpr\dtMISZ+\dtMHGT\relax>\dtOT \dtBLK=1 \fi
      \else
        \ifnum\numexpr\dtOB-\dtFLZ\relax>\PSTEP \dtBLK=1 \fi
      \fi
    \fi
  \fi
  %% end the scan early by running the loop counter out
  \ifnum\dtBLK=1 \dtSGI=\NSEG \fi}

%%% P_TryMove, with the cheap form of sliding: when the full move is blocked,
%%% keep whichever single axis is still free.  Against an axis-aligned wall
%%% this is exactly a projection onto the wall; against a diagonal it loses
%%% some speed but never sticks.
\newcount\dtNXF \newcount\dtNYF
\def\dtTryMove#1#2{%
  \dtNXF=\numexpr#1\relax \dtNYF=\numexpr#2\relax
  %% the player's own size, every time: a monster's move borrows this scan
  %% and leaves its own radius behind in these two registers
  \dtMRAD=\PRAD \dtMHGT=\PHEAD \dtMSELF=-1
  \dtFLZ=\dtFloorZ\relax
  \dtPX=\numexpr\PX/65536\relax \dtPY=\numexpr\PY/65536\relax
  \dtMoveTest{\dtPX}{\dtPY}{\dtNXF/65536}{\dtNYF/65536}%
  \ifnum\dtBLK=0
    \PX=\dtNXF \PY=\dtNYF
  \else
    \dtMoveTest{\dtPX}{\dtPY}{\dtNXF/65536}{\dtPY}%
    \ifnum\dtBLK=0
      \PX=\dtNXF
    \else
      \dtMoveTest{\dtPX}{\dtPY}{\dtPX}{\dtNYF/65536}%
      \ifnum\dtBLK=0 \PY=\dtNYF \fi
    \fi
  \fi
  %% Walkover triggers fire on the move that was actually taken, after the
  %% slide has decided what that is -- \dtPX/\dtPY still hold the origin,
  %% because \dtMoveTest writes \dtCX/\dtCY and never touches them.  Before
  %% \dtLocatePlayer, so that a floor this move just started is already
  %% moving when the eye height is settled against it.
  \dtCrossLines
  \dtLocatePlayer
  %% After \dtLocatePlayer, because the touch scan reads \dtPSX/\dtPSY and the
  %% floor under the player, and all three are what that call just settled.
  %% Doom reaches items from P_CheckPosition, i.e. on the move that was
  %% taken -- so a move blocked by a wall still collects what it slid past.
  \dtTouchThings
  \dtFollowFloor}

%%% One step along the facing angle: the movement primitive both the scripted
%%% walk and the key loop go through.
%%
%% P_DeathThink's first line: a dead player does not walk.  The guard is here
%% rather than in the two loops that call this, because \dtUse and \dtFire
%% each carry their own and a third copy in the caller would be a fourth
%% place to forget.
\def\dtStepFwd#1{%
  \ifnum\dtHEALTH>0
    \dtTryMove{\PX+\fmul{\dtcos{\the\PA}}{#1}}%
              {\PY+\fmul{\dtsin{\the\PA}}{#1}}%
  \fi}
%% ...and one step sideways, which is the same move a quarter turn round:
%% positive is to the player's left, where increasing angle already points.
\def\dtStepSide#1{%
  \ifnum\dtHEALTH>0
    \dtTryMove{\PX-\fmul{\dtsin{\the\PA}}{#1}}%
              {\PY+\fmul{\dtcos{\the\PA}}{#1}}%
  \fi}

%%% ---- one seg -------------------------------------------------------------
%%% Door surfaces are decided once per seg, not once per column: the two
%%% \ifcsname lookups would otherwise run for every pixel column the seg
%%% covers, and neither answer can change across a seg.
%%%
%%% In Doom a shut door is the UPPER texture of a two-sided line whose back
%%% sector has collapsed its ceiling onto its floor -- so the face is the
%%% upper span of any seg whose BACK sector is a door.  The track is the
%%% door sector's own one-sided walls, which is why it keys on the FRONT
%%% sector instead.  Both fall out of the dtDT flag \dtFindDoors already
%%% sets, so this needs no new map array.
%%
%% The seg-level constants come in two halves with \dtPlaneSeg between them,
%% and they are macros of their own because M19's masked pass has to rebuild
%% exactly this state for a seg it is revisiting after the wall pass is over
%% -- without re-running the plane setup, which has already had its say.  The
%% split is where \dtPlaneSeg already sat, so the order here is unchanged.
\newif\ifdt@solid \newif\ifdt@doorface \newif\ifdt@doortrack
\def\dtSegDraw#1{%
  \dt@segsect{#1}%
  \dtPlaneSeg
  \dt@segtex
  \dt@segproj}
\def\dt@segsect#1{%
  \dtG=#1\relax
  \dtLD=\csname dtGl\the\dtG\endcsname\relax
  \dtSDE=\csname dtGs\the\dtG\endcsname\relax
  \ifnum\dtSDE=0
    \dtFS=\csname dtLF\the\dtLD\endcsname\relax
    \dtBS=\csname dtLB\the\dtLD\endcsname\relax
  \else
    \dtFS=\csname dtLB\the\dtLD\endcsname\relax
    \dtBS=\csname dtLF\the\dtLD\endcsname\relax
  \fi
  \ifnum\dtBS<0 \dt@solidtrue \else\dt@solidfalse \fi
  \dt@doorfacefalse \dt@doortrackfalse
  \ifcsname dtDT\the\dtFS\endcsname \dt@doortracktrue \fi
  \ifnum\dtBS<0\else
    \ifcsname dtDT\the\dtBS\endcsname \dt@doorfacetrue \fi
  \fi}
\def\dt@segtex{%
  \dtVA=\csname dtGa\the\dtG\endcsname\relax
  \dtVB=\csname dtGb\the\dtG\endcsname\relax
  %% ---- textures for this seg's three possible surfaces ----
  %% seg -> linedef -> side -> sidedef, exactly as Doom reaches them.  All
  %% of this is constant along a seg, so it is decided once here rather than
  %% per column, and a map with no sidedefs at all (the built-in level)
  %% leaves every index at -1 and stays on the procedural tables.
  \dtSDI=-1
  \ifnum\dtSDE=0
    \ifcsname dtLR\the\dtLD\endcsname
      \dtSDI=\csname dtLR\the\dtLD\endcsname\relax\fi
  \else
    \ifcsname dtLL\the\dtLD\endcsname
      \dtSDI=\csname dtLL\the\dtLD\endcsname\relax\fi
  \fi
  \dtWTU=-1 \dtWTL=-1 \dtWTM=-1
  \dtUOFF=0 \dtVOFF=0 \dtPEGF=0
  \ifnum\dtSDI<0\else
    \dtWTU=\csname dtYu\the\dtSDI\endcsname\relax
    \dtWTL=\csname dtYl\the\dtSDI\endcsname\relax
    \dtWTM=\csname dtYm\the\dtSDI\endcsname\relax
    \dtUOFF=\csname dtYx\the\dtSDI\endcsname\relax
    \dtVOFF=\csname dtYy\the\dtSDI\endcsname\relax
    \ifcsname dtGo\the\dtG\endcsname
      \advance\dtUOFF by\csname dtGo\the\dtG\endcsname\relax\fi
    \ifcsname dtLG\the\dtLD\endcsname
      \dtPEGF=\csname dtLG\the\dtLD\endcsname\relax\fi
  \fi
  %% Doom's fake contrast: a wall running east-west is drawn one light level
  %% darker and one running north-south one lighter, so that two walls
  %% meeting at a corner are told apart even where the sector light is flat.
  \dtFAKE=0
  \ifnum\csname dtVy\the\dtVA\endcsname=\csname dtVy\the\dtVB\endcsname
    \dtFAKE=-1
  \else
    \ifnum\csname dtVx\the\dtVA\endcsname=\csname dtVx\the\dtVB\endcsname
      \dtFAKE=1 \fi
  \fi}
\def\dt@segproj{%
  \dtXform{\csname dtVx\the\dtVA\endcsname}{\csname dtVy\the\dtVA\endcsname}%
  \d@one=\ra \l@one=\rb
  \dtXform{\csname dtVx\the\dtVB\endcsname}{\csname dtVy\the\dtVB\endcsname}%
  \d@two=\ra \l@two=\rb
  \dtSKIP=0
  \ifnum\d@one<\NEARZ \ifnum\d@two<\NEARZ \dtSKIP=1 \fi\fi
  \ifnum\dtSKIP=0
    \u@one=0
    \u@two=\numexpr\csname dtGn\the\dtG\endcsname*65536\relax
    \ifnum\d@one<\NEARZ
      \rc=\numexpr\d@two-\d@one\relax
      \rd=\numexpr(\NEARZ-\d@one)*65536/\the\rc\relax
      \l@one=\numexpr\l@one+\fmul{\numexpr\l@two-\l@one\relax}{\the\rd}\relax
      \u@one=\numexpr\u@one+\fmul{\numexpr\u@two-\u@one\relax}{\the\rd}\relax
      \d@one=\NEARZ
    \fi
    \ifnum\d@two<\NEARZ
      \rc=\numexpr\d@one-\d@two\relax
      \rd=\numexpr(\NEARZ-\d@two)*65536/\the\rc\relax
      \l@two=\numexpr\l@two+\fmul{\numexpr\l@one-\l@two\relax}{\the\rd}\relax
      \u@two=\numexpr\u@two+\fmul{\numexpr\u@one-\u@two\relax}{\the\rd}\relax
      \d@two=\NEARZ
    \fi
    %% Lateral is measured to the viewer's LEFT (+y with the camera facing
    %% +x), so screen x decreases as lateral grows.  Getting this backwards
    %% only mirrors a z-buffered image, but it inverts the facing test below.
    \x@one=\numexpr\SCRW/2-\l@one*\PPLANE/\d@one\relax
    \x@two=\numexpr\SCRW/2-\l@two*\PPLANE/\d@two\relax
    %% a seg faces its own subsector only when it projects left to right
    \ifnum\x@one<\x@two \dt@segspans \fi
  \fi}

\def\dt@segspans{%
  \s@one=\dtScale{\the\d@one}\s@two=\dtScale{\the\d@two}%
  \re=\numexpr\fmul{\the\u@one}{\the\s@one}\relax
  \rf=\numexpr\fmul{\the\u@two}{\the\s@two}\relax
  \wid=\numexpr\x@two-\x@one\relax
  %% M19: a masked middle cannot be drawn now.  A column fills from the
  %% outside in, the middle of a portal is exactly the gap the traversal has
  %% not reached yet, and the texture is transparent -- so whatever comes
  %% later has to show through it.  Doom has the same problem and the same
  %% answer: R_StoreWallRange records a drawseg and R_DrawMasked composites
  %% it afterwards, in the pass that draws the sprites.  This records the
  %% six numbers that pass needs; nothing else about the seg is kept,
  %% because \dt@segsect and \dt@segtex can rebuild the rest from its index.
  \dt@maskrec
  %% These MUST be terminated explicitly.  A line ending after a control word
  %% leaves TeX skipping blanks, so no space token is produced and the number
  %% scan runs on into the next line -- where it expanded \dtWhile itself and
  %% ran the column loop from inside the assignment.
  \cx=\ifnum\x@one<0 0\else\the\x@one\fi\relax
  \rb=\ifnum\x@two>\SCRW \the\SCRW\else\the\x@two\fi\relax
  \dtWhile{\cx<\rb}{%
    \ct=\CT\cx \cb=\CB\cx
    \ifnum\ct<\cb \dt@column \fi
    \advance\cx by1\relax }}

\def\dt@column{%
  \rc=\numexpr(\cx-\x@one)*65536/\wid\relax
  \ra=\numexpr\s@one+\fmul{\numexpr\s@two-\s@one\relax}{\the\rc}\relax
  \ifnum\ra<1 \ra=1 \fi
  %% front sector floor and ceiling, projected to screen rows
  \ytr=\numexpr\HALFH-((\csname dtSc\the\dtFS\endcsname)*65536-\PZ)%
       *\ra/65536/65536\relax
  \ybr=\numexpr\HALFH-((\csname dtSf\the\dtFS\endcsname)*65536-\PZ)%
       *\ra/65536/65536\relax
  \yt=\ytr \ifnum\yt<\ct \yt=\ct \fi \ifnum\yt>\cb \yt=\cb \fi
  \yb=\ybr \ifnum\yb<\yt \yb=\yt \fi \ifnum\yb>\cb \yb=\cb \fi
  %% texture column and shading are constant all the way down a column
  \rd=\numexpr\re+\fmul{\numexpr\rf-\re\relax}{\the\rc}\relax
  \rd=\numexpr\rd*65536/\ra\relax
  %% u in plain map units, measured from the LINEDEF's start rather than the
  %% seg's: \dtUOFF carries the seg's offset along its linedef plus the
  %% sidedef's own x offset.  Adding it here, after the perspective divide,
  %% rather than to \u@one/\u@two before it, keeps \fmul's operands the size
  %% they were -- a 256-unit offset on a long seg would otherwise push the
  %% product past 2^31.  It is also exact: u is already map units by now.
  %% Wrapping to a texture width happens per surface, because an upper and a
  %% lower on one seg need not be the same width.
  \rd=\numexpr\rd/65536+\dtUOFF\relax
  \wlight=\numexpr14-\ra*80/\PPLANE/8192
     +(255-(\csname dtSl\the\dtFS\endcsname))/48\relax
  \ifnum\wlight>8 \wlight=8 \fi \ifnum\wlight<0 \wlight=0 \fi
  %% Doom's own light arithmetic for the real-texture path: lightnum is the
  %% sector light's top four bits, startmap is (15-lightnum)*4, and the
  %% distance term is the scale index halved.  The 320/\SCRW factor is
  %% Doom's own SCREENWIDTH/viewwidth, so a narrower view fades over the
  %% same world distance rather than over the same number of columns.
  \ifnum\dtCMN>0
    \dtWCM=\csname dtSl\the\dtFS\endcsname\relax
    \divide\dtWCM by16
    \advance\dtWCM by\dtFAKE
    \ifnum\dtWCM<0 \dtWCM=0 \fi \ifnum\dtWCM>15 \dtWCM=15 \fi
    \dtWCM=\numexpr(15-\dtWCM)*4\relax
    \dtTU=\numexpr\ra*320/\SCRW\relax
    \divide\dtTU by4096
    \ifnum\dtTU>47 \dtTU=47 \fi
    \divide\dtTU by2
    \advance\dtWCM by-\dtTU
    \ifnum\dtWCM<0 \dtWCM=0 \fi
    \ifnum\dtWCM>\numexpr\dtCMN-1\relax \dtWCM=\numexpr\dtCMN-1\relax \fi
    %% Doom's fixedcolormap: under the light visor or invulnerability every
    %% surface is drawn through ONE table whatever its light and distance.
    \ifnum\dtFIXCM<0 \else \dtWCM=\dtFIXCM \fi
    \edef\dt@cml{dtCM\the\dtWCM @}%
  \fi
  \vstepr=\numexpr65536*65536/\ra\relax
  %% The column half of a plane's world position.  Only fetched when a plane
  %% in this seg actually wants it, so a map with no flats pays one \if.
  \ifdt@pcol \edef\dt@uxq{\the\PU\cx}\edef\dt@uyq{\the\PV\cx}\fi
  \ifdt@psky
    \edef\dt@skyk{dtX\the\dtSKYT @\the\PK\cx @}%
    \expandafter\let\expandafter\dt@spick\csname dt@pg\the\PJ\cx\endcsname
  \fi
  \ifdt@solid \dt@solidcol \else\dt@portalcol \fi}

%%% Shade table and ramp base for one wall span.  \DOORRAMP and \TRAKRAMP
%%% are declared with the rest of the door machinery further down; they are
%%% only read here, at render time, so the forward reference is fine.
%%%
%%% Every surface setup ends by settling three things the span expansion
%%% reads: \dt@wrun (which run), \dt@vmod (where v wraps) and either
%%% \dt@uu/\dt@tx (procedural) or \dt@ux/\dt@cml (real).
\def\dt@wallstyle#1#2{%
  \dtTU=\rd \divide\dtTU by64 \multiply\dtTU by64
  \dtTU=\numexpr\rd-\dtTU\relax \ifnum\dtTU<0 \advance\dtTU by64 \fi
  \edef\dt@uu{\the\dtTU}%
  \def\dt@tx{#1}\wbase=\numexpr(#2)+\wlight\relax
  \def\dt@vmod{4194304}\let\dt@wrun\WallRun \dt@vsetup}
\def\dt@plainwall{\dt@wallstyle{dtT}{\csname dtSw\the\dtFS\endcsname}}
\def\dt@facewall{\dt@wallstyle{dtD}{\the\DOORRAMP}}
\def\dt@trackwall{\dt@wallstyle{dtD}{\the\TRAKRAMP}}

%%% A real composited texture: wrap u to its width, point \dt@ux at that
%%% column, and take the wrap for v from the texture's own height.
\def\dt@usetex#1{%
  \dtWT=#1\relax
  \dtTW=\csname dtXw\the\dtWT\endcsname\relax
  \dtTU=\rd \divide\dtTU by\dtTW \multiply\dtTU by\dtTW
  \dtTU=\numexpr\rd-\dtTU\relax \ifnum\dtTU<0 \advance\dtTU by\dtTW \fi
  %% u splits into its group of eight and its place within the group, and
  %% both are constant down the whole span.
  \dtTG=\dtTU \divide\dtTG by8
  \edef\dt@ux{dtX\the\dtWT @\the\dtTG @}%
  \multiply\dtTG by8
  \expandafter\let\expandafter\dt@pick
    \csname dt@pg\number\numexpr\dtTU-\dtTG\relax\endcsname
  \edef\dt@vmod{\csname dtXm\the\dtWT\endcsname}%
  \let\dt@wrun\WallTex \dt@vsetup}

%%% Is there real artwork for this surface?  A sidedef can name a texture
%%% the WAD has no TEXTURE1 entry for, which leaves the index allocated but
%%% the width at 0 -- so both tests are needed, and the answer is "fall back
%%% to the procedural table" rather than an error.
\def\dt@txok#1{\ifnum#1<0 \dt@realtexfalse\else
  \ifnum\csname dtXw#1\endcsname>0 \dt@realtextrue\else\dt@realtexfalse\fi\fi}

%%% v steps by \vstepr per screen row.  It has to be reduced modulo the
%%% surface's own wrap before the span runs, because the run only ever
%%% subtracts the modulus once per row.
\def\dt@vsetup{%
  \dtTV=\dt@vmod\relax
  \vstep=\vstepr
  \ifnum\vstep<\dtTV\else
    \tmpa=\vstep \divide\tmpa by\dtTV \multiply\tmpa by\dtTV
    \advance\vstep by-\tmpa
  \fi}

%%% ---- where a texture's v origin sits ------------------------------------
%%% Doom pegs a wall texture to one end of its surface and lets the other end
%%% run off, and the linedef's two pegging bits pick which.  The default for
%%% an upper is the BACK ceiling, which is what makes a door's slab travel
%%% upward with the ceiling instead of standing still while the gap opens.
%%% \dtTV holds the surface's wrap, which is also the texture's height in v,
%%% so "bottom aligned here" is just that value.
\def\dt@pegsolid{%
  \ifnum\dtPEGF>1
    \vstart=\numexpr\dtTV+(\yt-\ybr)*\vstep+\dtVOFF*65536\relax
  \else
    \vstart=\numexpr(\yt-\ytr)*\vstep+\dtVOFF*65536\relax
  \fi
  \dt@wrapv}
\def\dt@pegup{%
  \ifodd\dtPEGF
    \vstart=\numexpr(\yt-\ytr)*\vstep+\dtVOFF*65536\relax
  \else
    \vstart=\numexpr\dtTV+(\yt-\utr)*\vstep+\dtVOFF*65536\relax
  \fi
  \dt@wrapv}
\def\dt@peglow{%
  \ifnum\dtPEGF>1
    \vstart=\numexpr(\lb-\ytr)*\vstep+\dtVOFF*65536\relax
  \else
    \vstart=\numexpr(\lb-\lbr)*\vstep+\dtVOFF*65536\relax
  \fi
  \dt@wrapv}

\def\dt@solidcol{%
  \dt@txok{\the\dtWTM}%
  \ifdt@realtex \dt@usetex{\the\dtWTM}%
  \else \ifdt@doortrack \dt@trackwall \else \dt@plainwall \fi\fi
  \dt@pegsolid
  \dt@setceil
  \edef\dt@t{\dt@prun{\the\ct}{\the\yt}%
    \dt@wrun{\number\numexpr\yb-\yt\relax}{\the\vstart}{\the\vstep}%
            {\the\wbase}}%
  \dt@setfloor
  \edef\dt@b{\dt@prun{\the\yb}{\the\cb}}%
  \dtAppTop{\the\cx}{\dt@t}\dtPreBot{\the\cx}{\dt@b}%
  %% A solid wall takes the whole remaining opening, so its silhouette is the
  %% top clip meeting the bottom one.  \CD says the same thing more cheaply
  %% and is kept as the early-out for sprites wholly behind this wall.
  \dt@silT{\the\cx}{\the\ra}{\the\cb}%
  \CT\cx=\cb \CD\cx=\ra \advance\dtOPEN by-1 }

%%% A portal draws the step down from the front ceiling and the step up from
%%% the front floor, then leaves the gap between the back sector's own floor
%%% and ceiling open for whatever the traversal reaches next.
\def\dt@portalcol{%
  \utr=\numexpr\HALFH-((\csname dtSc\the\dtBS\endcsname)*65536-\PZ)%
       *\ra/65536/65536\relax
  \lbr=\numexpr\HALFH-((\csname dtSf\the\dtBS\endcsname)*65536-\PZ)%
       *\ra/65536/65536\relax
  \ut=\utr \ifnum\ut<\yt \ut=\yt \fi \ifnum\ut>\yb \ut=\yb \fi
  \lb=\lbr \ifnum\lb<\ut \lb=\ut \fi \ifnum\lb>\yb \lb=\yb \fi
  %% The upper span is the door itself when the back sector is a door: the
  %% shut ceiling has come down to the floor, so this one span is the whole
  %% visible slab.  It shrinks upward on its own as the ceiling rises,
  %% because \ut already tracks the animated dtSc.  With a real texture on
  %% the sidedef the structural door table is not needed -- the map names
  %% BIGDOOR1 or EXITDOOR there and that is what gets drawn.
  \dt@txok{\the\dtWTU}%
  \ifdt@realtex \dt@usetex{\the\dtWTU}%
  \else
    \ifdt@doorface \dt@facewall \else
      \ifdt@doortrack \dt@trackwall \else \dt@plainwall \fi\fi
  \fi
  \dt@pegup
  \dt@setceil
  \edef\dt@t{\dt@prun{\the\ct}{\the\yt}%
    \dt@wrun{\number\numexpr\ut-\yt\relax}{\the\vstart}{\the\vstep}%
            {\the\wbase}}%
  %% The lower span is the step up out of the doorway, not part of the door,
  %% so it goes back to the front sector's own wall.
  \dt@txok{\the\dtWTL}%
  \ifdt@realtex \dt@usetex{\the\dtWTL}%
  \else
    \ifdt@doorface \dt@plainwall \fi
  \fi
  \dt@peglow
  \dt@setfloor
  \edef\dt@b{\dt@wrun{\number\numexpr\yb-\lb\relax}{\the\vstart}{\the\vstep}%
      {\number\numexpr\wbase+1\relax}%
    \dt@prun{\the\yb}{\the\cb}}%
  \dtAppTop{\the\cx}{\dt@t}\dtPreBot{\the\cx}{\dt@b}%
  %% Whatever the upper and lower spans just covered is hidden from anything
  %% further away, which is what a sprite behind a part-open door or a raised
  %% floor has to be tested against.  Record only a real narrowing: a portal
  %% flush with the current opening contributes nothing here.
  \ifnum\ut>\ct \dt@silT{\the\cx}{\the\ra}{\the\ut}\fi
  \ifnum\lb<\cb \dt@silB{\the\cx}{\the\ra}{\the\lb}\fi
  \CT\cx=\ut \CB\cx=\lb
  %% A portal whose opening has collapsed -- a closed door -- is a wall, and
  %% must record its depth like one.  Leaving \CD alone here left the column
  %% reading as "never closed, infinitely far", so every sprite behind a shut
  %% door passed the depth test and drew straight through it.
  \ifnum\ut<\lb\else\CD\cx=\ra \advance\dtOPEN by-1 \fi}

%%% \dtTV is the current surface's wrap, settled by \dt@vsetup.
\def\dt@wrapv{%
  \ifnum\vstart<0
    \tmpa=\vstart \divide\tmpa by\dtTV \advance\tmpa by-1
    \multiply\tmpa by\dtTV \advance\vstart by-\tmpa
  \fi
  \ifnum\vstart<\dtTV\else
    \tmpa=\vstart \divide\tmpa by\dtTV \multiply\tmpa by\dtTV
    \advance\vstart by-\tmpa
  \fi}

%%% ==========================================================  THINGS (M4)
%%% Doom defers masked rendering: walls record per-column clip arrays and
%%% sprites are composited afterwards, back to front, tested against them.
%%% This renderer commits wall pixels into per-column strings during
%%% traversal, so the equivalent is to assemble each column first, then splice
%%% sprite bands into the middle of it.

%%% \dtSplice{column}{startrow}{rows}{data} -- fully expandable.  Emits the
%%% first #2 pixels of #1, then #4, then drops #3 pixels, then the rest.
%%% Two hex characters per pixel throughout.
\def\dtSplice#1#2#3#4{%
  \dt@spA{\number\numexpr#2\relax}{\number\numexpr#3\relax}{#4}#1\dt@end}
\def\dt@spA#1#2#3{%
  \ifnum#1>0 \expandafter\dt@spN\else\expandafter\dt@spI\fi{#1}{#2}{#3}}
\def\dt@spN#1#2#3#4#5{#4#5%
  \expandafter\dt@spA\expandafter{\number\numexpr#1-1\relax}{#2}{#3}}
\def\dt@spI#1#2#3{#3\dt@spD{#2}}
\def\dt@spD#1{%
  \ifnum#1>0 \expandafter\dt@spDN\else\expandafter\dt@spDE\fi{#1}}
\def\dt@spDN#1#2#3{\expandafter\dt@spD\expandafter{\number\numexpr#1-1\relax}}
%% the tail is grabbed in one delimited argument rather than walked
\def\dt@spDE#1#2\dt@end{#2}

%%% Splice into a stored column.  Only the column needs pre-expansion, so the
%%% single \expandafter idiom is safe here -- but it IS needed: \dt@spN grabs
%%% two character tokens, and macro arguments are not expanded while being
%%% grabbed, so an unexpanded \csname would be swallowed as pixel data.
\def\dtSpliceCol#1#2#3#4{%
  \edef\dt@sc{\csname dtFC#1\endcsname}%
  \expandafter\dt@scB\expandafter{\dt@sc}{#1}{#2}{#3}{#4}}
\def\dt@scB#1#2#3#4#5{%
  \expandafter\xdef\csname dtFC#2\endcsname{\dtSplice{#1}{#3}{#4}{#5}}}

%%% Assemble every column once the BSP pass is done: top spans, then whatever
%%% never closed, then bottom spans.  Sprites are spliced into these.
\def\dtBuildColumns{%
  \cx=0
  \dtWhile{\cx<\SCRW}{%
    \expandafter\xdef\csname dtFC\the\cx\endcsname{%
      \csname dtTS\the\cx\endcsname
      \Rep{\number\numexpr\CB\cx-\CT\cx\relax}{\hx{250}}%
      \csname dtBS\the\cx\endcsname}%
    \advance\cx by1\relax }}

%%% ---- thing classification -------------------------------------------------
%%% Two answers to "what does this thing look like".  \dtSpriteKey is the real
%%% one: Doom's mobjinfo maps a thing type to a sprite, and wad.tex reads that
%%% sprite's actual picture data.  \dtThingClass below is the fallback for
%%% everything that has no picture -- no WAD loaded, no sprite lump in the
%%% WAD, or a type outside this table -- and draws a coloured billboard
%%% instead.  The built-in map has no sprite data at all, so it is entirely on
%%% the fallback path.
%%%
%%% The key is four characters of sprite name plus a frame letter, and the
%%% loader appends the rotation.  Frame A is the first frame of the thing's
%%% spawn state for anything alive or collectable; corpses left lying around
%%% the map are a late frame of the monster's death animation, which is why
%%% the frame letter is part of the key rather than assumed.
\def\dtSpriteType#1#2{\expandafter\gdef\csname dtSN#1\endcsname{#2}}
\def\dtSpriteKey#1{\ifcsname dtSN#1\endcsname\csname dtSN#1\endcsname\fi}
\def\dt@empty{}
%% monsters
\dtSpriteType{3004}{POSSA}\dtSpriteType{9}{SPOSA}\dtSpriteType{65}{CPOSA}
\dtSpriteType{3001}{TROOA}\dtSpriteType{3002}{SARGA}\dtSpriteType{58}{SARGA}
\dtSpriteType{3005}{HEADA}\dtSpriteType{3006}{SKULA}\dtSpriteType{3003}{BOSSA}
\dtSpriteType{69}{BOS2A}\dtSpriteType{68}{BSPIA}\dtSpriteType{71}{PAINA}
\dtSpriteType{66}{SKELA}\dtSpriteType{67}{FATTA}\dtSpriteType{64}{VILEA}
\dtSpriteType{16}{CYBRA}\dtSpriteType{7}{SPIDA}\dtSpriteType{84}{SSWVA}
\dtSpriteType{72}{KEENA}\dtSpriteType{88}{BBRNA}
%% weapons and ammunition
\dtSpriteType{2001}{SHOTA}\dtSpriteType{82}{SGN2A}\dtSpriteType{2002}{MGUNA}
\dtSpriteType{2003}{LAUNA}\dtSpriteType{2004}{PLASA}\dtSpriteType{2005}{CSAWA}
\dtSpriteType{2006}{BFUGA}\dtSpriteType{2007}{CLIPA}\dtSpriteType{2048}{AMMOA}
\dtSpriteType{2008}{SHELA}\dtSpriteType{2049}{SBOXA}\dtSpriteType{2010}{ROCKA}
\dtSpriteType{2046}{BROKA}\dtSpriteType{2047}{CELLA}\dtSpriteType{17}{CELPA}
%% health, armour, powerups, keys
\dtSpriteType{2011}{STIMA}\dtSpriteType{2012}{MEDIA}\dtSpriteType{2014}{BON1A}
\dtSpriteType{2015}{BON2A}\dtSpriteType{2018}{ARM1A}\dtSpriteType{2019}{ARM2A}
\dtSpriteType{2013}{SOULA}\dtSpriteType{2022}{PINVA}\dtSpriteType{2023}{PSTRA}
\dtSpriteType{2024}{PINSA}\dtSpriteType{2025}{SUITA}\dtSpriteType{2026}{PMAPA}
\dtSpriteType{2045}{PVISA}\dtSpriteType{83}{MEGAA}\dtSpriteType{8}{BPAKA}
\dtSpriteType{5}{BKEYA}\dtSpriteType{40}{BSKUA}\dtSpriteType{13}{RKEYA}
\dtSpriteType{38}{RSKUA}\dtSpriteType{6}{YKEYA}\dtSpriteType{39}{YSKUA}
%% the barrel, and the scenery that fills out a room
\dtSpriteType{2035}{BAR1A}\dtSpriteType{2028}{COLUA}\dtSpriteType{30}{COL1A}
\dtSpriteType{31}{COL2A}\dtSpriteType{32}{COL3A}\dtSpriteType{33}{COL4A}
\dtSpriteType{36}{COL5A}\dtSpriteType{37}{COL6A}\dtSpriteType{41}{CEYEA}
\dtSpriteType{42}{FSKUA}\dtSpriteType{43}{TRE1A}\dtSpriteType{54}{TRE2A}
\dtSpriteType{44}{TBLUA}\dtSpriteType{45}{TGRNA}\dtSpriteType{46}{TREDA}
\dtSpriteType{55}{SMBTA}\dtSpriteType{56}{SMGTA}\dtSpriteType{57}{SMRTA}
\dtSpriteType{47}{SMITA}\dtSpriteType{48}{ELECA}\dtSpriteType{34}{CANDA}
\dtSpriteType{35}{CBRAA}\dtSpriteType{70}{FCANA}\dtSpriteType{49}{GOR1A}
\dtSpriteType{50}{GOR2A}\dtSpriteType{51}{GOR3A}\dtSpriteType{52}{GOR4A}
\dtSpriteType{53}{GOR5A}\dtSpriteType{25}{POL1A}\dtSpriteType{26}{POL6A}
\dtSpriteType{27}{POL4A}\dtSpriteType{28}{POL2A}\dtSpriteType{29}{POL3A}
\dtSpriteType{24}{POL5A}\dtSpriteType{73}{HDB1A}\dtSpriteType{74}{HDB2A}
\dtSpriteType{75}{HDB3A}\dtSpriteType{76}{HDB4A}\dtSpriteType{77}{HDB5A}
\dtSpriteType{78}{HDB6A}\dtSpriteType{79}{POB1A}\dtSpriteType{80}{POB2A}
\dtSpriteType{81}{BRS1A}
%% corpses: a late frame of the death animation, not frame A
\dtSpriteType{15}{PLAYN}\dtSpriteType{18}{POSSL}\dtSpriteType{19}{SPOSL}
\dtSpriteType{20}{TROOM}\dtSpriteType{21}{SARGN}\dtSpriteType{22}{HEADL}
\dtSpriteType{10}{PLAYW}\dtSpriteType{12}{PLAYW}

%%% ---- frame animation -----------------------------------------------------
%%% Doom's world is not still: a torch flickers, a bonus pulses, a zombie
%%% shifts its weight.  Each of those is a spawn state that loops over a few
%%% frames of one sprite, so all of it reduces to "which letter is showing
%%% now" -- the picture reader already keys on a frame letter.
%%%
%%%   \dtSpriteAnim{type}{letters}{tics}
%%%
%%% cycles that type through `letters', holding each for `tics' frames.  A
%%% type with no entry keeps the single frame in \dtSN and costs one
%%% \ifcsname per thing per frame.  The letters are written out rather than
%%% derived as base+n because Doom's sequences are not all runs: a bonus goes
%%% ABCDCB, which is a run up and back down again.
%%%
%%% The phase is the thing's own index, so two torches in a room do not
%%% flicker in lockstep.  Doom randomises it at spawn (P_SpawnMobj perturbs
%%% the first tic count); the index gives the same staggering, is
%%% deterministic -- which a regression harness needs -- and costs no array.
\newcount\dtANIM \dtANIM=0
\newcount\dtMFA \newcount\dtMFB \newcount\dtMFC
\def\dtSpriteAnim#1#2#3{%
  \dtMFC=0 \dt@ancount#2\relax
  \expandafter\gdef\csname dtMA#1\endcsname{#2}%
  \expandafter\xdef\csname dtMN#1\endcsname{\the\dtMFC}%
  \expandafter\gdef\csname dtMT#1\endcsname{#3}}
\def\dt@ancount#1{\ifx\relax#1\else
  \advance\dtMFC by1 \expandafter\dt@ancount\fi}

%% Which frame of type #1's sequence is showing.  \divide, not \numexpr:
%% \numexpr's division rounds, and this is modular arithmetic (INTERNALS 6).
%% Both inputs are non-negative, so the truncation has no sign case.
\def\dtAnimIndex#1{%
  \dtMFA=\dtANIM
  \divide\dtMFA by\csname dtMT#1\endcsname\relax
  \advance\dtMFA by\dtTI
  \dtMFB=\dtMFA
  \divide\dtMFB by\csname dtMN#1\endcsname\relax
  \multiply\dtMFB by\csname dtMN#1\endcsname\relax
  \advance\dtMFA by-\dtMFB}

%% The k-th letter of a sequence, 0-based.  Expandable, because it is used
%% inside the \edef that rebuilds the sprite key.
\def\dtAnimAt#1#2{\dt@anA{#1}#2\relax}
\def\dt@anA#1#2{\ifnum#1>0
  \expandafter\dt@anB\else\expandafter\dt@anC\fi{#1}{#2}}
\def\dt@anB#1#2{\expandafter\dt@anA\expandafter{\number\numexpr#1-1\relax}}
\def\dt@anC#1#2{#2\dt@angobble}
\def\dt@angobble#1\relax{}

%% A key is four characters of name and a frame letter, so swapping the frame
%% is dropping the fifth character and appending another.
\def\dt@dropf#1#2#3#4#5\relax{#1#2#3#4}
\def\dtAnimTick{\global\advance\dtANIM by1\relax}

%%% The spawn-state loops, from Doom's info.c.  \dtANIM advances once per
%%% rendered frame, so these hold for that many frames; the numbers are Doom's
%%% own tic counts.  Types absent here have a spawn state that does not loop
%%% -- the cacodemon, the pain elemental and the candelabra really do stand
%%% still -- and cost nothing but the \ifcsname that misses.
%% monsters: two frames of idle, ten tics each
\dtSpriteAnim{3004}{AB}{10}\dtSpriteAnim{9}{AB}{10}\dtSpriteAnim{65}{AB}{10}
\dtSpriteAnim{3001}{AB}{10}\dtSpriteAnim{3002}{AB}{10}\dtSpriteAnim{58}{AB}{10}
\dtSpriteAnim{3003}{AB}{10}\dtSpriteAnim{69}{AB}{10}\dtSpriteAnim{68}{AB}{10}
\dtSpriteAnim{66}{AB}{10}\dtSpriteAnim{67}{AB}{10}\dtSpriteAnim{64}{AB}{10}
\dtSpriteAnim{16}{AB}{10}\dtSpriteAnim{7}{AB}{10}\dtSpriteAnim{84}{AB}{10}
\dtSpriteAnim{3006}{AB}{10}
%% bonuses and powerups: a pulse, some of which runs up and back down
\dtSpriteAnim{2014}{ABCDCB}{6}\dtSpriteAnim{2015}{ABCDCB}{6}
\dtSpriteAnim{2013}{ABCDCB}{6}\dtSpriteAnim{2026}{ABCDCB}{6}
\dtSpriteAnim{2022}{ABCD}{6}\dtSpriteAnim{2024}{ABCD}{6}
\dtSpriteAnim{83}{ABCD}{6}\dtSpriteAnim{2045}{AB}{6}
\dtSpriteAnim{2018}{AB}{6}\dtSpriteAnim{2019}{AB}{6}
%% keys, both card and skull
\dtSpriteAnim{5}{AB}{10}\dtSpriteAnim{40}{AB}{10}\dtSpriteAnim{13}{AB}{10}
\dtSpriteAnim{38}{AB}{10}\dtSpriteAnim{6}{AB}{10}\dtSpriteAnim{39}{AB}{10}
%% scenery that burns or blinks
\dtSpriteAnim{2035}{AB}{6}\dtSpriteAnim{70}{ABC}{4}
\dtSpriteAnim{44}{ABCD}{4}\dtSpriteAnim{45}{ABCD}{4}\dtSpriteAnim{46}{ABCD}{4}
\dtSpriteAnim{55}{ABCD}{4}\dtSpriteAnim{56}{ABCD}{4}\dtSpriteAnim{57}{ABCD}{4}
\dtSpriteAnim{41}{ABCB}{6}\dtSpriteAnim{42}{ABC}{6}

%%% The billboard fallback.  This is a deliberate approximation of mobjinfo:
%%% enough to tell a monster from a pickup from a bit of scenery, sized and
%%% coloured accordingly.
\newcount\dtTKR \newcount\dtTKH \newcount\dtTKW \newcount\dtTKD
\def\dtThingClass#1{%
  \dtTKD=1 \dtTKR=208 \dtTKH=40 \dtTKW=8
  \ifnum#1<5 \dtTKD=0 \fi
  \ifnum#1=11 \dtTKD=0 \fi
  \ifnum#1=14 \dtTKD=0 \fi
  \ifnum#1=2035 \dtTKR=32 \dtTKH=32 \dtTKW=10 \else
    \ifnum#1>2999 \ifnum#1<3007 \dtTKR=64 \dtTKH=56 \dtTKW=20 \fi\fi
    \ifnum#1=9    \dtTKR=64 \dtTKH=56 \dtTKW=20 \fi
    \ifnum#1=16   \dtTKR=64 \dtTKH=56 \dtTKW=20 \fi
    \ifnum#1=58   \dtTKR=64 \dtTKH=56 \dtTKW=20 \fi
    \ifnum#1>63 \ifnum#1<72 \dtTKR=64 \dtTKH=56 \dtTKW=20 \fi\fi
    \ifnum#1>83 \ifnum#1<90 \dtTKR=64 \dtTKH=56 \dtTKW=20 \fi\fi
    \ifnum#1>1999 \ifnum#1<2030
      \dtTKR=112 \dtTKH=16 \dtTKW=10
      \ifnum#1=2011 \dtTKR=80 \fi
      \ifnum#1=2012 \dtTKR=80 \fi
      \ifnum#1=2014 \dtTKR=80 \fi
      \ifnum#1=2015 \dtTKR=96 \fi
      \ifnum#1=2018 \dtTKR=96 \fi
      \ifnum#1=2019 \dtTKR=96 \fi
    \fi\fi
  \fi}

%%% ---- projection ----------------------------------------------------------
\newcount\dtTI \newcount\dtNVIS \newcount\dtTSC \newcount\dtTXC
\newcount\dtTHW \newcount\dtTYT \newcount\dtTYB \newcount\dtTXA \newcount\dtTXB
\newcount\dtVK \newcount\dtVJ \newcount\dtBEST \newcount\dtBESTD
\newcount\dtSPSH \newcount\dtSPT \newcount\dtSPU \newcount\dtSPOK
\newcount\dtSPCT \newcount\dtSPCB \newcount\dtSPV
%% These were once \countdef'd from a hand-picked block at 300, on the premise
%% that plain TeX's \newcount allocator stops at 255.  It does not here: TeX
%% Live builds the pdftex format from etex.src, so \newcount runs well past
%% 255 (300 allocations reach 350 with no error).  The block was not merely
%% unnecessary, it was wrong -- under LaTeX the allocator starts around 250
%% and runs straight through 300..399, so \VIEWH and \dtHDA became the SAME
%% register and the HUD pass overwrote the view height every frame.  Nothing
%% in this file may pick a register number by hand; let the allocator do it.
\newcount\dtSRA \newcount\dtSRB \newcount\dtSRC
\newcount\dtSRD \newcount\dtSRE \newcount\dtSRW
\newcount\dtSRH \newcount\dtSRU \newcount\dtSRM
\newcount\dtSRV \newcount\dtSRT
\newcount\dtSRS

%% Geometry of thing \dtTI in view space.  \dtSPOK=0 means it is not on screen.
%% \dt@tk carries the sprite key and \ifdt@spr whether its picture is loaded;
%% both are read again by the drawing pass, which calls this a second time.
\newif\ifdt@spr \newif\ifdt@sprflip
\def\dtThingAngle{\ifcsname dtHa\the\dtTI\endcsname
  \csname dtHa\the\dtTI\endcsname\else0\fi}
\def\dt@thingkey{%
  \edef\dt@ty{\csname dtHt\the\dtTI\endcsname}%
  \edef\dt@tk{\dtSpriteKey{\dt@ty}}%
  \dt@sprfalse \dt@sprflipfalse
  \ifx\dt@tk\dt@empty\else
    %% Swap in the frame this thing is showing now, before the rotation digit
    %% goes on: the loader stored one picture per (name, frame, rotation), so
    %% the frame has to be settled first or the lookup misses and the thing
    %% falls back to rotation 0.
    %%
    %% A monster carries a state and the state owns the letter, so it does
    %% not go through the per-type spawn loop at all -- that loop is what
    %% every torch and bonus still uses.  The test is on dtOs rather than on
    %% the type, so a map with no monster of that type pays nothing.
    %% A missile is asked first because it is the one thing here whose
    %% four-character sprite NAME changes with its state -- plasma flies as
    %% PLSS and explodes as PLSE -- so it cannot go through the walkers
    %% below, which all assume the name was settled by \dtSN.
    \ifcsname dtJs\the\dtTI\endcsname
      \dt@misframe
    \else
    \ifcsname dtOs\the\dtTI\endcsname
      \dt@monframe
    \else
      \ifcsname dtMA\dt@ty\endcsname \dt@animframe \fi
    \fi\fi
    \dt@thingrot
  \fi}

\def\dt@animframe{%
  \dtAnimIndex{\dt@ty}%
  %% The sequence has to reach \dt@anA as literal characters.  Macro
  %% arguments are not expanded while being grabbed, so passing
  %% {\csname dtMA..\endcsname} hands \dt@anA the \csname token itself as the
  %% first "letter" -- the same trap \dtSpliceCol documents from the other
  %% side.  Building the call in an \edef expands the sequence in place,
  %% because \dt@anA is \noexpand'd there and so grabs nothing yet.
  \edef\dt@acall{\noexpand\dt@anA{\the\dtMFA}%
    \csname dtMA\dt@ty\endcsname\relax}%
  \edef\dt@tk{\expandafter\dt@dropf\dt@tk\relax\dt@acall}}

%% R_ProjectSprite: the rotation is the angle from the viewer to the thing
%% less the thing's own facing, biased by half a sector so rotation 1 is
%% centred on "facing you" rather than starting there.  A thing whose sprite
%% is the same from every side is stored as rotation 0 and answers for all
%% eight, which is what the fallback below picks up.
\def\dt@thingrot{%
  \dtAtan{\csname dtHx\the\dtTI\endcsname-\dtPSX}%
         {\csname dtHy\the\dtTI\endcsname-\dtPSY}%
  \dtnorm\dtAQ{\dtANG-\dtThingAngle+4608}%
  \divide\dtAQ by1024 \advance\dtAQ by1
  %% \dt@trot, not \dt@tr: play.tex's column-to-row transpose is a macro of
  %% that name, defined once at load, and this is rewritten once per thing per
  %% frame.  The engine's scratch used to overwrite it, so \dtAsciiFrame's
  %% `\expandafter\dt@tr\dt@col' expanded to a bare rotation digit and dropped
  %% the whole column into the enclosing mode as text -- one Overfull hbox per
  %% column, and only with a WAD, since without sprites this line never runs.
  %% Per-frame scratch must not share a name with anything that has to
  %% survive a frame.  See INTERNALS section 6, "register aliasing"; this is the
  %% same trap in the macro namespace.
  \edef\dt@trot{\the\dtAQ}%
  \ifcsname dtPw\dt@tk\dt@trot\endcsname
    \edef\dt@tk{\dt@tk\dt@trot}\dt@sprtrue
  \else
    \ifcsname dtPw\dt@tk0\endcsname \edef\dt@tk{\dt@tk0}\dt@sprtrue \fi
  \fi
  \ifdt@spr \ifcsname dtPf\dt@tk\endcsname \dt@sprfliptrue \fi\fi}

%% A thing that has been taken is gone from the world, and the cheapest place
%% to say so is here: \dtSPOK=0 is already "not on screen", and both the
%% collection pass and the draw pass go through this macro, so neither needs
%% to know that pickups exist.
\def\dt@thinggeom{%
  \dtSPOK=0
  \ifcsname dtHG\the\dtTI\endcsname\else
    %% A missile whose explosion has finished has left the world too, but it
    %% is not "taken" -- its slot is waiting to be handed to the next shot,
    %% so it cannot be marked with dtHG, which is permanent.  State 2 is that
    %% wait, and it is read here for the same reason dtHG is: everything that
    %% looks at a thing comes through this macro.
    \ifcsname dtJs\the\dtTI\endcsname
      \ifnum\csname dtJs\the\dtTI\endcsname<2 \dt@thinggeomB \fi
    \else \dt@thinggeomB \fi
  \fi}
\def\dt@thinggeomB{%
  \dt@thingkey
  \ifdt@spr \dtTKD=1
  \else \dtThingClass{\csname dtHt\the\dtTI\endcsname}\fi
  \ifnum\dtTKD=1
    \dtXform{\csname dtHx\the\dtTI\endcsname}{\csname dtHy\the\dtTI\endcsname}%
    \ifnum\ra>\NEARZ
      \dtTSC=\dtScale{\the\ra}%
      \dtTXC=\numexpr\SCRW/2-\rb*\PPLANE/\ra\relax
      \ifdt@spr \dt@sprbox \else \dt@billbox \fi
      \ifnum\dtTXB>0 \ifnum\dtTXA<\SCRW
        \ifnum\dtTYT<\dtTYB \dtSPOK=1 \fi
      \fi\fi
    \fi
  \fi}

%% Scale is PPLANE*65536/depth, so a world length becomes pixels with ONE
%% division by 65536.  Two divisions -- as the vertical formulas correctly
%% use, because their inputs are already 16.16 -- put every sprite under the
%% 1-pixel clamp no matter how close it was.
\def\dt@billbox{%
  \dtTHW=\numexpr\dtTKW*\dtTSC/65536\relax
  \ifnum\dtTHW<1 \dtTHW=1 \fi
  \dtTXA=\numexpr\dtTXC-\dtTHW\relax
  \dtTXB=\numexpr\dtTXC+\dtTHW\relax
  \dtTYB=\numexpr\HALFH-((\csname dtHz\the\dtTI\endcsname)*65536-\PZ)%
         *\dtTSC/65536/65536\relax
  \dtTYT=\numexpr\HALFH-((\csname dtHz\the\dtTI\endcsname+\dtTKH)%
         *65536-\PZ)*\dtTSC/65536/65536\relax}

%% A picture carries its own placement.  leftoffset is how far left of the
%% thing the picture starts and topoffset how far above it the top row sits,
%% so a torch stands on the floor and a corpse lies slightly below it without
%% either being special-cased.  Both boxes are built by adding the scaled
%% extent to the near edge rather than projecting both edges, which keeps
%% width and height in pixels exactly what the resampler is told they are.
\def\dt@sprbox{%
  \dtSRW=\csname dtPw\dt@tk\endcsname\relax
  \dtSRH=\csname dtPh\dt@tk\endcsname\relax
  \dtTXA=\numexpr\dtTXC-(\csname dtPl\dt@tk\endcsname)*\dtTSC/65536\relax
  \dtTXB=\numexpr\dtTXA+\dtSRW*\dtTSC/65536\relax
  \ifnum\dtTXB>\dtTXA\else \dtTXB=\numexpr\dtTXA+1\relax \fi
  \dtTYT=\numexpr\HALFH-((\csname dtHz\the\dtTI\endcsname
         +\csname dtPt\dt@tk\endcsname)*65536-\PZ)%
         *\dtTSC/65536/65536\relax
  \dtTYB=\numexpr\dtTYT+\dtSRH*\dtTSC/65536\relax
  \ifnum\dtTYB>\dtTYT\else \dtTYB=\numexpr\dtTYT+1\relax \fi}

%%% M19 puts masked segs through this same pass, so a map with no things at
%%% all can still have something to composite.
\def\dtDrawThings{\ifnum\numexpr\NTHING+\dtNMASK\relax>0 \dt@things \fi}

%%% Collect what is visible, then draw far to near.  Selection sort: the lists
%%% are short (only what is on screen), and it needs no auxiliary structure.
\def\dt@things{%
  \dtNVIS=0 \dtTI=0
  \dtWhile{\dtTI<\NTHING}{%
    \dt@thinggeom
    \ifnum\dtSPOK=1
      \expandafter\xdef\csname dtVI\the\dtNVIS\endcsname{\the\dtTI}%
      \expandafter\xdef\csname dtVD\the\dtNVIS\endcsname{\the\dtTSC}%
      \expandafter\xdef\csname dtVM\the\dtNVIS\endcsname{0}%
      \advance\dtNVIS by1
    \fi
    \advance\dtTI by1\relax }%
  \dt@maskvis
  \dtVK=0
  \dtWhile{\dtVK<\dtNVIS}{\dt@thingpick\advance\dtVK by1\relax }}

\def\dt@thingpick{%
  \dtBEST=-1 \dtBESTD=0 \dtVJ=0
  \dtWhile{\dtVJ<\dtNVIS}{%
    \ifnum\csname dtVD\the\dtVJ\endcsname>0
      \ifnum\dtBEST<0
        \dtBEST=\dtVJ \dtBESTD=\csname dtVD\the\dtVJ\endcsname\relax
      \else
        \ifnum\csname dtVD\the\dtVJ\endcsname<\dtBESTD
          \dtBEST=\dtVJ \dtBESTD=\csname dtVD\the\dtVJ\endcsname\relax
        \fi
      \fi
    \fi
    \advance\dtVJ by1\relax }%
  \ifnum\dtBEST<0\else
    \dtKI=\csname dtVI\the\dtBEST\endcsname\relax
    \dtKE=\csname dtVM\the\dtBEST\endcsname\relax
    \expandafter\xdef\csname dtVD\the\dtBEST\endcsname{-1}%
    \ifnum\dtKE=1
      \dt@maskdraw{\the\dtKI}%
    \else
      \dtTI=\dtKI
      \dt@thinggeom
      \ifnum\dtSPOK=1 \dt@thingcols \fi
    \fi
  \fi}

%%% The clip in force in column \cx at the current sprite's depth.  The lists
%%% run near to far, so walk while the record is nearer than the sprite and
%%% keep the last row seen; the sentinel pair stops the walk when the list
%%% runs out, and \dt@silstop swallows whatever is left.
%%%
%%% \dt@sl is expanded exactly once on the way in.  \dt@silscan grabs two
%%% brace groups and macro arguments are not expanded while being grabbed, so
%%% an unexpanded \csname would arrive as the scale.
\def\dt@silscan#1#2{\ifnum#1>\dtTSC \dtSPV=#2\relax
  \expandafter\dt@silscan\else\expandafter\dt@silstop\fi}
\def\dt@silstop#1\dt@silend{}
\def\dt@thingclip{%
  \dtSPV=0
  \edef\dt@sl{\csname dtZT\the\cx\endcsname}%
  \expandafter\dt@silscan\dt@sl{-1}{0}\dt@silend
  \dtSPCT=\dtSPV
  \dtSPV=\SCRH
  \edef\dt@sl{\csname dtZB\the\cx\endcsname}%
  \expandafter\dt@silscan\dt@sl{-1}{0}\dt@silend
  \dtSPCB=\dtSPV}

%%% One column band per screen column.  \CD rejects a sprite wholly behind
%%% the wall that closed the column without touching the lists at all; what
%%% survives that is clipped to the opening as it stood at the sprite's own
%%% depth, so a step, a doorframe or a door still rising cuts the band
%%% instead of being ignored.
\def\dt@thingcols{\ifdt@spr \dt@sprcols \else \dt@billcols \fi}

\def\dt@billcols{%
  %% distance shading, the same falloff the walls use
  \dtSPSH=\numexpr13-\dtTSC*80/\PPLANE/8192\relax
  \ifnum\dtSPSH>12 \dtSPSH=12 \fi \ifnum\dtSPSH<0 \dtSPSH=0 \fi
  \cx=\ifnum\dtTXA<0 0\else\the\dtTXA\fi\relax
  \rb=\ifnum\dtTXB>\SCRW \the\SCRW\else\the\dtTXB\fi\relax
  \dtWhile{\cx<\rb}{%
    \ifnum\dtTSC>\CD\cx \dt@thingclip \dt@thingband \fi
    \advance\cx by1\relax }}

%%% A flat billboard reads as a signpost, so taper the top toward the edges:
%%% the silhouette becomes a rounded column rather than a rectangle.  This is
%%% the one piece that real picture-format sprite data would replace.
\def\dt@thingband{%
  \dtSPT=\numexpr\cx-\dtTXC\relax
  \ifnum\dtSPT<0 \dtSPT=-\dtSPT \fi
  \dtSPT=\numexpr\dtSPT*65536/\dtTHW\relax
  \dtSPU=\numexpr\dtTYB-\dtTYT\relax
  \dtSPU=\numexpr\dtSPU*\dtSPT/65536*\dtSPT/65536/2\relax
  \yt=\numexpr\dtTYT+\dtSPU\relax
  \yb=\dtTYB
  %% These bounds start at 0 and \SCRH and only ever narrow, so clipping to
  %% the silhouette also clips to the screen.
  \ifnum\yt<\dtSPCT \yt=\dtSPCT \fi
  \ifnum\yb>\dtSPCB \yb=\dtSPCB \fi
  \ifnum\yt<\yb
    \edef\dt@sp{\Rep{\number\numexpr\yb-\yt\relax}%
      {\hx{\numexpr\dtTKR+\dtSPSH\relax}}}%
    \edef\dt@next{\noexpand\dtSpliceCol{\the\cx}{\the\yt}%
      {\number\numexpr\yb-\yt\relax}{\dt@sp}}%
    \dt@next
  \fi}

%%% ---- real sprite graphics ------------------------------------------------
%%% Doom's picture format is column oriented and transparency aware, which
%%% maps unusually well onto a renderer that already builds one column at a
%%% time and splices bands into them.  wad.tex has read each source column
%%% into a list of posts -- {topdelta}{length}{pixels} -- with the pixels
%%% still in the WAD's own bytes, which with a real PLAYPAL installed are
%%% already the palette indices the PDF wants.
%%%
%%% \numexpr rounds, and a resampler has to truncate.  For a non-negative
%%% numerator floor(a/b) is round((2a-b+1)/2b); the +1 is what stops an exact
%%% multiple rounding up to the next row.
\def\dtFloorDiv#1#2{\number\numexpr(2*(#1)-(#2)+1)/(2*(#2))\relax}

%%% The resampler, source driven: walk a post's pixels once and emit each into
%%% however many destination rows it covers.  That count is exactly what \Rep
%%% already produces, so a whole post is one expansion with no assignments and
%%% no random access into the run -- which matters, because the run is a plain
%%% string of bytes out of the file rather than an indexable array.
%%%
%%% Source driven costs O(source height) per screen column instead of
%%% O(screen height), and the column count already falls with distance, so a
%%% distant sprite is a few columns of 56 and a near one is no worse than the
%%% wall spans beside it.
%%%
%%% \dtSprRun{n}{f}{y} followed by 2n hex characters:
%%%   n  source pixels left in this post
%%%   f  16.16 destination offset of this pixel from the picture's top row
%%%   y  destination row it starts on
%%% \dtSRS, \dtSRT, \dtSPCT/\dtSPCB and \dtSRM are constant across a post and
%%% are read rather than passed.  They are registers, not macros, for a
%%% reason: `\ifnum\dtSRC>\dt@srH\relax' looks like it terminates the number
%%% scan, but \ifnum pushes that \relax back, so it arrives at the head of the
%%% branch and the branch value is never read.  A register needs no
%%% terminator at all.
\def\dt@srY#1{\number\numexpr\dtSRT+\dtFloorDiv{#1}{65536}\relax}
%% rows to emit for one source pixel, clipped to the column's silhouette.
%% Summed over a post this telescopes to min(y1,hi)-max(y0,lo) exactly, which
%% is what lets the splice below trust its own row count.
\def\dt@srK#1#2{\number\numexpr
  \ifnum#2>\dtSPCB \dtSPCB\else#2\fi
  -\ifnum#1<\dtSPCT \dtSPCT\else#1\fi\relax}
\def\dtSprRun#1#2#3{%
  \ifnum#1>0 \expandafter\dt@srA\else\expandafter\dt@srZ\fi{#1}{#2}{#3}}
\def\dt@srZ#1#2#3{}
%% Expand the next row boundary once and carry it: computing it for the clip
%% test and again for the recursion would double the divisions per pixel.
\def\dt@srA#1#2#3#4#5{%
  \expandafter\dt@srB\expandafter{\dt@srY{(#2)+\dtSRS}}{#1}{#2}{#3}{#4#5}}
\def\dt@srB#1#2#3#4#5{%
  \Rep{\dt@srK{#4}{#1}}{\dt@srPx#5}%
  \expandafter\dt@srD\expandafter{\number\numexpr(#3)+\dtSRS\relax}{#2}{#1}}
\def\dt@srD#1#2#3{%
  \expandafter\dtSprRun\expandafter{\number\numexpr(#2)-1\relax}{#1}{#3}}
%% One pixel.  The source byte's own two hex characters are the COLORMAP key,
%% so nothing is decoded going in or coming out.  Without a COLORMAP -- a
%% PWAD with no such lump -- the artwork goes through at full brightness.
\def\dt@srCM#1#2{\csname dtCM\the\dtSRM @#1#2\endcsname}
\def\dt@srRaw#1#2{#1#2}
\let\dt@srPx\dt@srRaw

%%% One post.  Its band follows from the same mapping the run uses, so the
%%% rows the splice replaces and the pixels the run emits are the same number
%%% by construction rather than by agreement.
\def\dt@onepost#1#2#3{%
  \dtSRA=\numexpr\dtSRS*(#1)\relax
  \dtSRB=\dt@srY{\the\dtSRA}\relax
  \dtSRC=\dt@srY{\numexpr\dtSRS*((#1)+(#2))\relax}\relax
  \dtSRD=\ifnum\dtSRB<\dtSPCT \dtSPCT\else\dtSRB\fi\relax
  \dtSRE=\ifnum\dtSRC>\dtSPCB \dtSPCB\else\dtSRC\fi\relax
  \ifnum\dtSRD<\dtSRE
    \edef\dt@sp{\dtSprRun{#2}{\the\dtSRA}{\the\dtSRB}#3}%
    \edef\dt@next{\noexpand\dtSpliceCol{\the\cx}{\the\dtSRD}%
      {\number\numexpr\dtSRE-\dtSRD\relax}{\dt@sp}}%
    \dt@next
  \fi}

%%% Transparency means one splice per post rather than one band per column,
%%% which is the whole reason the format is worth reading in the first place.
%%% The list ends with a {-1} sentinel, written when it was built.
\def\dt@postwalk#1#2#3{%
  \ifnum#1<0 \expandafter\dt@postend\else\expandafter\dt@postrun\fi
  {#1}{#2}{#3}}
\def\dt@postrun#1#2#3{\dt@onepost{#1}{#2}{#3}\dt@postwalk}
\def\dt@postend#1#2#3{}

\def\dt@sprcol{%
  \dtSRU=\dtFloorDiv{(\cx-\dtTXA)*\dtSRV}{65536}\relax
  \ifnum\dtSRU<0 \dtSRU=0 \fi
  \ifnum\dtSRU<\dtSRW\else \dtSRU=\numexpr\dtSRW-1\relax \fi
  %% Four of the eight rotations are stored as the mirror of another, so the
  %% picture is read once and the column order reverses here.  Nothing below
  %% this line knows about it.
  \ifdt@sprflip \dtSRU=\numexpr\dtSRW-1-\dtSRU\relax \fi
  \edef\dt@pl{\csname dtQ\dt@tk @\the\dtSRU\endcsname}%
  \expandafter\dt@postwalk\dt@pl}

%%% Doom picks a colormap from distance and sector light.  The falloff is the
%%% one the wall spans use, rescaled from the fourteen shades a ramp offers to
%%% COLORMAP's thirty-two, so artwork and geometry fade together.
%%
%% dtHl is the light a thing was PLACED in, which was the whole answer while
%% no sector's light ever changed.  A strobing room has to strobe the things
%% standing in it, so a thing that knows its sector reads that sector's light
%% as it is now; one that does not -- a missile, or anything dropped into the
%% built-in map by a fixture -- keeps the answer it always had.
\def\dtThingLight{\ifcsname dtHl\the\dtTI\endcsname
  \ifcsname dtHs\the\dtTI\endcsname
    \csname dtSl\csname dtHs\the\dtTI\endcsname\endcsname
  \else\csname dtHl\the\dtTI\endcsname\fi
  \else160\fi}
\def\dt@sprcols{%
  \dtSRM=\numexpr24-2*(\dtTSC*80/\PPLANE/8192)
        +(255-\dtThingLight)/32\relax
  \ifnum\dtSRM<0 \dtSRM=0 \fi \ifnum\dtSRM>31 \dtSRM=31 \fi
  %% A missile and its explosion are Doom's MF_SHADOW-free fullbright frames:
  %% info.c marks every one of them 32768, which R_ProjectSprite reads as
  %% "colormap 0 whatever the sector says".  A fireball that dimmed with
  %% distance would stop reading as something burning.
  \ifcsname dtJs\the\dtTI\endcsname \dtSRM=0 \fi
  %% ...and a fixed colormap outranks both, as it does in R_ProjectSprite.
  \ifnum\dtCMN>0 \ifnum\dtFIXCM<0 \else \dtSRM=\dtFIXCM \fi\fi
  \dtSRT=\dtTYT
  \dtSRS=\numexpr(\dtTYB-\dtTYT)*65536/\dtSRH\relax
  \dtSRV=\numexpr\dtSRW*65536/(\dtTXB-\dtTXA)\relax
  \cx=\ifnum\dtTXA<0 0\else\the\dtTXA\fi\relax
  \rb=\ifnum\dtTXB>\SCRW \the\SCRW\else\the\dtTXB\fi\relax
  \dtWhile{\cx<\rb}{%
    \ifnum\dtTSC>\CD\cx \dt@thingclip \dt@sprcol \fi
    \advance\cx by1\relax }}

%%% ------------------------------ masked middle textures (M19) --------------
%%% A middle texture on a two-sided line is the one wall surface Doom draws
%%% with its holes in: the room beyond shows through the bars.  It cannot be
%%% drawn during the wall pass, and not for want of trying -- a column fills
%%% from the OUTSIDE IN, the middle of a portal is exactly the gap the
%%% traversal has not reached yet, and writing into the middle of an
%%% unfinished column is the one thing this renderer's design does not do.
%%%
%%% Doom has the same problem and the same answer: a masked middle is not
%%% drawn with the walls at all.  R_StoreWallRange records a drawseg and
%%% R_DrawMasked composites it afterwards, against the drawseg's own clip
%%% arrays -- in the pass that draws the sprites.
%%%
%%% So a masked middle here is a sprite that happens to be a wall.  Every
%%% piece it needs already exists: \dtSplice writes into an assembled column,
%%% dtZT/dtZB record the per-column silhouette to clip against, \dt@thingpick
%%% sorts back to front, and \dtSprRun walks a chain of posts.  What the wall
%%% pass had to grow is the six numbers below, which \dt@column computes and
%%% throws away.
\newcount\dtNMASK  \dtNMASK=0
\newcount\dtMAXMASK \dtMAXMASK=64
\newcount\dtKB \newcount\dtKC \newcount\dtKD \newcount\dtKE
\newcount\dtKF \newcount\dtKG \newcount\dtKH \newcount\dtKI
\newcount\dtKQ \newcount\dtKS \newcount\dtKT \newcount\dtKU
\newcount\dtKV \newcount\dtKW
%%% dtKA<x> is the sky's per-column ray angle and is NOT one of these.

%% Recorded from \dt@segspans, which is the last point at which all six are
%% in hand.  Nothing else about the seg is kept: \dt@segsect and \dt@segtex
%% rebuild the sectors, the sidedef, the offsets and the pegging from its
%% index, which is why they are macros.
\def\dt@maskrec{%
  \ifnum\dtWTM<0 \else
  \ifnum\dtBS<0 \else
  \ifnum\dtNMASK<\dtMAXMASK
    \ifcsname dtXn\the\dtWTM\endcsname
      \expandafter\xdef\csname dtKm\the\dtNMASK\endcsname{%
        {\the\dtG}{\the\x@one}{\the\x@two}{\the\s@one}{\the\s@two}%
        {\the\re}{\the\rf}}%
      \global\advance\dtNMASK by1
    \fi
  \fi\fi\fi}

%% Sorted with the sprites, by the NEARER of its two ends.  A masked seg's
%% depth is a range where a sprite's is a point, so no single order is right
%% for both -- one end of the bars can be in front of a monster and the other
%% behind it.  Doom sorts drawsegs and sprites by their own rules and accepts
%% exactly the same limitation; the nearest scale is the cheap answer and it
%% is wrong in the cases Doom is wrong in.
%% \dt@mk is expanded exactly once on the way in, for the reason
%% \dt@thingclip gives about \dt@sl: \expandafter on the \csname builds the
%% control sequence but does not expand it, so the walker would be handed one
%% token where it wanted seven brace groups.
\def\dt@maskvis{%
  \dtKI=0
  \dtWhile{\dtKI<\dtNMASK}{%
    \edef\dt@mk{\csname dtKm\the\dtKI\endcsname}%
    \expandafter\dt@masknear\dt@mk
    \expandafter\xdef\csname dtVI\the\dtNVIS\endcsname{\the\dtKI}%
    \expandafter\xdef\csname dtVD\the\dtNVIS\endcsname{\the\dtKS}%
    \expandafter\xdef\csname dtVM\the\dtNVIS\endcsname{1}%
    \advance\dtNVIS by1
    \advance\dtKI by1\relax }}
\def\dt@masknear#1#2#3#4#5#6#7{\dtKS=\ifnum#4>#5 #4\else#5\fi\relax}

\def\dt@maskdraw#1{%
  \dtKI=\numexpr#1\relax
  \edef\dt@mk{\csname dtKm\the\dtKI\endcsname}%
  \expandafter\dt@maskgo\dt@mk}
\def\dt@maskgo#1#2#3#4#5#6#7{%
  \dt@segsect{#1}\dt@segtex
  \x@one=#2\relax \x@two=#3\relax
  \s@one=#4\relax \s@two=#5\relax
  \re=#6\relax \rf=#7\relax
  \wid=\numexpr\x@two-\x@one\relax
  \dt@masksurf
  \cx=\ifnum\x@one<0 0\else\the\x@one\fi\relax
  \rb=\ifnum\x@two>\SCRW \the\SCRW\else\the\x@two\fi\relax
  \dtWhile{\cx<\rb}{\dt@maskcol\advance\cx by1\relax }}

%% R_RenderMaskedSegRange's pegging.  DONTPEGBOTTOM stands the texture on the
%% bottom of the opening and otherwise it hangs from the top -- and either
%% way it is drawn ONCE rather than tiled, which is the visual difference
%% between a masked middle and a wall and the reason the wall path's \dt@wrapv
%% is not wanted here.
\def\dt@masksurf{%
  \dtKT=\csname dtSc\the\dtFS\endcsname\relax
  \ifnum\csname dtSc\the\dtBS\endcsname<\dtKT
    \dtKT=\csname dtSc\the\dtBS\endcsname\relax \fi
  \dtKB=\csname dtSf\the\dtFS\endcsname\relax
  \ifnum\csname dtSf\the\dtBS\endcsname>\dtKB
    \dtKB=\csname dtSf\the\dtBS\endcsname\relax \fi
  \dtKH=\csname dtXh\the\dtWTM\endcsname\relax
  \dtKW=\csname dtXw\the\dtWTM\endcsname\relax
  \ifnum\dtPEGF>1
    \dtKU=\numexpr\dtKB+\dtKH+\dtVOFF\relax
  \else
    \dtKU=\numexpr\dtKT+\dtVOFF\relax
  \fi
  \dtKV=\numexpr\dtKU-\dtKH\relax }

%% \CD rejects a column the wall pass closed nearer than this, exactly as it
%% does for a sprite.
\def\dt@maskcol{%
  \rc=\numexpr(\cx-\x@one)*65536/\wid\relax
  \ra=\numexpr\s@one+\fmul{\numexpr\s@two-\s@one\relax}{\the\rc}\relax
  \ifnum\ra<1 \ra=1 \fi
  \ifnum\ra>\CD\cx \dt@maskcolB \fi}

%% u, derived exactly as \dt@column derives it -- same interpolation, same
%% perspective divide, same \dtUOFF added after it for the same reason.  The
%% wrap is by \divide against the texture's own width, because widths are
%% powers of two and heights are not and nothing here may assume a mask.
\def\dt@maskcolB{%
  \rd=\numexpr\re+\fmul{\numexpr\rf-\re\relax}{\the\rc}\relax
  \rd=\numexpr\rd*65536/\ra\relax
  \rd=\numexpr\rd/65536+\dtUOFF\relax
  \dtKC=\rd \divide\dtKC by\dtKW \multiply\dtKC by\dtKW
  \dtKC=\numexpr\rd-\dtKC\relax \ifnum\dtKC<0 \advance\dtKC by\dtKW \fi
  \ifcsname dtQMK\the\dtWTM @\the\dtKC\endcsname \dt@maskcolC \fi}

\def\dt@maskcolC{%
  \dtSRT=\numexpr\HALFH-(\dtKU*65536-\PZ)*\ra/65536/65536\relax
  \dtKE=\numexpr\HALFH-(\dtKV*65536-\PZ)*\ra/65536/65536\relax
  \ifnum\dtKE>\dtSRT
    \dtSRS=\numexpr(\dtKE-\dtSRT)*65536/\dtKH\relax
    %% The silhouette at THIS column's depth, which is \dt@thingclip's whole
    %% job; the seg's own record sits at exactly \ra and \dt@silscan takes
    %% only what is strictly nearer, so the bars are not clipped away by the
    %% opening they are hanging in.
    \dtTSC=\ra
    \dt@thingclip
    %% ...and then the opening itself, so a texture taller than the gap is
    %% cut off by it rather than spilling through the floor above.
    \dtKF=\numexpr\HALFH-(\dtKT*65536-\PZ)*\ra/65536/65536\relax
    \dtKG=\numexpr\HALFH-(\dtKB*65536-\PZ)*\ra/65536/65536\relax
    \ifnum\dtSPCT<\dtKF \dtSPCT=\dtKF \fi
    \ifnum\dtSPCB>\dtKG \dtSPCB=\dtKG \fi
    \ifnum\dtSPCT<0 \dtSPCT=0 \fi
    \ifnum\dtSPCB>\SCRH \dtSPCB=\SCRH \fi
    \ifnum\dtSPCT<\dtSPCB
      \dt@maskshade
      \edef\dt@pl{\csname dtQMK\the\dtWTM @\the\dtKC\endcsname}%
      \expandafter\dt@postwalk\dt@pl
    \fi
  \fi}

%% \dt@column's colormap arithmetic, landing in \dtSRM because that is what
%% the post walker's \dt@srCM reads.  Same light, same fake contrast, same
%% halved scale index: the bars have to fade with the wall they are set in.
\def\dt@maskshade{%
  \dtSRM=0
  \ifnum\dtCMN>0
    \dtKD=\csname dtSl\the\dtFS\endcsname\relax
    \divide\dtKD by16
    \advance\dtKD by\dtFAKE
    \ifnum\dtKD<0 \dtKD=0 \fi \ifnum\dtKD>15 \dtKD=15 \fi
    \dtKD=\numexpr(15-\dtKD)*4\relax
    \dtKQ=\numexpr\ra*320/\SCRW\relax
    \divide\dtKQ by4096
    \ifnum\dtKQ>47 \dtKQ=47 \fi
    \divide\dtKQ by2
    \advance\dtKD by-\dtKQ
    \ifnum\dtKD<0 \dtKD=0 \fi
    \ifnum\dtKD>\numexpr\dtCMN-1\relax \dtKD=\numexpr\dtCMN-1\relax \fi
    \ifnum\dtFIXCM<0 \else \dtKD=\dtFIXCM \fi
    \dtSRM=\dtKD
  \fi}

%%% ==========================================  PLAYER STATE / PICKUPS (M14)
%%% Until now the player has been a camera: a position, a facing and an eye
%%% height.  Doom's player is also a row of numbers, and the status bar is
%%% those numbers rather than a picture -- so the bar cannot be finished
%%% without them, and the numbers are worth nothing until something in the
%%% world changes one.  Pickups are that something, and they are the first
%%% thing in this project that alters the map permanently: a thing that is
%%% taken stops being drawn and stops being collectable, for the rest of the
%%% run.
%%%
%%% P_TouchSpecialThing is a switch over sprite names.  The equivalent here
%%% is a table keyed by thing type, in the shape \dtSpriteType already uses:
%%% a kind and two payload fields.  A type with no entry is not collectable
%%% at all, which is how every monster, every corpse and every bit of scenery
%%% comes out of it, and how the built-in map's things stay where they are.

\newcount\dtHEALTH \dtHEALTH=100      % Doom's MAXHEALTH; bonuses go to 200
\newcount\dtARMOR  \dtARMOR=0
\newcount\dtARMCL  \dtARMCL=0         % armour class: 1 green, 2 blue
\newcount\dtWEAP   \dtWEAP=1          % readyweapon; 1 is the pistol
\newcount\dtBACKP  \dtBACKP=0
\newcount\dtNPICK  \dtNPICK=0         % how many things have been taken
%%% Did this touch actually give anything.  EVERY assignment to it ends in an
%%% explicit \relax, and that is not decoration: \fi, \else and \or are
%%% EXPANDABLE, so a bare `\dtGOT=1' at the end of a branch keeps scanning for
%%% digits, expands the \fi that should have closed the branch, runs on
%%% through the \or into the `\ifnum\dtGOT=1' below -- and evaluates that
%%% test while the assignment that would have made it true is still in flight.
%%% The result was a health bonus that raised health and was never collected,
%%% with no error and a plausible frame.  INTERNALS section 6's rule about
%%% terminating number scans applies to conditionals, not just to line ends.
\newcount\dtGOT

%%% Doom's blockmap test is a box, not a circle: |dx| and |dy| are each
%%% compared against the sum of the two radii.  An item is 20 and the player
%%% 16, so this is 36.  It is a count rather than a constant in the source so
%%% that a fixture can widen it, on the \DOORRAMP precedent.
\newcount\dtPICKR \dtPICKR=36

%%% Ammunition is four numbered slots rather than four registers, because the
%%% weapon table has to index it: 0 bullets, 1 shells, 2 cells, 3 rockets.
%%%   dtUA<k>  how much is held        dtUM<k>  the maximum
%%%   dtUW<n>  present if weapon n is owned
%%%   dtUT<n>  weapon n's ammo slot, -1 for none
%%%   dtUS<n>  weapon n's on-screen picture
%%%   dtUK<n>  present if key n is held (0 blue, 1 yellow, 2 red)
\def\dt@ammodef#1#2#3{%
  \expandafter\gdef\csname dtUA#1\endcsname{#2}%
  \expandafter\gdef\csname dtUM#1\endcsname{#3}}
\dt@ammodef{0}{50}{200}\dt@ammodef{1}{0}{50}
\dt@ammodef{2}{0}{300}\dt@ammodef{3}{0}{50}

%%% Doom's clipammo[]: what one "clip" of each slot is worth.  A found weapon
%%% gives two of them, a pickup gives the number in its table row.
\def\dtClipAmmo#1{\ifcase#1 10\or4\or20\or1\else0\fi}

\def\dt@weapdef#1#2#3{%
  \expandafter\gdef\csname dtUT#1\endcsname{#2}%
  \expandafter\gdef\csname dtUS#1\endcsname{#3}}
\dt@weapdef{0}{-1}{PUNGA0}\dt@weapdef{1}{0}{PISGA0}
\dt@weapdef{2}{1}{SHTGA0}\dt@weapdef{3}{0}{CHGGA0}
\dt@weapdef{4}{3}{MISGA0}\dt@weapdef{5}{2}{PLSGA0}
\dt@weapdef{6}{2}{BFGGA0}\dt@weapdef{7}{-1}{SAWGA0}
\dt@weapdef{8}{1}{SHT2A0}
%% the fist and the pistol, as Doom starts you
\expandafter\gdef\csname dtUW0\endcsname{1}
\expandafter\gdef\csname dtUW1\endcsname{1}

%%% ---- the item table ------------------------------------------------------
%%% \dtItem{thing type}{kind}{a}{b}.  The kinds are Doom's give functions:
%%%
%%%   1 body        a units of health, refused at 100 or over
%%%   2 bonus       a units of health, always taken, ceiling 200
%%%   3 armour      class a, refused if the armour held is already as good
%%%   4 armbonus    a points of armour, always taken, ceiling 200
%%%   5 ammo        b clips of slot a
%%%   6 weapon      weapon a, with two clips of its ammo
%%%   7 key         key a
%%%   8 backpack    doubles every maximum, then a clip of each
%%%   9 soul        a units of health, ceiling 200
%%%  10 mega        health 200 and blue armour
%%%  11 berserk     a units of health, out comes the fist, and it hits ten
%%%                times as hard for the rest of the level
%%%  12 power       timed power a: 0 invulnerability, 1 partial invisibility,
%%%                2 the radiation suit, 3 light amplification
%%%
%%% The computer map is the one powerup still absent, and for the reason all
%%% five used to be: there is no automap here for it to fill in, so
%%% collecting one would delete it from the level in exchange for nothing.
\def\dtItem#1#2#3#4{%
  \expandafter\gdef\csname dtIK#1\endcsname{#2}%
  \expandafter\gdef\csname dtIA#1\endcsname{#3}%
  \expandafter\gdef\csname dtIB#1\endcsname{#4}}
%% health and armour
\dtItem{2011}{1}{10}{0}\dtItem{2012}{1}{25}{0}
\dtItem{2014}{2}{1}{0}\dtItem{2015}{4}{1}{0}
\dtItem{2018}{3}{1}{0}\dtItem{2019}{3}{2}{0}
\dtItem{2013}{9}{100}{0}\dtItem{83}{10}{0}{0}\dtItem{2023}{11}{100}{0}
%% the timed powers
\dtItem{2022}{12}{0}{0}\dtItem{2024}{12}{1}{0}
\dtItem{2025}{12}{2}{0}\dtItem{2045}{12}{3}{0}
%% ammunition
\dtItem{2007}{5}{0}{1}\dtItem{2048}{5}{0}{5}
\dtItem{2008}{5}{1}{1}\dtItem{2049}{5}{1}{5}
\dtItem{2047}{5}{2}{1}\dtItem{17}{5}{2}{5}
\dtItem{2010}{5}{3}{1}\dtItem{2046}{5}{3}{5}
\dtItem{8}{8}{0}{0}
%% weapons
\dtItem{2001}{6}{2}{0}\dtItem{82}{6}{8}{0}\dtItem{2002}{6}{3}{0}
\dtItem{2003}{6}{4}{0}\dtItem{2004}{6}{5}{0}\dtItem{2006}{6}{6}{0}
\dtItem{2005}{6}{7}{0}
%% keys, card and skull alike
\dtItem{5}{7}{0}{0}\dtItem{40}{7}{0}{0}
\dtItem{6}{7}{1}{0}\dtItem{39}{7}{1}{0}
\dtItem{13}{7}{2}{0}\dtItem{38}{7}{2}{0}

%%% ---- giving --------------------------------------------------------------
%%% Every one of these sets \dtGOT when it actually gave something, and that
%%% is what decides whether the thing disappears.  It is \dtFIRED again, for
%%% the same reason: a medikit picked up at full health would otherwise
%%% vanish and heal nothing, which is exactly the bug Doom's return values
%%% exist to prevent.
%%%
%%% Own registers, not \tmpa-\tmpk and not the touch scan's: these run below
%%% a loop that is holding its own bounds, and INTERNALS section 6 records what
%%% sharing a register across a call into anything costs.
\newcount\dtGVA \newcount\dtGVB \newcount\dtGVC \newcount\dtGVI

\def\dtGiveBody#1{%
  \ifnum\dtHEALTH<100
    \advance\dtHEALTH by#1\relax
    \ifnum\dtHEALTH>100 \dtHEALTH=100 \fi
    \dtGOT=1\relax
  \fi}

\def\dtGiveArmor#1{%
  \ifnum\dtARMOR<\numexpr100*(#1)\relax
    \dtARMOR=\numexpr100*(#1)\relax
    \dtARMCL=#1\relax
    \dtGOT=1\relax
  \fi}

\def\dtGiveAmmo#1#2{%
  \dtGVA=\csname dtUA#1\endcsname\relax
  \dtGVB=\csname dtUM#1\endcsname\relax
  \ifnum\dtGVA<\dtGVB
    \dtGVC=\dtClipAmmo{#1}\relax
    \multiply\dtGVC by\numexpr#2\relax
    \advance\dtGVA by\dtGVC
    \ifnum\dtGVA>\dtGVB \dtGVA=\dtGVB \fi
    \expandafter\xdef\csname dtUA#1\endcsname{\the\dtGVA}%
    \dtGOT=1\relax
  \fi}

%%% P_GiveWeapon: two clips with a weapon found on the floor, and the weapon
%%% comes up only if it is new -- picking up a second shotgun is worth its
%%% shells and nothing else.
\def\dtGiveWeapon#1{%
  \ifnum\csname dtUT#1\endcsname<0 \else
    \dtGiveAmmo{\csname dtUT#1\endcsname}{2}%
  \fi
  \ifcsname dtUW#1\endcsname\else
    \expandafter\gdef\csname dtUW#1\endcsname{1}%
    \dtWEAP=#1\relax
    %% Doom's evil grin is a NEW weapon, not any weapon: walking over a
    %% second shotgun is worth its shells and no expression at all.  So it
    %% goes inside this branch, beside the two things that already know the
    %% difference.
    \dtGRINC=\dtEVILC
    \dtGOT=1\relax
  \fi}

\def\dtGiveKey#1{%
  \ifcsname dtUK#1\endcsname\else
    \expandafter\gdef\csname dtUK#1\endcsname{1}%
  \fi
  \dtGOT=1\relax}

\def\dt@gbody{\dtGiveBody{\dt@iva}}
\def\dt@gbonus{%
  \advance\dtHEALTH by\dt@iva\relax
  \ifnum\dtHEALTH>200 \dtHEALTH=200 \fi
  \dtGOT=1\relax}
\def\dt@garmour{\dtGiveArmor{\dt@iva}}
\def\dt@garmbonus{%
  \advance\dtARMOR by\dt@iva\relax
  \ifnum\dtARMOR>200 \dtARMOR=200 \fi
  \ifnum\dtARMCL=0 \dtARMCL=1 \fi
  \dtGOT=1\relax}
\def\dt@gammo{\dtGiveAmmo{\dt@iva}{\dt@ivb}}
\def\dt@gweapon{\dtGiveWeapon{\dt@iva}}
\def\dt@gkey{\dtGiveKey{\dt@iva}}
\def\dt@gsoul{%
  \advance\dtHEALTH by\dt@iva\relax
  \ifnum\dtHEALTH>200 \dtHEALTH=200 \fi
  \dtGOT=1\relax}
\def\dt@gmega{\dtHEALTH=200 \dtGiveArmor{2}\dtGOT=1\relax}
\def\dt@gbers{%
  \dtGiveBody{\dt@iva}%
  \dtWEAP=0
  \ifdt@pow \dtBERSERK=1 \fi
  \dtGOT=1\relax}

%%% ---- the timed powers (M23) ----------------------------------------------
%%% P_GivePower.  Each is a counter that \dtPlayerTick runs down, and each is
%%% read in exactly one place: invulnerability in \dtDamagePlayer, the suit
%%% in \dt@slime, the visor in \dt@powerview, and invisibility by whatever is
%%% trying to aim at you.  Counts, so a fixture can hand one out directly.
%%%
%%% Doom's lengths are 30, 60, 60 and 120 seconds.  They are turned into
%%% frames by \dtMISRATE's ratio -- two frames to five tics -- for the reason
%%% the weapon periods are: a power has to last as long as it does against
%%% the pace of the player wearing it.
%%%
%%% \DTNOPOW puts the four back on the floor where M22 left them, and takes
%%% the strength out of the berserk pack.
\newif\ifdt@pow \dt@powtrue
\newcount\dtINVUL \newcount\dtINVIS \newcount\dtSUIT \newcount\dtLAMP
\newcount\dtBERSERK
\dtINVUL=0 \dtINVIS=0 \dtSUIT=0 \dtLAMP=0 \dtBERSERK=0
\def\dtPowerTime#1{\ifcase#1 420\or840\or840\or1680\else0\fi}
\def\dt@gpower{%
  \ifdt@pow
    \ifcase\dt@iva\relax
        \dtINVUL=\dtPowerTime{0}\relax
      \or \dtINVIS=\dtPowerTime{1}\relax
      \or \dtSUIT=\dtPowerTime{2}\relax
      \or \dtLAMP=\dtPowerTime{3}\relax
    \fi
    \dtGOT=1\relax
  \fi}

%%% What the powers do to the VIEW, settled once a tick.  The visor draws
%%% everything through COLORMAP 1 and invulnerability through table 32, the
%%% inverse map -- and both blink through their last fifty frames, which is
%%% Doom's `powers > 4*32 || powers & 8' with the tics turned into frames.
%%% Invulnerability is asked second so that it wins, as it does in Doom.
%%%
%%% Table 32 exists only when wad.tex read a real COLORMAP, so without one
%%% the powers still work and simply cannot be seen.
\newcount\dtPWA \newcount\dtPWB \newcount\dtPWC
\def\dt@blink#1{%
  \dtPWA=1
  \ifnum#1<52
    \dtPWB=#1\relax \divide\dtPWB by3
    \ifodd\dtPWB \else \dtPWA=0 \fi
  \fi}
\def\dt@powerview{%
  \dtPWC=-1
  \ifnum\dtLAMP>0  \dt@blink{\dtLAMP}\ifnum\dtPWA=1 \dtPWC=1 \fi\fi
  \ifnum\dtINVUL>0 \dt@blink{\dtINVUL}\ifnum\dtPWA=1 \dtPWC=32 \fi\fi
  \ifcsname dtCM32@00\endcsname \else \dtPWC=-1 \fi
  \dt@setfixcm{\the\dtPWC}}
%%% Switching the fixed colormap on takes the distance term out of the plane
%%% shader, and switching it off puts back whatever was there -- which is not
%%% necessarily 20480, because a fixture is allowed to have moved it.
\def\dt@plitkeep{20480}
\def\dt@setfixcm#1{%
  \ifnum#1<0
    \ifnum\dtFIXCM<0 \else \let\dt@plit\dt@plitkeep \fi
  \else
    \ifnum\dtFIXCM<0 \let\dt@plitkeep\dt@plit \def\dt@plit{0}\fi
  \fi
  \dtFIXCM=#1\relax}

%%% The backpack doubles every maximum once, then gives a clip of each --
%%% including of a slot that was already full, which is why the doubling has
%%% to happen first.
\def\dt@gpack{%
  \ifnum\dtBACKP=0
    \dtBACKP=1 \dtGOT=1\relax
    \dtGVI=0
    \dtWhile{\dtGVI<4}{%
      \expandafter\xdef\csname dtUM\the\dtGVI\endcsname
        {\number\numexpr2*\csname dtUM\the\dtGVI\endcsname\relax}%
      \advance\dtGVI by1\relax }%
  \fi
  \dtGVI=0
  \dtWhile{\dtGVI<4}{\dtGiveAmmo{\the\dtGVI}{1}\advance\dtGVI by1\relax }}

%%% ---- the touch scan ------------------------------------------------------
%%% Doom reaches items through the blockmap from P_CheckPosition.  There is no
%%% blockmap here -- \dtCrossLines walks the whole seg list for the same
%%% reason -- so this walks the thing list, which on MAP01 is 61 entries
%%% against collision's 601 segs and costs correspondingly less.
%%%
%%% Only things whose type is in the item table are considered at all, and a
%%% thing already taken is skipped by the same test that stops it being
%%% drawn.  dtHG<i> is dtLU for things: a flag hung off a csname, absent
%%% until it means something.
\newcount\dtTUI \newcount\dtTUA \newcount\dtTUB \newcount\dtTUD

\def\dtTouchThings{\ifnum\NTHING>0 \dt@touchall \fi}
\def\dt@touchall{%
  \dtTUI=0
  \dtWhile{\dtTUI<\NTHING}{%
    \ifcsname dtHG\the\dtTUI\endcsname\else
      \ifcsname dtIK\csname dtHt\the\dtTUI\endcsname\endcsname \dt@touch \fi
    \fi
    \advance\dtTUI by1\relax }}

%%% The box test, then a height test.  Doom compares the two objects' z
%%% extents; things here carry the floor of the sector they stand on rather
%%% than a live z, so this compares that against the floor under the player.
%%% Without it a stimpack on a ledge would be collectable from underneath it.
\def\dt@touch{%
  \dtTUA=\numexpr\csname dtHx\the\dtTUI\endcsname-\dtPSX\relax
  \ifnum\dtTUA<0 \dtTUA=-\dtTUA \fi
  \ifnum\dtTUA<\dtPICKR
    \dtTUB=\numexpr\csname dtHy\the\dtTUI\endcsname-\dtPSY\relax
    \ifnum\dtTUB<0 \dtTUB=-\dtTUB \fi
    \ifnum\dtTUB<\dtPICKR
      \dtTUD=\numexpr\csname dtHz\the\dtTUI\endcsname-\dtFloorZ\relax
      \ifnum\dtTUD<0 \dtTUD=-\dtTUD \fi
      \ifnum\dtTUD<\PHEAD \dt@give \fi
    \fi
  \fi}

%%% The payload fields reach the give routines as literal digits.  A macro's
%%% arguments are not expanded while being grabbed, so handing \dtGiveBody a
%%% bare \csname would pass it the token rather than the number -- the trap
%%% \dtSpliceCol documents from the other side.  Two \edefs settle it once.
\def\dt@give{%
  \dtGOT=0
  \edef\dt@ity{\csname dtHt\the\dtTUI\endcsname}%
  \edef\dt@iva{\csname dtIA\dt@ity\endcsname}%
  \edef\dt@ivb{\csname dtIB\dt@ity\endcsname}%
  \ifcase\csname dtIK\dt@ity\endcsname\relax
    \or\dt@gbody
    \or\dt@gbonus
    \or\dt@garmour
    \or\dt@garmbonus
    \or\dt@gammo
    \or\dt@gweapon
    \or\dt@gkey
    \or\dt@gpack
    \or\dt@gsoul
    \or\dt@gmega
    \or\dt@gbers
    \or\dt@gpower
  \fi
  \ifnum\dtGOT=1
    \expandafter\gdef\csname dtHG\the\dtTUI\endcsname{1}%
    \global\advance\dtNPICK by1
  \fi}

%%% ---- choosing a weapon (M23) ---------------------------------------------
%%% Doom's number keys.  Until this the weapon in your hands was whichever
%%% you had picked up last, which is fine for a scripted walk and no use at
%%% all to somebody holding a rocket launcher in a corridor.  A key asks for
%%% a SLOT rather than a weapon, and two slots hold a pair: 1 is the fist
%%% and the chainsaw, 3 the shotgun and the super shotgun.  Pressing a
%%% pair's key again swaps between them, as it does in Doom II.
%%%
%%% A weapon you do not own is refused, and that is the only test: Doom
%%% also lets you raise an empty weapon and then lowers it again.
\def\dt@weapto#1{\ifcsname dtUW#1\endcsname \dtWEAP=#1\relax \fi}
\def\dt@weappair#1#2{%
  \ifnum\dtWEAP=#1\relax
    \dt@weapto{#2}%
  \else
    \ifcsname dtUW#1\endcsname \dtWEAP=#1\relax \else \dt@weapto{#2}\fi
  \fi}
\def\dtSelectWeapon#1{%
  \ifnum\dtHEALTH>0
    \ifcase#1\relax
      \or \dt@weappair{7}{0}%
      \or \dt@weapto{1}%
      \or \dt@weappair{2}{8}%
      \or \dt@weapto{3}%
      \or \dt@weapto{4}%
      \or \dt@weapto{5}%
      \or \dt@weapto{6}%
    \fi
  \fi}

%%% What the status bar and the frame log ask for: the ammunition the weapon
%%% in hand actually uses, or -1 for the fist and the chainsaw, which use
%%% none and show no number at all.
\def\dtAmmoNow{\ifnum\csname dtUT\the\dtWEAP\endcsname<0 -1\else
  \csname dtUA\csname dtUT\the\dtWEAP\endcsname\endcsname\fi}

%%% ================================================================  HUD (M9)
%%% The weapon and the status bar are pictures in the same format as a sprite,
%%% composited in screen space: no projection, no depth test, no silhouette.
%%% So this is the sprite path with the box set by hand and the clip opened to
%%% the whole column -- \dt@sprcol does not care where its numbers came from.
%%%
%%% Doom lays both out against a 320x200 screen whose last 32 rows are the
%%% status bar, so everything scales by \SCRW/320 rather than by depth.
\newcount\dtHDA \newcount\dtHDB \newcount\dtHDC
\newcount\dtHDD \newcount\dtHDN \newcount\dtHDX \newcount\dtHDW
\newcount\dtHDK \newcount\dtHDY \newcount\dtHDP
%% The COLORMAP table a blit is drawn through.  0 -- full brightness -- for
%% everything on the bar; the weapon takes the view's fixed colormap when
%% there is one, because Doom draws it as part of the view and not of the bar.
\newcount\dtBLM \dtBLM=0
%% The picture in your hands is the weapon you are holding, which is a table
%% lookup now rather than a constant.  Every one of the nine is an ordinary
%% sprite lump, so picking a shotgun up changes this with nothing else to do.
\def\dtWeapon{\csname dtUS\the\dtWEAP\endcsname}

%% \dtBlit{key}{x1}{x2}{ytop}{vstep}.  Horizontal scaling falls out of the
%% box, vertical is given, so the two need not agree -- and at the default
%% 200x120 they do not, because that is not 320x200's aspect.
\def\dtBlit#1#2#3#4#5{%
  \ifcsname dtPw#1\endcsname
    \edef\dt@tk{#1}\dt@sprflipfalse
    \dtSRW=\csname dtPw#1\endcsname\relax
    \dtSRH=\csname dtPh#1\endcsname\relax
    \dtTXA=\numexpr#2\relax \dtTXB=\numexpr#3\relax
    \ifnum\dtTXB>\dtTXA\else \dtTXB=\numexpr\dtTXA+1\relax \fi
    \dtSRT=\numexpr#4\relax
    \dtSRS=\numexpr#5\relax
    \dtSRV=\numexpr\dtSRW*65536/(\dtTXB-\dtTXA)\relax
    \dtSRM=\dtBLM
    \dtSPCT=0 \dtSPCB=\SCRH
    \cx=\ifnum\dtTXA<0 0\else\the\dtTXA\fi\relax
    \rb=\ifnum\dtTXB>\SCRW \the\SCRW\else\the\dtTXB\fi\relax
    \dtWhile{\cx<\rb}{\dt@sprcol\advance\cx by1\relax }%
  \fi}

%% The bar first, then the weapon standing on it.
%%
%% A weapon patch carries negative offsets that place it: Doom puts its left
%% edge at -leftoffset and its top row at -topoffset in a 320x200 view, which
%% for the pistol is 126 and 106 -- and 106 plus its 62 rows is exactly 168,
%% flush with the top of the bar.  So the offsets are the placement and there
%% is nothing to centre by hand.
%%
%% A status bar wider than 320 is a widescreen one, drawn with its middle 320
%% across the view and the extra falling off both sides, so it is anchored by
%% that overhang rather than stretched to fit.
%%% ---- numbers and a face on the bar ---------------------------------------
%%% Doom lays the bar out in the 320x200 screen and V_DrawPatch subtracts a
%%% patch's own offsets from the position it is given, so this does too: the
%%% face patches carry (-5,-2) and the digit `1' carries (-1,0), and without
%%% that subtraction the face sits five columns left of where Doom puts it
%%% and every 1 is a pixel out.
%%
%% \dtHDP rather than \dtHDX: \dt@huddig walks a cursor leftward in \dtHDX
%% and calls this for every digit, so sharing the register would leave the
%% cursor holding a scaled screen column and the second digit would land
%% somewhere else entirely.  Cheap to write, silent to get wrong.
\def\dt@hudpat#1#2#3{%
  \ifcsname dtPw#1\endcsname
    \dtHDP=\numexpr((#2)-(\csname dtPl#1\endcsname))*\dtHDA/65536\relax
    \dtBlit{#1}{\the\dtHDP}%
      {\numexpr\dtHDP+(\csname dtPw#1\endcsname)*\dtHDA/65536\relax}%
      {\numexpr((#3)-(\csname dtPt#1\endcsname))*\dtHDB/65536\relax}%
      {\the\dtHDB}%
  \fi}

%%% STlib_drawNum: the x it is given is the RIGHT edge, digits are emitted
%%% least significant first walking leftward, and the pitch is the width of
%%% the digit ZERO for all ten -- `1' is three pixels narrower than the rest
%%% and Doom still advances by 14, which is what keeps a column of numbers
%%% aligned.  Taking each digit's own width instead makes 111 too narrow.
%%%
%%% \divide, not \numexpr: this is modular arithmetic and \numexpr's division
%%% rounds, so 107 would show its last digit as 8 (INTERNALS section 6).
%%% The tall digits and the small ones down the right of the bar differ only
%%% in which font they come from, so the drawer takes a prefix.  \dt@hudnum
%%% keeps its name and its meaning: M14's two layout checks are written about
%%% that macro, and the pitch is still the width of whichever font's own
%%% ZERO -- which is the rule they pin, not a property of STTNUM.
\def\dt@hudnum#1#2#3{\dt@hudnumP{STTNUM}{#1}{#2}{#3}}
\def\dt@hudnumP#1#2#3#4{%
  \ifcsname dtPw#10\endcsname
    \def\dt@hnp{#1}%
    \dtHDN=\numexpr#2\relax
    \ifnum\dtHDN<0 \dtHDN=0 \fi
    \dtHDW=\csname dtPw#10\endcsname\relax
    \dtHDX=\numexpr#3\relax
    \dtHDK=\numexpr#4\relax
    \dt@huddig
  \fi}
%%% \dt@hudpat expands its argument -- it is a \csname key at both ends and
%%% \dtBlit \edef's it again -- so \dt@hnp goes in bare, with no \edef here to
%%% build the call.
\def\dt@huddig{%
  \dtHDD=\dtHDN \divide\dtHDD by10
  \advance\dtHDX by-\dtHDW
  \dt@hudpat{\dt@hnp\number\numexpr\dtHDN-10*\dtHDD\relax}%
    {\the\dtHDX}{\the\dtHDY}%
  \dtHDN=\dtHDD
  \advance\dtHDK by-1
  \ifnum\dtHDN>0 \ifnum\dtHDK>0 \dt@huddig \fi\fi}

%%% ST_calcPainOffset: five faces spread over 0..100 health, and the division
%%% is by 101 rather than 100 so that full health lands on face 0 rather than
%%% off the end.  The three looks within a pain level are Doom's random
%%% glance left, ahead and right, shuffled every ST_FACETIME tics; \dtANIM is
%%% the clock the sprite animation already runs on, which makes this
%%% deterministic for the same reason the per-thing phase is.
%%%
%%% Everything the face says beyond this -- the grin, the ouch, the turn
%%% toward whatever hit you -- is \dtFaceKey's business, further down.
%%
%% Which face is a separate question from where to put it, and it is the half
%% worth asking directly: the thresholds are arithmetic and a wrong one still
%% draws a perfectly good face.  \dtFaceKey names the picture, so a fixture
%% can print it for a set of healths instead of inferring it from pixels.
%%
%% \divide, not \numexpr, and the difference is visible at exactly one health:
%% 80 gives (100-80)*5 = 100, and 100/101 rounds to 1 but truncates to 0.  So
%% a rounding division puts a wounded face on a player who has taken a fifth
%% of one, one health point early.
\def\dt@facepick{%
  \dtHDD=\dtHEALTH
  \ifnum\dtHDD>100 \dtHDD=100 \fi
  \dtHDN=\numexpr(100-\dtHDD)*5\relax
  \divide\dtHDN by101
  \ifnum\dtHDN<0 \dtHDN=0 \fi
  \ifnum\dtHDN>4 \dtHDN=4 \fi
  \dtHDK=\dtANIM \divide\dtHDK by17
  \dtHDD=\dtHDK \divide\dtHDD by3 \multiply\dtHDD by3
  \advance\dtHDK by-\dtHDD}
%%% ST_updateFaceWidget is a priority list, and the order is the whole of it:
%%% a dead player does not grin, and a player who has just been shot does not
%%% glance about.  Doom's own order, less the states nothing here can reach:
%%%
%%%   dead            STFDEAD0
%%%   just hurt       STFOUCH<pain> for a big hit, otherwise STFTL/STFTR
%%%                   toward whatever did it
%%%   new weapon      STFEVL<pain>, the evil grin
%%%   holding fire    STFKILL<pain>, the rampage
%%%   invulnerable    STFGOD0, the gold eyes
%%%   otherwise       STFST<pain><glance>
%%%
%%% Every branch names a picture rather than drawing one, so a fixture can
%%% ask which face without going near a pixel -- the split \dt@facepick was
%%% written for, now that there is something to split.
\def\dtFaceKey{%
  \ifnum\dtHEALTH<1 STFDEAD0\else
    \ifnum\dtHURTC>0 \dt@facehit\else
      \ifnum\dtGRINC>0 STFEVL\the\dtHDN\else
        \ifnum\dtFIREC>0 STFKILL\the\dtHDN\else
          \ifnum\dtINVUL>0 STFGOD0\else
            STFST\the\dtHDN\the\dtHDK\fi\fi\fi\fi\fi}
%%% Doom's threshold is 20 points in one hit, and below it the face turns
%%% toward the attacker instead.  \dtATTL was settled when the damage landed,
%%% because that is when the geometry of it was known.
\def\dt@facehit{%
  \ifnum\dtHURT>20 STFOUCH\the\dtHDN\else
    \ifnum\dtATTL=1 STFTL\the\dtHDN0\else STFTR\the\dtHDN0\fi\fi}
\def\dt@hudface{\dt@facepick \dt@hudpat{\dtFaceKey}{143}{168}}

%%% Doom's own coordinates: ST_Y is 168, the three numbers sit three rows
%%% into the bar, and their right edges are at 44 (ammo), 90 (health) and
%%% 221 (armour), with the percent signs starting where the digits end.
%%% ---- the rest of the bar (M17) -------------------------------------------
%%% The arms panel, the keys and the eight small numbers down the right --
%%% all of them pictures in the format \dt@readpic already reads, at
%%% positions out of Doom's st_stuff.c, drawn by the \dt@hudpat that already
%%% draws the face.  No machinery, which is what the roadmap said it was.
%%%
%%%   arms background STARMS   104,168      40x32
%%%   arms numbers 2..7        111 + 12*(i mod 3), 172 + 10*(i div 3)
%%%   keys STKEYS0..2          239 at y 171, 181, 191
%%%   ammunition held          right edge 288, at y 173, 179, 185, 191
%%%   ammunition maximum       right edge 314, at the same four rows
%%%
%%% None of them touches the four boxes M14's checks measure: STARMS ends at
%%% 143 and the face box starts at 148, the keys sit at 239 against an armour
%%% box ending at 222, and the small numbers are out at 276..314.  That was
%%% arithmetic done against the WAD's own patch sizes, not a hope.
\newcount\dtHDI \newcount\dtHDJ

%%% Weapons 2..7 in a three-by-two grid, from STYSNUM when you own one and
%%% STGNUM when you do not.  dtUW is the subject, not \dtWEAP: the panel says
%%% what you are carrying, not what is in your hands.
%%
%% The prefix is chosen by an \ifcsname *inside* the \csname key \dt@hudpat
%% builds, which works because \ifcsname expands its own contents -- checked
%% in isolation before it went in, because "it should expand" is how most of
%% section 6 starts.
\def\dt@hudarms{%
  \dt@hudpat{STARMS}{104}{168}%
  \dtHDI=0
  \dtWhile{\dtHDI<6}{%
    \dtHDJ=\dtHDI \divide\dtHDJ by3
    \dt@hudpat{\ifcsname dtUW\number\numexpr\dtHDI+2\relax\endcsname
        STYSNUM\else STGNUM\fi\number\numexpr\dtHDI+2\relax}%
      {\number\numexpr111+12*(\dtHDI-3*\dtHDJ)\relax}%
      {\number\numexpr172+10*\dtHDJ\relax}%
    \advance\dtHDI by1\relax }}

%%% Three slots, not six.  Doom draws a card and its skull in the same place
%%% and keeps them apart in the widget; dtUK here cannot tell them apart at
%%% all, because \dtItem gives both spellings of a colour the same key
%%% number.  Recording that is better than pretending to six.
\def\dt@hudkeys{%
  \dtHDI=0
  \dtWhile{\dtHDI<3}{%
    \ifcsname dtUK\the\dtHDI\endcsname
      \dt@hudpat{STKEYS\the\dtHDI}{239}%
        {\number\numexpr171+10*\dtHDI\relax}%
    \fi
    \advance\dtHDI by1\relax }}

%%% Four rows of held-and-maximum.  The maxima are what a backpack doubles,
%%% so this is the only place on the bar that ever says a backpack did
%%% anything at all.
\def\dt@hudammo{%
  \dtHDI=0
  \dtWhile{\dtHDI<4}{%
    \dtHDY=\numexpr173+6*\dtHDI\relax
    \dt@hudnumP{STYSNUM}{\csname dtUA\the\dtHDI\endcsname}{288}{3}%
    \dt@hudnumP{STYSNUM}{\csname dtUM\the\dtHDI\endcsname}{314}{3}%
    \advance\dtHDI by1\relax }}

\def\dt@hudstate{%
  \dtHDY=171
  \dt@hudnum{\dtHEALTH}{90}{3}%
  \dt@hudpat{STTPRCNT}{90}{171}%
  \dt@hudnum{\dtARMOR}{221}{3}%
  \dt@hudpat{STTPRCNT}{221}{171}%
  %% The fist and the chainsaw draw no ammo count at all, which Doom does by
  %% pointing the widget at a sentinel its draw routine refuses.
  \ifnum\csname dtUT\the\dtWEAP\endcsname<0 \else
    \dt@hudnum{\dtAmmoNow}{44}{3}%
  \fi
  \dt@hudface
  \dt@hudarms
  \dt@hudkeys
  \dt@hudammo}

\def\dtDrawHUD{%
  \dtHDA=\numexpr\SCRW*65536/320\relax
  \dtHDB=\numexpr\SCRH*65536/200\relax
  \ifcsname dtPwSTBAR\endcsname
    %% \numexpr takes a leading minus only before a number, never
    %% before a bracket, so the subtraction goes the other way round.
    \dtHDC=\numexpr(320-(\csname dtPwSTBAR\endcsname))/2*\dtHDA/65536\relax
    \dtBlit{STBAR}{\the\dtHDC}%
      {\numexpr\dtHDC+(\csname dtPwSTBAR\endcsname)*\dtHDA/65536\relax}%
      {\numexpr\SCRH-(\csname dtPhSTBAR\endcsname)*\dtHDB/65536\relax}%
      {\the\dtHDB}%
  \fi
  %% The numbers go on the bar, so they are drawn after it.  They cost
  %% nothing when no WAD is loaded: \dt@hudpat guards on the picture.
  \dt@hudstate
  \ifcsname dtPw\dtWeapon\endcsname
    \dtHDC=\numexpr0-(\csname dtPl\dtWeapon\endcsname)*\dtHDA/65536\relax
    \ifnum\dtCMN>0 \ifnum\dtFIXCM<0 \else \dtBLM=\dtFIXCM \fi\fi
    \dtBlit{\dtWeapon}{\the\dtHDC}%
      {\numexpr\dtHDC+(\csname dtPw\dtWeapon\endcsname)*\dtHDA/65536\relax}%
      {\numexpr0-(\csname dtPt\dtWeapon\endcsname)*\dtHDB/65536\relax}%
      {\the\dtHDB}%
    \dtBLM=0
  \fi}

%%% ==============================================================  DOORS (M5)
%%% A door is a sector whose ceiling animates.  Nothing in the renderer or in
%%% collision needs to be told: both already read dtSc every frame, so a
%%% descending ceiling closes the view and the P_LineOpening headroom test at
%%% the same moment, and the depth fix in \dt@portalcol makes it occlude
%%% sprites the instant it touches the floor.
\newcount\dtNDOOR \dtNDOOR=0
\newcount\dtDI \newcount\dtDJ \newcount\dtDK \newcount\dtDL
\newcount\dtDSEC \newcount\dtDMIN \newcount\dtDTGT \newcount\dtDST
\newcount\dtDSPD
\newcount\DOORSPD  \DOORSPD=8     % ceiling units per frame
\newcount\DOORWAIT \DOORWAIT=20   % frames held open before closing again
%% Palette ramps for the two door surfaces.  Only rows that darken
%% monotonically across shades 0..13 may be used -- see the PLAYPAL note in
%% INTERNALS -- so keep these among 16, 32, 48, 64, 80, 96, 112, 128, 176,
%% 192, 208.  A caller can change them from \DTHOOK.
\newcount\DOORRAMP \DOORRAMP=128  % the door slab
\newcount\TRAKRAMP \TRAKRAMP=192  % the track it slides in

\def\dtAddLineSpec#1#2{\expandafter\xdef\csname dtLP#1\endcsname{#2}}
\def\dtLineSpec#1{\ifcsname dtLP#1\endcsname\csname dtLP#1\endcsname\else0\fi}

%%% ---- what a linedef special means ---------------------------------------
%%% Doom's table is a hundred-odd numbers, but each one is really six
%%% independent choices, and the numbering is only a shorthand for the
%%% combination.  Splitting them apart is what lets one state machine serve
%%% every door type and one serve every floor type:
%%%
%%%   kind    what moves: 1 = a ceiling (door), 2 = a floor (lift, mover);
%%%           M16 and M23 carry it on to 3 a ceiling that is not a door,
%%%           4 stairs, 5 a light, 6 the exit, 7 a teleporter, 8 a floor
%%%           that rises
%%%   fun     which motion, within that kind
%%%   act     what fires it: 1 = use, manual; 2 = use, switch; 3 = walk over;
%%%           4 = shoot it
%%%   rep     0 = once for the whole level, 1 = every time
%%%   bytag   0 = the line's own back sector, 1 = every sector with its tag
%%%   spd     1, or 2 for the types Doom calls "fast"
%%%
%%% Stored as one control sequence per declared special, which answers in a
%%% single lookup and costs one string each.  An \ifcase over a range that is
%%% sparse up to 123 would be a hundred branches to miscount, and a chain of
%%% \ifnum would run the whole chain on every miss -- and this is consulted
%%% once per seg by the use scan and once per seg again by the crossing scan.
\newcount\dtTKIND \newcount\dtTFUN \newcount\dtTACT
\newcount\dtTREP  \newcount\dtTBYTAG \newcount\dtTSPD
\def\dtActSet#1#2#3#4#5#6{%
  \dtTKIND=#1 \dtTFUN=#2 \dtTACT=#3 \dtTREP=#4 \dtTBYTAG=#5 \dtTSPD=#6 }
\def\dtActDef#1#2#3#4#5#6#7{%
  \expandafter\gdef\csname dtLA#1\endcsname{\dtActSet{#2}{#3}{#4}{#5}{#6}{#7}}}
%%% Anything not declared below leaves kind 0 and is inert, which is how the
%%% built-in map -- and every special this engine does not implement -- comes
%%% out of here.
\def\dtLineAction#1{%
  \dtTKIND=0 \dtTFUN=0 \dtTACT=0 \dtTREP=0 \dtTBYTAG=0 \dtTSPD=1
  \ifcsname dtLA#1\endcsname\csname dtLA#1\endcsname\fi}

%%% The key a special demands, as a SEVENTH field that is not in \dtActSet.
%%% Six of the sixty-odd rows above want one, so a seventh argument would be
%%% sixty edits to make one table wider -- and every one of them a chance to
%%% shift a column by one silently.  A sparse second table with "absent means
%%% none" is the shape \dtLineTag and \dtLineSpec already use for exactly
%%% this argument.
%%%
%%% 1 blue, 2 yellow, 3 red, i.e. dtUK's numbering plus one so that 0 can
%%% mean "no key".  Doom takes either the card or the skull of a colour, and
%%% \dtItem gives both spellings the same key number, so the conflation M17
%%% recorded as a gap in the status bar is exactly right here.
\def\dtActKey#1#2{\expandafter\gdef\csname dtLK#1\endcsname{#2}}
\def\dtLineKey#1{\ifcsname dtLK#1\endcsname\csname dtLK#1\endcsname\else0\fi}
%% p_doors.c, EV_VerticalDoor: 26 and 32 blue, 27 and 34 yellow, 28 and 33
%% red.  Note 33 is RED and 34 is YELLOW -- the manual doors do not run in
%% colour order and the comments on those two rows above used to say they
%% did, which cost nothing while no key was ever asked for.
\dtActKey{26}{1}\dtActKey{32}{1}
\dtActKey{27}{2}\dtActKey{34}{2}
\dtActKey{28}{3}\dtActKey{33}{3}
%% ...and EV_DoLockedDoor's six switches, which are new here: nothing before
%% M20 could tell them from the unlocked kind, so there was no reason to
%% carry them.
\dtActKey{99}{1}\dtActKey{133}{1}
\dtActKey{134}{3}\dtActKey{135}{3}
\dtActKey{136}{2}\dtActKey{137}{2}

%%% Doors on a ceiling.  fun: 1 open and stay, 2 open, wait and close,
%%% 3 close and stay.
%%                kind fun act rep tag spd
\dtActDef{1}    {1}  {2} {1} {1} {0} {1}   % DR open wait close
\dtActDef{26}   {1}  {2} {1} {1} {0} {1}   % DR blue   -- the key each of
\dtActDef{27}   {1}  {2} {1} {1} {0} {1}   % DR yellow    these wants is in
\dtActDef{28}   {1}  {2} {1} {1} {0} {1}   % DR red       \dtActKey below
\dtActDef{31}   {1}  {1} {1} {0} {0} {1}   % D1 open stay
\dtActDef{32}   {1}  {1} {1} {0} {0} {1}   % D1 blue
\dtActDef{33}   {1}  {1} {1} {0} {0} {1}   % D1 red
\dtActDef{34}   {1}  {1} {1} {0} {0} {1}   % D1 yellow
%% EV_DoLockedDoor's switches: open the tagged door and leave it, fast.
\dtActDef{99}   {1}  {1} {2} {1} {1} {2}   % SR blue fast
\dtActDef{133}  {1}  {1} {2} {0} {1} {2}   % S1 blue fast
\dtActDef{134}  {1}  {1} {2} {1} {1} {2}   % SR red fast
\dtActDef{135}  {1}  {1} {2} {0} {1} {2}   % S1 red fast
\dtActDef{136}  {1}  {1} {2} {1} {1} {2}   % SR yellow fast
\dtActDef{137}  {1}  {1} {2} {0} {1} {2}   % S1 yellow fast
\dtActDef{117}  {1}  {2} {1} {1} {0} {2}   % DR fast
\dtActDef{118}  {1}  {1} {1} {0} {0} {2}   % D1 fast
\dtActDef{2}    {1}  {1} {3} {0} {1} {1}   % W1 open stay
\dtActDef{3}    {1}  {3} {3} {0} {1} {1}   % W1 close
\dtActDef{4}    {1}  {2} {3} {0} {1} {1}   % W1 open wait close
\dtActDef{29}   {1}  {2} {2} {0} {1} {1}   % S1 open wait close
\dtActDef{42}   {1}  {3} {2} {1} {1} {1}   % SR close
\dtActDef{50}   {1}  {3} {2} {0} {1} {1}   % S1 close
\dtActDef{61}   {1}  {1} {2} {1} {1} {1}   % SR open stay
\dtActDef{63}   {1}  {2} {2} {1} {1} {1}   % SR open wait close
\dtActDef{75}   {1}  {3} {3} {1} {1} {1}   % WR close
\dtActDef{86}   {1}  {1} {3} {1} {1} {1}   % WR open stay
\dtActDef{90}   {1}  {2} {3} {1} {1} {1}   % WR open wait close
\dtActDef{103}  {1}  {1} {2} {0} {1} {1}   % S1 open stay
\dtActDef{105}  {1}  {2} {3} {1} {1} {2}   % WR fast open wait close
\dtActDef{106}  {1}  {1} {3} {1} {1} {2}   % WR fast open stay
\dtActDef{107}  {1}  {3} {3} {1} {1} {2}   % WR fast close
\dtActDef{108}  {1}  {2} {3} {0} {1} {2}   % W1 fast open wait close
\dtActDef{109}  {1}  {1} {3} {0} {1} {2}   % W1 fast open stay
\dtActDef{110}  {1}  {3} {3} {0} {1} {2}   % W1 fast close
\dtActDef{111}  {1}  {2} {2} {0} {1} {2}   % S1 fast open wait close
\dtActDef{112}  {1}  {1} {2} {0} {1} {2}   % S1 fast open stay
\dtActDef{113}  {1}  {3} {2} {0} {1} {2}   % S1 fast close
\dtActDef{114}  {1}  {2} {2} {1} {1} {2}   % SR fast open wait close
\dtActDef{115}  {1}  {1} {2} {1} {1} {2}   % SR fast open stay
\dtActDef{116}  {1}  {3} {2} {1} {1} {2}   % SR fast close
%%% Floors.  fun: 1 lift (lower, wait, raise again), 2 lower to the highest
%%% neighbouring floor, 3 lower to the lowest neighbouring floor.
\dtActDef{10}   {2}  {1} {3} {0} {1} {1}   % W1 lift
\dtActDef{21}   {2}  {1} {2} {0} {1} {1}   % S1 lift
\dtActDef{62}   {2}  {1} {2} {1} {1} {1}   % SR lift
\dtActDef{88}   {2}  {1} {3} {1} {1} {1}   % WR lift
\dtActDef{120}  {2}  {1} {3} {1} {1} {2}   % WR fast lift
\dtActDef{121}  {2}  {1} {3} {0} {1} {2}   % W1 fast lift
\dtActDef{122}  {2}  {1} {2} {0} {1} {2}   % S1 fast lift
\dtActDef{123}  {2}  {1} {2} {1} {1} {2}   % SR fast lift
\dtActDef{19}   {2}  {2} {3} {0} {1} {1}   % W1 floor lower to highest
\dtActDef{45}   {2}  {2} {2} {1} {1} {1}   % SR floor lower to highest
\dtActDef{83}   {2}  {2} {3} {1} {1} {1}   % WR floor lower to highest
\dtActDef{102}  {2}  {2} {2} {0} {1} {1}   % S1 floor lower to highest
\dtActDef{23}   {2}  {3} {2} {0} {1} {1}   % S1 floor lower to lowest
\dtActDef{38}   {2}  {3} {3} {0} {1} {1}   % W1 floor lower to lowest
\dtActDef{60}   {2}  {3} {2} {1} {1} {1}   % SR floor lower to lowest
\dtActDef{82}   {2}  {3} {3} {1} {1} {1}   % WR floor lower to lowest
%%% Ceilings.  fun: 1 crush -- down to floor + 8 and back up, for ever --
%%% 2 lower to the floor, 3 raise to the highest neighbouring ceiling,
%%% 4 stop whatever this sector's ceiling is doing.
\dtActDef{6}    {3}  {1} {3} {0} {1} {1}   % W1 crusher
\dtActDef{25}   {3}  {1} {3} {0} {1} {1}   % W1 crusher, slow
\dtActDef{73}   {3}  {1} {3} {1} {1} {1}   % WR crusher, slow
\dtActDef{77}   {3}  {1} {3} {1} {1} {2}   % WR crusher, fast
\dtActDef{141}  {3}  {1} {3} {0} {1} {1}   % W1 crusher, silent
\dtActDef{49}   {3}  {1} {2} {0} {1} {1}   % S1 crusher
\dtActDef{57}   {3}  {4} {3} {0} {1} {1}   % W1 stop a crusher
\dtActDef{74}   {3}  {4} {3} {1} {1} {1}   % WR stop a crusher
\dtActDef{41}   {3}  {2} {2} {0} {1} {1}   % S1 ceiling down to the floor
\dtActDef{43}   {3}  {2} {2} {1} {1} {1}   % SR ceiling down to the floor
\dtActDef{40}   {3}  {3} {3} {0} {1} {1}   % W1 ceiling up to the highest
%%% Stairs.  fun: 1 eight units a step, 2 sixteen and fast.
\dtActDef{8}    {4}  {1} {3} {0} {1} {1}   % W1 stairs
\dtActDef{7}    {4}  {1} {2} {0} {1} {1}   % S1 stairs
\dtActDef{100}  {4}  {2} {3} {0} {1} {2}   % W1 stairs, turbo
\dtActDef{127}  {4}  {2} {2} {0} {1} {2}   % S1 stairs, turbo
%%% Lights.  fun: 1 the brightest neighbour, 2 full, 3 down to 35,
%%% 4 the dimmest neighbour.  Doom's strobes and its glow are not in this
%%% table at all: they are properties of a SECTOR, not things a line does,
%%% and M23 gives them a thinker each.
\dtActDef{12}   {5}  {1} {3} {0} {1} {1}   % W1 light to the brightest
\dtActDef{80}   {5}  {1} {3} {1} {1} {1}   % WR light to the brightest
\dtActDef{13}   {5}  {2} {3} {0} {1} {1}   % W1 light to 255
\dtActDef{81}   {5}  {2} {3} {1} {1} {1}   % WR light to 255
\dtActDef{35}   {5}  {3} {3} {0} {1} {1}   % W1 light down to 35
\dtActDef{79}   {5}  {3} {3} {1} {1} {1}   % WR light down to 35
\dtActDef{104}  {5}  {4} {3} {0} {1} {1}   % W1 light to the dimmest
%%% The exit.  It acts on no sector at all, which is why its tag field is 0
%%% and why \dt@fireact has to take it before the tag scan.
\dtActDef{11}   {6}  {1} {2} {0} {0} {1}   % S1 exit
\dtActDef{52}   {6}  {1} {3} {0} {0} {1}   % W1 exit
\dtActDef{51}   {6}  {2} {2} {0} {0} {1}   % S1 secret exit
\dtActDef{124}  {6}  {2} {3} {0} {0} {1}   % W1 secret exit
%%% ---- M23: the rest of the table ------------------------------------------
%%% A fourth way to fire a line, which is `act' 4: shoot it.  Doom has three
%%% of these and two of them are below with the floors.
\dtActDef{46}   {1}  {1} {4} {1} {1} {1}   % GR open stay
%%% More floors that go DOWN, so they are marked at load like the rest of
%%% kind 2.  fun: 4 lower to eight above the highest neighbouring floor,
%%% 5 travel between the lowest and highest neighbours for ever, 6 stop one
%%% that is doing that.
\dtActDef{36}   {2}  {4} {3} {0} {1} {2}   % W1 turbo lower
\dtActDef{70}   {2}  {4} {2} {1} {1} {2}   % SR turbo lower
\dtActDef{71}   {2}  {4} {2} {0} {1} {2}   % S1 turbo lower
\dtActDef{98}   {2}  {4} {3} {1} {1} {2}   % WR turbo lower
\dtActDef{37}   {2}  {3} {3} {0} {1} {1}   % W1 lower to lowest (and change)
\dtActDef{84}   {2}  {3} {3} {1} {1} {1}   % WR lower to lowest (and change)
\dtActDef{53}   {2}  {5} {3} {0} {1} {1}   % W1 perpetual lift
\dtActDef{87}   {2}  {5} {3} {1} {1} {1}   % WR perpetual lift
\dtActDef{54}   {2}  {6} {3} {0} {1} {1}   % W1 stop a perpetual lift
\dtActDef{89}   {2}  {6} {3} {1} {1} {1}   % WR stop a perpetual lift
%%% Two more light switches, which M16's machine already knew how to obey.
\dtActDef{138}  {5}  {2} {2} {1} {1} {1}   % SR light to 255
\dtActDef{139}  {5}  {3} {2} {1} {1} {1}   % SR light down to 35
%%% Teleporters.  Kind 7 acts on the PLAYER rather than on a sector, and the
%%% tag names where to: a sector with a teleport destination standing in it.
\dtActDef{39}   {7}  {1} {3} {0} {1} {1}   % W1 teleport
\dtActDef{97}   {7}  {1} {3} {1} {1} {1}   % WR teleport
%%% Floors that go UP.  Kind 8, and unlike kind 2 its target is settled when
%%% the line fires -- "the next floor up" is a question about the sector as
%%% it stands now, not as the map was loaded.  fun: 1 the next higher
%%% neighbouring floor, 2 the lowest neighbouring ceiling, 3 that less eight
%%% (Doom's crushing floor), 4 by 24, 5 by 32, 6 by 512, 7 by the shortest
%%% lower texture around it.  Ten more is "and change": the floor takes the
%%% flat of the sector in front of the line that moved it.
\dtActDef{18}   {8}  {1} {2} {0} {1} {1}   % S1 raise to next floor
\dtActDef{69}   {8}  {1} {2} {1} {1} {1}   % SR raise to next floor
\dtActDef{119}  {8}  {1} {3} {0} {1} {1}   % W1 raise to next floor
\dtActDef{128}  {8}  {1} {3} {1} {1} {1}   % WR raise to next floor
\dtActDef{129}  {8}  {1} {3} {1} {1} {2}   % WR raise to next floor, turbo
\dtActDef{130}  {8}  {1} {3} {0} {1} {2}   % W1 raise to next floor, turbo
\dtActDef{131}  {8}  {1} {2} {0} {1} {2}   % S1 raise to next floor, turbo
\dtActDef{132}  {8}  {1} {2} {1} {1} {2}   % SR raise to next floor, turbo
\dtActDef{20}   {8} {11} {2} {0} {1} {1}   % S1 raise to next floor and change
\dtActDef{22}   {8} {11} {3} {0} {1} {1}   % W1 raise to next floor and change
\dtActDef{68}   {8} {11} {2} {1} {1} {1}   % SR raise to next floor and change
\dtActDef{95}   {8} {11} {3} {1} {1} {1}   % WR raise to next floor and change
\dtActDef{47}   {8} {11} {4} {0} {1} {1}   % G1 raise to next floor and change
\dtActDef{5}    {8}  {2} {3} {0} {1} {1}   % W1 raise to lowest ceiling
\dtActDef{91}   {8}  {2} {3} {1} {1} {1}   % WR raise to lowest ceiling
\dtActDef{101}  {8}  {2} {2} {0} {1} {1}   % S1 raise to lowest ceiling
\dtActDef{64}   {8}  {2} {2} {1} {1} {1}   % SR raise to lowest ceiling
\dtActDef{24}   {8}  {2} {4} {0} {1} {1}   % G1 raise to lowest ceiling
\dtActDef{55}   {8}  {3} {2} {0} {1} {1}   % S1 raise to ceiling less 8
\dtActDef{56}   {8}  {3} {3} {0} {1} {1}   % W1 raise to ceiling less 8
\dtActDef{65}   {8}  {3} {2} {1} {1} {1}   % SR raise to ceiling less 8
\dtActDef{94}   {8}  {3} {3} {1} {1} {1}   % WR raise to ceiling less 8
\dtActDef{58}   {8}  {4} {3} {0} {1} {1}   % W1 raise 24
\dtActDef{92}   {8}  {4} {3} {1} {1} {1}   % WR raise 24
\dtActDef{59}   {8} {14} {3} {0} {1} {1}   % W1 raise 24 and change
\dtActDef{93}   {8} {14} {3} {1} {1} {1}   % WR raise 24 and change
\dtActDef{15}   {8} {14} {2} {0} {1} {1}   % S1 raise 24 and change
\dtActDef{66}   {8} {14} {2} {1} {1} {1}   % SR raise 24 and change
\dtActDef{14}   {8} {15} {2} {0} {1} {1}   % S1 raise 32 and change
\dtActDef{67}   {8} {15} {2} {1} {1} {1}   % SR raise 32 and change
\dtActDef{140}  {8}  {6} {2} {0} {1} {1}   % S1 raise 512
\dtActDef{30}   {8}  {7} {3} {0} {1} {1}   % W1 raise by shortest lower
\dtActDef{96}   {8}  {7} {3} {1} {1} {1}   % WR raise by shortest lower

%%% One pass over the linedefs, marking every sector that any special can
%%% move.  A manual door acts on the linedef's BACK sector -- you always face
%%% a door from outside it -- and everything else acts on every sector
%%% carrying the line's tag, which may be nowhere near the line.
%%%
%%% The heights a sector moves between are worked out here rather than when
%%% the trigger fires, following the door precedent: they are a property of
%%% the map's shape, they cost a scan over every linedef each, and firing is
%%% then pure state assignment.  Doom recomputes them at activation, which
%%% differs only if two specials name one sector with different targets --
%%% the first one wins here, and it is recorded in the handoff.
\def\dtFindDoors{%
  \dtDI=0
  \dtWhile{\dtDI<\NLINE}{%
    \dtLineAction{\dtLineSpec{\the\dtDI}}%
    \ifnum\dtTKIND=0\else \dt@markline \fi
    \advance\dtDI by1\relax }%
  %% M23's two lists hang off the same pass, so a fixture that adds a special
  %% from \DTHOOK and calls this again gets them rebuilt along with the rest.
  \dtFindLists
  \ifnum\dtNDOOR>0 \message{[doom.tex] \the\dtNDOOR\space doors}\fi
  \ifnum\dtNPLAT>0 \message{[doom.tex] \the\dtNPLAT\space moving floors}\fi}

\def\dt@markline{%
  \ifnum\dtTBYTAG=0
    \ifnum\dtTKIND=1
      \dtDSEC=\csname dtLB\the\dtDI\endcsname\relax
      \ifnum\dtDSEC<0\else \dt@markdoor \fi
    \fi
  \else
    %% Kinds 4, 5 and 6 mark nothing.  A staircase's targets are a chain and
    %% cannot be known until a trigger names its first step; a light and an
    %% exit have no state to keep between frames at all.
    \def\dt@tagdo{\ifcase\dtTKIND \or
        \dtDSEC=\dtMVSEC \dt@markdoor \or
        \dt@markplat \or
        \dtCESEC=\dtMVSEC \dt@markceil \fi}%
    \dt@fortag{\dtLineTag{\the\dtDI}}%
  \fi}

%%% The shut height is the sector's FLOOR, which is where Doom's T_VerticalDoor
%%% drives a closing ceiling.  For a door that starts shut -- every door on
%%% MAP01, and every fixture here -- the floor and the ceiling are already the
%%% same number, so this is the value that was being recorded anyway.  It only
%%% differs for a sector that is standing open when it is marked, which is
%%% exactly the case a `close' type has to handle: recording the ceiling there
%%% would make closing the door a no-op that looks like a dead trigger.
\def\dt@markdoor{%
  \ifcsname dtDT\the\dtDSEC\endcsname\else
    \expandafter\xdef\csname dtDT\the\dtDSEC\endcsname{1}%
    \expandafter\xdef\csname dtDS\the\dtDSEC\endcsname{0}%
    \expandafter\xdef\csname dtDW\the\dtDSEC\endcsname{0}%
    \expandafter\xdef\csname dtDA\the\dtDSEC\endcsname{1}%
    \expandafter\xdef\csname dtDC\the\dtDSEC\endcsname
      {\csname dtSf\the\dtDSEC\endcsname}%
    \dt@doorhigh
    \expandafter\xdef\csname dtDO\the\dtDSEC\endcsname{\the\dtDMIN}%
    \expandafter\xdef\csname dtDX\the\dtNDOOR\endcsname{\the\dtDSEC}%
    \global\advance\dtNDOOR by1
  \fi}

%%% Doom opens a door to the lowest neighbouring ceiling less four, so the
%%% lintel stays visible.  A door with no neighbours falls back to a fixed
%%% lift of 72.
\def\dt@doorhigh{%
  \dtDMIN=32767 \dtDJ=0
  \dtWhile{\dtDJ<\NLINE}{%
    \dtDK=\csname dtLB\the\dtDJ\endcsname\relax
    \ifnum\dtDK<0\else
      \dtDL=\csname dtLF\the\dtDJ\endcsname\relax
      \ifnum\dtDK=\dtDSEC \dt@doorcand{\the\dtDL}\fi
      \ifnum\dtDL=\dtDSEC \dt@doorcand{\the\dtDK}\fi
    \fi
    \advance\dtDJ by1\relax }%
  \ifnum\dtDMIN>32000
    \dtDMIN=\numexpr\csname dtSc\the\dtDSEC\endcsname+72\relax
  \else
    \advance\dtDMIN by-4
  \fi}
\def\dt@doorcand#1{%
  \ifnum\csname dtSc#1\endcsname<\dtDMIN
    \dtDMIN=\csname dtSc#1\endcsname\relax
  \fi}

%%% ---- one tick ------------------------------------------------------------
\def\dtDoorTick{%
  \dtDI=0
  \dtWhile{\dtDI<\dtNDOOR}{%
    \dtDSEC=\csname dtDX\the\dtDI\endcsname\relax
    \dt@door
    \advance\dtDI by1\relax }}

%%% State 4 is "open and idle": a door raised by an open-and-stay trigger has
%%% no wait to run down and must not drift back to state 0, which the use key
%%% reads as "shut" and would answer by opening it again.
%%% Doom's "fast" door types run at twice the speed of the ordinary ones.
%%% dtDF marks a sector whose door was declared fast; the speed is resolved
%%% into a register here rather than expanded inline, because a macro that
%%% expands to \numexpr inside another \numexpr has to get its own \relax
%%% right and a register simply cannot.
\def\dt@door{%
  \dtDSPD=\DOORSPD
  \ifcsname dtDF\the\dtDSEC\endcsname \multiply\dtDSPD by2 \fi
  \dtDST=\csname dtDS\the\dtDSEC\endcsname\relax
  \ifcase\dtDST \or\dt@dooropen \or\dt@doorwait \or\dt@doorshut \or\fi}

%%% dtDA is read through \dtDoorMode rather than directly, and the difference
%%% matters: the door fixtures build a door by writing dtDT/dtDS/dtDW/dtDC/
%%% dtDO straight from \DTHOOK, so every door written before this flag
%%% existed has no dtDA at all.  \csname would bind that to \relax and the
%%% \ifnum would detonate three macros away from the cause.  Absent means
%%% "open, wait and close", which is what a door did before there was a
%%% choice, so those fixtures keep running unchanged and keep meaning what
%%% they meant when they were written.
\def\dtDoorMode#1{\ifcsname dtDA#1\endcsname\csname dtDA#1\endcsname\else1\fi}
\def\dtPlatRet#1{\ifcsname dtER#1\endcsname\csname dtER#1\endcsname\else1\fi}

\def\dt@dooropen{%
  \dtDTGT=\csname dtDO\the\dtDSEC\endcsname\relax
  \dtDJ=\numexpr\csname dtSc\the\dtDSEC\endcsname+\dtDSPD\relax
  \ifnum\dtDJ<\dtDTGT\else
    \dtDJ=\dtDTGT
    \ifnum\dtDoorMode{\the\dtDSEC}=1
      \expandafter\xdef\csname dtDS\the\dtDSEC\endcsname{2}%
      \expandafter\xdef\csname dtDW\the\dtDSEC\endcsname{\the\DOORWAIT}%
    \else
      \expandafter\xdef\csname dtDS\the\dtDSEC\endcsname{4}%
    \fi
  \fi
  \expandafter\xdef\csname dtSc\the\dtDSEC\endcsname{\the\dtDJ}}

\def\dt@doorwait{%
  \dtDJ=\numexpr\csname dtDW\the\dtDSEC\endcsname-1\relax
  \ifnum\dtDJ>0
    \expandafter\xdef\csname dtDW\the\dtDSEC\endcsname{\the\dtDJ}%
  \else
    \expandafter\xdef\csname dtDS\the\dtDSEC\endcsname{3}%
  \fi}

\def\dt@doorshut{%
  \dtDTGT=\csname dtDC\the\dtDSEC\endcsname\relax
  \dtDJ=\numexpr\csname dtSc\the\dtDSEC\endcsname-\dtDSPD\relax
  \ifnum\dtDJ>\dtDTGT\else
    \dtDJ=\dtDTGT
    \expandafter\xdef\csname dtDS\the\dtDSEC\endcsname{0}%
  \fi
  \expandafter\xdef\csname dtSc\the\dtDSEC\endcsname{\the\dtDJ}}

%%% ---- P_UseLines ----------------------------------------------------------
%%% Nearest door line within 64 units whose midpoint is in front of the
%%% player.  Uses the seg list, like collision, so it needs no new arrays.
%%% What is picked is now the nearest usable LINE rather than the nearest door
%%% sector, because a switch need not have a back sector at all -- two of
%%% MAP01's are one-sided walls -- and acts on sectors named by its tag,
%%% which may be on the other side of the map.
\newcount\dtULN \newcount\dtUOK
\def\dtUse{%
  %% \dtPointSide reads \dtPSX/\dtPSY and the facing test reads \vsin/\vcos,
  %% and \dtUse is reachable from \DTHOOK and from the key loop as well as
  %% from the frame loop.  Setting them here rather than trusting the caller
  %% to have run \dtViewSetup first costs four assignments a press.
  \dtPSX=\numexpr\PX/65536\relax
  \dtPSY=\numexpr\PY/65536\relax
  \vsin=\dtsin{\the\PA}\vcos=\dtcos{\the\PA}%
  \dtDI=0 \dtULN=-1 \dtDMIN=64
  \dtWhile{\dtDI<\NSEG}{\dt@useseg\advance\dtDI by1\relax }%
  \ifnum\dtULN<0\else \dtFireLine{\the\dtULN}\fi}

\def\dt@useseg{%
  \dtDJ=\csname dtGl\the\dtDI\endcsname\relax
  \dtLineAction{\dtLineSpec{\the\dtDJ}}%
  \ifnum\dtTKIND=0\else
    \ifnum\dtTACT<3 \dt@usenear \fi
  \fi}

\def\dt@usenear{%
  \dtCA=\csname dtGa\the\dtDI\endcsname\relax
  \dtCB=\csname dtGb\the\dtDI\endcsname\relax
  \dtAX=\numexpr(\csname dtVx\the\dtCA\endcsname
        +\csname dtVx\the\dtCB\endcsname)/2-\dtPSX\relax
  \dtAY=\numexpr(\csname dtVy\the\dtCA\endcsname
        +\csname dtVy\the\dtCB\endcsname)/2-\dtPSY\relax
  %% in front: the midpoint offset agrees with the facing direction
  \ifnum\numexpr\dtAX*\vcos/65536+\dtAY*\vsin/65536\relax>0
    %% \dtIsqrt assigns into its first argument; it does not expand to a
    %% number.  Written as \dtDK=\dtIsqrt{...} the register assignment ate
    %% \isq@n from the macro body, and the rest of that body -- starting with
    %% a bare "=" -- was typeset, which is why the failure surfaced as
    %% "You can't use \numexpr in horizontal mode".
    \dtIsqrt\dtDK{\dtAX*\dtAX+\dtAY*\dtAY}%
    \ifnum\dtDK<\dtDMIN
      %% A manual door still has to *be* a door -- its back sector must have
      %% been marked -- or a line whose special this build does not implement
      %% would shadow the real door behind it.  A switch carries its target
      %% in its tag and cannot be checked here.
      \dtUOK=0
      \ifnum\dtTACT=1
        \dtDL=\csname dtLB\the\dtDJ\endcsname\relax
        \ifnum\dtDL<0\else
          \ifcsname dtDT\the\dtDL\endcsname \dtUOK=1 \fi
        \fi
      \else
        \dtUOK=1
      \fi
      \ifnum\dtUOK=1 \dtDMIN=\dtDK \dtULN=\dtDJ \fi
    \fi
  \fi}

%%% Using a moving door reverses it, as in Doom.
%%%
%%% \DTUSE holds the use key rather than tapping it, and Doom's usedown latch
%%% means a held key fires P_UseLines once and then nothing.  Without the
%%% equivalent here, a scripted walk opens its door on one frame, reverses it
%%% on the next, and oscillates in front of it forever.  Under \dt@usehold
%%% only a fully shut door responds; the interactive e key goes through the
%%% unlatched path and can still reverse a door in motion.
\newif\ifdt@usehold
\def\dt@usedoor{%
  \dtDST=\csname dtDS\the\dtDSEC\endcsname\relax
  \ifdt@usehold \ifnum\dtDST>0 \dtDST=-1 \fi\fi
  \ifcase\dtDST
    \expandafter\xdef\csname dtDS\the\dtDSEC\endcsname{1}%
  \or
    \expandafter\xdef\csname dtDS\the\dtDSEC\endcsname{3}%
  \or
    \expandafter\xdef\csname dtDS\the\dtDSEC\endcsname{3}%
  \or
    \expandafter\xdef\csname dtDS\the\dtDSEC\endcsname{1}%
  \or
    %% state 4: standing open because something opened it for good.  Using it
    %% by hand closes it again, which is what Doom does with a sector that
    %% has no door thinker left on it.
    \expandafter\xdef\csname dtDS\the\dtDSEC\endcsname{3}%
  \fi}

%%% ==============================================  MOVING FLOORS (M13)
%%% A lift is a door lying down: an animating sector height, a per-frame tick
%%% and an activation scan, with dtSf where the door machine has dtSc.  It
%%% needs nothing from the renderer or from collision, because both already
%%% read dtSf every frame -- the floor a plane is drawn at, the height the
%%% step-up test measures against, and the surface \dtFollowFloor eases the
%%% eye toward are all the same number, so the player rides a lift down
%%% without a line of code saying so.
%%%
%%% States, and they are the door's own with the sense flipped:
%%%   0  idle, up      1  lowering      2  waiting, down
%%%   3  raising       4  idle, down (a mover that does not come back)
%%%   5  waiting, up -- a perpetual lift at the top of its travel
%%%   6  held by a stop line; dtEQ remembers which of the others it was in
\newcount\dtNPLAT \dtNPLAT=0
\newcount\dtMVI \newcount\dtMVJ \newcount\dtMVK \newcount\dtMVL
\newcount\dtMVSEC \newcount\dtMVST \newcount\dtMVTGT \newcount\dtMVSPD
\newcount\dtMVX
\newcount\PLATSPD  \PLATSPD=8    % floor units per frame
\newcount\PLATWAIT \PLATWAIT=20  % frames held down before rising again

%%% Every sector carrying tag #1, with \dtMVSEC set, through \dt@tagdo.  Tag
%%% 0 matches nothing: it is what an untagged linedef and an untagged sector
%%% both report, and letting those meet would attach every manual door in the
%%% map to every plain sector in it.
\newcount\dtTAGN
\def\dt@fortag#1{%
  \ifnum#1=0 \else
    \dtTAGN=0
    \dtWhile{\dtTAGN<\NSEC}{%
      \ifnum\dtSecTag{\the\dtTAGN}=#1\relax
        \dtMVSEC=\dtTAGN \dt@tagdo
      \fi
      \advance\dtTAGN by1\relax }%
  \fi}

%%
%% M23 added three funs.  4 is Doom's turbo lower: the highest neighbouring
%% floor and eight more, so the lip of the pit still shows.  5 is the
%% perpetual lift, the one mover here whose HIGH target is not where it
%% started -- it runs between the lowest and the highest floor around it, so
%% both ends are scanned and dtEP says not to stop at either.  6 is a stop
%% line, which must mark nothing at all: it has no targets of its own, and
%% "the first special to name a sector wins" would otherwise let a stop line
%% that happened to come first in the lump freeze the lift it was meant to
%% stop before it ever moved.
\def\dt@markplat{\ifnum\dtTFUN=6 \else \dt@markplatB \fi}
\def\dt@markplatB{%
  \ifcsname dtET\the\dtMVSEC\endcsname\else
    \ifcase\dtTFUN \or\dt@floorlow \or\dt@floorhigh \or\dt@floorlow
      \or\dt@floorhigh
         \ifnum\dtMVX=\csname dtSf\the\dtMVSEC\endcsname\relax \else
           \advance\dtMVX by8 \fi
      \or\dt@floorlow \fi
    %% A "lower" type whose target sits above the floor it would move must
    %% stay put rather than rise: Doom drives these downward only.
    \ifnum\dtMVX>\csname dtSf\the\dtMVSEC\endcsname
      \dtMVX=\csname dtSf\the\dtMVSEC\endcsname\relax\fi
    \expandafter\xdef\csname dtET\the\dtMVSEC\endcsname{1}%
    \expandafter\xdef\csname dtES\the\dtMVSEC\endcsname{0}%
    \expandafter\xdef\csname dtEW\the\dtMVSEC\endcsname{0}%
    \expandafter\xdef\csname dtEH\the\dtMVSEC\endcsname
      {\csname dtSf\the\dtMVSEC\endcsname}%
    \expandafter\xdef\csname dtEL\the\dtMVSEC\endcsname{\the\dtMVX}%
    \expandafter\xdef\csname dtER\the\dtMVSEC\endcsname
      {\ifnum\dtTFUN=1 1\else\ifnum\dtTFUN=5 1\else0\fi\fi}%
    \ifnum\dtTFUN=5
      \dt@floorhigh
      \ifnum\dtMVX>\csname dtSf\the\dtMVSEC\endcsname\relax
        \expandafter\xdef\csname dtEH\the\dtMVSEC\endcsname{\the\dtMVX}%
      \fi
      \expandafter\xdef\csname dtEP\the\dtMVSEC\endcsname{1}%
    \fi
    \ifnum\dtTSPD=2
      \expandafter\xdef\csname dtEF\the\dtMVSEC\endcsname{1}\fi
    \expandafter\xdef\csname dtEX\the\dtNPLAT\endcsname{\the\dtMVSEC}%
    \global\advance\dtNPLAT by1
  \fi}

%%% P_FindLowestFloorSurrounding / P_FindHighestFloorSurrounding, over the
%%% two-sided linedefs that touch this sector.  \dt@doorhigh walks the
%%% linedefs the same way and for the same reason: a sector does not store
%%% its neighbours, and the linedef list is where the adjacency lives.
%%%
%%% Both fall back to the sector's own floor when it has no two-sided
%%% neighbour at all.  Doom's high scan starts at -500 map units and would
%%% drop such a floor through the world.
\def\dt@floorlow{%
  \dtMVX=\csname dtSf\the\dtMVSEC\endcsname\relax
  \dtMVJ=0
  \dtWhile{\dtMVJ<\NLINE}{%
    \dtMVK=\csname dtLB\the\dtMVJ\endcsname\relax
    \ifnum\dtMVK<0\else
      \dtMVL=\csname dtLF\the\dtMVJ\endcsname\relax
      \ifnum\dtMVK=\dtMVSEC \dt@floorcandlo{\the\dtMVL}\fi
      \ifnum\dtMVL=\dtMVSEC \dt@floorcandlo{\the\dtMVK}\fi
    \fi
    \advance\dtMVJ by1\relax }}
\def\dt@floorcandlo#1{%
  \ifnum\csname dtSf#1\endcsname<\dtMVX \dtMVX=\csname dtSf#1\endcsname\relax\fi}

\def\dt@floorhigh{%
  \dtMVX=-32000 \dtMVJ=0
  \dtWhile{\dtMVJ<\NLINE}{%
    \dtMVK=\csname dtLB\the\dtMVJ\endcsname\relax
    \ifnum\dtMVK<0\else
      \dtMVL=\csname dtLF\the\dtMVJ\endcsname\relax
      \ifnum\dtMVK=\dtMVSEC \dt@floorcandhi{\the\dtMVL}\fi
      \ifnum\dtMVL=\dtMVSEC \dt@floorcandhi{\the\dtMVK}\fi
    \fi
    \advance\dtMVJ by1\relax }%
  \ifnum\dtMVX=-32000
    \dtMVX=\csname dtSf\the\dtMVSEC\endcsname\relax\fi}
\def\dt@floorcandhi#1{%
  \ifnum#1=\dtMVSEC\else
    \ifnum\csname dtSf#1\endcsname>\dtMVX
      \dtMVX=\csname dtSf#1\endcsname\relax\fi
  \fi}

%%% ---- one tick ------------------------------------------------------------
\def\dtPlatTick{%
  \dtMVI=0
  \dtWhile{\dtMVI<\dtNPLAT}{%
    \dtMVSEC=\csname dtEX\the\dtMVI\endcsname\relax
    \dt@plat
    \advance\dtMVI by1\relax }}

\def\dt@plat{%
  \dtMVSPD=\PLATSPD
  \ifcsname dtEF\the\dtMVSEC\endcsname \multiply\dtMVSPD by2 \fi
  \dtMVST=\csname dtES\the\dtMVSEC\endcsname\relax
  \ifcase\dtMVST \or\dt@platdown \or\dt@platwait \or\dt@platup \or
    \or\dt@plattop \fi}

\def\dt@platdown{%
  \dtMVTGT=\csname dtEL\the\dtMVSEC\endcsname\relax
  \dtMVJ=\numexpr\csname dtSf\the\dtMVSEC\endcsname-\dtMVSPD\relax
  \ifnum\dtMVJ>\dtMVTGT\else
    \dtMVJ=\dtMVTGT
    \ifnum\dtPlatRet{\the\dtMVSEC}=1
      \expandafter\xdef\csname dtES\the\dtMVSEC\endcsname{2}%
      \expandafter\xdef\csname dtEW\the\dtMVSEC\endcsname{\the\PLATWAIT}%
    \else
      \expandafter\xdef\csname dtES\the\dtMVSEC\endcsname{4}%
    \fi
  \fi
  \expandafter\xdef\csname dtSf\the\dtMVSEC\endcsname{\the\dtMVJ}%
  \dt@platride}

\def\dt@platwait{%
  \dtMVJ=\numexpr\csname dtEW\the\dtMVSEC\endcsname-1\relax
  \ifnum\dtMVJ>0
    \expandafter\xdef\csname dtEW\the\dtMVSEC\endcsname{\the\dtMVJ}%
  \else
    \expandafter\xdef\csname dtES\the\dtMVSEC\endcsname{3}%
  \fi}

\def\dt@platup{%
  \dtMVTGT=\csname dtEH\the\dtMVSEC\endcsname\relax
  \dtMVJ=\numexpr\csname dtSf\the\dtMVSEC\endcsname+\dtMVSPD\relax
  \ifnum\dtMVJ<\dtMVTGT\else
    \dtMVJ=\dtMVTGT
    %% A perpetual lift waits at the top as it waited at the bottom and then
    %% goes down again.  States 5 and 6 are M23's: 5 is that wait, and 6 is a
    %% lift a stop line is holding, which nothing ticks.
    \ifcsname dtEP\the\dtMVSEC\endcsname
      \expandafter\xdef\csname dtES\the\dtMVSEC\endcsname{5}%
      \expandafter\xdef\csname dtEW\the\dtMVSEC\endcsname{\the\PLATWAIT}%
    \else
      \expandafter\xdef\csname dtES\the\dtMVSEC\endcsname{0}%
    \fi
  \fi
  \expandafter\xdef\csname dtSf\the\dtMVSEC\endcsname{\the\dtMVJ}%
  \dt@platride}

\def\dt@plattop{%
  \dtMVJ=\numexpr\csname dtEW\the\dtMVSEC\endcsname-1\relax
  \ifnum\dtMVJ>0
    \expandafter\xdef\csname dtEW\the\dtMVSEC\endcsname{\the\dtMVJ}%
  \else
    \expandafter\xdef\csname dtES\the\dtMVSEC\endcsname{1}%
  \fi}

%%% Standing on a lift has to move the camera, and nothing else in the engine
%%% will do it: \dtFollowFloor is reached only from \dtTryMove, so a player
%%% who is not walking never settles against the floor.  A descending lift
%%% left the eye hanging at the height it started from, with the floor
%%% visibly dropping away underneath -- no error and a perfectly steady
%%% frame.  Doom calls P_CalcHeight every tic regardless of movement.
%%%
%%% Settling here, in the two branches that actually moved a floor, rather
%%% than once a frame in the main loop, is deliberate: \dtFollowFloor halves
%%% the remaining gap, so calling it a second time in a frame the player also
%%% walked would converge at a different rate and shift every view-height
%%% expectation in the harness.  This way a map with no moving floor cannot
%%% tell the difference, and \dtPlatTick loops zero times on one.
\def\dt@platride{%
  \ifnum\dtMVSEC=\dtSECN \dtFollowFloor \fi}

%%% Doors and floors are ticked together, so a caller that animates the world
%%% has one macro to call.  \dtDoorTick stays as it was: the door fixtures
%%% drive it directly, and a check that names it should keep meaning what it
%%% meant when it was written.
%%
%% \dtSpecTick is M23's: the lights that strobe, the walls that scroll and
%% the floors that burn.  Last, because the floor under the player is the
%% last thing the three movers before it can change.
\def\dtSectorTick{\dtDoorTick\dtPlatTick\dtCeilTick\dtSpecTick}

%%% =================================  CEILINGS, STAIRS, LIGHTS, EXIT (M16)
%%% Four more of Doom's linedef kinds, and three of them needed no new idea:
%%% M13's table was built to be extended, a crushing ceiling is a door with a
%%% different state machine, a staircase is a row of lifts told to go up
%%% instead of down, and a light special is an assignment.
%%%
%%%   kind 3   a ceiling: crush, lower, raise, or stop
%%%   kind 4   a staircase built out from the tagged sector
%%%   kind 5   a light level
%%%   kind 6   the level exit
%%%
%%% Ceiling movers get their own arrays rather than borrowing the door's,
%%% because a door and a crusher can want the same sector and Doom lets them:
%%%
%%%   dtRT<s>  is-mover flag     dtRS<s>  state: 0 idle, 1 down, 2 up, 3 done
%%%   dtRL<s>  low target        dtRH<s>  high target
%%%   dtRC<s>  present if this is a crusher, i.e. it comes back down for ever
%%%   dtRF<s>  fast              dtRX<k>  mover k's sector
\newcount\dtNCEIL \dtNCEIL=0
\newcount\dtCEI \newcount\dtCEJ \newcount\dtCEK \newcount\dtCEL
\newcount\dtCESEC \newcount\dtCEST \newcount\dtCETGT \newcount\dtCESPD
\newcount\dtCEX
\newcount\CEILSPD \CEILSPD=8      % ceiling units per frame
\newcount\CRUSHDMG \CRUSHDMG=10   % what a closing ceiling costs the player
%%% Doom's crushers stop eight units above the floor, which is the gap the
%%% blood shows through.
\newcount\CRUSHGAP \CRUSHGAP=8

%%% ---- marking -------------------------------------------------------------
%%% \dt@markdoor's shape, reached the same way: \dt@markline dispatches on
%%% \dtTKIND.  The targets are settled here rather than at activation for the
%%% reason the door's are, and with the same consequence -- two specials
%%% naming one sector with different targets, and the first one marked wins.
\def\dt@markceil{%
  \ifcsname dtRT\the\dtCESEC\endcsname\else
    \expandafter\xdef\csname dtRT\the\dtCESEC\endcsname{1}%
    \expandafter\xdef\csname dtRS\the\dtCESEC\endcsname{0}%
    \expandafter\xdef\csname dtRH\the\dtCESEC\endcsname
      {\csname dtSc\the\dtCESEC\endcsname}%
    \ifcase\dtTFUN
      \or \expandafter\xdef\csname dtRL\the\dtCESEC\endcsname
            {\number\numexpr\csname dtSf\the\dtCESEC\endcsname
              +\CRUSHGAP\relax}%
          \expandafter\xdef\csname dtRC\the\dtCESEC\endcsname{1}%
      \or \expandafter\xdef\csname dtRL\the\dtCESEC\endcsname
            {\csname dtSf\the\dtCESEC\endcsname}%
      \or \dt@ceilhigh
          \expandafter\xdef\csname dtRH\the\dtCESEC\endcsname{\the\dtCEX}%
          \expandafter\xdef\csname dtRL\the\dtCESEC\endcsname
            {\csname dtSc\the\dtCESEC\endcsname}%
    \fi
    %% A "stop" line marks a sector nobody else did.  Give it a low target
    %% of its own ceiling so the arrays are complete either way: an absent
    %% dtRL is \relax at the point \dt@ceildown reads it.
    \ifcsname dtRL\the\dtCESEC\endcsname\else
      \expandafter\xdef\csname dtRL\the\dtCESEC\endcsname
        {\csname dtSc\the\dtCESEC\endcsname}%
    \fi
    \ifnum\dtTSPD=2
      \expandafter\xdef\csname dtRF\the\dtCESEC\endcsname{1}\fi
    \expandafter\xdef\csname dtRX\the\dtNCEIL\endcsname{\the\dtCESEC}%
    \global\advance\dtNCEIL by1
  \fi}

%%% P_FindHighestCeilingSurrounding, which is \dt@floorhigh's walk with dtSc
%%% where it has dtSf.  A sector with no two-sided neighbour keeps its own.
\def\dt@ceilhigh{%
  \dtCEX=-32000 \dtCEJ=0
  \dtWhile{\dtCEJ<\NLINE}{%
    \dtCEK=\csname dtLB\the\dtCEJ\endcsname\relax
    \ifnum\dtCEK<0\else
      \dtCEL=\csname dtLF\the\dtCEJ\endcsname\relax
      \ifnum\dtCEK=\dtCESEC \dt@ceilcand{\the\dtCEL}\fi
      \ifnum\dtCEL=\dtCESEC \dt@ceilcand{\the\dtCEK}\fi
    \fi
    \advance\dtCEJ by1\relax }%
  \ifnum\dtCEX=-32000
    \dtCEX=\csname dtSc\the\dtCESEC\endcsname\relax\fi}
\def\dt@ceilcand#1{%
  \ifnum#1=\dtCESEC\else
    \ifnum\csname dtSc#1\endcsname>\dtCEX
      \dtCEX=\csname dtSc#1\endcsname\relax\fi
  \fi}

%%% ---- one tick ------------------------------------------------------------
\def\dtCeilTick{%
  \dtCEI=0
  \dtWhile{\dtCEI<\dtNCEIL}{%
    \dtCESEC=\csname dtRX\the\dtCEI\endcsname\relax
    \dt@ceil
    \advance\dtCEI by1\relax }}

\def\dt@ceil{%
  \dtCESPD=\CEILSPD
  \ifcsname dtRF\the\dtCESEC\endcsname \multiply\dtCESPD by2 \fi
  \dtCEST=\csname dtRS\the\dtCESEC\endcsname\relax
  \ifcase\dtCEST \or\dt@ceildown \or\dt@ceilup \or\fi}

\def\dt@ceildown{%
  \dtCETGT=\csname dtRL\the\dtCESEC\endcsname\relax
  \dtCEJ=\numexpr\csname dtSc\the\dtCESEC\endcsname-\dtCESPD\relax
  \ifnum\dtCEJ>\dtCETGT\else
    \dtCEJ=\dtCETGT
    \ifcsname dtRC\the\dtCESEC\endcsname
      \expandafter\xdef\csname dtRS\the\dtCESEC\endcsname{2}%
    \else
      \expandafter\xdef\csname dtRS\the\dtCESEC\endcsname{3}%
    \fi
  \fi
  \expandafter\xdef\csname dtSc\the\dtCESEC\endcsname{\the\dtCEJ}%
  \dt@crushride}

\def\dt@ceilup{%
  \dtCETGT=\csname dtRH\the\dtCESEC\endcsname\relax
  \dtCEJ=\numexpr\csname dtSc\the\dtCESEC\endcsname+\dtCESPD\relax
  \ifnum\dtCEJ<\dtCETGT\else
    \dtCEJ=\dtCETGT
    \ifcsname dtRC\the\dtCESEC\endcsname
      \expandafter\xdef\csname dtRS\the\dtCESEC\endcsname{1}%
    \else
      \expandafter\xdef\csname dtRS\the\dtCESEC\endcsname{3}%
    \fi
  \fi
  \expandafter\xdef\csname dtSc\the\dtCESEC\endcsname{\the\dtCEJ}}

%%% \dt@platride's argument in the ceiling's direction, and the same reason
%%% for living here rather than in the frame loop: only the sector the player
%%% is standing in can crush them, and \dtSECN is what says which that is.
%%% The test is collision's own headroom test, so a ceiling that has not yet
%%% come down to them costs nothing.
\def\dt@crushride{%
  \ifnum\dtCESEC=\dtSECN
    \ifnum\numexpr\csname dtSc\the\dtCESEC\endcsname
          -\csname dtSf\the\dtCESEC\endcsname\relax<\PHEAD
      \dtDamagePlayer{\the\CRUSHDMG}{\the\dtPSX}{\the\dtPSY}%
    \fi
  \fi}

%%% ---- stairs --------------------------------------------------------------
%%% EV_BuildStairs, and the one thing here whose targets CANNOT be settled at
%%% marking time: a staircase is a chain, and which sector is the third step
%%% depends on which one the trigger named as the first.  So this walks the
%%% chain when it fires and hands each step to the moving-floor machine M13
%%% already built -- a stair is a lift told to go up and stay there, which is
%%% dtES=3 with dtEH set to that step's own target.
%%%
%%% The chain rule is Doom's: from the sector just raised, find a two-sided
%%% line whose other side carries the same floor flat, and raise that one
%%% another step.  Same flat is what stops a staircase running away through a
%%% whole level.
\newcount\dtSTPI \newcount\dtSTPJ \newcount\dtSTPK \newcount\dtSTPL
\newcount\dtSTPN \newcount\dtSTPH \newcount\dtSTPS \newcount\dtSTPF
\newcount\dtSTPG
\newcount\STAIRMAX \STAIRMAX=64    % a chain longer than this is a bad map

%%% dtAf is ABSENT on a map with no flats -- the built-in one, and any WAD
%%% render under \DTNOTEX -- and \csname would bind that to \relax and
%%% detonate two macros away.  Read it through this, the way dtLT and dtSt
%%% are read through \dtLineTag and \dtSecTag.  Everything reporting -1 means
%%% every neighbour matches, so a staircase on the built-in map spreads
%%% through all three rooms -- which is the right answer for a map with no
%%% flats to tell its floors apart with.
\def\dtSecFlat#1{\ifcsname dtAf#1\endcsname\csname dtAf#1\endcsname\else-1\fi}

\def\dt@firestair{%
  \ifcsname dtSU\the\dtMVSEC\endcsname\else
    \dtSTPS=\dtMVSEC
    \dtSTPH=\csname dtSf\the\dtSTPS\endcsname\relax
    \dtSTPF=\dtSecFlat{\the\dtSTPS}\relax
    \dtSTPN=8 \ifnum\dtTFUN=2 \dtSTPN=16 \fi
    \dtSTPG=0
    \dtWhile{\dtSTPG<\STAIRMAX}{%
      \advance\dtSTPH by\dtSTPN
      \dt@stairstep
      \dt@stairnext
      \advance\dtSTPG by1\relax }%
    \dtFIRED=1
  \fi}

%%% One step: mark the sector as a mover if it is not one already, point it
%%% at its new height, and start it rising.  dtSU is "already part of a
%%% staircase", which is dtLU for sectors -- without it a chain that loops
%%% back on itself runs until \STAIRMAX and builds a tower.
\def\dt@stairstep{%
  \expandafter\xdef\csname dtSU\the\dtSTPS\endcsname{1}%
  \ifcsname dtET\the\dtSTPS\endcsname\else
    \expandafter\xdef\csname dtET\the\dtSTPS\endcsname{1}%
    \expandafter\xdef\csname dtEW\the\dtSTPS\endcsname{0}%
    \expandafter\xdef\csname dtER\the\dtSTPS\endcsname{0}%
    \expandafter\xdef\csname dtEX\the\dtNPLAT\endcsname{\the\dtSTPS}%
    \global\advance\dtNPLAT by1
  \fi
  \expandafter\xdef\csname dtEL\the\dtSTPS\endcsname
    {\csname dtSf\the\dtSTPS\endcsname}%
  \expandafter\xdef\csname dtEH\the\dtSTPS\endcsname{\the\dtSTPH}%
  \expandafter\xdef\csname dtES\the\dtSTPS\endcsname{3}}

%%% The next step in the chain, or none: a two-sided line touching this
%%% sector whose other side carries the same floor flat and is not already a
%%% stair.  With none, \dtSTPG is run out and the loop ends.
\def\dt@stairnext{%
  \dtSTPK=-1 \dtSTPI=0
  \dtWhile{\dtSTPI<\NLINE}{%
    \dtSTPJ=\csname dtLB\the\dtSTPI\endcsname\relax
    \ifnum\dtSTPJ<0\else
      \dtSTPL=\csname dtLF\the\dtSTPI\endcsname\relax
      \ifnum\dtSTPL=\dtSTPS \dt@staircand{\the\dtSTPJ}\fi
      \ifnum\dtSTPJ=\dtSTPS \dt@staircand{\the\dtSTPL}\fi
    \fi
    \advance\dtSTPI by1\relax }%
  \ifnum\dtSTPK<0
    \dtSTPG=\STAIRMAX
  \else
    \dtSTPS=\dtSTPK
  \fi}
%%% Numbers arrive as \the of a register rather than as a bare \csname, which
%%% is what \dt@floorcandlo and \dt@doorcand already do: a macro's arguments
%%% are not expanded while being grabbed, so an unexpanded lookup would
%%% travel into the \ifcsname key below.
\def\dt@staircand#1{%
  \ifnum\dtSTPK<0
    \ifcsname dtSU#1\endcsname\else
      \ifnum\dtSecFlat{#1}=\dtSTPF \dtSTPK=#1\relax\fi
    \fi
  \fi}

%%% ---- lights --------------------------------------------------------------
%%% No state machine and no tick: Doom's non-strobing light specials are
%%% assignments, and the renderer reads dtSl every frame already.  fun:
%%% 1 the brightest neighbour, 2 full, 3 down to 35, 4 the dimmest.
\newcount\dtLTI \newcount\dtLTJ \newcount\dtLTK \newcount\dtLTV
%%% \dtLTV starts at the sector's own level rather than at whatever the last
%%% call left in it, so an undeclared fun writes the level back unchanged
%%% instead of blacking the room out.  Every assignment ending an \ifcase arm
%%% is terminated: \or is expandable (INTERNALS section 6).
\def\dt@firelight{%
  \dtLTV=\csname dtSl\the\dtMVSEC\endcsname\relax
  \ifcase\dtTFUN
    \or \dt@lightnear{1}%
    \or \dtLTV=255\relax
    \or \dtLTV=35\relax
    \or \dt@lightnear{0}%
  \fi
  \expandafter\xdef\csname dtSl\the\dtMVSEC\endcsname{\the\dtLTV}%
  \dtFIRED=1 }

%%% P_FindMaxSurroundingLight and its opposite, which is \dt@floorhigh's walk
%%% again with dtSl in place of dtSf.
\def\dt@lightnear#1{%
  \dtLTV=\ifnum#1=1 0\else255\fi\relax
  \dtLTI=0
  \dtWhile{\dtLTI<\NLINE}{%
    \dtLTJ=\csname dtLB\the\dtLTI\endcsname\relax
    \ifnum\dtLTJ<0\else
      \dtLTK=\csname dtLF\the\dtLTI\endcsname\relax
      \ifnum\dtLTJ=\dtMVSEC \dt@lightcand{\the\dtLTK}{#1}\fi
      \ifnum\dtLTK=\dtMVSEC \dt@lightcand{\the\dtLTJ}{#1}\fi
    \fi
    \advance\dtLTI by1\relax }}
\def\dt@lightcand#1#2{%
  \ifnum#1=\dtMVSEC\else
    \ifnum#2=1
      \ifnum\csname dtSl#1\endcsname>\dtLTV
        \dtLTV=\csname dtSl#1\endcsname\relax\fi
    \else
      \ifnum\csname dtSl#1\endcsname<\dtLTV
        \dtLTV=\csname dtSl#1\endcsname\relax\fi
    \fi
  \fi}

%%% ---- firing a ceiling ----------------------------------------------------
%%% \dt@fireplat's shape: act only on an idle sector, so firing one twice in a
%%% frame is the same as firing it once and the crossing scan can ignore a
%%% linedef the BSP split.  "Stop" is the exception and has to be, because it
%%% is the one type whose whole job is to act on a sector that is NOT idle.
\def\dt@fireceil{%
  \ifcsname dtRT\the\dtCESEC\endcsname
    \dtCEST=\csname dtRS\the\dtCESEC\endcsname\relax
    \ifnum\dtTFUN=4
      \ifnum\dtCEST>0
        \ifnum\dtCEST<3
          \expandafter\xdef\csname dtRS\the\dtCESEC\endcsname{3}%
          \dtFIRED=1
        \fi
      \fi
    \else
      \ifnum\dtCEST=0
        %% fun 3 is the only one that starts by going UP
        \expandafter\xdef\csname dtRS\the\dtCESEC\endcsname
          {\ifnum\dtTFUN=3 2\else1\fi}%
        \dtFIRED=1
      \fi
    \fi
  \fi}

%%% ---- the exit ------------------------------------------------------------
%%% Doom leaves the level.  There is nowhere to go here, so this records that
%%% it happened and the frame loop stops: a run that ends because the player
%%% reached the exit is a different thing from one that ran out of frames,
%%% and the log has to say which.
\newcount\dtEXIT \dtEXIT=0
\def\dt@fireexit{%
  \ifnum\dtEXIT=0
    \dtEXIT=\ifnum\dtTFUN=2 2\else1\fi\relax
    \message{^^J[doom.tex] exit reached on frame \the\dtFRAME
      \ifnum\dtEXIT=2 \space(the secret one)\fi}%
  \fi
  \dtFIRED=1 }

%%% ==============================  THE REST OF THE SECTOR MACHINERY (M23)
%%% Everything M16 left for later, and it comes in two halves.
%%%
%%% The first half is the SECTOR special -- byte 22 of a sector, which is not
%%% the linedef special that names it by tag.  Nobody fires one: a sector
%%% simply IS a strobing light, or a pool of nukage, or a secret, for as long
%%% as the level lasts.  So these are thinkers, spawned once at load and
%%% ticked once a frame, which is the one shape the table of linedef
%%% specials could not give M16's four lights.
%%%
%%%   1  light flickers off at random       9  secret
%%%   2  strobe, fast                      11  20 damage, and the level ends
%%%   3  strobe, slow                          when it has nearly killed you
%%%   4  strobe, fast, and 20 damage       12  strobe, slow, synchronised
%%%   5  10 damage                         13  strobe, fast, synchronised
%%%   7  5 damage                          16  20 damage
%%%   8  light glows                       17  light flickers like fire
%%%
%%%   dtSs<s>  the special; ABSENT means none (\dtSecSpec)
%%%   dtBT<s>  the light thinker it spawned: 0 none, 1 flash, 2 strobe,
%%%            3 glow, 4 fire.  Present for every sector that HAD a special,
%%%            which is what makes a second \dtSpawnSpecials harmless
%%%   dtBH<s> dtBL<s>   the two levels it moves between
%%%   dtBC<s>  frames left before it next changes
%%%   dtBD<s>  a strobe's dark time; a glow's direction
%%%   dtBX<k>  thinker k's sector
%%%
%%% Doom's two timed doors, 10 and 14, are the ones left out.
%%%
%%% \DTNOSPEC is the off switch, and it takes the whole first half with it:
%%% nothing strobes, no floor hurts and no secret is counted, which is the
%%% map every check before M23 was written on.
\newif\ifdt@spec \dt@spectrue
\newcount\dtNLITE   \dtNLITE=0      % sectors carrying a light thinker
\newcount\dtNSECRET \dtNSECRET=0    % secret sectors on this map
\newcount\dtSECRETS \dtSECRETS=0    % ...and how many have been walked into
\newcount\dtLEVELT  \dtLEVELT=0     % frames since the level began
\newcount\dtZLI \newcount\dtZLS \newcount\dtZLK \newcount\dtZLN
\newcount\dtZLV \newcount\dtZLX \newcount\dtZLC \newcount\dtZLR
\newcount\dtZLQ \newcount\dtZFA \newcount\dtZFB
%%% Doom's timings in tics, as frames on \dtMISRATE's ratio of two to five:
%%% STROBEBRIGHT 5, FASTDARK 15, SLOWDARK 35, and GLOWSPEED's 8 a tic.
%%% Counts, so a fixture can move them.
\newcount\LITEBRIGHT \LITEBRIGHT=2
\newcount\LITEFAST   \LITEFAST=6
\newcount\LITESLOW   \LITESLOW=14
\newcount\LITEGLOW   \LITEGLOW=20
%%% A burning floor hurts on every thirty-second tic of the LEVEL's clock
%%% rather than of the time spent standing in it, so a quick enough dash
%%% across a strip of nukage costs nothing.  Thirteen frames.
\newcount\SLIMETIME  \SLIMETIME=13

\def\dt@litekind#1{\ifcase#1 0\or1\or2\or2\or2\or0\or0\or0\or3\or0\or0\or
  0\or2\or2\or0\or0\or0\or4\else0\fi}
\def\dt@slimedmg#1{\ifcase#1 0\or0\or0\or0\or20\or10\or0\or5\or0\or0\or0\or
  20\or0\or0\or0\or0\or20\else0\fi}

%%% P_Random's stand-in for a sector, on \dt@monrand's argument: the
%%% animation clock, the sector's own index and a salt.  0..255.
\def\dt@specrand#1#2{%
  \dtZLR=\numexpr\dtANIM*73+(#1)*41+(#1)*(#1)*11+(#2)*19\relax
  \dtZLQ=\dtZLR \divide\dtZLQ by256 \multiply\dtZLQ by256
  \advance\dtZLR by-\dtZLQ}

%%% ---- spawning ------------------------------------------------------------
%%% Once, after the map loads, on the \dtFindDoors precedent, and guarded on
%%% dtBT the way \dt@markdoor is guarded on dtDT -- so a fixture can give the
%%% built-in map a special from \DTHOOK and call this again.  The \DTNOSPEC
%%% guard is in here rather than round the engine's own call for the same
%%% fixture's sake: under the off switch its call has to do nothing too.
\def\dtSpawnSpecials{\ifdt@spec \dt@spawnspecials \fi}
\def\dt@spawnspecials{%
  \dtZLI=0
  \dtWhile{\dtZLI<\NSEC}{%
    \ifcsname dtSs\the\dtZLI\endcsname
      \ifcsname dtBT\the\dtZLI\endcsname\else \dt@specone \fi
    \fi
    \advance\dtZLI by1\relax }%
  \ifnum\dtNLITE>0
    \message{[doom.tex] \the\dtNLITE\space light thinkers}\fi
  \ifnum\dtNSECRET>0
    \message{[doom.tex] \the\dtNSECRET\space secret sectors}\fi}
\def\dt@specone{%
  \dtZLK=\dtSecSpec{\the\dtZLI}\relax
  \ifnum\dtZLK=9 \global\advance\dtNSECRET by1 \fi
  \dtZLN=\dt@litekind{\the\dtZLK}\relax
  \expandafter\xdef\csname dtBT\the\dtZLI\endcsname{\the\dtZLN}%
  \ifnum\dtZLN>0 \dt@litespawn \fi}

%%% Every light thinker moves between the sector's own level and the dimmest
%%% of its neighbours -- P_FindMinSurroundingLight, which never answers
%%% brighter than the sector it was asked about.  \dt@lightnear is M16's walk
%%% and wants its sector in \dtMVSEC; that register belongs to the movers, so
%%% it is borrowed and put back.
\def\dt@litespawn{%
  \dtZLV=\csname dtSl\the\dtZLI\endcsname\relax
  \dtZLQ=\dtMVSEC \dtMVSEC=\dtZLI
  \dt@lightnear{0}%
  \dtMVSEC=\dtZLQ
  \dtZLX=\dtLTV \ifnum\dtZLX>\dtZLV \dtZLX=\dtZLV \fi
  \dtZLC=1
  \ifcase\dtZLN
  \or
  \or
    %% A strobe with nothing dimmer beside it goes all the way to black, and
    %% the two synchronised kinds start in step where the others do not.
    \ifnum\dtZLX=\dtZLV \dtZLX=0 \fi
    \expandafter\xdef\csname dtBD\the\dtZLI\endcsname{\the\LITEFAST}%
    \ifnum\dtZLK=3
      \expandafter\xdef\csname dtBD\the\dtZLI\endcsname{\the\LITESLOW}\fi
    \ifnum\dtZLK=12
      \expandafter\xdef\csname dtBD\the\dtZLI\endcsname{\the\LITESLOW}\fi
    \ifnum\dtZLK<12
      \dt@specrand{\dtZLI}{0}%
      \dtZLQ=\dtZLR \divide\dtZLQ by4 \multiply\dtZLQ by4
      \dtZLC=\numexpr\dtZLR-\dtZLQ+1\relax
    \fi
  \or
    \expandafter\xdef\csname dtBD\the\dtZLI\endcsname{-1}%
  \or
    \advance\dtZLX by16
    \ifnum\dtZLX>\dtZLV \dtZLX=\dtZLV \fi
  \fi
  \expandafter\xdef\csname dtBH\the\dtZLI\endcsname{\the\dtZLV}%
  \expandafter\xdef\csname dtBL\the\dtZLI\endcsname{\the\dtZLX}%
  \expandafter\xdef\csname dtBC\the\dtZLI\endcsname{\the\dtZLC}%
  \expandafter\xdef\csname dtBX\the\dtNLITE\endcsname{\the\dtZLI}%
  \global\advance\dtNLITE by1 }

%%% ---- one tick ------------------------------------------------------------
\def\dtSpecTick{%
  \ifdt@spec
    \advance\dtLEVELT by1
    \dtLightTick
    \dtScrollTick
    \dtFloorTick
  \fi}

\def\dtLightTick{%
  \dtZLI=0
  \dtWhile{\dtZLI<\dtNLITE}{%
    \dtZLS=\csname dtBX\the\dtZLI\endcsname\relax
    \dtZLC=\numexpr\csname dtBC\the\dtZLS\endcsname-1\relax
    \ifcase\csname dtBT\the\dtZLS\endcsname\relax
      \or\dt@liteflash \or\dt@litestrobe \or\dt@liteglow \or\dt@litefire \fi
    \expandafter\xdef\csname dtBC\the\dtZLS\endcsname{\the\dtZLC}%
    \advance\dtZLI by1\relax }}
\def\dt@liteset#1{\expandafter\xdef\csname dtSl\the\dtZLS\endcsname{#1}}

%%% T_LightFlash.  Doom's bright time is `(P_Random & 64) + 1', and that mask
%%% is a single bit: the light is on for one tic or for sixty-five, never
%%% anything between, which is what makes it read as a loose connection
%%% rather than as a strobe.  The dark time is one to eight tics.
\def\dt@liteflash{%
  \ifnum\dtZLC<1
    \dt@specrand{\dtZLS}{1}%
    \ifnum\csname dtSl\the\dtZLS\endcsname=\csname dtBH\the\dtZLS\endcsname
      \relax
      \dt@liteset{\csname dtBL\the\dtZLS\endcsname}%
      \dtZLQ=\dtZLR \divide\dtZLQ by3 \multiply\dtZLQ by3
      \dtZLC=\numexpr\dtZLR-\dtZLQ+1\relax
    \else
      \dt@liteset{\csname dtBH\the\dtZLS\endcsname}%
      \dtZLC=\ifnum\dtZLR<128 1\else26\fi\relax
    \fi
  \fi}

%%% T_StrobeFlash: dark for one time, bright for another, for ever.
\def\dt@litestrobe{%
  \ifnum\dtZLC<1
    \ifnum\csname dtSl\the\dtZLS\endcsname=\csname dtBL\the\dtZLS\endcsname
      \relax
      \dt@liteset{\csname dtBH\the\dtZLS\endcsname}%
      \dtZLC=\LITEBRIGHT
    \else
      \dt@liteset{\csname dtBL\the\dtZLS\endcsname}%
      \dtZLC=\csname dtBD\the\dtZLS\endcsname\relax
    \fi
  \fi}

%%% T_Glow.  Doom steps eight a tic and BOUNCES -- it takes the step back
%%% when it overshoots.  At twenty a frame a bounce would pin any glow
%%% narrower than twenty at one level for good, so this clamps to the end it
%%% reached and turns round there instead.
\def\dt@liteglow{%
  \dtZLV=\csname dtSl\the\dtZLS\endcsname\relax
  \ifnum\csname dtBD\the\dtZLS\endcsname<0
    \advance\dtZLV by-\LITEGLOW
    \ifnum\dtZLV>\csname dtBL\the\dtZLS\endcsname\relax \else
      \dtZLV=\csname dtBL\the\dtZLS\endcsname\relax
      \expandafter\xdef\csname dtBD\the\dtZLS\endcsname{1}%
    \fi
  \else
    \advance\dtZLV by\LITEGLOW
    \ifnum\dtZLV<\csname dtBH\the\dtZLS\endcsname\relax \else
      \dtZLV=\csname dtBH\the\dtZLS\endcsname\relax
      \expandafter\xdef\csname dtBD\the\dtZLS\endcsname{-1}%
    \fi
  \fi
  \dt@liteset{\the\dtZLV}%
  \dtZLC=1 }

%%% T_FireFlicker: every fourth tic, the full level less a random 0, 16, 32
%%% or 48 -- but never under the floor \dt@litespawn gave it.
\def\dt@litefire{%
  \ifnum\dtZLC<1
    \dt@specrand{\dtZLS}{2}%
    \dtZLQ=\dtZLR \divide\dtZLQ by4 \multiply\dtZLQ by4
    \dtZLV=\numexpr\csname dtBH\the\dtZLS\endcsname-(\dtZLR-\dtZLQ)*16\relax
    \ifnum\dtZLV<\csname dtBL\the\dtZLS\endcsname\relax
      \dtZLV=\csname dtBL\the\dtZLS\endcsname\relax
    \fi
    \dt@liteset{\the\dtZLV}%
    \dtZLC=2
  \fi}

%%% ---- the floor under the player ------------------------------------------
%%% P_PlayerInSpecialSector.  \dtSECN is the player's own sector here, for
%%% the reason \dt@crushride can trust it: every tick that moves it puts it
%%% back before it returns.
\def\dtFloorTick{%
  \ifnum\dtHEALTH>0
    \ifcsname dtSs\the\dtSECN\endcsname \dt@floorspec \fi
  \fi}
\def\dt@floorspec{%
  \dtZFA=\dtSecSpec{\the\dtSECN}\relax
  \ifnum\dtZFA=9
    %% A secret is found once.  Doom clears the special as it counts it, and
    %% so does this -- to 0 rather than by forgetting it, since dtBT's guard
    %% reads "this sector had one".
    \global\advance\dtSECRETS by1
    \expandafter\xdef\csname dtSs\the\dtSECN\endcsname{0}%
    \message{^^J[doom.tex] secret \the\dtSECRETS\space of \the\dtNSECRET}%
  \else
    \dtZFB=\dt@slimedmg{\the\dtZFA}\relax
    \ifnum\dtZFB>0 \dt@slime \fi
  \fi}

%%% The radiation suit stops the two weak floors outright and the strong ones
%%% almost always -- Doom lets them through on `P_Random < 5', which is the
%%% suit failing about once in fifty.
%%%
%%% Special 11 is the end of Knee-Deep in the Dead: it turns invulnerability
%%% off, cannot quite kill you, and ends the level once you are down to ten.
\def\dt@slime{%
  \dtZLQ=\dtLEVELT \divide\dtZLQ by\SLIMETIME \multiply\dtZLQ by\SLIMETIME
  \ifnum\dtZLQ=\dtLEVELT
    \ifnum\dtZFA=11
      \dtINVUL=0
      \ifnum\dtZFB<\dtHEALTH \else \dtZFB=\numexpr\dtHEALTH-1\relax \fi
    \else
      \ifnum\dtSUIT>0
        \ifnum\dtZFB<20
          \dtZFB=0
        \else
          \dt@specrand{\dtSECN}{3}%
          \ifnum\dtZLR<5 \else \dtZFB=0 \fi
        \fi
      \fi
    \fi
    \ifnum\dtZFB>0
      \dtDamagePlayer{\the\dtZFB}{\the\dtPSX}{\the\dtPSY}%
      \global\advance\dtBURNS by1
    \fi
    \ifnum\dtZFA=11
      \ifnum\dtHEALTH<11 \dtTFUN=1 \dt@fireexit \fi
    \fi
  \fi}
\newcount\dtBURNS \dtBURNS=0        % hits from a floor, for the harness

%%% ---- walls that scroll ---------------------------------------------------
%%% Linedef special 48 is not fired by anything either: the texture on the
%%% line's front slides left one unit a tic for the whole level.  That is the
%%% sidedef's x offset, which \dt@segtex reads for every seg every frame, so
%%% the renderer needs telling nothing.  Two and a half units a frame is two
%%% and three alternately.
%%%
%%%   dtBW<k>  scroller k's sidedef
\newcount\dtNSCROLL \dtNSCROLL=0
\newcount\dtZWI \newcount\dtZWS
\def\dtScrollTick{%
  \dtZWI=0
  \dtWhile{\dtZWI<\dtNSCROLL}{%
    \dtZWS=\csname dtBW\the\dtZWI\endcsname\relax
    \expandafter\xdef\csname dtYx\the\dtZWS\endcsname
      {\number\numexpr\csname dtYx\the\dtZWS\endcsname
        +\ifodd\dtLEVELT 3\else2\fi\relax}%
    \advance\dtZWI by1\relax }}

%%% ---- the two lists a line cannot be found without ------------------------
%%% The scrollers, and the segs of every line that is fired by SHOOTING it.
%%% Both are rebuilt from nothing on every call, so calling this twice is the
%%% same as calling it once -- \dtFindDoors does, and a fixture that adds a
%%% special from \DTHOOK calls \dtFindDoors again.
%%%
%%% The shot scan runs once per pellet, and a super shotgun has twenty of
%%% them: asking every seg on the map what its special means would be twelve
%%% thousand table lookups a blast.  On a map with no such line this list is
%%% empty and a shot costs one \ifnum for it.
%%%
%%%   dtBG<k>  seg k of a line with an `act' of 4
\newcount\dtNGUNL \dtNGUNL=0
\def\dtFindLists{%
  \global\dtNSCROLL=0 \global\dtNGUNL=0
  \dtZWI=0
  \dtWhile{\dtZWI<\NLINE}{%
    \ifnum\dtLineSpec{\the\dtZWI}=48
      \ifcsname dtLR\the\dtZWI\endcsname
        \ifnum\csname dtLR\the\dtZWI\endcsname<0 \else
          \expandafter\xdef\csname dtBW\the\dtNSCROLL\endcsname
            {\csname dtLR\the\dtZWI\endcsname}%
          \global\advance\dtNSCROLL by1
        \fi
      \fi
    \fi
    \advance\dtZWI by1\relax }%
  \dtZWI=0
  \dtWhile{\dtZWI<\NSEG}{%
    \ifcsname dtLP\csname dtGl\the\dtZWI\endcsname\endcsname
      \dtLineAction{\dtLineSpec{\csname dtGl\the\dtZWI\endcsname}}%
      \ifnum\dtTACT=4
        \expandafter\xdef\csname dtBG\the\dtNGUNL\endcsname{\the\dtZWI}%
        \global\advance\dtNGUNL by1
      \fi
    \fi
    \advance\dtZWI by1\relax }}

%%% ---- teleporters ---------------------------------------------------------
%%% EV_Teleport.  The line's tag names a sector, and what it is looking for
%%% in that sector is a THING: type 14, the teleport destination, which is
%%% the only reason that thing is on the map at all.  Where it stands and
%%% which way it faces are where the player arrives and which way they look.
%%%
%%% This is the second half of the section and the first kind here that
%%% moves the player rather than a sector.  It is reached from \dtCrossLines,
%%% i.e. from inside \dtTryMove after the move was committed -- so writing
%%% \PX and \PY here simply replaces where that move ended, and the
%%% \dtLocatePlayer, the touch scan and the floor follow that \dtTryMove was
%%% about to run all happen at the far end without being asked.  \PZ is set
%%% outright rather than eased: you do not climb out of a teleporter.
%%%
%%% Whatever is standing on the pad when you arrive dies of it, which is
%%% Doom's PIT_StompThing.
\newcount\dtTELEPORTS \dtTELEPORTS=0   % for the harness
\newcount\dtTPI \newcount\dtTPS \newcount\dtTPF
\newcount\dtTPA \newcount\dtTPB \newcount\dtTPR
\def\dt@fireteleport{%
  \ifnum\dtFIRED=0 \ifnum\dtXSIDE=0 \ifnum\dtTPDONE=0
    \dtTPI=0
    \dtWhile{\dtTPI<\NTHING}{%
      \ifnum\csname dtHt\the\dtTPI\endcsname=14 \dt@teledest \fi
      \advance\dtTPI by1\relax }%
  \fi\fi\fi}

%% Which sector the destination stands in.  wad.tex records it for every
%% thing it places; one dropped into the built-in map by a fixture has none
%% and is located here, which is safe for the reason given above -- the
%% caller is about to locate the player anyway.
\def\dt@teledest{%
  \ifcsname dtHs\the\dtTPI\endcsname
    \dtTPS=\csname dtHs\the\dtTPI\endcsname\relax
  \else
    \dtPSX=\csname dtHx\the\dtTPI\endcsname\relax
    \dtPSY=\csname dtHy\the\dtTPI\endcsname\relax
    \dtPointInSSec
    \dtTPS=\dtSECN
  \fi
  \ifnum\dtTPS=\dtMVSEC
    \PX=\numexpr\csname dtHx\the\dtTPI\endcsname*65536\relax
    \PY=\numexpr\csname dtHy\the\dtTPI\endcsname*65536\relax
    \PZ=\numexpr(\csname dtSf\the\dtTPS\endcsname+\VIEWH)*65536\relax
    \ifcsname dtHa\the\dtTPI\endcsname
      \dtnorm\PA{\csname dtHa\the\dtTPI\endcsname}%
    \fi
    \dt@telefrag
    \global\advance\dtTELEPORTS by1
    \dtFIRED=1 \dtTPDONE=1
    %% the first destination wins, as it does in Doom
    \dtTPI=\NTHING
  \fi}

%% \dtHurtThing reads \dtMBI.  The monster tick is not running -- a move is
%% made between ticks -- so it is borrowed, and explicitly, for the reason
%% \dt@hitscan gives.
\def\dt@telefrag{%
  \ifdt@mon
    \dtTPA=\numexpr\PX/65536\relax \dtTPB=\numexpr\PY/65536\relax
    \dtTPF=0
    \dtWhile{\dtTPF<\NTHING}{%
      \ifcsname dtOs\the\dtTPF\endcsname
        \ifnum\csname dtOs\the\dtTPF\endcsname<4 \dt@telefragone \fi
      \fi
      \advance\dtTPF by1\relax }%
  \fi}
\def\dt@telefragone{%
  \dtTPR=\numexpr\PRAD
    +\csname dtOr\csname dtHt\the\dtTPF\endcsname\endcsname\relax
  \dtTPS=\numexpr\csname dtHx\the\dtTPF\endcsname-\dtTPA\relax
  \ifnum\dtTPS<0 \dtTPS=-\dtTPS \fi
  \ifnum\dtTPS<\dtTPR
    \dtTPS=\numexpr\csname dtHy\the\dtTPF\endcsname-\dtTPB\relax
    \ifnum\dtTPS<0 \dtTPS=-\dtTPS \fi
    \ifnum\dtTPS<\dtTPR
      \dtHURTBY=-1 \dtMBI=\dtTPF
      \dtHurtThing{10000}%
    \fi
  \fi}

%%% ---- floors that rise ----------------------------------------------------
%%% EV_DoFloor's upward half, and M16's staircase had already built the
%%% machine for it: a stair is a lift told to go up and stay there, which is
%%% dtES=3 with dtEH set to wherever it is going.  So this section is only
%%% the question of WHERE, asked when the line fires rather than at load --
%%% "the next floor up from here" depends on where here has got to.
%%%
%%% One walk over the sector's neighbours serves every answer; \dt@nbrdo is
%%% \let to whichever question is being asked.  \dtRZI is the line it is
%%% standing on, which the texture-height question needs.
\newcount\dtRZI \newcount\dtRZJ \newcount\dtRZK \newcount\dtRZH
\newcount\dtRZX \newcount\dtRZF \newcount\dtRZT \newcount\dtRZU
\def\dt@nbrwalk{%
  \dtRZI=0
  \dtWhile{\dtRZI<\NLINE}{%
    \dtRZJ=\csname dtLB\the\dtRZI\endcsname\relax
    \ifnum\dtRZJ<0 \else
      \dtRZK=\csname dtLF\the\dtRZI\endcsname\relax
      \ifnum\dtRZJ=\dtRZK \else
        \ifnum\dtRZJ=\dtMVSEC \dt@nbrdo{\the\dtRZK}\fi
        \ifnum\dtRZK=\dtMVSEC \dt@nbrdo{\the\dtRZJ}\fi
      \fi
    \fi
    \advance\dtRZI by1\relax }}

%% P_FindNextHighestFloor: the lowest neighbouring floor that is still above
%% this one.  With none, the floor stays where it is.
\def\dt@raisenext{%
  \dtRZH=32767
  \let\dt@nbrdo\dt@nbrnext \dt@nbrwalk
  \ifnum\dtRZH=32767 \dtRZH=\dtRZX \fi}
\def\dt@nbrnext#1{%
  \ifnum\csname dtSf#1\endcsname>\dtRZX
    \ifnum\csname dtSf#1\endcsname<\dtRZH
      \dtRZH=\csname dtSf#1\endcsname\relax
    \fi
  \fi}

%% P_FindLowestCeilingSurrounding, and never past the sector's own ceiling.
%% #1 is the eight units Doom's crushing floor stops short by.
\def\dt@raiseceil#1{%
  \dtRZH=32767
  \let\dt@nbrdo\dt@nbrceil \dt@nbrwalk
  \ifnum\dtRZH>\csname dtSc\the\dtMVSEC\endcsname\relax
    \dtRZH=\csname dtSc\the\dtMVSEC\endcsname\relax
  \fi
  \advance\dtRZH by-#1\relax}
\def\dt@nbrceil#1{%
  \ifnum\csname dtSc#1\endcsname<\dtRZH
    \dtRZH=\csname dtSc#1\endcsname\relax
  \fi}

%% raiseToTexture: by the height of the shortest lower texture on any line
%% around the sector.  With no texture to measure -- the built-in map, or
%% \DTNOTEX -- there is no answer and the floor stays put.
\def\dt@raisetex{%
  \dtRZH=32767
  \let\dt@nbrdo\dt@nbrtex \dt@nbrwalk
  \ifnum\dtRZH=32767 \dtRZH=\dtRZX \else \advance\dtRZH by\dtRZX \fi}
\def\dt@nbrtex#1{\dt@nbrtexside{dtLR}\dt@nbrtexside{dtLL}}
\def\dt@nbrtexside#1{%
  \ifcsname #1\the\dtRZI\endcsname
    \ifnum\csname #1\the\dtRZI\endcsname<0 \else
      \dtRZU=\csname dtYl\csname #1\the\dtRZI\endcsname\endcsname\relax
      \ifnum\dtRZU<0 \else
        \ifnum\csname dtXh\the\dtRZU\endcsname>0
          \ifnum\csname dtXh\the\dtRZU\endcsname<\dtRZH
            \dtRZH=\csname dtXh\the\dtRZU\endcsname\relax
          \fi
        \fi
      \fi
    \fi
  \fi}

%%% Firing one.  Only a floor that is standing still, which is EV_DoFloor
%%% refusing a sector that already has a thinker and is what makes firing
%%% twice in a frame the same as firing once.
\def\dt@fireraise{%
  \dtRZT=0
  \ifcsname dtET\the\dtMVSEC\endcsname
    \dtRZT=\csname dtES\the\dtMVSEC\endcsname\relax
    \ifnum\dtRZT=4 \dtRZT=0 \fi
  \fi
  \ifnum\dtRZT=0
    \dtRZX=\csname dtSf\the\dtMVSEC\endcsname\relax
    \dtRZF=\dtTFUN \ifnum\dtRZF>10 \advance\dtRZF by-10 \fi
    \dtRZH=\dtRZX
    \ifcase\dtRZF
      \or \dt@raisenext
      \or \dt@raiseceil{0}%
      \or \dt@raiseceil{8}%
      \or \dtRZH=\numexpr\dtRZX+24\relax
      \or \dtRZH=\numexpr\dtRZX+32\relax
      \or \dtRZH=\numexpr\dtRZX+512\relax
      \or \dt@raisetex
    \fi
    \ifnum\dtRZH>\dtRZX \dt@raisego \fi
  \fi}

%% \dt@stairstep's body, less the staircase.  The sector becomes a mover if
%% it was not one, and its targets are overwritten if it was -- so a floor
%% that is ALSO somebody's lift comes back from this with the lift's own low
%% target gone, which is the price of settling a lift's targets at load.
\def\dt@raisego{%
  \ifcsname dtET\the\dtMVSEC\endcsname\else
    \expandafter\xdef\csname dtET\the\dtMVSEC\endcsname{1}%
    \expandafter\xdef\csname dtEW\the\dtMVSEC\endcsname{0}%
    \expandafter\xdef\csname dtEX\the\dtNPLAT\endcsname{\the\dtMVSEC}%
    \global\advance\dtNPLAT by1
  \fi
  \expandafter\xdef\csname dtER\the\dtMVSEC\endcsname{0}%
  \expandafter\xdef\csname dtEL\the\dtMVSEC\endcsname{\the\dtRZX}%
  \expandafter\xdef\csname dtEH\the\dtMVSEC\endcsname{\the\dtRZH}%
  \expandafter\xdef\csname dtES\the\dtMVSEC\endcsname{3}%
  \ifnum\dtTSPD=2
    \expandafter\xdef\csname dtEF\the\dtMVSEC\endcsname{1}\fi
  \ifnum\dtTFUN>10 \dt@raisechange \fi
  \dtFIRED=1 }

%% "...and change": the floor takes the look of the sector in FRONT of the
%% line that raised it -- its flat, and its procedural ramp so that the
%% built-in map shows it too.  Doom's raise-to-nearest kind also clears the
%% sector's special, which is how a pool of nukage is paved over.
\def\dt@raisechange{%
  \dtRZU=\csname dtLF\the\dtFLINE\endcsname\relax
  \ifnum\dtRZU<0 \else
    \ifcsname dtAf\the\dtRZU\endcsname
      \expandafter\xdef\csname dtAf\the\dtMVSEC\endcsname
        {\csname dtAf\the\dtRZU\endcsname}%
    \fi
    \expandafter\xdef\csname dtSp\the\dtMVSEC\endcsname
      {\csname dtSp\the\dtRZU\endcsname}%
    \ifnum\dtTFUN=11
      \ifcsname dtSs\the\dtMVSEC\endcsname
        \expandafter\xdef\csname dtSs\the\dtMVSEC\endcsname{0}\fi
    \fi
  \fi}

%%% ---- lines that are shot -------------------------------------------------
%%% P_ShootSpecialLine.  A hitscan that CROSSES one of these fires it,
%%% whether or not the line is a wall -- Doom calls it from the trace before
%%% asking about the opening -- provided nothing stopped the shot first.  So
%%% this wants the nearest listed seg the ray crosses, nearer than whatever
%%% the aim scan found, and then one sight test to just short of it.
%%%
%%% Everything is measured the way \dt@aimone measures a monster: along the
%%% ray and across it, in map units, from \dtPLX/\dtPLY with the ray's own
%%% \vsin/\vcos.  A seg is crossed when its two ends are on opposite sides,
%%% and the crossing is the point where the across measure reaches zero.
\newcount\dtGZI \newcount\dtGZS \newcount\dtGZL \newcount\dtGZD
\newcount\dtGZA \newcount\dtGZB \newcount\dtGZC \newcount\dtGZE
\newcount\dtGZF \newcount\dtGZG \newcount\dtGZT \newcount\dtGZK
\def\dt@shootlines{%
  \dtGZD=\ifnum\dtSHI<0 \dtGUNRANGE\else\dtSHD\fi\relax
  \ifnum\dtGunMelee{\dtWEAP}=1
    \ifnum\dtGZD>\dtMELEE \dtGZD=\dtMELEE \fi
  \fi
  \dtGZL=-1 \dtGZI=0
  \dtWhile{\dtGZI<\dtNGUNL}{%
    \dtGZS=\csname dtBG\the\dtGZI\endcsname\relax
    \ifcsname dtLU\csname dtGl\the\dtGZS\endcsname\endcsname\else
      \dt@shootseg
    \fi
    \advance\dtGZI by1\relax }%
  \ifnum\dtGZL<0 \else
    \dtSight{\the\dtPLX}{\the\dtPLY}%
      {\dtPLX+\vcos*(\dtGZD-4)/65536}{\dtPLY+\vsin*(\dtGZD-4)/65536}%
    \ifnum\dtSEE=1 \dtXSIDE=0 \dtFireLine{\the\dtGZL}\fi
  \fi}
\def\dt@shootseg{%
  \dtGZK=\csname dtGa\the\dtGZS\endcsname\relax
  \dtGZA=\numexpr\csname dtVx\the\dtGZK\endcsname-\dtPLX\relax
  \dtGZB=\numexpr\csname dtVy\the\dtGZK\endcsname-\dtPLY\relax
  \dtGZK=\csname dtGb\the\dtGZS\endcsname\relax
  \dtGZC=\numexpr\csname dtVx\the\dtGZK\endcsname-\dtPLX\relax
  \dtGZE=\numexpr\csname dtVy\the\dtGZK\endcsname-\dtPLY\relax
  %% across the ray, for each end
  \dtGZF=\numexpr(\dtGZB*\vcos-\dtGZA*\vsin)/65536\relax
  \dtGZG=\numexpr(\dtGZE*\vcos-\dtGZC*\vsin)/65536\relax
  \dtGZT=0
  \ifnum\dtGZF<0 \ifnum\dtGZG>0 \dtGZT=1 \fi\fi
  \ifnum\dtGZF>0 \ifnum\dtGZG<0 \dtGZT=1 \fi\fi
  \ifnum\dtGZT=1
    %% along the ray, for each end, and then for the crossing itself
    \dtGZA=\numexpr(\dtGZA*\vcos+\dtGZB*\vsin)/65536\relax
    \dtGZC=\numexpr(\dtGZC*\vcos+\dtGZE*\vsin)/65536\relax
    \dtGZT=\numexpr\dtGZA+(\dtGZC-\dtGZA)*\dtGZF/(\dtGZF-\dtGZG)\relax
    \ifnum\dtGZT>0
      \ifnum\dtGZT<\dtGZD
        \dtGZD=\dtGZT
        \dtGZL=\csname dtGl\the\dtGZS\endcsname\relax
      \fi
    \fi
  \fi}

%%% ==========================================  TRIGGERS: FIRING A LINE (M13)
%%% One entry point for every activation route.  A line that has already been
%%% spent -- Doom's once-only W1 and S1 types -- is recorded in dtLU and
%%% never fires again, which is also what stops a repeatable *walkover* line
%%% firing twice in one move when the BSP happened to split it into two segs.
%%% For the repeatable ones that is handled instead by making a trigger act
%%% only on a sector that is standing idle, so firing it twice in a frame is
%%% the same as firing it once.  The manual use key keeps its own behaviour:
%%% pressing a door in motion reverses it, and that is what a player expects.
%%% The key test goes HERE, at the one entry point every route goes through,
%%% and before anything else happens: a refused door must not spend a
%%% once-only line and must not turn a switch face, and both of those are
%%% inside \dt@fireact.  Doom refuses in the same place for the same reason.
%%%
%%% \dtLOCKED is a register rather than a conditional because check.py reads
%%% it: "the door did not open" and "the door was refused for want of a key"
%%% are different statements and a fixture that cannot tell them apart would
%%% pass on a door that was simply out of reach.
\newcount\dtFLINE \newcount\dtLOCKED \newcount\dtNEEDKEY
\newcount\dtLOCKN \dtLOCKN=0        % how many refusals, for the harness
\newif\ifdt@lock \dt@locktrue       % \DTNOLOCK opens everything regardless
\def\dtFireLine#1{%
  \dtFLINE=#1\relax
  \dtLineAction{\dtLineSpec{\the\dtFLINE}}%
  \ifnum\dtTKIND=0\else
    \ifcsname dtLU\the\dtFLINE\endcsname\else
      \dt@keytest
      \ifnum\dtLOCKED=0 \dt@fireact \fi
    \fi
  \fi}
\def\dt@keytest{%
  \dtLOCKED=0
  \ifdt@lock
    \dtNEEDKEY=\dtLineKey{\dtLineSpec{\the\dtFLINE}}\relax
    \ifnum\dtNEEDKEY>0
      \ifcsname dtUK\number\numexpr\dtNEEDKEY-1\relax\endcsname\else
        \dtLOCKED=1 \global\advance\dtLOCKN by1
      \fi
    \fi
  \fi}

%%% Doom's EV_ functions return whether they found anything to do, and the
%%% switch only changes face and the once-only line is only spent if they
%%% did.  \dtFIRED is that return value: a tag naming no sector, or naming
%%% one that is already in motion, leaves the trigger armed rather than
%%% quietly consuming it.  It is also what keeps a held use key from
%%% flickering a repeatable switch on and off once a frame.
\newcount\dtFIRED
\def\dt@fireact{%
  \dtFIRED=0
  %% The exit acts on no sector, so it is taken before the tag scan rather
  %% than inside it -- its tag field is 0 and \dt@fortag would match nothing.
  \ifnum\dtTKIND=6
    \dt@fireexit
  \else
  \ifnum\dtTBYTAG=0
    \dtDSEC=\csname dtLB\the\dtFLINE\endcsname\relax
    \ifnum\dtDSEC<0\else \dt@firedoor \fi
  \else
    \def\dt@tagdo{\ifcase\dtTKIND \or
        \dtDSEC=\dtMVSEC \dt@firedoor \or
        \dt@fireplat \or
        \dtCESEC=\dtMVSEC \dt@fireceil \or
        \dt@firestair \or
        \dt@firelight \or \or
        \dt@fireteleport \or
        \dt@fireraise \fi}%
    \dt@fortag{\dtLineTag{\the\dtFLINE}}%
  \fi
  \fi
  \ifnum\dtFIRED=1
    \ifnum\dtTREP=0
      \expandafter\xdef\csname dtLU\the\dtFLINE\endcsname{1}%
    \fi
    \ifnum\dtTACT=2 \dt@swline{\the\dtFLINE}\fi
    %% a switch that is shot changes face exactly as one that is pressed
    \ifnum\dtTACT=4 \dt@swline{\the\dtFLINE}\fi
  \fi}

%%% The mode a door runs in belongs to the line that fired it, not to the
%%% sector: the same door can be opened for good by one trigger and opened
%%% briefly by another, so dtDA and dtDF are written here rather than at
%%% marking time.
\def\dt@firedoor{%
  \ifcsname dtDT\the\dtDSEC\endcsname
    \expandafter\xdef\csname dtDA\the\dtDSEC\endcsname
      {\ifnum\dtTFUN=2 1\else0\fi}%
    \ifnum\dtTSPD=2
      \expandafter\xdef\csname dtDF\the\dtDSEC\endcsname{1}\fi
    \ifnum\dtTACT=1 \dtFIRED=1 \dt@usedoor \else \dt@doortrig \fi
  \fi}

%%% A triggered door acts only on an idle sector, as Doom's EV_DoDoor does by
%%% refusing a sector that already carries a thinker.  That is what makes
%%% firing one twice in a frame the same as firing it once, which in turn is
%%% what lets the crossing scan ignore a linedef the BSP split into two segs.
\def\dt@doortrig{%
  \dtDST=\csname dtDS\the\dtDSEC\endcsname\relax
  \ifnum\dtTFUN=3
    \ifnum\dtDST=4 \dt@doorset{3}\fi
    \ifnum\dtDST=2 \dt@doorset{3}\fi
  \else
    \ifnum\dtDST=0 \dt@doorset{1}\fi
  \fi}
\def\dt@doorset#1{%
  \expandafter\xdef\csname dtDS\the\dtDSEC\endcsname{#1}\dtFIRED=1 }

\def\dt@fireplat{%
  \ifcsname dtET\the\dtMVSEC\endcsname
    \dtMVST=\csname dtES\the\dtMVSEC\endcsname\relax
    \ifnum\dtTFUN=6
      %% EV_StopPlat: hold a lift that is doing something, and remember what,
      %% so that firing its own line again takes up where it left off.
      \ifnum\dtMVST>0 \ifnum\dtMVST<6 \ifnum\dtMVST=4 \else
        \expandafter\xdef\csname dtEQ\the\dtMVSEC\endcsname{\the\dtMVST}%
        \dt@platset{6}%
      \fi\fi\fi
    \else
      \ifnum\dtMVST=0 \dt@platset{1}\fi
      %% a lift called back up from the bottom
      \ifnum\dtMVST=4
        \ifnum\dtPlatRet{\the\dtMVSEC}=1 \dt@platset{3}\fi
      \fi
      \ifnum\dtMVST=6
        \dt@platset{\csname dtEQ\the\dtMVSEC\endcsname}%
      \fi
    \fi
  \fi}
\def\dt@platset#1{%
  \expandafter\xdef\csname dtES\the\dtMVSEC\endcsname{#1}\dtFIRED=1 }

%%% ---- switch faces --------------------------------------------------------
%%% Pressing a switch turns SW1SOMETHING into SW2SOMETHING.  The pairing is
%%% not in the WAD in Doom's own data -- the engine holds the list -- but the
%%% naming convention is, so the counterpart is found by name at load and
%%% recorded in dtXs, both ways round, by wad.tex.  Here it is only a matter
%%% of writing the other index into the sidedef the seg already reads, so
%%% the next frame draws the pressed face with nothing else told anything.
%%%
%%% Doom tries top, then middle, then bottom, and stops at the first surface
%%% that is a switch; a line whose bottom is SW1COMP and whose top is an
%%% ordinary wall must not have its top swapped for something arbitrary.
\newcount\dtSWS \newcount\dtSWX \newcount\dtSWY \newcount\dtSWT
\def\dt@swline#1{%
  \ifcsname dtLR#1\endcsname
    \dtSWS=\csname dtLR#1\endcsname\relax
    \ifnum\dtSWS<0\else
      \dtSWT=0
      \dt@swtry{dtYu}%
      \ifnum\dtSWT=0 \dt@swtry{dtYm}\fi
      \ifnum\dtSWT=0 \dt@swtry{dtYl}\fi
    \fi
  \fi}
\def\dt@swtry#1{%
  \ifcsname #1\the\dtSWS\endcsname
    \dtSWX=\csname #1\the\dtSWS\endcsname\relax
    \ifnum\dtSWX<0\else
      \ifcsname dtXs\the\dtSWX\endcsname
        \dtSWY=\csname dtXs\the\dtSWX\endcsname\relax
        %% only if the counterpart actually composited; an unresolved name
        %% falls back to the procedural bricks, and swapping a switch for
        %% those would read as the wall being replaced rather than pressed
        \ifnum\csname dtXw\the\dtSWY\endcsname>0
          \expandafter\xdef\csname #1\the\dtSWS\endcsname{\the\dtSWY}%
          \dtSWT=1
        \fi
      \fi
    \fi
  \fi}

%%% ---- walking over a line -------------------------------------------------
%%% Doom collects the lines a move crossed in P_TryMove and fires them after
%%% the move succeeds.  The crossing test is the one collision already uses,
%%% \dt@segcross, which is exactly the question a walkover asks; running it
%%% here rather than inside \dtBlocked keeps it off the rejected candidate
%%% moves, of which there are two for every blocked step.
\newcount\dtXGI \newcount\dtXGL \newcount\dtXGF
\newcount\dtXSIDE \dtXSIDE=0
%% Set by a teleport and cleared with the scan.  The BSP splits a linedef
%% into segs and the scan below meets every one of them, measured from where
%% the move began -- so a repeatable teleporter cut in two would otherwise
%% fire twice in one stride.  The doors' answer, acting only on a sector that
%% is standing idle, has nothing to hold on to here.
\newcount\dtTPDONE \dtTPDONE=0
\def\dtCrossLines{%
  \dtCX=\numexpr\PX/65536\relax \dtCY=\numexpr\PY/65536\relax
  \dtOX=\numexpr\dtPX-\dtCX\relax \dtOY=\numexpr\dtPY-\dtCY\relax
  \dtXGF=0
  \ifnum\dtOX=0\else\dtXGF=1 \fi
  \ifnum\dtOY=0\else\dtXGF=1 \fi
  \dtTPDONE=0
  \ifnum\dtXGF=1 \dt@crossscan \fi
  %% every other way of firing a line counts as its front
  \dtXSIDE=0 \dtTPDONE=0 }

\def\dt@crossscan{%
  \dtBXL=\ifnum\dtOX<0 \dtOX\else0\fi\relax \advance\dtBXL by-\PRAD
  \dtBXH=\ifnum\dtOX>0 \dtOX\else0\fi\relax \advance\dtBXH by\PRAD
  \dtBYL=\ifnum\dtOY<0 \dtOY\else0\fi\relax \advance\dtBYL by-\PRAD
  \dtBYH=\ifnum\dtOY>0 \dtOY\else0\fi\relax \advance\dtBYH by\PRAD
  \dtXGI=0
  \dtWhile{\dtXGI<\NSEG}{\dt@crossseg\advance\dtXGI by1\relax }}

%%% The \ifcsname on dtLP misses for all but a handful of lines on a real map
%%% -- 12 of MAP01's 370 -- so the scan costs one failed lookup per seg and
%%% the crossing arithmetic runs almost nowhere.
\def\dt@crossseg{%
  \dtXGL=\csname dtGl\the\dtXGI\endcsname\relax
  \ifcsname dtLP\the\dtXGL\endcsname
    \ifcsname dtLU\the\dtXGL\endcsname\else
      \dt@crosstest
    \fi
  \fi}

\def\dt@crosstest{%
  \dtCA=\csname dtGa\the\dtXGI\endcsname\relax
  \dtCB=\csname dtGb\the\dtXGI\endcsname\relax
  \dtAX=\numexpr\csname dtVx\the\dtCA\endcsname-\dtCX\relax
  \dtAY=\numexpr\csname dtVy\the\dtCA\endcsname-\dtCY\relax
  \dtBX=\numexpr\csname dtVx\the\dtCB\endcsname-\dtCX\relax
  \dtBY=\numexpr\csname dtVy\the\dtCB\endcsname-\dtCY\relax
  \dtRJ=0
  \ifnum\dtAX>\dtBXH \ifnum\dtBX>\dtBXH \dtRJ=1 \fi\fi
  \ifnum\dtAX<\dtBXL \ifnum\dtBX<\dtBXL \dtRJ=1 \fi\fi
  \ifnum\dtAY>\dtBYH \ifnum\dtBY>\dtBYH \dtRJ=1 \fi\fi
  \ifnum\dtAY<\dtBYL \ifnum\dtBY<\dtBYL \dtRJ=1 \fi\fi
  \ifnum\dtRJ=0
    \dtHIT=0
    \dt@segcross
    \ifnum\dtHIT=1
      %% Which side the move STARTED on, which a teleporter has to know: it
      %% takes you only when you walk in through its front, or stepping off
      %% the pad at the far end would send you straight back.  \dt@segcross
      %% left it behind in \dtDA -- the move's origin against the seg -- and a
      %% linedef's front is on its right, so a negative cross product is the
      %% front of a seg that runs with its line and the back of one that
      %% runs against it.
      \dtXSIDE=\ifnum\dtDA<0 0\else1\fi\relax
      \ifnum\csname dtGs\the\dtXGI\endcsname=1
        \dtXSIDE=\numexpr1-\dtXSIDE\relax
      \fi
      \dt@crossfire
    \fi
  \fi}

\def\dt@crossfire{%
  \dtLineAction{\dtLineSpec{\the\dtXGL}}%
  \ifnum\dtTACT=3 \dtFireLine{\the\dtXGL}\fi}

%%% ============================================================  ENEMIES (M15)
%%% The first thing here that acts on its own.  Everything before this moves
%%% because the player moved: a door because it was pressed, a lift because a
%%% switch was thrown, a torch because the frame counter advanced.  A monster
%%% decides.
%%%
%%% Doom's mobj is a state machine and the states are the whole of it: an
%%% idle loop until it sees you, a walk cycle while it closes, an attack, a
%%% flinch when hit, and a death animation ending in the corpse it leaves
%%% behind.  Each state names a run of sprite frames, which is the mechanism
%%% \dtSpriteAnim already drives -- the difference is that the sequence is
%%% chosen per THING rather than per type, because two zombiemen in a room
%%% are not doing the same thing.
%%%
%%%   dtOs<i>   state: 0 idle, 1 chasing, 2 attacking, 3 pain, 4 dying,
%%%             5 dead -- a corpse, and never ticked again
%%%   dtOh<i>   health remaining
%%%   dtOc<i>   tics left on the frame showing now
%%%   dtOf<i>   which letter of the state's sequence that is
%%%
%%% and per TYPE, out of Doom's mobjinfo:
%%%
%%%   dtOq<t>  spawn health   dtOr<t>  radius       dtOp<t>  speed
%%%   dtOw<t>  painchance     dtOe<t>  melee dice   dtOm<t>  melee multiplier
%%%   dtOg<t>  hitscan shots per attack; a zero multiplier means no melee
%%%   dtOa<t>@<s>  dtOn<t>@<s>  dtOd<t>@<s>  state s's letters, count, tics
%%%
%%% \DTNOMON leaves every monster asleep in its spawn loop and solid to
%%% nothing, which is what the world did before this section existed.  That
%%% is not decoration: a dozen checks in check.py are written about a map
%%% that moves only when the player moves, and they have to go on meaning
%%% that.  \ifdt@mon and the block scan it guards are up with the collision
%%% code, which is the first thing to ask.
\newcount\dtNMON  \dtNMON=0        % monsters this map spawned
\newcount\dtKILLS \dtKILLS=0
\newcount\dtMBI \newcount\dtMBS \newcount\dtMBC \newcount\dtMBF
\newcount\dtMBH \newcount\dtMBN \newcount\dtMBW \newcount\dtMBX
\newcount\dtMBY \newcount\dtMBA \newcount\dtMBB \newcount\dtMBD
\newcount\dtMBR \newcount\dtMBP \newcount\dtMBQ \newcount\dtMBZ
\newcount\dtMBV \newcount\dtMBK \newcount\dtMBE \newcount\dtMBL
\newcount\dtMBG \newcount\dtMBJ
%%% The player's position in map units, taken once at the top of the tick.
%%% \dtPSX/\dtPSY cannot be used for it: a monster that moves has to be
%%% located in the map, and the only point locator here writes exactly those
%%% two registers -- so the second monster in the list would otherwise walk
%%% toward wherever the first one ended up.
\newcount\dtPLX \newcount\dtPLY

%%% Doom's MELEERANGE less the 20 P_CheckMeleeRange takes off, plus the
%%% player's own radius.  Counts, so a fixture can move them.
\newcount\dtMELEE \dtMELEE=60
%%% MISSILERANGE.  Doom's P_CheckMissileRange is a chain of special cases on
%%% top of the sight test; this is the part that matters -- do not shoot at
%%% something you cannot reach.
\newcount\dtMRANGE \dtMRANGE=2048

\def\dtMonster#1#2#3#4#5#6#7#8{%
  \expandafter\gdef\csname dtOq#1\endcsname{#2}%
  \expandafter\gdef\csname dtOr#1\endcsname{#3}%
  \expandafter\gdef\csname dtOp#1\endcsname{#4}%
  \expandafter\gdef\csname dtOw#1\endcsname{#5}%
  \expandafter\gdef\csname dtOe#1\endcsname{#6}%
  \expandafter\gdef\csname dtOm#1\endcsname{#7}%
  \expandafter\gdef\csname dtOg#1\endcsname{#8}}

%%% \dtMonAnim{type}{state}{letters}{tics}.  \dt@ancount is \dtSpriteAnim's
%%% own counter and both of these run at load, so there is no second walker.
%%
%% State 5 is declared here as a copy of state 4, because a corpse IS the
%% death animation stopped on its last frame and has no letters of its own.
%% Leaving it undeclared is not a missing picture, it is an expansion
%% accident: \csname binds the absent sequence to \relax, the frame walker
%% is handed a control sequence where it wanted a letter, and the run ends
%% with a hundred errors every one of which names something several
%% expansions away from the cause.
\def\dtMonAnim#1#2#3#4{%
  \dtMFC=0 \dt@ancount#3\relax
  \expandafter\gdef\csname dtOa#1@#2\endcsname{#3}%
  \expandafter\xdef\csname dtOn#1@#2\endcsname{\the\dtMFC}%
  \expandafter\gdef\csname dtOd#1@#2\endcsname{#4}%
  \ifnum#2=4
    \expandafter\gdef\csname dtOa#1@5\endcsname{#3}%
    \expandafter\xdef\csname dtOn#1@5\endcsname{\the\dtMFC}%
    \expandafter\gdef\csname dtOd#1@5\endcsname{#4}%
  \fi}

%%% ---- the table -----------------------------------------------------------
%%% Doom's own numbers, from info.c.
%%%
%%% Speed is applied once per rendered frame, and that needs saying because
%%% it is not quite what Doom does.  Doom moves a monster `speed' units per
%%% call of A_Chase, which sits on a walk frame lasting three or four tics --
%%% about two units a tic against a running player's sixteen.  This engine's
%%% frame is a whole player step of forty units, so moving a zombieman eight
%%% units in it leaves it at a fifth of the player's pace: nearer Doom's real
%%% ratio than moving it once every fourth frame would be, and it keeps the
%%% tick to one rule.  \dtSpriteAnim reads Doom's tic counts as frames for
%%% the same reason.
%%%
%%% A monster whose real attack is a PROJECTILE -- the imp's fireball, the
%%% cacodemon's, the baron's -- has a zero in the `shots' column and gets its
%%% missile from \dtMisFrom, in the projectile section.  The melee dice in
%%% the same row are the other arm of the same attack: A_TroopAttack uses its
%%% fists at arm's length and throws past it, never both.  Under \DTNOMIS
%%% nothing is thrown and the fists are all such a monster has.
%%              type  hp  rad spd pain dice mult shots
\dtMonster{3004}  {20} {20} {8} {200} {0}  {0}  {1}   % zombieman
\dtMonster{9}     {30} {20} {8} {170} {0}  {0}  {3}   % shotgun guy
\dtMonster{65}    {70} {20} {8} {170} {0}  {0}  {2}   % chaingunner
\dtMonster{84}    {50} {20} {8} {170} {0}  {0}  {3}   % Wolfenstein SS
\dtMonster{3001}  {60} {20} {8} {200} {8}  {3}  {0}   % imp
\dtMonster{3002} {150} {30}{10} {180} {10} {4}  {0}   % demon
\dtMonster{58}   {150} {30}{10} {180} {10} {4}  {0}   % spectre
\dtMonster{3005} {400} {31} {8} {128} {6}  {10} {0}   % cacodemon
\dtMonster{3003}{1000} {24} {8}  {50} {8}  {10} {0}   % baron of hell
\dtMonster{69}   {500} {24} {8}  {50} {8}  {10} {0}   % hell knight
\dtMonster{66}   {300} {20}{10} {100} {6}  {10} {0}   % revenant
\dtMonster{3006} {100} {16} {8} {256} {8}  {3}  {0}   % lost soul
%% M23's four, which are rows and nothing else: a thing that walks, takes
%% damage, and either fires a hitscan or throws a missile the table below
%% already knows how to fly.  Doom's two that are NOT rows -- the arch-vile,
%% which raises the dead, and the pain elemental, which spawns lost souls --
%% are still scenery here, and so are Keen and the boss brain.
\dtMonster{67}   {600} {48} {8}  {80} {0}  {0}  {0}   % mancubus
\dtMonster{68}   {500} {64}{12} {128} {0}  {0}  {0}   % arachnotron
\dtMonster{16}  {4000} {40}{16}  {20} {0}  {0}  {0}   % cyberdemon
\dtMonster{7}   {3000}{128}{12}  {40} {0}  {0}  {3}   % spider mastermind

%%% The sequences, also from info.c.  The letters are written out for the
%%% reason \dtSpriteAnim records: Doom's runs are not all runs, and a picker
%%% that assumed base+n would ask for a lump that is not there.
%%
%% A wrong letter here does not crash.  \dtWantSprite asks for a picture that
%% does not exist, \dt@thingrot's rotation-0 fallback misses too, and the
%% thing drops to the coloured billboard -- wrong, visibly, and quietly.
%% That is why check.py counts the pictures the loader actually resolved
%% instead of trusting this table.
%% zombieman
\dtMonAnim{3004}{0}{AB}{10}     \dtMonAnim{3004}{1}{AABBCCDD}{4}
\dtMonAnim{3004}{2}{EF}{8}      \dtMonAnim{3004}{3}{G}{6}
\dtMonAnim{3004}{4}{HIJKL}{5}
%% shotgun guy
\dtMonAnim{9}{0}{AB}{10}        \dtMonAnim{9}{1}{AABBCCDD}{3}
\dtMonAnim{9}{2}{EF}{10}        \dtMonAnim{9}{3}{G}{6}
\dtMonAnim{9}{4}{HIJKL}{5}
%% chaingunner
\dtMonAnim{65}{0}{AB}{10}       \dtMonAnim{65}{1}{AABBCCDD}{3}
\dtMonAnim{65}{2}{EFEF}{4}      \dtMonAnim{65}{3}{G}{6}
\dtMonAnim{65}{4}{HIJKLMN}{5}
%% Wolfenstein SS
\dtMonAnim{84}{0}{AB}{10}       \dtMonAnim{84}{1}{AABBCCDD}{3}
\dtMonAnim{84}{2}{EFGF}{10}     \dtMonAnim{84}{3}{H}{6}
\dtMonAnim{84}{4}{IJKLM}{5}
%% imp
\dtMonAnim{3001}{0}{AB}{10}     \dtMonAnim{3001}{1}{AABBCCDD}{3}
\dtMonAnim{3001}{2}{EFG}{8}     \dtMonAnim{3001}{3}{H}{4}
\dtMonAnim{3001}{4}{IJKLM}{8}
%% demon and spectre share every frame
\dtMonAnim{3002}{0}{AB}{10}     \dtMonAnim{3002}{1}{AABBCCDD}{2}
\dtMonAnim{3002}{2}{EFG}{8}     \dtMonAnim{3002}{3}{H}{2}
\dtMonAnim{3002}{4}{IJKLMN}{8}
\dtMonAnim{58}{0}{AB}{10}       \dtMonAnim{58}{1}{AABBCCDD}{2}
\dtMonAnim{58}{2}{EFG}{8}       \dtMonAnim{58}{3}{H}{2}
\dtMonAnim{58}{4}{IJKLMN}{8}
%% cacodemon
\dtMonAnim{3005}{0}{A}{10}      \dtMonAnim{3005}{1}{A}{3}
\dtMonAnim{3005}{2}{BCD}{5}     \dtMonAnim{3005}{3}{EF}{3}
\dtMonAnim{3005}{4}{GHIJKL}{8}
%% baron and knight
\dtMonAnim{3003}{0}{AB}{10}     \dtMonAnim{3003}{1}{AABBCCDD}{3}
\dtMonAnim{3003}{2}{EFG}{8}     \dtMonAnim{3003}{3}{H}{2}
\dtMonAnim{3003}{4}{IJKLMNO}{8}
\dtMonAnim{69}{0}{AB}{10}       \dtMonAnim{69}{1}{AABBCCDD}{3}
\dtMonAnim{69}{2}{EFG}{8}       \dtMonAnim{69}{3}{H}{2}
\dtMonAnim{69}{4}{IJKLMNO}{8}
%% revenant
\dtMonAnim{66}{0}{AB}{10}       \dtMonAnim{66}{1}{AABBCCDD}{2}
\dtMonAnim{66}{2}{GH}{6}        \dtMonAnim{66}{3}{L}{5}
\dtMonAnim{66}{4}{MNOPQ}{7}
%% lost soul
\dtMonAnim{3006}{0}{AB}{10}     \dtMonAnim{3006}{1}{AB}{6}
\dtMonAnim{3006}{2}{CD}{6}      \dtMonAnim{3006}{3}{E}{3}
\dtMonAnim{3006}{4}{FGHIJ}{6}
%% mancubus.  Doom's fires three pairs of fireballs to a volley and this one
%% throws a single ball each time its attack comes round.
\dtMonAnim{67}{0}{AB}{15}       \dtMonAnim{67}{1}{AABBCCDDEEFF}{4}
\dtMonAnim{67}{2}{GHI}{8}       \dtMonAnim{67}{3}{J}{6}
\dtMonAnim{67}{4}{KLMNOPQRST}{6}
%% arachnotron
\dtMonAnim{68}{0}{AB}{10}       \dtMonAnim{68}{1}{AABBCCDDEEFF}{3}
\dtMonAnim{68}{2}{AGH}{4}       \dtMonAnim{68}{3}{I}{6}
\dtMonAnim{68}{4}{JKLMNOP}{7}
%% cyberdemon
\dtMonAnim{16}{0}{AB}{10}       \dtMonAnim{16}{1}{AABBCCDD}{3}
\dtMonAnim{16}{2}{EF}{9}        \dtMonAnim{16}{3}{G}{10}
\dtMonAnim{16}{4}{HIJKLMNOP}{10}
%% spider mastermind
\dtMonAnim{7}{0}{AB}{10}        \dtMonAnim{7}{1}{AABBCCDDEEFF}{3}
\dtMonAnim{7}{2}{AGH}{4}        \dtMonAnim{7}{3}{I}{6}
\dtMonAnim{7}{4}{JKLMNOPQRS}{10}

%%% ---- line of sight -------------------------------------------------------
%%% P_CheckSight, without the REJECT lump and without the BSP: does a
%%% straight line from (#1,#2) to (#3,#4) cross anything that stops sight?
%%%
%%% The arithmetic is \dt@segcross's with the two segments the other way
%%% round, and it has its own register set for the reason INTERNALS section 6
%%% gives: it is called from inside the thing loop, which is holding several
%%% of its own, and from inside the shot scan, which is holding more.
%%
%% A seg stops sight when its linedef is one-sided, or when the opening
%% between its two sectors has closed -- a shut door, a lift at the top --
%% or, since M21, when the sight line does not pass THROUGH that opening at
%% the height it arrives at.  That last test needs the z of both ends and is
%% set up just below.
\newcount\dtSEE
\newcount\dtSTI \newcount\dtSTA \newcount\dtSTB \newcount\dtSTC
\newcount\dtSTD \newcount\dtSTX \newcount\dtSTY \newcount\dtSTU
\newcount\dtSTV \newcount\dtSTP \newcount\dtSTQ \newcount\dtSTR
\newcount\dtSTS \newcount\dtSTK \newcount\dtSTH \newcount\dtSTL
\newcount\dtSTM \newcount\dtSTN \newcount\dtSTW
\newcount\dtSTZ \newcount\dtSTG \newcount\dtSTJ

%%% ---- the z half of P_CheckSight (M21) ------------------------------------
%%% The two ends of the sight line in world z.  Registers rather than two
%%% more arguments because \dtSight is called from five places, three of
%%% which are already carrying a tick's worth of state -- and because a
%%% SENTINEL is what makes adding them safe: \dtNOZ means "this caller has
%%% not said", and a caller that has not said gets exactly the M20 behaviour
%%% rather than a ray at z=0 that every raised floor in the map would stop.
%%% \dtSight puts the sentinel back when it is done, so a value cannot leak
%%% from one scan into the next.
\newcount\dtNOZ \dtNOZ=-2147483
\newcount\dtSEYE \newcount\dtSTGT
\dtSEYE=\dtNOZ \dtSTGT=\dtNOZ
\newif\ifdt@zsight \dt@zsighttrue
%%% Where a sight line starts and what it aims at, above the floor.  Doom
%%% uses `z + (height>>1) + 8' for the looker, which is about forty units up
%%% a fifty-six unit monster, and checks a RANGE of the target's body; one
%%% ray at the middle of it is the cheap reading of the same thing.  Counts,
%%% so a fixture can move them.
\newcount\dtMONEYE \dtMONEYE=40
\newcount\dtBODY   \dtBODY=28

\def\dtSight#1#2#3#4{%
  \dtSTX=\numexpr#1\relax \dtSTY=\numexpr#2\relax
  \dtSTU=\numexpr(#3)-\dtSTX\relax \dtSTV=\numexpr(#4)-\dtSTY\relax
  %% the sight segment's own bounding box, in coordinates relative to its
  %% first end.  That is what keeps every product below well inside 2^31 and
  %% rejects almost every seg in four comparisons.
  \dtSTP=\ifnum\dtSTU<0 \dtSTU\else0\fi\relax
  \dtSTQ=\ifnum\dtSTU>0 \dtSTU\else0\fi\relax
  \dtSTR=\ifnum\dtSTV<0 \dtSTV\else0\fi\relax
  \dtSTS=\ifnum\dtSTV>0 \dtSTV\else0\fi\relax
  \dtSEE=1 \dtSTI=0
  \dtWhile{\dtSTI<\NSEG}{\dt@sightseg\advance\dtSTI by1\relax }%
  \dtSEYE=\dtNOZ \dtSTGT=\dtNOZ}

\def\dt@sightseg{%
  \dtSTK=\csname dtGa\the\dtSTI\endcsname\relax
  \dtSTA=\numexpr\csname dtVx\the\dtSTK\endcsname-\dtSTX\relax
  \dtSTB=\numexpr\csname dtVy\the\dtSTK\endcsname-\dtSTY\relax
  \dtSTK=\csname dtGb\the\dtSTI\endcsname\relax
  \dtSTC=\numexpr\csname dtVx\the\dtSTK\endcsname-\dtSTX\relax
  \dtSTD=\numexpr\csname dtVy\the\dtSTK\endcsname-\dtSTY\relax
  \dtSTL=0
  \ifnum\dtSTA>\dtSTQ \ifnum\dtSTC>\dtSTQ \dtSTL=1 \fi\fi
  \ifnum\dtSTA<\dtSTP \ifnum\dtSTC<\dtSTP \dtSTL=1 \fi\fi
  \ifnum\dtSTB>\dtSTS \ifnum\dtSTD>\dtSTS \dtSTL=1 \fi\fi
  \ifnum\dtSTB<\dtSTR \ifnum\dtSTD<\dtSTR \dtSTL=1 \fi\fi
  \ifnum\dtSTL=0 \dt@sightcross \fi}

%%% The seg's two ends on opposite sides of the sight line, and then the
%%% sight line's two ends on opposite sides of the seg.  Both are needed:
%%% either alone is satisfied by segs the line merely passes near.  Touching
%%% and collinear cases fall through as "not crossed", which errs toward
%%% letting a monster see you and never toward a wall that is not there.
\def\dt@sightcross{%
  \dtSTK=\numexpr\dtSTU*\dtSTB-\dtSTV*\dtSTA\relax
  \dtSTL=\numexpr\dtSTU*\dtSTD-\dtSTV*\dtSTC\relax
  \ifnum\dtSTK=0 \else\ifnum\dtSTL=0 \else
    \dtSTH=0
    \ifnum\dtSTK<0
      \ifnum\dtSTL>0 \dtSTH=1 \fi
    \else
      \ifnum\dtSTL<0 \dtSTH=1 \fi
    \fi
    \ifnum\dtSTH=1 \dt@sightcrossB \fi
  \fi\fi}
\def\dt@sightcrossB{%
  \dtSTK=\numexpr(\dtSTC-\dtSTA)*(0-\dtSTB)-(\dtSTD-\dtSTB)*(0-\dtSTA)\relax
  \dtSTL=\numexpr(\dtSTC-\dtSTA)*(\dtSTV-\dtSTB)%
        -(\dtSTD-\dtSTB)*(\dtSTU-\dtSTA)\relax
  \ifnum\dtSTK=0 \else\ifnum\dtSTL=0 \else
    \dtSTH=0
    \ifnum\dtSTK<0
      \ifnum\dtSTL>0 \dtSTH=1 \fi
    \else
      \ifnum\dtSTL<0 \dtSTH=1 \fi
    \fi
    %% Kept for the z test: these two are the sight line's own ends measured
    %% against the seg, and their ratio is exactly how far along the sight
    %% line the crossing happens.  \dt@sightblock overwrites \dtSTK on its
    %% first line, so they are copied out here rather than recomputed.
    \dtSTZ=\dtSTK \dtSTG=\dtSTL
    \ifnum\dtSTH=1 \dt@sightblock \fi
  \fi\fi}

\def\dt@sightblock{%
  \dtSTK=\csname dtGl\the\dtSTI\endcsname\relax
  \dtSTM=\csname dtLF\the\dtSTK\endcsname\relax
  \dtSTN=\csname dtLB\the\dtSTK\endcsname\relax
  \ifnum\dtSTN<0
    \dtSEE=0
  \else
    \ifnum\dtSTM<0
      \dtSEE=0
    \else
      \dtSTW=\csname dtSc\the\dtSTM\endcsname\relax
      \ifnum\csname dtSc\the\dtSTN\endcsname<\dtSTW
        \dtSTW=\csname dtSc\the\dtSTN\endcsname\relax\fi
      \dtSTH=\csname dtSf\the\dtSTM\endcsname\relax
      \ifnum\csname dtSf\the\dtSTN\endcsname>\dtSTH
        \dtSTH=\csname dtSf\the\dtSTN\endcsname\relax\fi
      \ifnum\dtSTW>\dtSTH\else \dtSEE=0 \fi
      %% ...and Doom's other half, which was the piece M15 left out: the
      %% sight line has to pass through that opening AT THE HEIGHT IT
      %% ARRIVES AT.  Without it a monster sees you over a parapet, which is
      %% invisible on a map whose openings all reach the floor -- MAP01's do,
      %% which is why this went unnoticed for two milestones.
      %%
      %% \dtSTZ/(\dtSTZ-\dtSTG) is the fraction along the sight line, and
      %% both scalings are in \numexpr's a*b/c form so the intermediates are
      %% 64-bit: \dtSTZ alone can reach a map extent squared.
      \ifdt@zsight
        \ifnum\dtSEYE>\dtNOZ
          \ifnum\dtSEE=1
            \dtSTJ=\numexpr\dtSTZ*65536/(\dtSTZ-\dtSTG)\relax
            \dtSTJ=\numexpr\dtSEYE+(\dtSTGT-\dtSEYE)*\dtSTJ/65536\relax
            \ifnum\dtSTJ<\dtSTH \dtSEE=0 \fi
            \ifnum\dtSTJ>\dtSTW \dtSEE=0 \fi
          \fi
        \fi
      \fi
    \fi
  \fi
  %% run the loop counter out rather than testing for it, exactly as
  %% \dt@segopen ends the collision scan
  \ifnum\dtSEE=0 \dtSTI=\NSEG \fi}

%%% ---- spawning ------------------------------------------------------------
%%% Once, after the map loads, on the \dtFindDoors precedent.  A thing whose
%%% type is not in the table gets no dtOs and is never looked at again, which
%%% keeps every barrel, corpse and lamp out of the tick without a single
%%% exclusion -- the argument \dtItem's absent entry already makes.
%%
%% Guarded on dtOs so a second call is harmless, exactly as \dt@markdoor is
%% guarded on dtDT.  That is what lets a fixture drop a monster into the
%% built-in map from \DTHOOK and spawn it -- \DTHOOK runs after this does --
%% which is how usewalk.tex installs a door.
\def\dtSpawnMonsters{%
  \dtMBI=0
  \dtWhile{\dtMBI<\NTHING}{%
    \edef\dt@mty{\csname dtHt\the\dtMBI\endcsname}%
    \ifcsname dtOs\the\dtMBI\endcsname \else
    \ifcsname dtOq\dt@mty\endcsname
      \expandafter\xdef\csname dtOs\the\dtMBI\endcsname{0}%
      \expandafter\xdef\csname dtOf\the\dtMBI\endcsname{0}%
      \expandafter\xdef\csname dtOh\the\dtMBI\endcsname
        {\csname dtOq\dt@mty\endcsname}%
      %% Staggered by the thing's own index, exactly as \dtAnimIndex is:
      %% Doom perturbs the first tic count at spawn so two torches in a room
      %% do not flicker in lockstep, and the same argument holds for two
      %% zombiemen shifting their weight.  It also spreads A_Look's sight
      %% scans across frames instead of running every one of them together.
      \dtMBC=\dtMBI
      \divide\dtMBC by\csname dtOd\dt@mty @0\endcsname\relax
      \multiply\dtMBC by\csname dtOd\dt@mty @0\endcsname\relax
      \dtMBC=\numexpr\dtMBI-\dtMBC+1\relax
      \expandafter\xdef\csname dtOc\the\dtMBI\endcsname{\the\dtMBC}%
      %% Which sector it is standing in, which M20's noise flood needs and
      %% nothing before it did.  Settled here rather than in the loader so
      %% that a monster dropped into the built-in map from \DTHOOK gets one
      %% too -- one BSP descent per monster, once, and \dt@monstep keeps it
      %% up to date from then on.
      \dtPSX=\csname dtHx\the\dtMBI\endcsname\relax
      \dtPSY=\csname dtHy\the\dtMBI\endcsname\relax
      \dtPointInSSec
      \expandafter\xdef\csname dtHs\the\dtMBI\endcsname{\the\dtSECN}%
      \global\advance\dtNMON by1
    \fi\fi
    \advance\dtMBI by1\relax }%
  %% The spawn loop now runs a point locator per monster, and that writes the
  %% registers the PLAYER's position lives in -- \dt@monall's hazard, one
  %% level up and at load time.  It is invisible in the engine's own call,
  %% which is followed by \dtLocatePlayer anyway; it is not invisible to the
  %% dozen fixtures that call \dtSpawnMonsters from \DTHOOK, which runs after
  %% that call and would leave \dtSECN pointing at the last monster's room
  %% for the first \dtTryMove to measure a step against.
  \dtLocatePlayer
  \ifnum\dtNMON>0
    \dt@checkstates
    \message{[doom.tex] \the\dtNMON\space monsters}\fi}

%%% A monster that reaches a state with no declared sequence takes the frame
%%% picker apart, and the error it raises names \dt@anA rather than the table
%%% row that is short a line.  So the table is checked once, at load, against
%%% the types this map actually spawned -- one line naming the type and the
%%% state, in place of a hundred naming neither.
\def\dt@checkstates{%
  \dtMBI=0
  \dtWhile{\dtMBI<\NTHING}{%
    \ifcsname dtOs\the\dtMBI\endcsname
      \edef\dt@mty{\csname dtHt\the\dtMBI\endcsname}%
      \dtMBN=0
      \dtWhile{\dtMBN<5}{%
        \ifcsname dtOa\dt@mty @\the\dtMBN\endcsname\else
          \message{^^J[doom.tex] WARNING: monster type \dt@mty\space has no
            state \the\dtMBN\space sequence}%
        \fi
        \advance\dtMBN by1\relax }%
    \fi
    \advance\dtMBI by1\relax }}

%%% ---- which picture a monster is showing ----------------------------------
%%% \dt@thingkey calls this instead of \dt@animframe for anything carrying a
%%% state, because the state owns the frame letter.  Same trap as
%%% \dt@animframe and the same answer: the sequence has to reach \dt@anA as
%%% literal characters, so the call is built in an \edef with the walker
%%% \noexpand'd.  The space after \dt@ty is required and costs nothing -- @
%%% is a letter here, so `dtOa\dt@ty@0' would be one control word, and a
%%% space after a control word produces no token at all.
\def\dt@monframe{%
  \edef\dt@acall{\noexpand\dt@anA{\csname dtOf\the\dtTI\endcsname}%
    \csname dtOa\dt@ty @\csname dtOs\the\dtTI\endcsname\endcsname\relax}%
  \edef\dt@tk{\expandafter\dt@dropf\dt@tk\relax\dt@acall}}

%%% ---- the tick ------------------------------------------------------------
\def\dtThingTick{\ifdt@mon \ifnum\dtNMON>0 \dt@monall \fi\fi}

%%% \dtLocatePlayer at the end is not tidiness.  \dt@monstep locates the
%%% monster it has just moved, and the only point locator here writes
%%% \dtPSX/\dtPSY/\dtSSEC/\dtSECN -- the registers the player's own position
%%% lives in.  \dtTryMove reads \dtFloorZ, i.e. dtSf<\dtSECN>, on its first
%%% line, so leaving \dtSECN pointing at whichever room the last monster
%%% walked into would measure the player's next step up against a floor on
%%% the other side of the map.  One BSP descent puts it back.
\def\dt@monall{%
  \dtPLX=\numexpr\PX/65536\relax
  \dtPLY=\numexpr\PY/65536\relax
  \dtMBI=0
  \dtWhile{\dtMBI<\NTHING}{%
    \ifcsname dtOs\the\dtMBI\endcsname \dt@montick \fi
    \advance\dtMBI by1\relax }%
  \dtLocatePlayer}

%%% ---- who a monster is fighting (M20) -------------------------------------
%%% Until M20 there was one target in the world and every monster had it, so
%%% the tick read \dtPLX/\dtPLY throughout.  Doom's mobj carries a `target'
%%% pointer, and P_DamageMobj repoints it at whatever just hurt you -- which
%%% is the whole of infighting.  dtOt<i> is that pointer; absent or -1 is the
%%% player, which is exactly what every monster was doing implicitly.
%%%
%%% \dtMTX/\dtMTY are then what the chase, the sight test, the bearing and
%%% the attack all read.  They are NOT \dtPLX/\dtPLY and NOT \dtPSX/\dtPSY,
%%% and this file now has three pairs of position registers for three
%%% different questions -- see the note on \dt@monall for what the second one
%%% cost.
\newcount\dtMTX \newcount\dtMTY \newcount\dtMTI \newcount\dtMSAVE
\def\dt@montarget{%
  \dtMTI=-1
  \ifcsname dtOt\the\dtMBI\endcsname
    \dtMTI=\csname dtOt\the\dtMBI\endcsname\relax
  \fi
  %% A target that has died goes back to the player.  Doom's A_Chase does
  %% the same test at the top of every chase frame, and without it a monster
  %% spends the rest of the level walking toward a corpse.
  \ifnum\dtMTI>-1
    \ifcsname dtOs\the\dtMTI\endcsname
      \ifnum\csname dtOs\the\dtMTI\endcsname>3 \dtMTI=-1 \fi
    \else \dtMTI=-1 \fi
  \fi
  \expandafter\xdef\csname dtOt\the\dtMBI\endcsname{\the\dtMTI}%
  \ifnum\dtMTI<0
    \dtMTX=\dtPLX \dtMTY=\dtPLY
  \else
    \dtMTX=\csname dtHx\the\dtMTI\endcsname\relax
    \dtMTY=\csname dtHy\the\dtMTI\endcsname\relax
  \fi}

%% Is there anything left to fight?  A dead player stops the world; a monster
%% target has already been checked for by \dt@montarget.
\def\dt@montgtlive{\ifnum\dtMTI<0 \ifnum\dtHEALTH>0 1\else0\fi\else1\fi}

%% The target's FEET, which is where P_SpawnMissile aims.  The player's is
%% \PZ less the eye height; a monster's is the floor it is standing on, which
%% is what dtHz has held since \dt@monstep started writing it.
\def\dt@montgtz{\number\numexpr
  \ifnum\dtMTI<0 \PZ/65536-\VIEWH\else\csname dtHz\the\dtMTI\endcsname\fi
  \relax}

%% One attack's worth of damage, at whichever of the two it is aimed at.
%% \dtHurtThing loops on \dtMBI and rewrites \dt@mty, and this tick is
%% standing in both, so both are saved across the call rather than the call
%% being trusted -- the argument \dt@hitscan makes about borrowing \dtMBI,
%% from inside the loop that owns it.
\def\dt@monhit#1{%
  \ifnum\dtMTI<0
    \dtDamagePlayer{#1}{\the\dtMBX}{\the\dtMBY}%
  \else
    \dtMSAVE=\dtMBI \edef\dt@savety{\dt@mty}%
    \dtHURTBY=\dtMBI \dtMBI=\dtMTI
    \dtHurtThing{#1}%
    \dtHURTBY=-1 \dtMBI=\dtMSAVE \let\dt@mty\dt@savety
  \fi}

\def\dt@montick{%
  \dtMBS=\csname dtOs\the\dtMBI\endcsname\relax
  \ifnum\dtMBS>4 \else
    \edef\dt@mty{\csname dtHt\the\dtMBI\endcsname}%
    \dtMBX=\csname dtHx\the\dtMBI\endcsname\relax
    \dtMBY=\csname dtHy\the\dtMBI\endcsname\relax
    \dt@montarget
    \dt@monadv
    \ifcase\dtMBS \dt@monidle \or\dt@monchase \or\dt@monatk
      \or\dt@monhurt \or\dt@mondie \fi
  \fi}

%%% One tic off the frame showing now, and the next letter when it runs out.
%%% \dtMBG says the letter changed, which is the cadence Doom's A_Look runs
%%% at; \dtMBW says the sequence came round, which is all the one-shot states
%%% -- attack, pain, death -- need to know.
\def\dt@monadv{%
  \dtMBW=0 \dtMBG=0
  \dtMBC=\numexpr\csname dtOc\the\dtMBI\endcsname-1\relax
  \ifnum\dtMBC>0
    \expandafter\xdef\csname dtOc\the\dtMBI\endcsname{\the\dtMBC}%
  \else
    \dtMBG=1
    \dtMBN=\csname dtOn\dt@mty @\the\dtMBS\endcsname\relax
    \dtMBF=\numexpr\csname dtOf\the\dtMBI\endcsname+1\relax
    \ifnum\dtMBF<\dtMBN\else \dtMBF=0 \dtMBW=1 \fi
    \expandafter\xdef\csname dtOf\the\dtMBI\endcsname{\the\dtMBF}%
    \expandafter\xdef\csname dtOc\the\dtMBI\endcsname
      {\csname dtOd\dt@mty @\the\dtMBS\endcsname}%
  \fi}

\def\dt@monset#1{%
  \expandafter\xdef\csname dtOs\the\dtMBI\endcsname{#1}%
  \expandafter\xdef\csname dtOf\the\dtMBI\endcsname{0}%
  \expandafter\xdef\csname dtOc\the\dtMBI\endcsname
    {\csname dtOd\dt@mty @#1\endcsname}}

%%% Doom perturbs almost everything a monster does with P_Random.  \dtANIM,
%%% the thing's own index and a salt stand in for it, which is the argument
%%% the per-thing animation phase and the face's glance both make:
%%% deterministic, which the harness needs, and no table.  0..255, as Doom's
%%% is.  The salt is what lets one monster roll three different shells in one
%%% frame without touching the loop counter it is standing in.
\def\dt@monrand#1{%
  \dtMBR=\numexpr\dtANIM*97+\dtMBI*57+\dtMBI*\dtMBI*13+(#1)*29\relax
  \ifnum\dtMBR<0 \dtMBR=-\dtMBR \fi
  \dtMBQ=\dtMBR \divide\dtMBQ by256 \multiply\dtMBQ by256
  \advance\dtMBR by-\dtMBQ}

%%% P_AproxDistance: the larger axis plus half the smaller, within about 12%
%%% of the true distance and needing no square root.  Doom uses it for every
%%% range test in the game, so this is not a shortcut taken for TeX's sake.
%%% \divide, because \numexpr's division rounds.
\def\dt@mondist{%
  \dtMBA=\numexpr\dtMTX-\dtMBX\relax \ifnum\dtMBA<0 \dtMBA=-\dtMBA \fi
  \dtMBB=\numexpr\dtMTY-\dtMBY\relax \ifnum\dtMBB<0 \dtMBB=-\dtMBB \fi
  \ifnum\dtMBA<\dtMBB \dtMBD=\dtMBA \dtMBA=\dtMBB \dtMBB=\dtMBD \fi
  \dtMBD=\dtMBB \divide\dtMBD by2 \advance\dtMBD by\dtMBA}

%%% A_Look.  A monster is asleep until it can see the player.  Doom also
%%% wakes one that hears a gunshot; that half is \dtNoiseAlert, which reaches
%%% in from the firing code and does not come through here at all.
%%
%% Only on a frame change, which is where Doom's A_Look sits too: the spawn
%% loop holds each letter ten tics, so this is one sight scan per monster per
%% ten frames rather than a scan of 601 segs per monster per frame.
\def\dt@monidle{%
  \ifnum\dtMBG=1
    \ifnum\dtHEALTH>0
      \dtSEYE=\numexpr\csname dtHz\the\dtMBI\endcsname+\dtMONEYE\relax
      \dtSTGT=\numexpr\PZ/65536-\VIEWH+\dtBODY\relax
      \dtSight{\the\dtMBX}{\the\dtMBY}{\the\dtPLX}{\the\dtPLY}%
      \ifnum\dtSEE=1 \dt@monset{1}\fi
    \fi
  \fi}

%%% A_Chase: close, and attack when you can.  Melee first, as Doom does --
%%% something with fists at arm's length uses them rather than shooting.
\def\dt@monchase{%
  \ifnum\dt@montgtlive=1
    \dt@mondist
    \dtMBV=0
    \ifnum\csname dtOm\dt@mty\endcsname>0
      \ifnum\dtMBD<\dtMELEE \dtMBV=1 \fi
    \fi
    \ifnum\dtMBV=0 \dt@monshoot \fi
    \ifnum\dtMBV=1 \dt@monset{2}\else \dt@monmove \fi
  \fi}

%%% Should it open fire?  A hitscan monster, in range, with a clear line, and
%%% then a die roll -- Doom rolls too, or every zombieman in a room would
%%% fire on the same tic for ever.
%%
%% "Has a ranged attack" was "has a hitscan" until M18; a monster that throws
%% something has one too, and the decision either way is the same die and the
%% same sight test.  \dtMBK is free here: \dt@monmove sets its own copy after
%% this returns, and \dt@monbullets is two frames away.
\def\dt@monshoot{%
  \dtMBK=0
  \ifnum\csname dtOg\dt@mty\endcsname>0 \dtMBK=1 \fi
  \ifdt@misok \ifcsname dtJm\dt@mty\endcsname \dtMBK=1 \fi\fi
  \ifnum\dtMBK=1
    \ifnum\dtMBD<\dtMRANGE
      \dt@monrand{0}%
      \ifnum\dtMBR<40
        \dtSEYE=\numexpr\csname dtHz\the\dtMBI\endcsname+\dtMONEYE\relax
        \dtSTGT=\numexpr\dt@montgtz+\dtBODY\relax
        \dtSight{\the\dtMBX}{\the\dtMBY}{\the\dtMTX}{\the\dtMTY}%
        \ifnum\dtSEE=1 \dtMBV=1 \fi
      \fi
    \fi
  \fi}

%%% P_Move, with the direct bearing rather than Doom's eight compass
%%% directions: the picture is the same and it needs no per-thing movedir.
%%% The move goes through the same swept test the player's does, with the
%%% monster's own radius and the floor it is standing on -- which is what
%%% \dtMRAD, \dtMHGT and \dtMSELF exist for.
%%
%% The bearing is taken first and turned into a destination, and the angle is
%% copied out of \dtANG before the move test: \dtANG is \dtAtan's result
%% register and anything the move fires is free to ask for another bearing.
\def\dt@monmove{%
  \dtAtan{\dtMTX-\dtMBX}{\dtMTY-\dtMBY}%
  \dtMBJ=\dtANG
  \dtMBP=\csname dtOp\dt@mty\endcsname\relax
  \dtMBA=\numexpr\dtMBX+(\dtcos{\the\dtMBJ})*\dtMBP/65536\relax
  \dtMBB=\numexpr\dtMBY+(\dtsin{\the\dtMBJ})*\dtMBP/65536\relax
  \dtMBE=\dtMRAD \dtMBL=\dtMHGT \dtMBK=\dtFLZ
  \dtMRAD=\csname dtOr\dt@mty\endcsname\relax
  \dtMHGT=\PHEAD
  \dtFLZ=\csname dtHz\the\dtMBI\endcsname\relax
  \dtMSELF=\dtMBI
  \dtMoveTest{\the\dtMBX}{\the\dtMBY}{\the\dtMBA}{\the\dtMBB}%
  \dtMRAD=\dtMBE \dtMHGT=\dtMBL \dtFLZ=\dtMBK \dtMSELF=-1
  \ifnum\dtBLK=0 \dt@monstep \fi
  %% face the way it meant to go whether or not the step was taken, so one
  %% wedged in a doorway still turns to look at you
  \expandafter\xdef\csname dtHa\the\dtMBI\endcsname{\the\dtMBJ}}

%%% Committing the step also has to put the monster on the floor of wherever
%%% it now is, or one that walks up MAP01's staircase sinks into it.  This is
%%% the one place dtHz stops being the load-time constant INTERNALS.md records.
\def\dt@monstep{%
  \expandafter\xdef\csname dtHx\the\dtMBI\endcsname{\the\dtMBA}%
  \expandafter\xdef\csname dtHy\the\dtMBI\endcsname{\the\dtMBB}%
  \dtPSX=\dtMBA \dtPSY=\dtMBB
  \dtPointInSSec
  \expandafter\xdef\csname dtHz\the\dtMBI\endcsname
    {\csname dtSf\the\dtSECN\endcsname}%
  \expandafter\xdef\csname dtHl\the\dtMBI\endcsname
    {\csname dtSl\the\dtSECN\endcsname}%
  \expandafter\xdef\csname dtHs\the\dtMBI\endcsname{\the\dtSECN}}

%%% Doom fires on a NAMED frame of the attack rather than at the end of it:
%%% the zombieman's A_PosAttack sits on the second of its three.  Firing when
%%% the sequence reaches its last letter is that, counted from the other end
%%% -- the shot lands part way through the animation and what follows is the
%%% recovery.  Waiting for the wrap instead put the whole attack's worth of
%%% frames between the decision and the damage, which reads as a monster
%%% that raises its rifle and never pulls the trigger.
%%%
%%% \dtMBF is only written when the letter changed, so it is read only under
%%% \dtMBG -- otherwise it would still hold whichever monster last advanced.
\def\dt@monatk{%
  \ifnum\dtMBW=1
    \dt@monset{1}%
  \else
    \ifnum\dtMBG=1
      \ifnum\dtMBF=\numexpr\csname dtOn\dt@mty @2\endcsname-1\relax
        \dt@monfire
      \fi
    \fi
  \fi}

\def\dt@monfire{%
  \dt@mondist
  \ifnum\csname dtOm\dt@mty\endcsname>0
    \ifnum\dtMBD<\dtMELEE \dt@monmelee \fi
  \fi
  \ifnum\csname dtOg\dt@mty\endcsname>0
    \ifnum\dtMBD<\dtMRANGE \dt@monbullets \fi
  \fi
  %% A_TroopAttack and its relatives are one function with two arms: the
  %% fists inside melee range and the fireball outside it, never both.  That
  %% \else is the whole of the exclusion, and it is why the imp's melee dice
  %% were already in M15's table with nothing to sit beside them.
  \ifdt@misok
    \ifcsname dtJm\dt@mty\endcsname
      \ifnum\dtMBD<\dtMRANGE
        \ifnum\dtMBD<\dtMELEE \else \dt@monmissile \fi
      \fi
    \fi
  \fi}

%%% Doom spawns a monster's missile at its own position plus thirty-two and
%%% aims it at the target's FEET, so a fireball thrown from a ledge comes
%%% down at you.
\def\dt@monmissile{%
  \dtMBA=\dtMTX \dtMBB=\dtMTY
  \ifnum\dtMTI<0 \ifnum\dtINVIS>0 \dt@monblurmis \fi\fi
  \dtSpawnMissile{\csname dtJm\dt@mty\endcsname}%
    {\the\dtMBX}{\the\dtMBY}%
    {\number\numexpr\csname dtHz\the\dtMBI\endcsname+32\relax}%
    {\the\dtMBA}{\the\dtMBB}{\dt@montgtz}{\the\dtMBI}}

%%% ---- aiming at somebody who is partly invisible (M23) --------------------
%%% Doom's whole treatment of the blur sphere is one line in each of two
%%% places: A_FaceTarget and P_SpawnMissile both add (P_Random-P_Random)<<21
%%% or <<20 to the angle when the target carries MF_SHADOW.  In this engine's
%%% 8192 to a turn those shifts are a difference of two rolls times four, for
%%% a hitscan, and times two, for something thrown.
%%
%% A thrown thing is simply thrown the wrong way: the aim point is swung
%% round the thrower by that much, at the range the target is actually at.
%% \dtMBD is that range -- \dt@monfire took it on its first line.
\def\dt@monblurmis{%
  \dt@monrand{11}\dtMBP=\dtMBR
  \dt@monrand{12}\advance\dtMBP by-\dtMBR
  \dtAtan{\dtMTX-\dtMBX}{\dtMTY-\dtMBY}%
  \dtnorm\dtMBP{\dtANG+2*\dtMBP}%
  \dtMBA=\numexpr\dtMBX+(\dtcos{\the\dtMBP})*\dtMBD/65536\relax
  \dtMBB=\numexpr\dtMBY+(\dtsin{\the\dtMBP})*\dtMBD/65536\relax}
%% A hitscan here has no direction to get wrong -- one that passes the sight
%% test hits -- so the same swing is turned into how far wide of the player
%% it would have passed: range times the angle in radians, and 8192 units to
%% a turn is 1304 to the radian.  Inside the player's radius it still lands.
\def\dt@monblur{%
  \dt@monrand{11}\dtMBP=\dtMBR
  \dt@monrand{12}\advance\dtMBP by-\dtMBR
  \ifnum\dtMBP<0 \dtMBP=-\dtMBP \fi
  \ifnum\numexpr\dtMBD*\dtMBP*4/1304\relax<\PRAD \else \dtSEE=0 \fi}

%%% A_TroopAttack and its relatives: (P_Random % dice + 1) * mult.
\def\dt@monmelee{%
  \dt@monrand{1}%
  \dtMBQ=\dtMBR \divide\dtMBQ by\csname dtOe\dt@mty\endcsname\relax
  \multiply\dtMBQ by\csname dtOe\dt@mty\endcsname\relax
  \dtMBQ=\numexpr\dtMBR-\dtMBQ+1\relax
  \multiply\dtMBQ by\csname dtOm\dt@mty\endcsname\relax
  \dt@monhit{\the\dtMBQ}}

%%% A_PosAttack: each shot is ((P_Random % 5) + 1) * 3, and Doom sprays them
%%% about the aim.  There is no spread here -- a hitscan that has passed the
%%% sight test hits -- so a shotgun guy's three shells are three rolls of the
%%% same die, salted by the shot number.
%%
%% The sight test runs again here rather than being taken on trust from
%% \dt@monshoot two frames ago: the player has had the whole attack animation
%% to step behind something, and Doom re-checks for the same reason.
\def\dt@monbullets{%
  \dtSEYE=\numexpr\csname dtHz\the\dtMBI\endcsname+\dtMONEYE\relax
  \dtSTGT=\numexpr\dt@montgtz+\dtBODY\relax
  \dtSight{\the\dtMBX}{\the\dtMBY}{\the\dtMTX}{\the\dtMTY}%
  \ifnum\dtSEE=1 \ifnum\dtMTI<0 \ifnum\dtINVIS>0 \dt@monblur \fi\fi\fi
  \ifnum\dtSEE=1
    \dtMBZ=0 \dtMBK=\csname dtOg\dt@mty\endcsname\relax
    \dtMBN=0
    \dtWhile{\dtMBN<\dtMBK}{%
      \dt@monrand{\the\dtMBN}%
      \dtMBQ=\dtMBR \divide\dtMBQ by5 \multiply\dtMBQ by5
      \dtMBQ=\numexpr(\dtMBR-\dtMBQ+1)*3\relax
      \advance\dtMBZ by\dtMBQ
      \advance\dtMBN by1\relax }%
    \dt@monhit{\the\dtMBZ}%
  \fi}

%%% A flinch runs its frames and hands back to the chase.  Doom keeps the
%%% target through it and so does this: dtOt is not touched.
\def\dt@monhurt{\ifnum\dtMBW=1 \dt@monset{1}\fi}

%%% The death animation stops on its last frame rather than looping, and that
%%% frame is the corpse the map keeps.  State 5 is never ticked again, so a
%%% corpse costs one \ifcsname a frame and nothing else.
\def\dt@mondie{%
  \ifnum\dtMBW=1
    \expandafter\xdef\csname dtOs\the\dtMBI\endcsname{5}%
    \expandafter\xdef\csname dtOf\the\dtMBI\endcsname
      {\number\numexpr\csname dtOn\dt@mty @4\endcsname-1\relax}%
  \fi}

%%% ---- hurting a monster ---------------------------------------------------
%%% P_DamageMobj's monster half.  The caller has already set \dtMBI, because
%%% everything that can damage a monster is walking the thing list anyway.
%%%
%%% \dtHURTBY is who did it: -1 the player, otherwise a thing index.  Every
%%% caller sets it, and every caller sets it back, because the default has to
%%% stay "the player" for the four routes that are always the player.  It is
%%% what turns a monster caught in another's crossfire around, which is the
%%% whole of Doom's infighting -- there is no separate mechanism.
\newcount\dtHURTBY \dtHURTBY=-1
\def\dtHurtThing#1{%
  \edef\dt@mty{\csname dtHt\the\dtMBI\endcsname}%
  \dtMBH=\numexpr\csname dtOh\the\dtMBI\endcsname-(#1)\relax
  \ifnum\dtMBH>0
    \expandafter\xdef\csname dtOh\the\dtMBI\endcsname{\the\dtMBH}%
    %% Waking is not the same as flinching, and Doom does both here:
    %% "target = source" happens whatever the painchance roll says, so a
    %% monster shot in the back turns round even when it does not stagger.
    %%
    %% ...and "target = source" is literally that line: whatever hurt this
    %% thing is what it now wants.  A monster already fighting another one
    %% turns on the player the moment the player shoots it, which is Doom.
    \ifnum\dtHURTBY=\dtMBI \else
      \expandafter\xdef\csname dtOt\the\dtMBI\endcsname{\the\dtHURTBY}%
    \fi
    \ifnum\csname dtOs\the\dtMBI\endcsname=0 \dt@monset{1}\fi
    \dt@monrand{7}%
    \ifnum\dtMBR<\csname dtOw\dt@mty\endcsname\relax \dt@monset{3}\fi
  \else
    \expandafter\xdef\csname dtOh\the\dtMBI\endcsname{0}%
    \dt@monset{4}%
    \global\advance\dtKILLS by1
  \fi}

%%% ====================================================  DAMAGE TO THE PLAYER
%%% P_DamageMobj's player half.  Armour absorbs a third of a hit in a green
%%% jacket and a half in a blue one, and stops being armour when it runs out
%%% -- which is why the class is cleared here and not merely the points.
%%%
%%% \dtHURT is the size of the last hit and \dtHURTC how long the face goes
%%% on reacting to it; \dtATTL says the blow came from the player's left.
%%% All three exist for the status bar, and all three are what finally make
%%% the rest of the face reachable.
\newcount\dtHURT  \newcount\dtHURTC \newcount\dtATTL
\newcount\dtGRINC \newcount\dtFIREC \newcount\dtDEADC
\newcount\dtDMA \newcount\dtDMB \newcount\dtDMC
%% Doom's ST_TURNCOUNT, ST_EVILGRINCOUNT and ST_RAMPAGEDELAY.
\newcount\dtTURNC \dtTURNC=15
\newcount\dtEVILC \dtEVILC=70
\newcount\dtRAGEC \dtRAGEC=70

%%
%% Invulnerability is asked first and is the whole of the power: nothing is
%% absorbed, nothing is subtracted and the face does not flinch.  It is a
%% wrapper rather than one more \ifnum inside, so that the body below is the
%% body M15 wrote.
\def\dtDamagePlayer#1#2#3{%
  \ifnum\dtINVUL>0 \else \dt@damageplayer{#1}{#2}{#3}\fi}
\def\dt@damageplayer#1#2#3{%
  \ifnum\dtHEALTH>0
    \dtDMA=\numexpr#1\relax
    \ifnum\dtARMCL>0
      \dtDMB=2 \ifnum\dtARMCL=1 \dtDMB=3 \fi
      \dtDMC=\dtDMA \divide\dtDMC by\dtDMB
      \ifnum\dtDMC>\dtARMOR \dtDMC=\dtARMOR \fi
      \advance\dtARMOR by-\dtDMC
      \ifnum\dtARMOR<1 \dtARMCL=0 \fi
      \advance\dtDMA by-\dtDMC
    \fi
    \advance\dtHEALTH by-\dtDMA
    \ifnum\dtHEALTH<0 \dtHEALTH=0 \fi
    \dtHURT=\dtDMA \dtHURTC=\dtTURNC
    %% Which way to look.  \dtPLX/\dtPLY are set here rather than trusted
    %% from the caller for the reason \dtUse sets its own: this is reachable
    %% from a fixture as well as from the tick.  Increasing angle is
    %% anticlockwise and lateral is measured to the viewer's left, so a
    %% bearing less than half a turn ahead of the facing is on the left.
    \dtPLX=\numexpr\PX/65536\relax
    \dtPLY=\numexpr\PY/65536\relax
    \dtAtan{(#2)-\dtPLX}{(#3)-\dtPLY}%
    \dtnorm\dtAQ{\dtANG-\PA}%
    \dtATTL=0 \ifnum\dtAQ<4096 \dtATTL=1 \fi
  \fi}

%%% One tic of the player's own state: three face timers running down, and
%%% how long they have been dead.  Separate from \dtAnimTick because that one
%%% is a clock and this one is the player, and a fixture that freezes the
%%% world still wants to tell them apart.
%%% P_DeathThink's other half: the view sinks to six units above the floor.
%%% \dtFollowFloor already eases toward floor + \VIEWH and converges in two
%%% or three frames, so dying is that macro with a different eye height --
%%% and borrowing it rather than writing a second easing keeps a corpse's
%%% view settling at exactly the rate a living one does.
\newcount\dtDEADH \dtDEADH=6
\def\dtPlayerTick{%
  %% The weapon comes round, and the burst ends the frame the trigger is not
  %% pulled.  \dtHELD is cleared here rather than by whatever set it, because
  %% "nobody called \dtFire this frame" has no other place to be noticed.
  \ifnum\dtWCOOL>0 \advance\dtWCOOL by-1 \fi
  \ifnum\dtHELD=0 \dtREFIRE=0 \else \advance\dtREFIRE by1 \fi
  \dtHELD=0
  \ifnum\dtHURTC>0 \advance\dtHURTC by-1 \fi
  \ifnum\dtGRINC>0 \advance\dtGRINC by-1 \fi
  \ifnum\dtFIREC>0 \advance\dtFIREC by-1 \fi
  %% The four timed powers run down, and what they do to the view is settled
  %% for the frame about to be drawn.
  \ifnum\dtINVUL>0 \advance\dtINVUL by-1 \fi
  \ifnum\dtINVIS>0 \advance\dtINVIS by-1 \fi
  \ifnum\dtSUIT>0 \advance\dtSUIT by-1 \fi
  \ifnum\dtLAMP>0 \advance\dtLAMP by-1 \fi
  \dt@powerview
  \ifnum\dtHEALTH<1
    \advance\dtDEADC by1
    \dtDMB=\VIEWH \VIEWH=\dtDEADH \dtFollowFloor \VIEWH=\dtDMB
  \fi}

%%% =========================================================  PROJECTILES (M18)
%%% The first thing in this engine that exists without being on the map.
%%% Everything with a position so far was placed by THINGS and outlives the
%%% run; a missile is created by something that decided to throw it, flies for
%%% a few frames and is gone.  That is one new idea, and the rest of this
%%% section is the consequences of it.
%%%
%%% A missile IS a thing -- it goes in the same dtHx/dtHy/dtHz/dtHt arrays, so
%%% projection, sorting, occlusion and the sprite reader need to know nothing
%%% about it.  What it needs beyond a thing is a velocity and a state:
%%%
%%%   dtJs<i>   0 flying, 1 exploding, 2 spent -- the slot is free
%%%   dtJf<i>   which letter of the state's sequence is showing
%%%   dtJc<i>   tics left on that letter
%%%   dtJu<i> dtJv<i> dtJw<i>   velocity, map units per FRAME
%%%   dtJo<i>   who fired it: -1 the player, else a thing index
%%%
%%% and per TYPE, out of Doom's info.c:
%%%
%%%   dtJp<t>  speed (Doom's per-TIC number)   dtJd<t>  damage multiplier
%%%   dtJr<t>  radius        dtJb<t>  splash damage, 0 = none
%%%   dtJy<t>  1 if it sprays on impact (the BFG, and only the BFG)
%%%   dtJm<t>  what THING TYPE t puts into the world -- see below
%%%   dtJa<t>@<s> dtJn<t>@<s> dtJt<t>@<s>  state s's letters, count, tics
%%%   dtJk<t>@<s>  and the four-character sprite NAME of that state
%%%
%%% dtJk is the one array here with no counterpart in M15, and it is why a
%%% missile cannot borrow \dt@monframe: a monster wears one sprite name
%%% through all five of its states, and plasma flies as PLSS and explodes as
%%% PLSE.  Doom stores a sprite per state and never noticed the difference.
%%%
%%% Slots are RECYCLED.  \NTHING grows once, to \dtMAXMIS past the map's own
%%% thing count, and a spent missile's slot is handed to the next shot --
%%% otherwise a two-minute firefight would leave a thousand dead fireballs in
%%% every loop that walks the thing list.  dtHG cannot say "spent", because
%%% dtHG is permanent and this has to be undone; state 2 is that wait.
%%%
%%% \DTNOMIS puts the world back exactly as M15 left it: nothing is thrown,
%%% the monsters that would have thrown something use their fists as they did
%%% before, the three missile weapons go back to being inert, and no missile
%%% picture is read.  Every check written between M15 and M18 goes on meaning
%%% what it meant, which is the argument \DTNOMON already makes.
\newif\ifdt@misok \dt@misoktrue
\newcount\dtNMIS  \dtNMIS=0        % missiles in the air
%%% The map's own thing count, i.e. the first slot a missile may have.  It is
%%% \dtMISBASE and not \dtNTH0 because a control sequence is letters only:
%%% \dtNTH0 parses as \dtNTH followed by the character 0, so the declaration
%%% names one register, the assignment sets another value into it, and the
%%% -1 lands on the floor.  wad.tex records the same trap about \dtU16.
\newcount\dtMISBASE  \dtMISBASE=-1
\newcount\dtBOOMS \dtBOOMS=0       % explosions, for the harness
\newcount\dtMAXMIS \dtMAXMIS=24
%%% Doom splits a move longer than half MAXMOVE into steps so that a fast
%%% mobj cannot be on one side of a wall in one sample and the far side in
%%% the next.  Here the need is sharper: a frame is a whole 40-unit player
%%% step, so a plasma bolt covers sixty units in one and would cross most of
%%% MAP01's walls between two tests of where it is.
\newcount\dtMISSTEP \dtMISSTEP=16
%%% Doom's missile speeds are per TIC.  A frame here is a 40-unit player step
%%% against Doom's running 16 units a tic, so a frame is two and a half tics
%%% and this is the numerator of that ratio over two.  M15's note applies
%%% unchanged: what has to look like Doom is the RATIO of a missile's pace to
%%% the player's, not the number of units in the table -- and unlike a
%%% monster's speed, a missile's cannot simply be read as per-frame, because
%%% a fireball that travelled ten units a frame would never catch anybody.
\newcount\dtMISRATE \dtMISRATE=5
%%% Every missile in Doom is eight units tall, so this is a constant rather
%%% than a column of 8s in the table.
\newcount\dtMISHGT \dtMISHGT=8

\newcount\dtJI \newcount\dtJS \newcount\dtJF \newcount\dtJC
\newcount\dtJN \newcount\dtJA \newcount\dtJB \newcount\dtJD
\newcount\dtJE \newcount\dtJG \newcount\dtJH \newcount\dtJK
\newcount\dtJL \newcount\dtJM \newcount\dtJP \newcount\dtJQ
\newcount\dtJR \newcount\dtJU \newcount\dtJV \newcount\dtJW
\newcount\dtJX \newcount\dtJY \newcount\dtJZ \newcount\dtJO
\newcount\dtJT \newcount\dtJJ
%%% \dtJ -- no letter -- is wad.tex's thing-loop bound and is NOT one of
%%% these.  It runs at load and these run per frame, so they do not overlap
%%% today; the names are kept apart anyway, because INTERNALS section 6 records
%%% what one shared register cost.
%%%
%%% Three of these are load bearing across a whole tick and are worth naming:
%%% \dtJI is the missile being ticked, \dtJK the loop that found it, and
%%% \dtJJ every scan an explosion runs from inside that loop.  An explosion
%%% walks the thing list while the tick is walking the thing list, so the two
%%% cannot share -- the same collision \dt@monall documents between the
%%% player's position registers and a monster's.

%%% \dtMissile{type}{speed}{damage}{radius}{splash}{spray}
\def\dtMissile#1#2#3#4#5#6{%
  \expandafter\gdef\csname dtJp#1\endcsname{#2}%
  \expandafter\gdef\csname dtJd#1\endcsname{#3}%
  \expandafter\gdef\csname dtJr#1\endcsname{#4}%
  \expandafter\gdef\csname dtJb#1\endcsname{#5}%
  \expandafter\gdef\csname dtJy#1\endcsname{#6}}

%%% \dtMisAnim{type}{state}{sprite}{letters}{tics}.  State 0 also declares the
%%% type's dtSN, so the flight frame is the fallback the rest of the engine
%%% reaches for and the two can never disagree.
\def\dtMisAnim#1#2#3#4#5{%
  \dtMFC=0 \dt@ancount#4\relax
  \expandafter\gdef\csname dtJa#1@#2\endcsname{#4}%
  \expandafter\xdef\csname dtJn#1@#2\endcsname{\the\dtMFC}%
  \expandafter\gdef\csname dtJt#1@#2\endcsname{#5}%
  \expandafter\gdef\csname dtJk#1@#2\endcsname{#3}%
  \ifnum#2=0
    \edef\dt@next{\noexpand\dtSpriteType{#1}{#3\dt@anA{0}#4\relax}}\dt@next
  \fi}

%%% \dtMisFrom{thing type}{missile type} -- what this type of thing puts into
%%% the world.  A monster because it fires it; a WEAPON PICKUP because picking
%%% it up means the player will.  One array for both, because the loader wants
%%% exactly one question answered -- "if this thing is on the map, which
%%% missile pictures will be needed" -- and \dt@monfire only ever asks it
%%% about something carrying a state machine.
\def\dtMisFrom#1#2{\expandafter\gdef\csname dtJm#1\endcsname{#2}}

%%% ---- the table -----------------------------------------------------------
%%% Doom's own numbers.  The doomednum space runs to 89 and 2001..2049 with
%%% 3001..3006 for the monsters, so 9001 up is clear -- these types are never
%%% read out of a map and exist only to be handed to \dtSpawnMissile.
%%          type   spd dmg rad splash spray
\dtMissile{9001}  {10}  {3}  {6}   {0} {0}  % imp fireball        BAL1
\dtMissile{9002}  {10}  {5}  {6}   {0} {0}  % cacodemon           BAL2
\dtMissile{9003}  {15}  {8}  {6}   {0} {0}  % baron / knight      BAL7
\dtMissile{9004}  {10} {10} {11}   {0} {0}  % revenant            FATB
\dtMissile{9005}  {20} {20} {11} {128} {0}  % rocket              MISL
\dtMissile{9006}  {25}  {5} {13}   {0} {0}  % plasma              PLSS
\dtMissile{9007}  {25}{100} {13}   {0} {1}  % BFG ball            BFS1
\dtMissile{9008}  {20}  {8}  {6}   {0} {0}  % mancubus fireball   MANF
\dtMissile{9009}  {25}  {5} {13}   {0} {0}  % arachnotron plasma  APLS

\dtMisAnim{9001}{0}{BAL1}{AB}{4}    \dtMisAnim{9001}{1}{BAL1}{CDE}{6}
\dtMisAnim{9002}{0}{BAL2}{AB}{4}    \dtMisAnim{9002}{1}{BAL2}{CDE}{6}
\dtMisAnim{9003}{0}{BAL7}{AB}{4}    \dtMisAnim{9003}{1}{BAL7}{CDE}{6}
\dtMisAnim{9004}{0}{FATB}{AB}{2}    \dtMisAnim{9004}{1}{FBXP}{ABC}{6}
\dtMisAnim{9005}{0}{MISL}{A}{1}     \dtMisAnim{9005}{1}{MISL}{BCD}{6}
\dtMisAnim{9006}{0}{PLSS}{AB}{6}    \dtMisAnim{9006}{1}{PLSE}{ABCDE}{4}
\dtMisAnim{9007}{0}{BFS1}{AB}{4}    \dtMisAnim{9007}{1}{BFE1}{ABCDEF}{8}
\dtMisAnim{9008}{0}{MANF}{AB}{4}    \dtMisAnim{9008}{1}{MISL}{BCD}{6}
\dtMisAnim{9009}{0}{APLS}{AB}{5}    \dtMisAnim{9009}{1}{APBX}{ABCDE}{5}

%% monsters that throw one.  The imp, the cacodemon and the two bruisers keep
%% their melee dice as well -- Doom's A_TroopAttack uses the fists at melee
%% range and the fireball past it, and both halves were already in M15's
%% table waiting for this.
\dtMisFrom{3001}{9001}      % imp
\dtMisFrom{3005}{9002}      % cacodemon
\dtMisFrom{3003}{9003}      % baron of hell
\dtMisFrom{69}{9003}        % hell knight
\dtMisFrom{66}{9004}        % revenant -- and it homes; see \dt@mishome
\dtMisFrom{67}{9008}        % mancubus
\dtMisFrom{68}{9009}        % arachnotron
\dtMisFrom{16}{9005}        % cyberdemon: the player's own rocket
%% ...and the weapon pickups, which is the loader's half of the same question
\dtMisFrom{2003}{9005}      % rocket launcher
\dtMisFrom{2004}{9006}      % plasma rifle
\dtMisFrom{2006}{9007}      % BFG 9000

%%% The three weapons that throw rather than trace.  \dtGunPellets says zero
%%% for all three, which is what has kept them inert; this is the other half
%%% of that answer and the reason the zero can stay.
\def\dtGunMissile#1{\ifcase#1 0\or0\or0\or0\or9005\or9006\or9007\else0\fi}

%%% ---- which picture a missile is showing ----------------------------------
%%% \dt@monframe keeps the first four characters of the key and swaps the
%%% fifth, because a monster's sprite name never changes.  A missile's does,
%%% so this builds the whole key instead: the state's own name out of dtJk,
%%% then the letter its frame counter is on.  Same \edef discipline as
%%% \dt@monframe and for the same reason -- the sequence has to reach
%%% \dt@anA as literal characters, and the space after \dt@ty is what stops
%%% `dtJa\dt@ty@0' being read as a single control word.
\def\dt@misframe{%
  \edef\dt@acall{\noexpand\dt@anA{\csname dtJf\the\dtTI\endcsname}%
    \csname dtJa\dt@ty @\csname dtJs\the\dtTI\endcsname\endcsname\relax}%
  \edef\dt@tk{\csname dtJk\dt@ty @\csname dtJs\the\dtTI\endcsname\endcsname
    \dt@acall}}

%%% A missile's deterministic die, on the argument \dt@monrand makes: the
%%% animation clock, the missile's own slot, and a salt so one explosion can
%%% roll a different number for each thing it catches.
\def\dt@misrand#1{%
  \dtJR=\numexpr\dtANIM*89+\dtJI*43+\dtJI*\dtJI*17+(#1)*23\relax
  \ifnum\dtJR<0 \dtJR=-\dtJR \fi
  \dtJQ=\dtJR \divide\dtJQ by256 \multiply\dtJQ by256
  \advance\dtJR by-\dtJQ}

%%% ---- spawning ------------------------------------------------------------
%%% A free slot, or a new one.  \dtMISBASE is settled the first time something
%%% is thrown rather than at load, because \DTHOOK runs after the map is loaded
%%% and several fixtures use it to drop a thing into the built-in map -- a
%%% count taken before that would hand the first missile a monster's slot.
\def\dt@misslot{%
  \ifnum\dtMISBASE<0 \dtMISBASE=\NTHING \fi
  \dtJI=-1 \dtJJ=\dtMISBASE
  \dtWhile{\dtJJ<\NTHING}{%
    \ifnum\dtJI<0
      \ifnum\csname dtJs\the\dtJJ\endcsname=2 \dtJI=\dtJJ \fi
    \fi
    \advance\dtJJ by1\relax }%
  \ifnum\dtJI<0
    \ifnum\numexpr\NTHING-\dtMISBASE\relax<\dtMAXMIS
      \dtJI=\NTHING
      \dtAddThing{0}{0}{0}{0}%
      \expandafter\xdef\csname dtJs\the\dtJI\endcsname{2}%
    \fi
  \fi}

%%% P_SpawnMissile: from (#2,#3,#4) toward (#5,#6,#7), thrown by #8.
%%% The bearing is taken first and kept in \dtJZ, for the reason \dt@monmove
%%% records: \dtANG is only good until the next call of \dtAtan.
\def\dtSpawnMissile#1#2#3#4#5#6#7#8{%
  \ifdt@misok
    \dt@misslot
    \ifnum\dtJI<0 \else \dt@misspawn{#1}{#2}{#3}{#4}{#5}{#6}{#7}{#8}\fi
  \fi}
\def\dt@misspawn#1#2#3#4#5#6#7#8{%
  \dtAtan{(#5)-(#2)}{(#6)-(#3)}%
  \dtJZ=\dtANG
  %% Doom's per-tic speed as units per frame.  \divide, because this is the
  %% one place the ratio is applied and \numexpr would round it.
  \dtJP=\csname dtJp#1\endcsname\relax
  \multiply\dtJP by\dtMISRATE \divide\dtJP by2
  \dtJU=\numexpr(\dtcos{\the\dtJZ})*\dtJP/65536\relax
  \dtJV=\numexpr(\dtsin{\the\dtJZ})*\dtJP/65536\relax
  %% momz: rise over the run, at the speed it is flying.  P_AproxDistance
  %% again, and a floor of one so a target directly overhead cannot divide
  %% by zero.
  \dtJA=\numexpr(#5)-(#2)\relax \ifnum\dtJA<0 \dtJA=-\dtJA \fi
  \dtJB=\numexpr(#6)-(#3)\relax \ifnum\dtJB<0 \dtJB=-\dtJB \fi
  \ifnum\dtJA<\dtJB \dtJD=\dtJA \dtJA=\dtJB \dtJB=\dtJD \fi
  \dtJD=\dtJB \divide\dtJD by2 \advance\dtJD by\dtJA
  \ifnum\dtJD<1 \dtJD=1 \fi
  \dtJW=\numexpr((#7)-(#4))*\dtJP/\dtJD\relax
  %% Clear of whoever threw it.  P_CheckMissileSpawn nudges the missile half
  %% a tic forward and explodes it on the spot if that is already inside
  %% something; starting it outside the shooter's own radius is the same
  %% intent with no second move test.
  \dtJX=\numexpr(#2)+(\dtcos{\the\dtJZ})*\dt@misclear{#8}{#1}/65536\relax
  \dtJY=\numexpr(#3)+(\dtsin{\the\dtJZ})*\dt@misclear{#8}{#1}/65536\relax
  \expandafter\xdef\csname dtHx\the\dtJI\endcsname{\the\dtJX}%
  \expandafter\xdef\csname dtHy\the\dtJI\endcsname{\the\dtJY}%
  \expandafter\xdef\csname dtHz\the\dtJI\endcsname{\number\numexpr#4\relax}%
  \expandafter\xdef\csname dtHt\the\dtJI\endcsname{#1}%
  \expandafter\xdef\csname dtHa\the\dtJI\endcsname{\the\dtJZ}%
  \expandafter\xdef\csname dtHl\the\dtJI\endcsname{255}%
  \expandafter\xdef\csname dtJu\the\dtJI\endcsname{\the\dtJU}%
  \expandafter\xdef\csname dtJv\the\dtJI\endcsname{\the\dtJV}%
  \expandafter\xdef\csname dtJw\the\dtJI\endcsname{\the\dtJW}%
  \expandafter\xdef\csname dtJo\the\dtJI\endcsname{\number\numexpr#8\relax}%
  %% What it is following, if it follows anything.  Written for every missile
  %% rather than only the ones that home, because slots are recycled and a
  %% rocket must not inherit the target of the tracer that had its slot last.
  \expandafter\xdef\csname dtJg\the\dtJI\endcsname{\dt@mistracer{#1}{#8}}%
  \expandafter\xdef\csname dtJs\the\dtJI\endcsname{0}%
  \expandafter\xdef\csname dtJf\the\dtJI\endcsname{0}%
  \expandafter\xdef\csname dtJc\the\dtJI\endcsname{\csname dtJt#1@0\endcsname}%
  \global\advance\dtNMIS by1 }

%%% ---- the tracer (M23) ----------------------------------------------------
%%% A_Tracer.  The revenant's missile is the one thing in Doom that changes
%%% its mind in flight: every fourth tic it turns toward what it was thrown
%%% at, by at most 16.875 degrees, and leans its climb the same way.  At two
%%% and a half tics to a frame that is five eighths of a turn a frame, which
%%% is the 240 below.
%%%
%%%   dtJh<t>  present if missile type t homes
%%%   dtJg<i>  what missile i is following: -1 the player, a thing index, or
%%%            -2 for nothing at all
%%%
%%% The target is whatever the thrower was fighting when it threw, settled
%%% once -- Doom's `tracer = actor->target' -- so a revenant that turns on an
%%% imp afterwards does not take its missile with it.  \DTNOHOME leaves every
%%% dtJg at -2, which is the straight line M18 flew.
\newif\ifdt@home \dt@hometrue
\newcount\TRACERTURN \TRACERTURN=240
\def\dtMisHomes#1{\expandafter\gdef\csname dtJh#1\endcsname{1}}
\dtMisHomes{9004}
%% Expandable, because it is read inside the \xdef that stores it.  A
%% missile the player threw follows nothing: no weapon of Doom's homes.
\def\dt@mistracer#1#2{%
  \ifdt@home
    \ifcsname dtJh#1\endcsname
      \ifnum#2<0 -2\else
        \ifcsname dtOt#2\endcsname\csname dtOt#2\endcsname\else-1\fi
      \fi
    \else-2\fi
  \else-2\fi}

%% Where the target is now, or \dtJT=0 if it is dead -- in which case the
%% missile flies on the way it was going, as Doom's does.  Everything here
%% is scratch that \dt@misfly assigns afresh as soon as this returns, which
%% is why it runs first.
\def\dt@mishome{%
  \dtJH=\csname dtJg\the\dtJI\endcsname\relax
  \dtJT=1
  \ifnum\dtJH<0
    \ifnum\dtHEALTH<1 \dtJT=0 \fi
    \dtJA=\dtPLX \dtJB=\dtPLY
    \dtJD=\numexpr\PZ/65536-\VIEWH\relax
  \else
    \ifnum\csname dtOs\the\dtJH\endcsname>3 \dtJT=0 \fi
    \dtJA=\csname dtHx\the\dtJH\endcsname\relax
    \dtJB=\csname dtHy\the\dtJH\endcsname\relax
    \dtJD=\csname dtHz\the\dtJH\endcsname\relax
  \fi
  \ifnum\dtJT=1 \dt@misturn \fi}

%% Ahead of the move test, so the velocity that test reads is the new one.
\def\dt@misturn{%
  \dtJA=\numexpr\dtJA-\csname dtHx\the\dtJI\endcsname\relax
  \dtJB=\numexpr\dtJB-\csname dtHy\the\dtJI\endcsname\relax
  \dtAtan{\dtJA}{\dtJB}%
  \dtJZ=\csname dtHa\the\dtJI\endcsname\relax
  \dtnorm\dtJQ{\dtANG-\dtJZ}%
  \ifnum\dtJQ>4096 \advance\dtJQ by-8192 \fi
  \ifnum\dtJQ>\TRACERTURN \dtJQ=\TRACERTURN \fi
  \ifnum\dtJQ<-\TRACERTURN \dtJQ=-\TRACERTURN \fi
  \dtnorm\dtJZ{\dtJZ+\dtJQ}%
  \dtJP=\csname dtJp\dt@jty\endcsname\relax
  \multiply\dtJP by\dtMISRATE \divide\dtJP by2
  \expandafter\xdef\csname dtHa\the\dtJI\endcsname{\the\dtJZ}%
  \expandafter\xdef\csname dtJu\the\dtJI\endcsname
    {\number\numexpr(\dtcos{\the\dtJZ})*\dtJP/65536\relax}%
  \expandafter\xdef\csname dtJv\the\dtJI\endcsname
    {\number\numexpr(\dtsin{\the\dtJZ})*\dtJP/65536\relax}%
  %% ...and the climb: Doom aims forty units up the target and moves momz an
  %% eighth of a unit toward that slope each time.  One unit a frame here,
  %% over P_AproxDistance, with a floor of one so it cannot divide by zero.
  \ifnum\dtJA<0 \dtJA=-\dtJA \fi
  \ifnum\dtJB<0 \dtJB=-\dtJB \fi
  \ifnum\dtJA<\dtJB \dtJH=\dtJA \dtJA=\dtJB \dtJB=\dtJH \fi
  \dtJH=\dtJB \divide\dtJH by2 \advance\dtJH by\dtJA
  \ifnum\dtJH<1 \dtJH=1 \fi
  \dtJD=\numexpr(\dtJD+40-\csname dtHz\the\dtJI\endcsname)*\dtJP/\dtJH\relax
  \dtJQ=\csname dtJw\the\dtJI\endcsname\relax
  \ifnum\dtJD>\dtJQ \advance\dtJQ by1 \fi
  \ifnum\dtJD<\dtJQ \advance\dtJQ by-1 \fi
  \expandafter\xdef\csname dtJw\the\dtJI\endcsname{\the\dtJQ}}

%% The shooter's radius plus the missile's plus four.  The player is \PRAD and
%% has no dtOr; a monster's is the table's.
\def\dt@misclear#1#2{\number\numexpr
  \ifnum#1<0 \PRAD\else\csname dtOr\csname dtHt#1\endcsname\endcsname\fi
  +\csname dtJr#2\endcsname+4\relax}

%%% ---- the tick ------------------------------------------------------------
%%% \dtLocatePlayer at the end for exactly the reason \dt@monall carries one:
%%% putting a missile on its new sector runs the only point locator here, and
%%% that writes the registers the player's own position lives in.
\def\dtMissileTick{\ifnum\dtNMIS>0 \dt@misall \fi}
\def\dt@misall{%
  \dtPLX=\numexpr\PX/65536\relax
  \dtPLY=\numexpr\PY/65536\relax
  \dtJK=\dtMISBASE
  \dtWhile{\dtJK<\NTHING}{%
    \dtJI=\dtJK
    \dtJS=\csname dtJs\the\dtJI\endcsname\relax
    \ifcase\dtJS \dt@misfly \or\dt@misdying \fi
    \advance\dtJK by1\relax }%
  \dtLocatePlayer}

%%% One tic off the frame showing now.  \dtJM says the sequence came round,
%%% which is all the explosion needs to know; the flight frames simply loop.
\def\dt@misadv{%
  \dtJM=0
  \dtJC=\numexpr\csname dtJc\the\dtJI\endcsname-1\relax
  \ifnum\dtJC>0
    \expandafter\xdef\csname dtJc\the\dtJI\endcsname{\the\dtJC}%
  \else
    \dtJN=\csname dtJn\dt@jty @\the\dtJS\endcsname\relax
    \dtJF=\numexpr\csname dtJf\the\dtJI\endcsname+1\relax
    \ifnum\dtJF<\dtJN\else \dtJF=0 \dtJM=1 \fi
    \expandafter\xdef\csname dtJf\the\dtJI\endcsname{\the\dtJF}%
    \expandafter\xdef\csname dtJc\the\dtJI\endcsname
      {\csname dtJt\dt@jty @\the\dtJS\endcsname}%
  \fi}

%%% P_XYMovement and P_ZMovement, substepped.  The missile is walked toward
%%% its destination in steps of at most \dtMISSTEP, and each step is the same
%%% swept test the player's move goes through -- with \ifdt@mis on, so
%%% \dt@segopen asks whether the missile fits through the opening at its own
%%% height rather than whether it could climb the step.
\def\dt@misfly{%
  \edef\dt@jty{\csname dtHt\the\dtJI\endcsname}%
  \dt@misadv
  %% M23: a tracer turns before it moves, so the velocity read below is
  %% already the one it turned to.
  \ifnum\csname dtJg\the\dtJI\endcsname>-2 \dt@mishome \fi
  \dtJX=\csname dtHx\the\dtJI\endcsname\relax
  \dtJY=\csname dtHy\the\dtJI\endcsname\relax
  \dtJZ=\csname dtHz\the\dtJI\endcsname\relax
  \dtJU=\csname dtJu\the\dtJI\endcsname\relax
  \dtJV=\csname dtJv\the\dtJI\endcsname\relax
  \dtJW=\csname dtJw\the\dtJI\endcsname\relax
  \dtJA=\dtJU \ifnum\dtJA<0 \dtJA=-\dtJA \fi
  \dtJB=\dtJV \ifnum\dtJB<0 \dtJB=-\dtJB \fi
  \ifnum\dtJA<\dtJB \dtJA=\dtJB \fi
  \dtJN=\dtJA \divide\dtJN by\dtMISSTEP \advance\dtJN by1
  %% the mover's size and identity, saved and put back exactly as
  %% \dt@monmove does -- \dtMoveTest is shared with the player
  \dtJE=\dtMRAD \dtJG=\dtMHGT \dtJL=\dtMSELF
  \dtMRAD=\csname dtJr\dt@jty\endcsname\relax
  \dtMHGT=\dtMISHGT
  \dtMSELF=\csname dtJo\the\dtJI\endcsname\relax
  \dt@mistrue
  \dtJT=0 \dtJO=1 \dtJP=\numexpr\dtJN+1\relax
  \dtWhile{\dtJO<\dtJP}{\dt@misstep\advance\dtJO by1\relax }%
  \dt@misfalse
  \dtMRAD=\dtJE \dtMHGT=\dtJG \dtMSELF=\dtJL
  \expandafter\xdef\csname dtHx\the\dtJI\endcsname{\the\dtJX}%
  \expandafter\xdef\csname dtHy\the\dtJI\endcsname{\the\dtJY}%
  \expandafter\xdef\csname dtHz\the\dtJI\endcsname{\the\dtJZ}%
  %% \dtJT is 0 clear, 1 stopped by the world, 2 stopped by the player.  The
  %% floor and ceiling test only runs on a missile that is still flying, and
  %% it is another way of reaching 1.
  \ifnum\dtJT=0 \dt@missector \fi
  \ifnum\dtJT>0 \dt@misboom \fi}

%% One substep.  All THREE destinations are computed from where this frame
%% started rather than accumulated, so N substeps land exactly where one step
%% of the whole velocity would.  That is not tidiness: \numexpr rounds, so a
%% z velocity of -3 split over two substeps is -2 and -2, and an imp's
%% fireball descends a third faster than it was aimed -- it reaches the floor
%% a stride in front of the player and explodes against it, every time, with
%% nothing in the arithmetic obviously wrong.
\def\dt@misstep{%
  \ifnum\dtJT=0
    \dtJA=\numexpr\csname dtHx\the\dtJI\endcsname+\dtJU*\dtJO/\dtJN\relax
    \dtJB=\numexpr\csname dtHy\the\dtJI\endcsname+\dtJV*\dtJO/\dtJN\relax
    \dtJD=\numexpr\csname dtHz\the\dtJI\endcsname+\dtJW*\dtJO/\dtJN\relax
    \dtMISZ=\dtJD
    \dtMHIT=-1
    \dtMoveTest{\the\dtJX}{\the\dtJY}{\the\dtJA}{\the\dtJB}%
    \ifnum\dtBLK=1
      \dtJT=1
    \else
      \dt@mishitpl
      \ifnum\dtJT=0 \dtJX=\dtJA \dtJY=\dtJB \dtJZ=\dtJD \fi
    \fi
  \fi}

%% The player is not in \dt@mobjblock's scan -- that walks the things
%% carrying a state machine, and the player is not one of them.  Doom's box
%% test, on both axes and then on height, which is \dt@touch's shape with the
%% radii of a missile and a player rather than a pickup's reach.
%%
%% \dtJA/\dtJB/\dtJD are the substep's DESTINATION and the caller commits all
%% three the moment this returns, so nothing here may touch them; the missile
%% z it wants is in \dtMISZ, which is that destination already.  \dtJM/\dtJH/
%% \dtJQ are the three registers idle between \dt@misadv and the explosion.
\def\dt@mishitpl{%
  \ifnum\csname dtJo\the\dtJI\endcsname<0 \else
    \ifnum\dtHEALTH>0
      \dtJM=\numexpr\dtMRAD+\PRAD\relax
      \dtJH=\numexpr\dtPLX-\dtJA\relax \ifnum\dtJH<0 \dtJH=-\dtJH \fi
      \ifnum\dtJH<\dtJM
        \dtJQ=\numexpr\dtPLY-\dtJB\relax \ifnum\dtJQ<0 \dtJQ=-\dtJQ \fi
        \ifnum\dtJQ<\dtJM
          %% the player's feet, which is \PZ less the eye height -- the same
          %% pair \dt@touch measures a pickup against
          \dtJH=\numexpr\PZ/65536-\VIEWH\relax
          \ifnum\dtMISZ<\numexpr\dtJH+\PHEAD\relax
            \ifnum\numexpr\dtMISZ+\dtMISHGT\relax>\dtJH \dtJT=2 \fi
          \fi
        \fi
      \fi
    \fi
  \fi}

%% P_ZMovement's other half: the floor and the ceiling of wherever it now is.
%% One BSP descent per missile per frame, which also settles the sector it is
%% standing in -- and that is the descent \dtLocatePlayer above undoes.
\def\dt@missector{%
  \dtPSX=\dtJX \dtPSY=\dtJY
  \dtPointInSSec
  \ifnum\dtJZ<\csname dtSf\the\dtSECN\endcsname\relax \dtJT=1 \fi
  \ifnum\numexpr\dtJZ+\dtMISHGT\relax>\csname dtSc\the\dtSECN\endcsname\relax
    \dtJT=1 \fi
  \ifnum\dtJT=1
    \dtJZ=\csname dtSf\the\dtSECN\endcsname\relax
    \expandafter\xdef\csname dtHz\the\dtJI\endcsname{\the\dtJZ}%
  \fi}

%%% P_ExplodeMissile, with P_DamageMobj on whatever stopped it first.  Doom
%%% rolls (P_Random & 7) + 1 and multiplies by the type's damage, so a
%%% fireball is 3..24 and a rocket 20..160.
\def\dt@misboom{%
  \dt@misrand{0}%
  \dtJQ=\dtJR \divide\dtJQ by8 \multiply\dtJQ by8
  \dtJD=\numexpr(\dtJR-\dtJQ+1)*\csname dtJd\dt@jty\endcsname\relax
  %% Everything this explosion hurts was hurt BY whoever threw the missile,
  %% which is how a fireball that catches a zombieman in the back starts a
  %% fight the player is not in.  Set once here and put back at the end.
  \dtHURTBY=\csname dtJo\the\dtJI\endcsname\relax
  \ifnum\dtMHIT<0 \else
    %% \dtHurtThing reads \dtMBI, which is the monster tick's own loop
    %% counter.  That tick is not running here -- \dtMissileTick follows it
    %% -- so borrowing it is safe, and it is borrowed explicitly for the
    %% reason \dt@hitscan gives.
    \dtMBI=\dtMHIT
    \dtHurtThing{\the\dtJD}%
  \fi
  \ifnum\dtJT=2 \dtDamagePlayer{\the\dtJD}{\the\dtJX}{\the\dtJY}\fi
  \ifnum\csname dtJb\dt@jty\endcsname>0 \dt@misblast \fi
  \ifnum\csname dtJy\dt@jty\endcsname>0 \dt@misspray \fi
  \dtHURTBY=-1
  \expandafter\xdef\csname dtJu\the\dtJI\endcsname{0}%
  \expandafter\xdef\csname dtJv\the\dtJI\endcsname{0}%
  \expandafter\xdef\csname dtJw\the\dtJI\endcsname{0}%
  \expandafter\xdef\csname dtJs\the\dtJI\endcsname{1}%
  \expandafter\xdef\csname dtJf\the\dtJI\endcsname{0}%
  \expandafter\xdef\csname dtJc\the\dtJI\endcsname
    {\csname dtJt\dt@jty @1\endcsname}%
  \global\advance\dtBOOMS by1 }

%%% P_RadiusAttack.  Doom's falloff is linear in P_AproxDistance less the
%%% victim's radius, and it catches the player too -- which is what makes a
%%% rocket fired at the wall in front of you a mistake rather than a tactic.
\def\dt@misblast{%
  \dtJH=\csname dtJb\dt@jty\endcsname\relax
  \dtJJ=0
  \dtWhile{\dtJJ<\NTHING}{%
    \ifcsname dtOs\the\dtJJ\endcsname \dt@misblastone \fi
    \advance\dtJJ by1\relax }%
  %% and the player, at their own radius
  \dtJA=\numexpr\dtPLX-\dtJX\relax \ifnum\dtJA<0 \dtJA=-\dtJA \fi
  \dtJB=\numexpr\dtPLY-\dtJY\relax \ifnum\dtJB<0 \dtJB=-\dtJB \fi
  \ifnum\dtJA<\dtJB \dtJD=\dtJA \dtJA=\dtJB \dtJB=\dtJD \fi
  \dtJD=\dtJB \divide\dtJD by2 \advance\dtJD by\dtJA
  \advance\dtJD by-\PRAD
  \ifnum\dtJD<\dtJH
    \dtDamagePlayer{\number\numexpr\dtJH-\dtJD\relax}{\the\dtJX}{\the\dtJY}%
  \fi}
\def\dt@misblastone{%
  \ifnum\csname dtOs\the\dtJJ\endcsname>3 \else
    \dtJA=\numexpr\csname dtHx\the\dtJJ\endcsname-\dtJX\relax
    \ifnum\dtJA<0 \dtJA=-\dtJA \fi
    \dtJB=\numexpr\csname dtHy\the\dtJJ\endcsname-\dtJY\relax
    \ifnum\dtJB<0 \dtJB=-\dtJB \fi
    \ifnum\dtJA<\dtJB \dtJD=\dtJA \dtJA=\dtJB \dtJB=\dtJD \fi
    \dtJD=\dtJB \divide\dtJD by2 \advance\dtJD by\dtJA
    \advance\dtJD by-\csname dtOr\csname dtHt\the\dtJJ\endcsname\endcsname
    \ifnum\dtJD<\dtJH
      \dtMBI=\dtJJ
      \dtHurtThing{\number\numexpr\dtJH-\dtJD\relax}%
    \fi
  \fi}

%%% A_BFGSpray.  Doom fires forty tracers from the PLAYER, fanned across a
%%% quarter turn about the direction the ball was flying, and each one
%%% damages the first thing it meets for fifteen rolls of (P_Random&7)+1 --
%%% so the ball's own impact is the small half of a BFG shot, and standing
%%% behind cover saves you from one that went off in the open beside you.
%%%
%%% Until M23 this was an arc: everything inside it was caught, which on a
%%% crowd was more generous than Doom.  It is forty rays now, each one the
%%% aim scan a bullet already runs -- the nearest live thing the ray passes
%%% inside, then one sight test -- over Doom's own reach of 1024 units.  A
%%% monster standing behind another is shielded by it, and a wide one close
%%% up takes several rays, both of which are Doom.
%%%
%%% \dt@aim reads \dtPLX/\dtPLY and \vsin/\vcos and loops on \dtSHN; nothing
%%% is firing while a missile is being ticked, so all of that is free, and
%%% \dtJJ -- this section's own "scan run from inside the tick" -- counts the
%%% rays.  \dtGUNRANGE is borrowed for the reach and put back.
\newcount\dtBFGARC  \dtBFGARC=1024    % half the fan
\newcount\dtBFGRAYS \dtBFGRAYS=40
\newcount\dtBFGHITS \dtBFGHITS=0      % rays that found something, for the harness
\def\dt@misspray{%
  \dtJH=\csname dtHa\the\dtJI\endcsname\relax
  \dtJA=\dtGUNRANGE \dtGUNRANGE=1024
  \dtJJ=0
  \dtWhile{\dtJJ<\dtBFGRAYS}{%
    \dtnorm\dtJB{\dtJH-\dtBFGARC+\dtJJ*2*\dtBFGARC/\dtBFGRAYS}%
    \vsin=\dtsin{\the\dtJB}\vcos=\dtcos{\the\dtJB}%
    \dt@aim
    \ifnum\dtSHI<0 \else \dt@missprayhit \fi
    \advance\dtJJ by1\relax }%
  \dtGUNRANGE=\dtJA }
%% Fifteen rolls, salted by the ray and the roll so no two are the same
%% draw.  \dt@misrand writes \dtJR and \dtJQ, so the sum and the counter are
%% \dtJD and \dtJB -- the angle in \dtJB has been spent by now.
\def\dt@missprayhit{%
  \dtJD=0 \dtJB=0
  \dtWhile{\dtJB<15}{%
    \dt@misrand{\numexpr100+\dtJJ*15+\dtJB\relax}%
    \dtJQ=\dtJR \divide\dtJQ by8 \multiply\dtJQ by8
    \advance\dtJD by\numexpr\dtJR-\dtJQ+1\relax
    \advance\dtJB by1\relax }%
  \dtMBI=\dtSHI
  \dtHurtThing{\the\dtJD}%
  \global\advance\dtBFGHITS by1 }

%%% The explosion runs its frames and then the slot is free.  State 2 rather
%%% than dtHG, and the counter comes down here rather than at the impact, so
%%% \dtNMIS is "missiles the tick still has work for".
\def\dt@misdying{%
  \edef\dt@jty{\csname dtHt\the\dtJI\endcsname}%
  \dt@misadv
  \ifnum\dtJM=1
    \expandafter\xdef\csname dtJs\the\dtJI\endcsname{2}%
    \global\advance\dtNMIS by-1
  \fi}

%%% ====================================================  SOUND (M20)
%%% P_NoiseAlert.  Until now a monster woke on sight alone, which is half of
%%% Doom: firing a gun wakes everything the noise can reach, and the noise
%%% goes round corners a sight line cannot.
%%%
%%% P_RecursiveSound floods outward through the sectors, crossing any
%%% two-sided line that still has a gap in it -- so a shut door stops the
%%% sound exactly as it stops the shot -- and crossing at most one line
%%% marked ML_SOUNDBLOCK, which is the flag a mapper uses to keep a room
%%% quiet.  There is no distance limit at all; on MAP01 one shot wakes most
%%% of the level, and that is what Doom does.
%%%
%%%   dtSJ<s>  every two-sided line touching sector s, as {line}{other}
%%%            pairs ending {-1}{-1}.  Built once, because the flood is a
%%%            per-shot cost and walking every linedef per sector per shot
%%%            would be 59 x 370 on MAP01.
%%%   dtSA<s>  0 unvisited, else soundtraversed: 1 direct, 2 past a block
%%%   dtSQ/dtSR<n>  the queue: sector, and how many blocks it cost to reach
%%%   dtHs<i>  the sector thing i is standing in
\newif\ifdt@snd \dt@sndtrue
\newcount\dtNOISE \dtNOISE=0        % monsters woken by noise, for the harness
\newcount\dtNDH \newcount\dtNDT \newcount\dtNDI \newcount\dtNDS
\newcount\dtNDB \newcount\dtNDC \newcount\dtNDK \newcount\dtNDO
\newcount\dtNDP
\def\dtAddLineSound#1{\expandafter\gdef\csname dtLQ#1\endcsname{1}}

%%% Once, after the map loads, on the \dtFindDoors precedent.  A one-sided
%%% line joins nothing and a line whose two sides name the same sector is a
%%% self-loop, and both are dropped here rather than in the flood.
\def\dtBuildSound{%
  \dtNDI=0
  \dtWhile{\dtNDI<\NLINE}{%
    \dtNDS=\csname dtLF\the\dtNDI\endcsname\relax
    \dtNDO=\csname dtLB\the\dtNDI\endcsname\relax
    \ifnum\dtNDS<0 \else\ifnum\dtNDO<0 \else\ifnum\dtNDS=\dtNDO \else
      \dt@sndjoin{\the\dtNDS}{\the\dtNDO}%
      \dt@sndjoin{\the\dtNDO}{\the\dtNDS}%
    \fi\fi\fi
    \advance\dtNDI by1\relax }%
  %% the {-1}{-1} sentinel, written once each rather than carried through
  %% every append -- the same shape \dt@readpic ends a post chain with
  \dtNDI=0
  \dtWhile{\dtNDI<\NSEC}{%
    \edef\dt@sj{\ifcsname dtSJ\the\dtNDI\endcsname
      \csname dtSJ\the\dtNDI\endcsname\fi}%
    \expandafter\xdef\csname dtSJ\the\dtNDI\endcsname{\dt@sj{-1}{-1}}%
    \advance\dtNDI by1\relax }}
\def\dt@sndjoin#1#2{%
  \edef\dt@sj{\ifcsname dtSJ#1\endcsname\csname dtSJ#1\endcsname\fi}%
  \expandafter\xdef\csname dtSJ#1\endcsname{\dt@sj{\the\dtNDI}{#2}}}

%%% One shot's worth of noise.  \dtLocatePlayer at the end for the reason
%%% \dt@monall carries one: the flood starts by locating the player, and the
%%% only point locator here writes the registers the player's position lives
%%% in -- which is fine, except that it also wants \dtSSEC put back.
\def\dtNoiseAlert{%
  \ifdt@snd \ifnum\NSEC>0 \dt@sndgo \fi\fi}
\def\dt@sndgo{%
  \dtNDI=0
  \dtWhile{\dtNDI<\NSEC}{%
    \expandafter\xdef\csname dtSA\the\dtNDI\endcsname{0}%
    \advance\dtNDI by1\relax }%
  \dtPSX=\numexpr\PX/65536\relax \dtPSY=\numexpr\PY/65536\relax
  \dtPointInSSec
  \dtNDH=0 \dtNDT=0
  \dt@sndpush{\the\dtSECN}{1}%
  \dtWhile{\dtNDH<\dtNDT}{\dt@sndpop\advance\dtNDH by1\relax }%
  \dt@sndwake
  \dtLocatePlayer}

%% Doom's test is `validcount' plus soundtraversed <= soundblocks+1, which is
%% "not seen, or seen the long way round" -- a sector reached past a blocking
%% line and then reached directly has to be walked again, or the rooms beyond
%% it are never told.
\def\dt@sndpush#1#2{%
  \dtNDK=#1\relax \dtNDC=#2\relax
  \ifnum\dtNDK<0 \else
    \ifnum\csname dtSA\the\dtNDK\endcsname=0
      \dt@sndenq
    \else
      \ifnum\csname dtSA\the\dtNDK\endcsname>\dtNDC \dt@sndenq \fi
    \fi
  \fi}
\def\dt@sndenq{%
  \expandafter\xdef\csname dtSA\the\dtNDK\endcsname{\the\dtNDC}%
  \expandafter\xdef\csname dtSQ\the\dtNDT\endcsname{\the\dtNDK}%
  \expandafter\xdef\csname dtSR\the\dtNDT\endcsname{\the\dtNDC}%
  \advance\dtNDT by1\relax }

%% \dt@sj is expanded exactly once on the way in, for the reason
%% \dt@thingclip gives about \dt@sl.
\def\dt@sndpop{%
  \dtNDS=\csname dtSQ\the\dtNDH\endcsname\relax
  \dtNDB=\csname dtSR\the\dtNDH\endcsname\relax
  \ifcsname dtSJ\the\dtNDS\endcsname
    \edef\dt@sj{\csname dtSJ\the\dtNDS\endcsname}%
    \expandafter\dt@sndwalk\dt@sj
  \fi}
\def\dt@sndwalk#1#2{%
  \ifnum#1<0 \expandafter\dt@sndstop\else\expandafter\dt@sndone\fi{#1}{#2}}
\def\dt@sndstop#1#2{}
\def\dt@sndone#1#2{\dt@sndline{#1}{#2}\dt@sndwalk}

%% P_LineOpening, asked only whether there is a gap at all.  A shut door and
%% a lift at the top both close it, and both stop the sound -- which is the
%% one thing that makes this different from "wake the whole map".
\def\dt@sndline#1#2{%
  \dtNDO=\csname dtSc\csname dtLF#1\endcsname\endcsname\relax
  \ifnum\csname dtSc\csname dtLB#1\endcsname\endcsname<\dtNDO
    \dtNDO=\csname dtSc\csname dtLB#1\endcsname\endcsname\relax\fi
  \dtNDP=\csname dtSf\csname dtLF#1\endcsname\endcsname\relax
  \ifnum\csname dtSf\csname dtLB#1\endcsname\endcsname>\dtNDP
    \dtNDP=\csname dtSf\csname dtLB#1\endcsname\endcsname\relax\fi
  \ifnum\dtNDO>\dtNDP
    \ifcsname dtLQ#1\endcsname
      \ifnum\dtNDB<2 \dt@sndpush{#2}{2}\fi
    \else
      \dt@sndpush{#2}{\the\dtNDB}%
    \fi
  \fi}

%% Only a monster still in its spawn loop is woken -- one already chasing has
%% a target it chose and Doom does not take it away.  The target it wakes
%% with is the player, because the player is what made the noise.
\def\dt@sndwake{%
  \dtNDI=0
  \dtWhile{\dtNDI<\NTHING}{%
    \ifcsname dtOs\the\dtNDI\endcsname
      \ifnum\csname dtOs\the\dtNDI\endcsname=0
        \ifcsname dtHs\the\dtNDI\endcsname
          \ifnum\csname dtSA\csname dtHs\the\dtNDI\endcsname\endcsname>0
            \dtMBI=\dtNDI
            \edef\dt@mty{\csname dtHt\the\dtNDI\endcsname}%
            \expandafter\xdef\csname dtOt\the\dtNDI\endcsname{-1}%
            \dt@monset{1}%
            \global\advance\dtNOISE by1
          \fi
        \fi
      \fi
    \fi
    \advance\dtNDI by1\relax }}

%%% ==============================================================  FIRING (M15)
%%% P_LineAttack, with Doom's autoaim reduced to what a hitscan needs: a ray
%%% along the facing, the nearest live monster whose radius that ray passes
%%% inside, and one sight test to it so a wall between you stops the shot.
%%%
%%% This section is the HITSCAN weapons and the trigger all nine share.  The
%%% rocket launcher, the plasma rifle and the BFG throw things that fly:
%%% \dt@fireweap sends those to \dt@firemis, which spawns a missile and
%%% leaves the rest to the projectile section.
\newcount\dtSHI \newcount\dtSHA \newcount\dtSHB \newcount\dtSHC
\newcount\dtSHD \newcount\dtSHT \newcount\dtSHR \newcount\dtSHN
\newcount\dtSHK \newcount\dtSHP \newcount\dtSHQ \newcount\dtSHZ
\newcount\dtSHE \newcount\dtSHF \newcount\dtSHG \newcount\dtSHM
\newcount\dtSHV \newcount\dtSHW \newcount\dtSHY \newcount\dtSHL
\newcount\dtSHOTS \dtSHOTS=0
\newcount\dtHITS  \dtHITS=0
%%% Doom's MISSILERANGE again, this time as the length of the trace.
\newcount\dtGUNRANGE \dtGUNRANGE=2048
%%% Did the trigger actually fire?  \dtFIRED's argument in a new place: a
%%% shot refused for want of ammunition must not count, or the rampage face
%%% comes up on an empty pistol.  Its own register rather than \dtFIRED,
%%% which belongs to the linedef machinery and is written by \dtUse.
\newcount\dtSHOT

%%% Pellets and ammunition per shot, by weapon.  Doom's numbers: one bullet
%%% from the pistol and the chaingun, seven from the shotgun, twenty from the
%%% super shotgun for two shells.  The fist and the chainsaw swing rather
%%% than shoot, which is one "pellet" at arm's length.
%%% The rocket launcher, the plasma rifle and the BFG still say zero pellets,
%%% because they trace nothing; what they spend is real, and M18 is what made
%%% it worth spending.  Forty cells for a BFG shot is Doom's number and it is
%%% why the weapon has to be checked against the ammunition it needs rather
%%% than against one.
\def\dtGunPellets#1{\ifcase#1 1\or1\or7\or1\or0\or0\or0\or1\or20\else0\fi}
\def\dtGunUse#1{\ifcase#1 0\or1\or1\or1\or1\or1\or40\or0\or2\else0\fi}
\def\dtGunMelee#1{\ifcase#1 1\or0\or0\or0\or0\or0\or0\or1\else0\fi}

%%% ---- the weapon state machine (M21) --------------------------------------
%%% Doom's weapon is a state machine like everything else: P_FireWeapon puts
%%% it in its attack state and A_ReFire is what decides whether to go round
%%% again, so the trigger's rate is the length of that sequence.  Until M21
%%% there was none of this and \DTFIRE was one shot a frame, which made the
%%% chaingun and the BFG the same weapon.
%%%
%%% These are Doom's refire periods in TICS -- how often a held trigger comes
%%% round -- and they are turned into frames by \dtMISRATE's ratio, for M18's
%%% reason: what has to look like Doom is the rate against the player's own
%%% pace, and the player covers forty units a frame here against Doom's
%%% sixteen a tic.  Reading tics as frames directly would make every weapon
%%% two and a half times slower than the man carrying it.
\def\dtGunTics#1{\ifcase#1 14\or14\or30\or4\or20\or3\or40\or4\or57\else14\fi}
\def\dtGunDelay#1{\number\numexpr
  \ifnum\numexpr(\dtGunTics{#1})*2/5\relax<1 1\else(\dtGunTics{#1})*2/5\fi
  \relax}
%%% Doom's A_FireShotgun sprays its pellets about the aim with
%%% (P_Random - P_Random) << 18, which in this engine's 8192-to-a-turn is
%%% half that difference; the super shotgun uses << 19, i.e. twice as wide.
%%% The pistol and the chaingun spread too, but only once the trigger has
%%% already been round -- P_GunShot's `accurate' argument is !player->refire,
%%% which is why the first shot of a burst goes where you pointed it.
\def\dtGunSpread#1{\ifcase#1 0\or1\or1\or1\or0\or0\or0\or0\or2\else0\fi}
\def\dtGunAcc#1{\ifcase#1 0\or1\or0\or1\or0\or0\or0\or0\or0\else0\fi}
%%% Frames left before the weapon comes round, and how long the trigger has
%%% been held.  \DTNOWAIT is the off switch: one shot a frame, which is what
%%% every check written between M15 and M21 was measuring.
\newcount\dtWCOOL  \dtWCOOL=0
\newcount\dtREFIRE \dtREFIRE=0
\newcount\dtHELD   \dtHELD=0
\newif\ifdt@wsm \dt@wsmtrue

%%% A shot's deterministic die: the animation clock, how many shots have been
%%% fired, and which pellet.  \dtANIM alone would give every pellet of a
%%% shotgun blast identical damage; \dtSHOTS alone would make two shots in
%%% one frame identical.
\def\dt@gunrand#1{%
  \dtSHR=\numexpr\dtANIM*61+\dtSHOTS*151+(#1)*31\relax
  \ifnum\dtSHR<0 \dtSHR=-\dtSHR \fi
  \dtSHQ=\dtSHR \divide\dtSHQ by256 \multiply\dtSHQ by256
  \advance\dtSHR by-\dtSHQ}

%%% The trigger.  \dtHELD says it was pulled at all this frame, which is what
%%% \dtPlayerTick reads to decide whether the burst is still running -- a
%%% frame on cooldown is still a held trigger, and letting go is the only
%%% thing that makes the next shot accurate again.
\def\dtFire{%
  \dtSHOT=0 \dtHELD=1
  \ifnum\dtHEALTH>0
    \ifnum\dtWCOOL>0 \else \dt@fireweap \fi
  \fi
  \ifnum\dtSHOT=1 \ifdt@wsm
    \dtWCOOL=\dtGunDelay{\dtWEAP}\relax
  \fi\fi}

%%% Two kinds of weapon now, and the ammunition test is the half they share.
%%% A weapon that throws has no pellets at all, so the two branches are
%%% exclusive by construction rather than by an order of tests.
\def\dt@fireweap{%
  \dtSHK=\dtGunPellets{\dtWEAP}\relax
  \dtSHM=\dtGunMissile{\dtWEAP}\relax
  \dtSHT=\csname dtUT\the\dtWEAP\endcsname\relax
  \dtSHP=\dtGunUse{\dtWEAP}\relax
  \dtSHZ=1
  \ifnum\dtSHT<0 \else
    \ifnum\csname dtUA\the\dtSHT\endcsname<\dtSHP \dtSHZ=0 \fi
  \fi
  \ifnum\dtSHZ=1
    \ifnum\dtSHM>0
      \ifdt@misok \dt@firemis \fi
    \else
      \ifnum\dtSHK>0 \dt@firegun \fi
    \fi
  \fi}

%%% P_SpawnPlayerMissile.  Doom autoaims a player's missile VERTICALLY and
%%% leaves the bearing on the facing: you cannot look up, so the rocket has
%%% to climb to the imp on the ledge by itself.  The aim scan a bullet runs
%%% already answers "what is in front of me" -- over Doom's 1024 units for
%%% this, not a bullet's 2048 -- so the missile is thrown at a point that far
%%% along the facing and as high as the middle of whatever was found.  With
%%% nothing found, or under \DTNOAIM, it flies level as it did before M23.
%%%
%%% The bookkeeping is \dt@firegun's, line for line, and it is repeated
%%% rather than shared because what sits between the two halves -- a trace
%%% against a spawn -- is the whole of the difference.
%%
%% \dt@aim writes \dtSHA, \dtSHB and \dtSHC, so the launch height is worked
%% out after it and what it found is carried in \dtSHE and \dtSHF -- the two
%% dice \dt@hitscan sets afresh on every call and nothing here is rolling.
\newif\ifdt@aim \dt@aimtrue
\def\dt@firemis{%
  \expandafter\xdef\csname dtUA\the\dtSHT\endcsname
    {\number\numexpr\csname dtUA\the\dtSHT\endcsname-\dtSHP\relax}%
  \global\advance\dtSHOTS by1
  \dtSHOT=1 \dtFIREC=\dtRAGEC
  \dtNoiseAlert
  \dtPLX=\numexpr\PX/65536\relax
  \dtPLY=\numexpr\PY/65536\relax
  \vsin=\dtsin{\the\PA}\vcos=\dtcos{\the\PA}%
  \dtSHE=1024 \dtSHF=\dtNOZ
  \ifdt@aim
    \dtSHG=\dtGUNRANGE \dtGUNRANGE=1024
    \dt@aim
    \dtGUNRANGE=\dtSHG
    \ifnum\dtSHI<0 \else
      \dtSHE=\dtSHD
      \dtSHF=\numexpr\csname dtHz\the\dtSHI\endcsname+\dtBODY\relax
    \fi
  \fi
  \dtSHA=\numexpr\PZ/65536-\VIEWH+32\relax
  \ifnum\dtSHF=\dtNOZ \dtSHF=\dtSHA \fi
  \dtSpawnMissile{\the\dtSHM}{\the\dtPLX}{\the\dtPLY}{\the\dtSHA}%
    {\number\numexpr\dtPLX+\vcos*\dtSHE/65536\relax}%
    {\number\numexpr\dtPLY+\vsin*\dtSHE/65536\relax}%
    {\the\dtSHF}{-1}}

\def\dt@firegun{%
  \ifnum\dtSHT<0 \else
    \expandafter\xdef\csname dtUA\the\dtSHT\endcsname
      {\number\numexpr\csname dtUA\the\dtSHT\endcsname-\dtSHP\relax}%
  \fi
  \global\advance\dtSHOTS by1
  \dtSHOT=1 \dtFIREC=\dtRAGEC
  %% P_NoiseAlert, before the shot rather than after it: Doom's P_FireWeapon
  %% makes the noise as the weapon comes up, and a monster woken this frame
  %% has not moved yet either way.  It is here and not in \dtFire so that a
  %% trigger pulled with no ammunition makes no noise, which is the argument
  %% \dtSHOT already makes about the rampage face.
  \dtNoiseAlert
  %% \dtPLX/\dtPLY and the facing are settled here rather than trusted from
  %% \dtViewSetup, because firing is reachable from the key loop and from
  %% \DTHOOK as well as from the frame loop -- the rule \dtUse follows.
  %% AFTER \dtNoiseAlert, which runs \dtLocatePlayer of its own.
  \dtPLX=\numexpr\PX/65536\relax
  \dtPLY=\numexpr\PY/65536\relax
  \dtSHY=\PA
  \dt@hitscan}

%%% The aim scan.  Along and across the facing, both in map units: a monster
%%% is a candidate when it is in front, within range, and the ray passes
%%% inside its radius.  Nearest wins, and only then is the sight line tested
%%% -- one sight scan per shot rather than one per monster.
\def\dt@aim{%
  \dtSHI=-1 \dtSHD=\dtGUNRANGE
  \dtSHN=0
  \dtWhile{\dtSHN<\NTHING}{%
    \ifcsname dtOs\the\dtSHN\endcsname \dt@aimone \fi
    \advance\dtSHN by1\relax }%
  \ifnum\dtSHI<0 \else
    \dtSEYE=\numexpr\PZ/65536\relax
    \dtSTGT=\numexpr\csname dtHz\the\dtSHI\endcsname+\dtBODY\relax
    \dtSight{\the\dtPLX}{\the\dtPLY}%
      {\csname dtHx\the\dtSHI\endcsname}{\csname dtHy\the\dtSHI\endcsname}%
    \ifnum\dtSEE=0 \dtSHI=-1 \fi
  \fi}

\def\dt@aimone{%
  \ifnum\csname dtOs\the\dtSHN\endcsname>3
  \else
    \dtSHA=\numexpr\csname dtHx\the\dtSHN\endcsname-\dtPLX\relax
    \dtSHB=\numexpr\csname dtHy\the\dtSHN\endcsname-\dtPLY\relax
    \dtSHC=\numexpr(\dtSHA*\vcos+\dtSHB*\vsin)/65536\relax
    \ifnum\dtSHC>0
      \ifnum\dtSHC<\dtSHD
        \dtSHQ=\numexpr(\dtSHB*\vcos-\dtSHA*\vsin)/65536\relax
        \ifnum\dtSHQ<0 \dtSHQ=-\dtSHQ \fi
        %% \relax, because the bound is a macro expanding to digits and the
        %% line ends right after it: no space token, so the scan would run
        %% on into the branch.  \ifnum pushes the \relax back at the head of
        %% the branch, which is harmless here -- the branch is an assignment
        %% rather than a value being read (INTERNALS section 6).
        \ifnum\dtSHQ<\csname dtOr\csname dtHt\the\dtSHN\endcsname\endcsname
          \relax
          \dtSHI=\dtSHN \dtSHD=\dtSHC
        \fi
      \fi
    \fi
  \fi}

%%% Every pellet rolls its own damage, as P_GunShot does: 5*(P_Random%3+1)
%%% for a bullet and 2*(P_Random%10+1) for a fist.  And every pellet is aimed
%%% SEPARATELY, as Doom traces them: until M21 they all landed on the one
%%% thing a single aim scan picked, so the count and the damage were Doom's
%%% and where they went was not.  Seven pellets is seven aim scans and seven
%%% sight tests, and that is the price of the feature -- a shotgun costs
%%% seven times a pistol, and the pistol is what fires most of the time.
\def\dt@hitscan{%
  \dtSHE=3 \dtSHF=5
  \ifnum\dtGunMelee{\dtWEAP}=1 \dtSHE=10 \dtSHF=2 \fi
  %% \dtSHL, and NOT \dtSHN: \dt@aim loops on \dtSHN over the thing list, and
  %% it now runs once per pellet from inside this loop.  Sharing them is not
  %% a wrong answer, it is a run that never ends -- the aim scan resets the
  %% pellet counter to zero and leaves it at \NTHING, and with more pellets
  %% than things the bound is never reached.  INTERNALS section 6's register
  %% aliasing, in the one place a loop was nested inside another this week.
  \dtSHL=0
  \dtWhile{\dtSHL<\dtSHK}{\dt@pellet\advance\dtSHL by1\relax }}

%% Doom's own spread, and its own rule about when it applies.  \dt@gunrand is
%% salted twice per pellet because a single roll subtracted from itself is
%% zero: (P_Random - P_Random) needs two draws, and this is the salt argument
%% doing the job P_Random's own sequence does there.
\def\dt@pellet{%
  \dtSHW=0
  \ifnum\dtGunSpread{\dtWEAP}>0
    \dtSHZ=1
    \ifnum\dtGunAcc{\dtWEAP}=1 \ifnum\dtREFIRE=0 \dtSHZ=0 \fi\fi
    \ifnum\dtSHZ=1
      \dt@gunrand{\numexpr100+\dtSHL\relax}\dtSHW=\dtSHR
      \dt@gunrand{\numexpr200+\dtSHL\relax}\advance\dtSHW by-\dtSHR
      \dtSHW=\numexpr\dtSHW*\dtGunSpread{\dtWEAP}/2\relax
    \fi
  \fi
  \dtnorm\dtSHV{\dtSHY+\dtSHW}%
  \vsin=\dtsin{\the\dtSHV}\vcos=\dtcos{\the\dtSHV}%
  \dt@aim
  \ifnum\dtSHI<0 \else \dt@pellethit \fi
  %% M23: a line that is fired by shooting it, if the map has one at all.
  \ifnum\dtNGUNL>0 \dt@shootlines \fi}

\def\dt@pellethit{%
  \dtSHZ=1
  \ifnum\dtGunMelee{\dtWEAP}=1
    \ifnum\dtSHD<\dtMELEE\else \dtSHZ=0 \fi
  \fi
  \ifnum\dtSHZ=1
    \dt@gunrand{\the\dtSHL}%
    \dtSHQ=\dtSHR \divide\dtSHQ by\dtSHE \multiply\dtSHQ by\dtSHE
    \dtSHQ=\numexpr(\dtSHR-\dtSHQ+1)*\dtSHF\relax
    %% A_Punch under the berserk pack: ten times the damage, and the fist
    %% only -- the chainsaw is a different function and gets nothing.
    \ifnum\dtWEAP=0 \ifnum\dtBERSERK>0 \multiply\dtSHQ by10 \fi\fi
    %% \dtHurtThing reads \dtMBI.  The tick is not running here, so borrowing
    %% it is safe -- but it is borrowed explicitly rather than assumed,
    %% because that is the register the whole tick loops on.
    %%
    %% \dtHURTBY stays -1: this is the player shooting, and a monster the
    %% player shoots turns on the player whatever it was fighting before.
    \dtHURTBY=-1
    \dtMBI=\dtSHI
    \dtHurtThing{\the\dtSHQ}%
    \global\advance\dtHITS by1
  \fi}

%%% =========================================================  FRAME / PDF OUT
%%% This file uses TeX as a program, not a typesetter.  Any stray material
%%% reaching the main vertical list would trigger the page builder and emit
%%% blank pages between frames, so the output routine is neutralised and only
%%% explicit \shipout produces a page.
\output={\deadcycles=0 \setbox0=\box255 }
\topskip=0pt \maxdeadcycles=10000

\def\dtSetupPDF{%
  \pdfpagewidth=\numexpr\SCRW*\PIXSZ\relax bp
  \pdfpageheight=\numexpr\SCRH*\PIXSZ\relax bp
  \pdfhorigin=0bp \pdfvorigin=0bp
  \pdfcompresslevel=9
  \immediate\pdfobj{[/Indexed /DeviceRGB 255 <\DTPAL>]}%
  \edef\DTCS{\the\pdflastobj}%
  \edef\dt@res{/ColorSpace << /DoomPal \DTCS\space 0 R >>}%
  \pdfpageresources\expandafter{\dt@res}%
  \PPLANE=\numexpr\SCRW/2\relax
  \HALFH=\numexpr\SCRH/2\relax
  \dtPlaneShade
  \dtPlaneTables}

%%% "page" places the literal in absolute page coordinates.  Without it the
%%% box origin lands on a fractional coordinate and the per-column images
%%% leave unpainted seams between them.
\def\dtColumn#1#2{%
  \pdfliteral page{q \the\PIXSZ\space 0 0 \number\numexpr\SCRH*\PIXSZ\relax
    \space\number\numexpr(#1)*\PIXSZ\relax\space 0 cm
    BI /W 1 /H \the\SCRH\space /CS /DoomPal /BPC 8 /F /AHx ID
    #2> EI Q}}

%%% Columns are assembled by \dtBuildColumns and then spliced by
%%% \dtDrawThings, so emission is just a handoff.
\def\dtEmitColumn#1{\dtColumn{#1}{\csname dtFC#1\endcsname}}

%%% Shared by \dtRenderFrame and play.tex's key loop -- these two used to carry
%%% duplicate copies of this preamble, which is exactly how they would drift.
\def\dtViewSetup{%
  \dtnorm\PA{\PA}%
  \vsin=\dtsin{\the\PA}\vcos=\dtcos{\the\PA}%
  \dtPSX=\numexpr\PX/65536\relax
  \dtPSY=\numexpr\PY/65536\relax
  \dtPlaneFrame
  %% M19's drawseg list is rebuilt by the BSP walk, so it is emptied with the
  %% clip arrays and for the same reason.
  \global\dtNMASK=0
  \dtClearClip}

%%% The whole scene, short of putting it anywhere: walls, then assembly, then
%%% sprites.  Shared with the interactive loop.
\def\dtRenderScene{%
  \dtViewSetup
  \dtRenderBSP{\the\DTROOT}%
  \dtBuildColumns
  \dtDrawThings \dtDrawHUD}

\def\dtRenderFrame{%
  \dtRenderScene
  \setbox0=\hbox{}%
  \setbox0=\hbox{\cx=0
    \dtWhile{\cx<\SCRW}{\dtEmitColumn{\the\cx}\advance\cx by1\relax }}%
  \wd0=0pt \ht0=0pt \dp0=0pt
  \shipout\hbox{\box0}}

%%% ============================================================  LEVEL / MAIN
\PZ=\numexpr41*65536\relax
\PX=\numexpr100*65536\relax
\PY=\numexpr256*65536\relax
\PA=0

\expandafter\ifx\csname DTWAD\endcsname\relax
  \dtBuiltinMap
  \message{[doom.tex] built-in level: \the\NSEG\space segs,
    \the\NSSEC\space subsectors, \the\NNODE\space nodes}
\else
  \input wad.tex
  \expandafter\ifx\csname DTMAP\endcsname\relax\def\DTMAP{E1M1}\fi
  \expandafter\expandafter\expandafter\dtLoadWAD
    \expandafter\expandafter\expandafter{\expandafter\DTWAD\expandafter}%
    \expandafter{\DTMAP}%
\fi
%% wad.tex ends with \catcode`\@=12, so from here down \dt@foo would parse as
%% \dt followed by the characters @foo.  Everything below this line is in the
%% same file that set the letter catcode in the first place, so restore it.
%% The failure only appears with a WAD loaded, which is what makes it worth
%% closing at the source rather than by avoiding @ in the main section.
\catcode`\@=11

\dtFindDoors

%% \DTNOMON leaves every monster asleep and solid to nothing, i.e. the world
%% exactly as it was before M15.  It has to be read here rather than at the
%% top, because \dt@montrue is set with the rest of the enemy code.
%% Nothing is spawned at all under it, rather than spawned and left asleep:
%% a thing with no dtOs goes back through the per-type spawn loop in
%% \dt@thingkey, which is the animation every check written before M15 was
%% watching.  Spawning and freezing would have frozen those too.
\expandafter\ifx\csname DTNOMON\endcsname\relax\else\dt@monfalse\fi
\ifdt@mon \dtSpawnMonsters \fi

%% \DTNOMIS is M18's, on exactly the \DTNOMON precedent and read here for the
%% same reason: \dt@misoktrue is set with the rest of the projectile code.
%% Nothing is thrown under it, so a monster that would have thrown falls back
%% to the fists M15 gave it and the three missile weapons go back to being
%% inert -- the world every check between M15 and M18 was written against.
%% wad.tex reads it too, so the pictures are not read either.
\expandafter\ifx\csname DTNOMIS\endcsname\relax\else\dt@misokfalse\fi

%% \DTNOLOCK is M20's, and it is the smallest of these off switches because
%% the world it restores is the one every map without a locked door was
%% already in: six specials that opened for anybody now want a key, and this
%% is how a fixture written before they did goes on meaning what it meant.
\expandafter\ifx\csname DTNOLOCK\endcsname\relax\else\dt@lockfalse\fi

%% ...and \DTNOSND, which puts a monster back to waking on sight alone.
%% \dtBuildSound is skipped with it, so a fixture that sets it pays nothing
%% for a graph nothing will walk.
\expandafter\ifx\csname DTNOSND\endcsname\relax\else\dt@sndfalse\fi
\ifdt@snd \dtBuildSound \fi

%% M21's two.  \DTNOWAIT is one shot a frame, which is what \DTFIRE meant
%% from M15 to M20 and what a dozen checks are counting bullets against;
%% \DTNOZ drops the height half of the sight test, i.e. a monster that can
%% see you over a parapet, which is the world those same checks were
%% written in.
\expandafter\ifx\csname DTNOWAIT\endcsname\relax\else\dt@wsmfalse\fi
\expandafter\ifx\csname DTNOZ\endcsname\relax\else\dt@zsightfalse\fi

%% M23's four, each of them the world as M22 left it in one respect.
%% \DTNOSPEC: no sector special does anything -- no light strobes, no floor
%% burns, no wall scrolls and no secret is counted.  \DTNOPOW: the timed
%% powers stay on the floor and the berserk pack is a medikit with a fist.
%% \DTNOAIM: a player's missile flies level.  \DTNOHOME: a revenant's flies
%% straight.
\expandafter\ifx\csname DTNOSPEC\endcsname\relax\else\dt@specfalse\fi
\expandafter\ifx\csname DTNOPOW\endcsname\relax\else\dt@powfalse\fi
\expandafter\ifx\csname DTNOAIM\endcsname\relax\else\dt@aimfalse\fi
\expandafter\ifx\csname DTNOHOME\endcsname\relax\else\dt@homefalse\fi
\dtSpawnSpecials

%% Optional overrides, applied after the map load so they beat the WAD's
%% player start.  \DTX/\DTY are map units, \DTANG is in 8192ths of a turn.
\expandafter\ifx\csname DTX\endcsname\relax\else
  \PX=\numexpr\DTX*65536\relax\fi
\expandafter\ifx\csname DTY\endcsname\relax\else
  \PY=\numexpr\DTY*65536\relax\fi
\expandafter\ifx\csname DTANG\endcsname\relax\else
  \dtnorm\PA{\DTANG}\fi

\dtSetupPDF
%% Stand on the floor of the sector actually containing the start position.
%% This used to be sector 0's floor, which is only right by luck: on a real
%% map the player start is rarely in sector 0.  (MAP01 of doom2.wad starts in
%% sector 39, floor 0, while sector 0's floor is 104.)
\dtLocatePlayer
\dtSnapFloor
\message{[doom.tex] start: sector \the\dtSECN, floor \dtFloorZ,
  eye \number\numexpr\PZ/65536\relax}

%% \DTUSE presses "use" every frame, so a scripted walk opens the doors it
%% reaches.  Interactive mode has the e key instead, which taps rather than
%% holds -- see the usedown note on \dt@usedoor.
\expandafter\ifx\csname DTUSE\endcsname\relax
  \def\dtUseStep{}
\else
  \def\dtUseStep{\dt@useholdtrue\dtUse\dt@useholdfalse}
\fi

%% \DTFIRE holds the trigger down, which is what a scripted walk needs to
%% shoot anything.  How often that actually fires is the weapon's own
%% business -- \dtWCOOL, out of Doom's refire tics -- unless \DTNOWAIT has
%% put it back to one shot a frame.
\expandafter\ifx\csname DTFIRE\endcsname\relax
  \def\dtFireStep{}
\else
  \def\dtFireStep{\dtFire}
\fi

%% \DTHOOK runs once, after the map and the engine are fully loaded and before
%% any frame is rendered.  Anything after \input doom.tex is dead code -- the
%% file ends with \end -- so this is the only place a caller can get in.
\expandafter\ifx\csname DTHOOK\endcsname\relax\else\DTHOOK\fi

%% \DTTURN rotates by that many angle units after every step.  Combined with
%% the scripted walk this drags the player around every wall in the map at all
%% angles, which is how the tunnelling bug below was pinned down.
\expandafter\ifx\csname DTTURN\endcsname\relax
  \def\dtTurnStep{}
\else
  \def\dtTurnStep{\dtnorm\PA{\PA+\DTTURN}}
\fi

\expandafter\ifx\csname DTPLAY\endcsname\relax
  \message{[doom.tex] rendering \the\NFRAMES\space frames at
    \the\SCRW x\the\SCRH...}
  \dtFRAME=0
  \dtWhile{\dtFRAME<\NFRAMES}{%
    \dtRenderFrame
    \dtUseStep
    \dtFireStep
    \dtSectorTick
    \dtThingTick
    %% After the monsters, because a missile thrown this frame should not
    %% also move this frame -- Doom's thinker list has the same property for
    %% anything spawned during a tick -- and because \dt@misboom borrows
    %% \dtMBI, which the monster tick above is the owner of.
    \dtMissileTick
    \dtPlayerTick
    \dtAnimTick
    \dtStepFwd{2621440}%
    \dtTurnStep
    %% Reaching the exit ends the run, which is what makes it visible: a
    %% level that finished is a different thing from one that ran out of
    %% frames, and the frame count in the log is the difference.
    \ifnum\dtEXIT>0 \dtFRAME=\NFRAMES \fi
    \advance\dtFRAME by1
    \message{frame \the\dtFRAME\space at \number\numexpr\PX/65536\relax,%
      \number\numexpr\PY/65536\relax\space sec \the\dtSECN\space
      floor \dtFloorZ\space eye \number\numexpr\PZ/65536\relax\space
      hp \the\dtHEALTH\space ar \the\dtARMOR\space am \dtAmmoNow\space
      wp \the\dtWEAP\space got \the\dtNPICK\space
      kill \the\dtKILLS\space of \the\dtNMON}}
\else
  \input play.tex
  \dtKeyLoop
\fi
%% What the level came to, in the one line Doom's intermission screen says
%% it in.  play.tex ends with \catcode`\@=12 for the reason wad.tex does, so
%% nothing below this point may name a control sequence with an @ in it.
\message{^^J[doom.tex] tally: kills \the\dtKILLS/\the\dtNMON,
  items \the\dtNPICK, secrets \the\dtSECRETS/\the\dtNSECRET
  \ifnum\dtEXIT>0 , level complete\fi
  \ifnum\dtHEALTH<1 , and you are dead\fi}
\message{[doom.tex] done.}
\csname @@end\endcsname\end
